\documentclass[11pt,oneside]{report}
\usepackage[a4paper,margin=24mm]{geometry}
\usepackage[T1]{fontenc}
\usepackage{lmodern}
\usepackage{amsmath,amssymb,amsthm,mathtools,booktabs,longtable,array,microtype}
\usepackage{hyperref,xurl}
\hypersetup{hidelinks,pdftitle={Written proofs for J05, J06, J07, J21, J10, J15, and J14}}
\newtheorem{theorem}{Theorem}[section]
\newtheorem{lemma}[theorem]{Lemma}
\newtheorem{proposition}[theorem]{Proposition}
\newtheorem{corollary}[theorem]{Corollary}
\theoremstyle{definition}
\newtheorem{definition}[theorem]{Definition}
\numberwithin{equation}{section}
\newcommand{\F}{\mathbb F_2}
\newcommand{\C}{\mathbb C}
\newcommand{\Z}{\mathbb Z}
\newcommand{\nat}{\mathbb N}
\newcommand{\decl}[1]{\textnormal{\nolinkurl{#1}}}
\newcommand{\NE}{\operatorname{NE}}
\newcommand{\SW}{\operatorname{SW}}
\newcommand{\CP}{\operatorname{CP}}
\newcommand{\seamcardop}{\operatorname{card}}
\newcommand{\seamcard}[1]{\lvert #1\rvert}
\newcommand{\Lp}{L_+}\newcommand{\Lm}{L_-}
\newcommand{\Hp}{H_+}\newcommand{\Hm}{H_-}
\newcommand{\cP}{\mathcal P}
\newcommand{\cR}{\mathcal R}
\newcommand{\sswap}{\mathbin{\star}}
\newcommand{\sgn}{\operatorname{sgn}}
\newcommand{\Paths}{\mathcal P}
\newcommand{\Mov}{\mathrel{\stackrel{\mathrm{moves}}{\sim}}}
\newcommand{\OO}{\mathrm O}
\newcommand{\XX}{\mathrm X}
\newcommand{\ord}{\operatorname{ord}}
\newcommand{\Disc}{\operatorname{Disc}}
\newcommand{\Res}{\operatorname{Res}}
\newcommand{\wind}{\operatorname{wind}}
\newcommand{\Exp}{\operatorname{exp}}
\newcommand{\length}{\operatorname{length}}
\newcommand{\ii}{\mathrm i}
\newcommand{\Conf}{\operatorname{Conf}}
\newcommand{\Ree}{\operatorname{Re}}
\newcommand{\Imm}{\operatorname{Im}}
\setlength{\emergencystretch}{2em}
\setlength{\LTpre}{8pt}
\setlength{\LTpost}{8pt}
\title{Written proofs for seven native obligations\\
\large J05, J06, J07, J21, J10, J15, and J14\\
Updated consolidated document}
\author{}
\date{1 November 2026}
\begin{document}
\hypersetup{pageanchor=false}
\maketitle
\hypersetup{pageanchor=true}
\chapter*{Exact scope, completed proofs, and source limits}
\addcontentsline{toc}{chapter}{Exact scope, completed proofs, and source limits}

This standalone LaTeX source contains the complete written arguments for
the seven requested jobs, with their exact hypotheses, proofs, source
ledgers, and exceptional cases. It supersedes the mathematical status
of the separately preserved earlier snapshot. That snapshot has not
been modified. No Lean code is included.

The principal new seam result is stronger than the requested parity:
for every pair of seam sizes and every pair of retained states,
\[
 \decl{SeamCorrectionPath}\ k\ p\ q=\varnothing
 \qquad(k\geq3).
\]
If both retained factors change, the path type is empty at every
internal length, including zero. The proof uses signed rectangle
domains, nonnegative factor projections, and a reduction to at most
five movable points. The finite table after that proved reduction
lists every possible transition and its blocking points. No finite
diagnostic is extrapolated to arbitrary sizes.

The J21 chapter retains the completed explicit second representative,
including the correction chain, every boundary coefficient, the exact
anchor-class identification, and the second-vertex support bound.
The J14 and J15 chapters now include the missing root-covering
construction and the comparison with every positive full-twist word
whose closure is isotopic to the curve link.

Two distinctions must remain explicit. First, the displayed M405 API
omits the multiplicity/transversality hypothesis in the user's
geometric statement and is false. The J14 chapter proves a
counterexample to that exact API and proves the full word theorem
under the user's stated transverse hypotheses. Second, the literal
J10 selected-grid instantiation requires two field equations and the
primitive move definitions whose source files are absent from the
combined attachment. Its complete switch-and-move proof is supplied
for the exact marking equations that are reproduced, and the missing
native implications are stated individually. They are not silently
assumed.

\begin{longtable}{@{}>{\raggedright\arraybackslash}p{0.09\linewidth}>{\raggedright\arraybackslash}p{0.52\linewidth}>{\raggedright\arraybackslash}p{0.33\linewidth}@{}}
\textbf{Job}&\textbf{Statement established}&\textbf{Status or precise limit}\\\hline
\endfirsthead
\textbf{Job}&\textbf{Statement established}&\textbf{Status or precise limit}\\\hline
\endhead
J05&Every fixed-second correction of internal length at least three is
absent. The long count is the integer zero.&Complete. M602--M604 are
proved, with the arbitrary-endpoint M603 also covered by the combined
truncation theorem.\\[5pt]
J06&The one- and two-transition paths are in explicit bijection with
the remaining second-factor rectangles. All longer paths are absent.
The count equality holds in natural numbers.&Complete. M605 is proved.\\[5pt]
J07&Every correction path changing both factors is absent, at every
length.&Complete. M606 is proved. The synthesis at the end of this
chapter proves M601 and M607.\\[5pt]
J21&An explicit correction produces a degree-zero cycle in the exact
anchor class with the full second-vertex bound.&Complete. M901--M903
are proved.\\[5pt]
J10&One adjacent X switch in the two specified letter columns,
unchanged O markings, and an explicit finite move sequence to every
admissible shorter-word grid.&The selected-grid incidence and order
equations, and primitive native move recognition, remain source
checks. Their defining files are absent.\\[5pt]
J15&The exponent equals the discriminant order and polar intersection.
The branch and reduced-word length formulas hold. The full-twist
comparison also proves M407's broader positive-word quantifier.
&Complete as a written word-length proof. The reduced native sphere
form retains the audit's M403(b) geometric realization input.\\[5pt]
J14&For every reduced germ of multiplicity $S$ with transverse
projection, the root braid has the required positive full-twist
prefix and exactly the branch cycles. No normalization input is
left unproved.&The uncorrected M405 API is false, with an explicit
counterexample. The native sphere realization retains the stated
M403(b) and M419(B) dependencies.\\\hline
\end{longtable}

All Floer chains, auxiliary cyclic-covering modules, and matrix
coefficients use $\mathbb F_2$. Integer rectangle multiplicities,+braid exponents, analytic orders, and intersection multiplicities
retain their ordinary integer meanings. Each new lemma is labelled
as new and proved in the text. Existing Lean declarations are cited
by their native names. Separate job chapters repeat some elementary
domain lemmas so that their individual documents remain usable.
The proof dependency is acyclic: J14 uses only the independently
proved preparation and local argument-principle lemmas of J15,
whereas J15's reduced full-twist conclusion then uses J14.
\tableofcontents


\chapter{J05: First-sector seam parity}
\label{chapter:J05}



\section{Exact scope and status}

All coefficients in this document are in $\F$. Every cardinality before reduction
modulo two is an ordinary nonnegative integer. Let $n,m\in\mathbb N$, put
$N=n+m+1$, and use the native grid states, rectangles, retained indices, and
matching sources in the supplied seam listings. Write
\[
 \CP_k(p,q)=\decl{SeamCorrectionPath}\ k\ p\ q.
\]
The integer $k$ counts internal transitions. An entry rectangle and an exit
rectangle are also part of each such path. They are not counted in $k$.

The acceptance statement of J05 requires that, for every $n,m\in\mathbb N$,
every $x,y\in\operatorname{GridState}(n+1)$, and every
$z\in\operatorname{GridState}(m+1)$,
\begin{equation}\label{J05:eq:acceptance}
 \decl{seamLongCorrectionCount}((x,z),(y,z))=0\quad\text{in }\F.
\end{equation}
This document proves \eqref{J05:eq:acceptance} for all $n,m$.
In fact, Theorem~\ref{J05:thm:seam-first-long-empty} proves that every
fixed-second correction path with at least three internal transitions is
absent. The proof uses a signed domain of index one, two nonnegative factor
projections, and deletion of fixed state points. Every putative long path
compresses to at most five points, where an explicit order certificate
excludes it. The following flat-path theorem handles the thinner compressed
sizes and also proves M603 directly.

\begin{theorem}[New lemma: exclusion of completely flat correction paths]
\label{J05:thm:flat}
Let $n,m,k\in\mathbb N$ with $k\geq2$, and let $p,q$ be native retained indices.
There is no element of $\CP_k(p,q)$ whose every matching source has exactly one
northeast column.
\end{theorem}

\begin{corollary}[New lemma: M603]\label{J05:cor:three}
Let $n,m\in\mathbb N$ and let $p,q$ be retained indices. No element of
$\CP_3(p,q)$ has northeast count one at both interior matching sources.
Consequently, for every $x,y\in\decl{GridState}(n+1)$ and
$z\in\decl{GridState}(m+1)$, the set $\CP_3((x,z),(y,z))$ is empty.
\end{corollary}

\begin{corollary}[New lemma: J05 for a factor of size at most two]
\label{J05:cor:thin}
If $\min(n,m)\leq1$, then \eqref{J05:eq:acceptance} holds for these $n,m$.
In fact, every correction path with at least two internal transitions is absent.
\end{corollary}

Theorem~\ref{J05:thm:flat} concerns arbitrary retained endpoints.
The later domain argument handles all northeast counts for the fixed-second
acceptance statement. The failing module
\decl{GridConnectedSumSeamFlatThreeStepBarrier.lean} is not used as an
existing result.

\section{Native coordinates and rectangle conventions}

Identify $\operatorname{Fin}(N)$ with $\{0,\ldots,N-1\}$, preserving its order.
A grid state is a permutation $s$ from columns to rows. Its northeast and
southwest columns are
\[
 \NE(s)=\{a<n:s(a)>n\},\qquad
 \SW(s)=\{P\geq n:s(P)\leq n\}.
\]
These are the native sets characterized by
\decl{mem_seamNortheastColumns} and \decl{mem_seamSouthwestColumns}, as used in
the supplied proofs. The native identity \decl{seamSouthwestColumns_card} gives
$\seamcardop\SW(s)=\seamcardop\NE(s)+1$.

For a matching source with northeast count one, write its two southwest
columns as $P<Q$. Write
\[
 s(P)=p,\qquad s(Q)=q,\qquad p<q\leq n.
\]
The last inequality is the native increasing-source condition. The matching
partner swaps the rows at $P,Q$.

For distinct indices $a,b$, let $(a,b)_{\rm cyc}$ be the open cyclic interval.
It is $(a,b)$ if $a<b$ and
$\{c:c>a\text{ or }c<b\}$ if $a>b$. An oriented rectangle with left and right
columns $l,r$ has bottom and top rows $s(l),s(r)$. Its emptiness means
\begin{equation}\label{J05:eq:empty}
 \text{there is no column }c\in(l,r)_{\rm cyc}
 \text{ with }s(c)\in(s(l),s(r))_{\rm cyc}.
\end{equation}
This is \decl{isEmpty_iff_forall_notMem_cIoo}, or equivalently the supplied
\decl{flatEmpty_source}. Thus a row interval from a larger row to a smaller
row wraps. This orientation is essential below.

Suppose a correction path is completely flat. By \decl{flatEntry_data}, its
first matching source has a column $A<n$ and row $h=s(A)>n$. Each flat step
preserves this point by \decl{flatApply_high}. The same $A,h$ therefore serves
at every matching source. Because the northeast count is one, this is its
unique northeast point.

The native entry lemma \decl{flatEntry_data} says that one of the following
two rectangles is empty, expressed in the first matching source's coordinates:
\begin{equation}\label{J05:eq:entry}
 E_L=(A,P)\times(q,h),\qquad
 E_R=(Q,A)_{\rm cyc}\times(h,p)_{\rm cyc}.
\end{equation}
The native exit lemma \decl{flatExit_data} says that one of the following is
empty in the last matching source:
\begin{equation}\label{J05:eq:exit}
 F_L=(A,P)\times(h,p)_{\rm cyc},\qquad
 F_R=(Q,A)_{\rm cyc}\times(q,h).
\end{equation}
The letters $P,Q,p,q$ in \eqref{J05:eq:entry} and \eqref{J05:eq:exit} belong to their
respective endpoint matching sources. Formula \eqref{J05:eq:entry} already
accounts for the inverse matching move. It is not the rectangle obtained by
using the source rows $s(A),s(P)$ for the entry.

\section{Classification and alternation of flat steps}

The existing declarations \decl{flatStep_cases} and \decl{flatStep_apply}
classify every flat internal transition. Table~\ref{J05:tab:types} writes their
conclusions without Lean notation. Here $C$ is the exchanged column,
$c=s(C)$, and $s'$ is the next matching source. Values at columns not displayed
remain unchanged. In the last column, the ordered pair gives the new southwest
columns.

\begin{table}[ht]
\centering\small
\begin{tabular}{c|l|c|c|c}
Type & Required range & Rectangle sides & $(s'(P),s'(Q),s'(C))$ & New columns\\\hline
$\Lp$ & $C<n$, $q<c\leq n$ & $(C,P)$ & $(q,c,p)$ & $(P,Q)$\\
$\Lm$ & $C<n$, $c<p$ & $(Q,C)$ & $(c,p,q)$ & $(P,Q)$\\
$\Hp$ & $C>Q$, $c>n$ & $(C,P)$ & $(c,p,q)$ & $(Q,C)$\\
$\Hm$ & $n\leq C<P$, $c>n$ & $(Q,C)$ & $(q,c,p)$ & $(C,P)$
\end{tabular}
\caption{The native four flat transitions.}\label{J05:tab:types}
\end{table}

\begin{lemma}[New lemma: the northeast point separates the four types]
\label{J05:lem:separate}
In the notation above, an empty flat transition satisfies the following
additional inequalities:
\[
 \Lp:\ A<C,\qquad
 \Lm:\ C<A,\qquad
 \Hp:\ h<c,\qquad
 \Hm:\ c<h.
\]
\end{lemma}
\begin{proof}
For $\Lp$, the row interval is $(c,p)_{\rm cyc}$ and contains $h$ because
$h>n\geq c$. If $C<A$, the column $A$ lies between $C$ and $P$, contrary to
\eqref{J05:eq:empty}. The equality $C=A$ contradicts $c\leq n<h$.

For $\Lm$, the column interval is $(Q,C)_{\rm cyc}$ and the row interval
$(q,c)_{\rm cyc}$ contains $h$. If $A<C$, the point at $A$ is inside the
rectangle. Equality is again excluded by its row. Thus $C<A$.

For $\Hp$, $A<P<C$, so $A\in(C,P)_{\rm cyc}$. The row interval
$(c,p)_{\rm cyc}$ contains $h$ if $h>c$. Emptiness and injectivity therefore
give $h<c$.

For $\Hm$, $A<C<Q$, so $A\in(Q,C)_{\rm cyc}$. The row interval $(q,c)$
contains $h$ if $h<c$. Emptiness and injectivity give $c<h$.
\end{proof}

\begin{lemma}[New lemma: successive flat types alternate]
\label{J05:lem:alternate}
For two consecutive empty flat transitions, the ordered pair of types is one
of
\[
 (\Lp,\Lm),\quad(\Lm,\Lp),\quad
 (\Hp,\Hm),\quad(\Hm,\Hp).
\]
\end{lemma}
\begin{proof}
Use $C,c$ for the first exchange and $D,d$ for the second exchange. In this
proof $d$ is the row of $D$ immediately before the second transition. Whenever
the proof tests $D$ against the first rectangle, $D$ differs from $P,Q,C$ and
its row has not changed. For equal low types, $D\ne C$ follows from the row
range in Table~\ref{J05:tab:types}. For equal high types it follows from the
column range. For a low type followed by a high type, or the reverse, it
follows from the quadrant restrictions and Lemma~\ref{J05:lem:separate}.

Suppose the first type is $\Lp$. A second $\Lp$ has $d>c$.
The first rectangle excludes $D>C$, and equality is impossible, so $D<C$.
The second rectangle then contains the point at $C$, whose row is $p<q$.
A second $\Hp$ has left endpoint $D>Q$ and right endpoint $P$.
Its row interval ends at $q$, and it contains the point at $C$ of row $p$.
A second $\Hm$ uses a column $n\leq D<P$ of high row. That point already
lies inside the first rectangle $(C,P)\times(c,p)_{\rm cyc}$.
These contradictions leave only $\Lm$.

Suppose the first type is $\Lm$. A second $\Lm$ has $d<c$.
The first rectangle excludes $D<C$, so $D>C$. The second rectangle contains
the point at $C$, now of row $q>p$.
A second $\Hp$ uses a column $D>Q$ of high row. That point is inside the
first rectangle $(Q,C)_{\rm cyc}\times(q,c)_{\rm cyc}$.
A second $\Hm$ has sides $(Q,D)$ in the unchanged southwest columns
$P,Q$. Its row interval starts at $p$ and ends at a row above $n$. The point
at $C<D$ has row $q$, so it lies inside this rectangle.
Only $\Lp$ remains.

Suppose the first type is $\Hp$. The new southwest columns are $Q,C$,
with rows $p,q$, and the point at the old $P$ has row $c>h$.
A second $\Hp$ uses $D>C$. The first rectangle forces $d<c$.
The second rectangle has sides $(D,Q)$ and row interval $(d,p)_{\rm cyc}$.
It contains the point at the old $P$, whose row is $c>d$.
A second $\Lp$ has sides $(D,Q)$ and starts at a row $d\leq n$.
It again contains the point at the old $P$, whose row exceeds $n$.
A second $\Lm$ uses a low column with row $d<p$. That point is inside the
first rectangle $(C,P)_{\rm cyc}\times(c,p)_{\rm cyc}$.
Only $\Hm$ remains.

Suppose the first type is $\Hm$. The new southwest columns are $C,P$,
with rows $p,q$, and the point at the old $Q$ has row $n<c<h$.
A second $\Hm$ uses $D<C$. The first rectangle forces $d>c$.
The second rectangle has sides $(P,D)$ and row interval $(q,d)$, so it
contains the point at the old $Q$ of row $c$.
A second $\Lp$ uses a low column of row $d>q$. That point is inside the
first rectangle $(Q,C)_{\rm cyc}\times(q,c)$.
A second $\Lm$ has sides $(P,D)$ and wraps vertically from $q$ to a row
below $p$. It contains the point at the old $Q$, whose row exceeds $n$.
Only $\Hp$ remains. This proves all four assertions.
\end{proof}

\section{Persistent obstructions to the exit}

\begin{proof}[Proof of Theorem~\ref{J05:thm:flat}]
Assume a completely flat correction path has $k\geq2$ internal transitions.
Use the fixed northeast point $(A,h)$, and write $P<Q$ and $p<q$ for the
initial southwest columns and rows. By Lemma~\ref{J05:lem:alternate}, there are
four possible alternating sequences, determined by the first type. We treat
each sequence and prove all inequalities used in the obstruction.

\subsection*{First type $\Lp$}

Write $C_j,c_j$ for the column and entering row of the $j$th $\Lp$ exchange,
and $D_j,d_j$ for those of the $j$th $\Lm$ exchange. Only indices corresponding
to transitions present in the finite path are used. Put $d_0=p$.
The southwest columns stay equal to $P,Q$. Table~\ref{J05:tab:types} gives
\[
 \begin{array}{c|c|c}
 \text{last transition}&(s(P),s(Q))&\text{row left at the exchanged column}\\\hline
 \Lp(C_j)&(q,c_j)&s(C_j)=d_{j-1}\\
 \Lm(D_j)&(d_j,q)&s(D_j)=c_j.
 \end{array}
\]
Lemma~\ref{J05:lem:separate} gives $A<C_j<n$, $D_j<A$, $d_j<q<c_j$.
The first entry cannot be $E_L$, because the point at $C_1$ has row
$q<c_1<h$ and column $A<C_1<P$. Hence it is $E_R$.
The second exchange has $D_1<A$ and row $d_1<q$.
If $d_1<p$, this point is inside $E_R$. Equality is excluded by injectivity.
Therefore
\begin{equation}\label{J05:eq:Lplusinitial}
 p<d_1<q<c_1.
\end{equation}

The subsequent exchanges satisfy
\begin{equation}\label{J05:eq:Lplusrec}
 C_{j+1}>C_j,\quad c_{j+1}<c_j,\qquad
 D_{j+1}<D_j,\quad d_{j+1}>d_j,
\end{equation}
whenever the relevant transitions occur. To prove the first two inequalities,
suppose the last two transitions were $\Lp(C_j),\Lm(D_j)$ and
$d_{j-1}<d_j$. The next $\Lp$ rectangle goes from $C_{j+1}$ to $P$ and wraps
vertically from $c_{j+1}$ to $d_j$. The point at $C_j$ has row
$d_{j-1}<d_j$. Its emptiness forces $C_{j+1}>C_j$.
Test $C_{j+1}$ against the earlier $\Lp(C_j)$ rectangle. Its row has not
changed between that earlier source and its exchange, since the only
intervening changed columns were $P,Q,C_j,D_j$, all different from
$C_{j+1}$. The earlier rectangle forces $c_{j+1}<c_j$.

For the next $\Lm$ rectangle, the point at $D_j$ has row $c_j>c_{j+1}$.
If $D_j<D_{j+1}$, this point lies in its wrapping column and row intervals.
Thus $D_{j+1}<D_j$. Test $D_{j+1}$ against the earlier $\Lm(D_j)$ rectangle.
The intervening $\Lp$ transition did not change its row. That earlier rectangle
forces $d_{j+1}>d_j$. These arguments start with \eqref{J05:eq:Lplusinitial}
and prove \eqref{J05:eq:Lplusrec} by induction. Equalities are excluded either
by the current row range or by injectivity of the earlier state.

In particular, the original exchanged columns $C_1,D_1$ are never exchanged
again. Their rows remain $p,c_1$, respectively. At every source after at least
two transitions, the current southwest rows $p',q'$ satisfy
$p<p'$ and $q'<c_1$. The point at $C_1$ lies in $F_L$, and the point at
$D_1$ lies in $F_R$. This contradicts \eqref{J05:eq:exit}.

\subsection*{First type $\Lm$}

Write $C_j,c_j$ for the successive $\Lm$ exchanges and $D_j,d_j$ for the
successive $\Lp$ exchanges. Put $d_0=q$. The southwest columns stay $P,Q$,
and the row formulas are
\[
 \begin{array}{c|c|c}
 \Lm(C_j)&(c_j,p)&s(C_j)=d_{j-1}\\
 \Lp(D_j)&(p,d_j)&s(D_j)=c_j.
 \end{array}
\]
Here $C_j<A<D_j<n$ and $c_j<p<d_j$.
The point at $C_1$ blocks $E_R$, since its row is below $p$. Thus the entry is
$E_L$. Its emptiness at $D_1$ forces $d_1<q$. Consequently
\[
 c_1<p<d_1<q.
\]
The inductive inequalities are
\[
 C_{j+1}<C_j,\quad c_{j+1}>c_j,\qquad
 D_{j+1}>D_j,\quad d_{j+1}<d_j.
\]
For their proof, immediately before $\Lm(C_{j+1})$ the point at $C_j$ has
row $d_{j-1}>d_j$. The new rectangle has sides $(Q,C_{j+1})$ and wraps
vertically from $d_j$ to $c_{j+1}$. Testing $C_j$ forces
$C_{j+1}<C_j$. Testing $C_{j+1}$ against the preceding $\Lm(C_j)$ rectangle
then forces $c_{j+1}>c_j$.
Immediately before $\Lp(D_{j+1})$, the point at $D_j$ has row
$c_j<c_{j+1}$. Testing $D_j$ against the new rectangle from $D_{j+1}$ to
$P$ forces $D_{j+1}>D_j$. Testing $D_{j+1}$ against the preceding
$\Lp(D_j)$ rectangle forces $d_{j+1}<d_j$. Each test uses an unchanged row
at a column distinct from the intervening exchanged column, as in the first
case. These arguments prove the induction starting at $d_1<d_0=q$.

The original $C_1$ retains row $q$, and $D_1$ retains row $c_1$.
After at least two transitions, the current southwest rows satisfy
$c_1<p'$ and $q'<q$. The point at $D_1$ blocks $F_L$, and the point at $C_1$
blocks $F_R$. Neither exit is empty.

\subsection*{First type $\Hp$}

In this case the southwest rows stay equal to the original $p,q$.
Write $P_0=P$, $Q_0=Q$. Write $C_j,c_j$ for the successive $\Hp$ exchanges
and $D_j,d_j$ for the successive $\Hm$ exchanges, with $D_0=P_0$.
The southwest columns and the row left behind are
\[
 \begin{array}{c|c|c}
 \Hp(C_j)&(Q_0,C_j)&s(D_{j-1})=c_j\\
 \Hm(D_j)&(D_j,Q_0)&s(C_j)=d_j.
 \end{array}
\]
The type inequalities give $D_j<Q_0<C_j$ and $n<d_j<h<c_j$.
The point at $C_1$ blocks $E_R$, so the entry is $E_L$.
Testing $D_1$ against this entry forces $P_0<D_1$. Thus
\[
 P_0<D_1<Q_0<C_1.
\]
The inductive inequalities are
\[
 C_{j+1}<C_j,\quad c_{j+1}>c_j,\qquad
 D_{j+1}>D_j,\quad d_{j+1}<d_j.
\]
Suppose first that $D_{j-1}<D_j$.
The rectangle of $\Hp(C_{j+1})$ has sides $(C_{j+1},D_j)$ and wraps
vertically from $c_{j+1}$ to $p$. The point at $D_{j-1}$ has row $c_j$ and
lies in its column interval. Emptiness forces $c_{j+1}>c_j$.
Testing $C_{j+1}$ against the preceding $\Hp(C_j)$ rectangle now forces
$C_{j+1}<C_j$. The tested row is unchanged because $C_{j+1}>Q_0$ and
$C_{j+1}\ne C_j$, while the intervening transition changed only
$Q_0,C_j,D_j$.

The rectangle of $\Hm(D_{j+1})$ has sides $(C_{j+1},D_{j+1})$ and row
interval $(q,d_{j+1})$. The point at $C_j>C_{j+1}$ has row $d_j$.
Its emptiness forces $d_{j+1}<d_j$. Testing $D_{j+1}$ against the preceding
$\Hm(D_j)$ rectangle forces $D_{j+1}>D_j$. Again its row is unchanged
between these tests. This proves the induction.

The original $P_0$ retains row $c_1>h$, and $C_1$ retains row $d_1<h$.
Every later southwest pair $P',Q'$ after at least two transitions satisfies
$P_0<P'$ and $Q'<C_1$. The point at $P_0$ blocks $F_L$, and the point at
$C_1$ blocks $F_R$.

\subsection*{First type $\Hm$}

Keep $P_0=P$, $Q_0=Q$, and write $C_j,c_j$ for the successive $\Hm$
exchanges and $D_j,d_j$ for the successive $\Hp$ exchanges, with
$D_0=Q_0$. The southwest rows stay $p,q$, and the formulas are
\[
 \begin{array}{c|c|c}
 \Hm(C_j)&(C_j,P_0)&s(D_{j-1})=c_j\\
 \Hp(D_j)&(P_0,D_j)&s(C_j)=d_j.
 \end{array}
\]
Here $C_j<P_0<D_j$ and $n<c_j<h<d_j$.
The point at $C_1$ blocks $E_L$, so the entry is $E_R$.
Testing $D_1$ against $E_R$ forces $D_1<Q_0$. Hence
\[
 C_1<P_0<D_1<Q_0.
\]
The inductive inequalities are
\[
 C_{j+1}>C_j,\quad c_{j+1}<c_j,\qquad
 D_{j+1}<D_j,\quad d_{j+1}>d_j.
\]
Suppose that $D_j<D_{j-1}$. The rectangle of $\Hm(C_{j+1})$ has sides
$(D_j,C_{j+1})$ and row interval $(q,c_{j+1})$. The point at
$D_{j-1}$ has row $c_j$ and belongs to its column interval.
Emptiness forces $c_{j+1}<c_j$. Testing $C_{j+1}$ against the preceding
$\Hm(C_j)$ rectangle then forces $C_{j+1}>C_j$.
The rectangle of $\Hp(D_{j+1})$ has sides $(D_{j+1},C_{j+1})$ and wraps
vertically from $d_{j+1}$ to $p$. The point at $C_j<C_{j+1}$ has row $d_j$,
so emptiness forces $d_{j+1}>d_j$. Testing $D_{j+1}$ against the preceding
$\Hp(D_j)$ rectangle forces $D_{j+1}<D_j$. In both earlier-rectangle tests,
the tested column is different from all columns changed in the intervening
transition. Its row is therefore the same in the two tests.

The original $Q_0$ retains row $c_1<h$, and $C_1$ retains row $d_1>h$.
After at least two transitions, every southwest pair satisfies
$C_1<P'$ and $Q'<Q_0$. The point at $C_1$ blocks $F_L$, and the point at
$Q_0$ blocks $F_R$.

All four possible alternating sequences contradict the required empty exit.
This proves the theorem.
\end{proof}

\section{Consequences and exceptional cases}

\begin{proof}[Proof of Corollary~\ref{J05:cor:three}]
The declarations \decl{seamMatchingEntry_northeast_card} and
\decl{seamMatchingExit_northeast_card} give northeast count one at the first
and last matching sources of any correction path. For a three-transition
path whose two interior sources also have count one, Theorem~\ref{J05:thm:flat}
applies. For fixed second factor, the existing
\decl{seamCorrectionThree_middle_eq_one} supplies exactly these two interior
equalities. No proved three-step exclusion is repeated here.
\end{proof}

\begin{lemma}[New lemma: cardinality bound for the northeast set]
\label{J05:lem:countbound}
For every state $s$ of size $n+m+1$,
\[
 \seamcardop\NE(s)\leq\min(n,m).
\]
\end{lemma}
\begin{proof}
The inclusion into the $n$ columns $0,\ldots,n-1$ bounds the cardinality by
$n$. The restriction of the permutation $s$ to $\NE(s)$ is an injection into
the $m$ rows $n+1,\ldots,n+m$, which bounds the cardinality by $m$.
\end{proof}

\begin{proof}[Proof of Corollary~\ref{J05:cor:thin}]
Every matching source has nonempty northeast set by
\decl{seamNortheastColumns_nonempty_iff}. If $n=0$ or $m=0$,
Lemma~\ref{J05:lem:countbound} rules out every matching source, and therefore
every correction path. If $\min(n,m)=1$, every matching source has northeast
count one. Theorem~\ref{J05:thm:flat} excludes every correction path with
$k\geq2$. In both cases every summand in the native
\decl{seamLongCorrectionCount} is zero, proving the corollary.
\end{proof}

The sizes $(n+1,m+1)=(1,1),(1,2),(2,1),(2,2)$ are covered by this proof.
For the smallest nonempty complement, $n=m=1$, there is one matching source,
namely the state $(2,0,1)$ on columns $0,1,2$. Its southwest columns are
$1,2$ with rows $0,1$, and its partner is $(2,1,0)$. The positively oriented
matching rectangle has sides $(1,2)$ and row interval $(0,1)$.
There is no internal transition because distinct matching sources are
required. This checks the orientation and the strict versus weak inequalities
at the shared index $n$ directly. It is an exceptional-case check, not the
argument for arbitrary sizes.

More generally, if the number $M$ of matching sources is at most three, the
native long sum is empty by its definition. This is already recorded by the
audit and does not require a new parity argument.

\section{Projected integral domains of a seam correction}

All path coefficients remain in $\mathbb F_2$. Integral domains in this
section record the multiplicities of grid squares before any parity is
taken. Their use does not change the coefficient field of the chain
complex.

For a grid of size $h$, an integral domain is an integer-valued function
$D(c,r)$ on the $h^2$ squares. All cell and vertex indices in the next
two lemmas are cyclic modulo $h$. Define its mixed difference and corner
average by
\[
 \begin{split}
 (\partial D)(c,r)
 &=D(c,r)+D(c-1,r-1)-D(c-1,r)-D(c,r-1),\\
 (\mathcal A D)(c,r)
 &=\tfrac14\bigl(D(c,r)+D(c-1,r-1)+D(c-1,r)+D(c,r-1)\bigr).
 \end{split}
\]
For a state $u$, let $[u]$ be the indicator function of its points.
A domain from $u$ to $v$ satisfies $\partial D=[u]-[v]$. Its index is
\[
 \iota(D;u,v)=\sum_{c=0}^{h-1}
 \bigl((\mathcal A D)(c,u(c))+(\mathcal A D)(c,v(c))\bigr).
\]
This is a definition by finite sums. In particular it makes sense for
domains with negative coefficients.

\begin{lemma}[New lemma: domain index identities]
\label{J05:lem:seam-domain-index}
The following statements hold for integral domains on a toroidal grid.
\begin{enumerate}
\item If $D$ is a domain from $u$ to $v$ and $E$ is a domain from
$v$ to $w$, then
\[
 \iota(D+E;u,w)=\iota(D;u,v)+\iota(E;v,w).
\]
\item A domain with mixed difference zero has the form
$D(c,r)=a_c+b_r$, with integers $a_c,b_r$. Its index between any equal
endpoint states is $2\sum_c a_c+2\sum_r b_r$.
\item An empty native rectangle has index one. For any auxiliary
grid with marking permutation $O$ and any domain from $u$ to $v$,
\begin{equation}\label{J05:eq:seam-domain-maslov}
 \iota(D;u,v)=M_O(u)-M_O(v)+2\sum_c D(c,O(c)).
\end{equation}
Here $M_O$ is the native integer Maslov grading.
\end{enumerate}
\end{lemma}

\begin{proof}
Finite cyclic reindexing gives
\[
 \sum_{c,r}(\partial D)(c,r)(\mathcal A E)(c,r)
 =\sum_{c,r}(\partial E)(c,r)(\mathcal A D)(c,r).
\]
For completeness, after moving the differences from $D$ to $E$, the
coefficient of $D(c,r)$ on the left is
\[
 \tfrac14\bigl(E(c+1,r+1)-E(c+1,r-1)
             -E(c-1,r+1)+E(c-1,r-1)\bigr).
\]
Its coefficients are unchanged when both offsets are negated. This
proves the displayed symmetry. Substituting
$\partial D=[u]-[v]$ and $\partial E=[v]-[w]$ proves the first assertion
by expanding the definition of $\iota$.

If $\partial D=0$, the difference $D(c,r)-D(c-1,r)$ is independent
of $r$. Fixing one row and one column and summing these differences
gives $D(c,r)=a_c+b_r$. Each state occupies every row and every column
once. Substitution into the corner average therefore gives index
$2\sum a_c+2\sum b_r$.

The domain of a rectangle with sides $(a,b)$ and source $u$ is the
indicator of
\[
 [a,b)_{\mathrm{cyc}}\times[u(a),u(b))_{\mathrm{cyc}}.
\]
Its mixed difference is $[u]-[u\circ(a\ b)]$ by the native declaration
\decl{GridRectangleBetween.target_eq_swapColumns}.
If the rectangle is empty, its four endpoint corners each contribute
$1/4$ to the index and every other state point contributes zero.
The native declaration \decl{laurentMaslovDegree_add_card_OColumns}
then proves \eqref{J05:eq:seam-domain-maslov} for an empty rectangle.

To obtain the formula for an arbitrary domain, connect $u$ to $v$ by
adjacent-column transpositions. Such a sequence exists by moving the
desired value into the first incorrect column and repeating on the
remaining columns. Use the rectangle of width one for each transposition.
It is empty because there is no interior column. Their sum $D_0$ has
the required boundary, and the formula for $D_0$ follows by additivity
and telescoping of the native grading differences. The difference
$D-D_0$ has zero mixed difference. For a row or column annulus the
index is two and the sum of the $O$ multiplicities is one.
The second assertion, with arbitrary integer coefficients, now gives
the formula for $D-D_0$ and hence for $D$.
\end{proof}

\begin{lemma}[New lemma: nonnegative domains of index one]
\label{J05:lem:seam-index-one-rectangle}
Let $D$ be a nonnegative integral domain from distinct states $u$ to
$v$. Then $\iota(D;u,v)\ge1$. If its index is one, then $D$ is the
indicator domain of an empty native rectangle from $u$ to $v$.
\end{lemma}

\begin{proof}
Suppose that the states differ in $e$ columns. Then $e\ge2$, and the
mixed difference has a coefficient $1$ or $-1$ at each of their $2e$
distinct changed points. At each such point at least one adjacent cell
has a positive integer coefficient. Its corner average is therefore
at least $1/4$. All remaining contributions are nonnegative, so
$\iota(D;u,v)\ge e/2\ge1$.

If the index is one, this bound gives $e=2$. Write the changed columns
as $A,B$, and put $R=u(A)$ and $S=u(B)$. The target exchanges these
two rows. Let $V$ and $W$ be the nonempty proper cyclic cell intervals
$[A,B)$ and $[R,S)$, and let $D_0$ be the indicator of $V\times W$.
By Lemma~\ref{J05:lem:seam-domain-index},
\[
 D(c,r)=\mathbf1_V(c)\mathbf1_W(r)+a_c+b_r.
\]
Replace $a_c$ by its minimum separately on $V$ and its complement,
and replace $b_r$ by its minimum separately on $W$ and its complement.
The resulting domain $D_b$ is nonnegative, because on each of the four
products of intervals its value is the minimum of $D$ there.
The difference $D-D_b$ is a sum of row and column annuli with
nonnegative integer coefficients. The domain $D_b$ has the same
boundary, so its index is at least one by the preceding paragraph.
Each annulus has index two. Additivity and $\iota(D)=1$ force every
annulus coefficient to be zero. Thus $D=D_b$ is constant on each of
the four products of intervals.

Let the four constants be $a,b,c,d$, with $a$ on $V\times W$ and
$d$ on its opposite product. They are nonnegative integers and satisfy
$a+d-b-c=1$. The four changed points contribute exactly
$a+b+c+d$ to the index. Hence this sum is at most one, and the
only possibilities are $(a,b,c,d)=(1,0,0,0)$ and $(0,0,0,1)$.
These are the two oriented rectangle domains from $u$ to $v$.
The four corners already account for the full index. Every other
source-state corner average is zero, so there is no source point in
the rectangle's interior. This is the native emptiness condition
\decl{isEmpty_iff_forall_notMem_cIoo}.
\end{proof}

Now let $n,m\ge1$ and $N=n+m+1$. For a correction path from retained
state $p=\decl{seamState}(x,z)$ to retained state
$q=\decl{seamState}(y,w)$, define $D$ to be the sum of its entry,
internal, and exit rectangle domains, minus the domains of the selected
matching rectangles at all its matching sources.

\begin{lemma}[New lemma: the signed correction domain]
\label{J05:lem:seam-signed-domain}
The signed domain $D$ of a correction path satisfies
$\partial D=[p]-[q]$ and
\begin{equation}\label{J05:eq:seam-path-index}
 \iota(D;p,q)=1.
\end{equation}
\end{lemma}
\begin{proof}
The matching rectangle at a source with southwest columns $P<Q$ and
rows $a<b\le n$ is $[P,Q)\times[a,b)$. It is empty, because every
column strictly between the first two southwest columns has a high row.
Telescoping the rectangle boundaries gives $\partial D=[p]-[q]$.
There is one more positive rectangle than subtracted matching rectangles.
Each has index one, so additivity in Lemma~\ref{J05:lem:seam-domain-index}
gives \eqref{J05:eq:seam-path-index}.
\end{proof}

Define two restrictions, with the indicated factor indices, by
\begin{equation}\label{J05:eq:seam-domain-projections}
 \begin{aligned}
 D_1(i,j)&=D(i,j) &&(0\le i<n,\ 0\le j\le n),\\
 D_1(n,j)&=D(N-1,j) &&(0\le j\le n),\\
 D_2(i,j)&=D(n+i,n+j) &&(0\le i,j\le m).
 \end{aligned}
\end{equation}
Put $d=D(N-1,n)$.

\begin{lemma}[New lemma: positive projections and the index relation]
\label{J05:lem:seam-positive-projections}
For every correction path as above, $D_1$ is a nonnegative domain from
$x$ to $y$, and $D_2$ is a nonnegative domain from $z$ to $w$.
Their indices satisfy
\begin{equation}\label{J05:eq:seam-projected-index}
 \iota(D_1;x,y)+\iota(D_2;z,w)=1+2d.
\end{equation}
\end{lemma}

\begin{proof}
A matching rectangle uses only columns from $n$ through $N-2$ and
rows from $0$ through $n-1$. It therefore misses both restrictions in
\eqref{J05:eq:seam-domain-projections}. All remaining summands of $D$ are
nonnegative rectangle indicators, proving positivity.

For the boundary assertion, restricting a domain to cyclic lists of
cell columns $c_i$ and cell rows $r_j$ replaces its mixed difference
at $(i,j)$ by
\[
 \sum_{c\in(c_{i-1},c_i]_{\mathrm{cyc}}}
 \sum_{r\in(r_{j-1},r_j]_{\mathrm{cyc}}}(\partial D)(c,r).
\]
This formula follows by telescoping in both coordinates.
For $D_1$, the last column collects the original columns
$n,\ldots,N-1$, and row zero collects the original rows
$n+1,\ldots,N-1,0$. Under these collections the retained state $p$
becomes $[x]+m[(n,0)]$. The extra $m$ copies are the high points in
the second columns. The retained state $q$ gives
$[y]+m[(n,0)]$, so the extra terms cancel.
For $D_2$, column zero collects columns $0,\ldots,n$, and row zero
collects rows $0,\ldots,n$. The retained states become
$[z]+n[(0,0)]$ and $[w]+n[(0,0)]$. Again the extra terms cancel.
These statements use precisely the three supplied coordinate formulas
\decl{seamState_left_castSucc}, \decl{seamState_right_of_zero}, and
\decl{seamState_right_of_ne_zero}.

We prove the index relation by explicit auxiliary markings.
Let $O_1(i)=i$ on the first factor, let
$O_2(j)=j+1\pmod{m+1}$ on the second factor, and let
\[
 O(c)=c\quad(c<n),\qquad
 O(n+j)=n+O_2(j)\quad(0\le j\le m).
\]
Each is a permutation and can be made into a valid auxiliary grid by
choosing $X$ to be its row successor. The Maslov grading depends only on
the specified $O$ markings. The full marking set is the union of the
two projected marking sets with their common square $(N-1,n)$ counted
once. Consequently
\begin{equation}\label{J05:eq:seam-marking-domain-count}
 \sum_c D(c,O(c))=
 \sum_iD_1(i,O_1(i))+\sum_jD_2(j,O_2(j))-d.
\end{equation}

The native Maslov point-count formula gives
\begin{equation}\label{J05:eq:seam-auxiliary-maslov-sum}
 M_O(\decl{seamState}(x,z))=M_{O_1}(x)+M_{O_2}(z).
\end{equation}
Here is a direct verification, so no condition on the auxiliary
$X$ markings is being suppressed. For finite planar point sets $A,B$,
write $I(A,B)$ for the number of ordered pairs that increase strictly
in both coordinates. State points have integral coordinates and marks
have both coordinates shifted by $1/2$. The supplied formula
\decl{GridDiagram.maslovOℤ_eq_card} is
\[
 M_O(s)=I(s,s)-I(s,O)-I(O,s)+I(O,O)+1.
\]
Splitting each point set into its first $n$ and last $m+1$ columns gives
\[
\begin{array}{rcl}
 I(p,p)&=&I(x,x)+I(z,z)+nm,\\
 I(O,O)&=&I(O_1,O_1)+I(O_2,O_2)+nm,\\
 I(p,O)&=&I(x,O_1)+I(z,O_2)+nm-1,\\
 I(O,p)&=&I(O_1,x)+I(O_2,z)+nm.
\end{array}
\]
To check the first equality, put $t=x(n)$. Removing the last column
from $x$ removes $t$ increasing pairs. The cross pairs with second-column
points contribute $nm+t$. The second equality is the same computation
for the marking permutations.
For the third equality, removing the last marking of $O_1$, whose row
is $n$, removes $n+1$ pairs. All $n(m+1)$ pairs from first-column state
points to second-column markings contribute. The pairs within the second
columns agree with $I(z,O_2)$ because the low state row $t$ is at most
the lowest marking row $n+1/2$.
For the fourth equality, removing the last state point of $x$ removes
$t$ pairs, whereas the cross pairs contribute $nm+t$.
Substituting the four identities into the native formula proves
\eqref{J05:eq:seam-auxiliary-maslov-sum}.

Apply \eqref{J05:eq:seam-domain-maslov} to the three domains. Equations
\eqref{J05:eq:seam-path-index}, \eqref{J05:eq:seam-marking-domain-count}, and
\eqref{J05:eq:seam-auxiliary-maslov-sum} give
\eqref{J05:eq:seam-projected-index}.
\end{proof}


\section{A finite order certificate on the five-point seam grid}

\begin{lemma}[New lemma: the five-point exclusion]
\label{J05:lem:seam-size-five-truncation}
Let $n=m=2$ and let $p,q$ be arbitrary retained indices.
Every native complete correction path from $p$ to $q$ has at most two
internal transitions.
\end{lemma}
\begin{proof}
A state is written as its five rows, in column order. For example
$03142$ denotes the permutation $(0,3,1,4,2)$.
All digits in this proof are single row indices, rather than an integer
written in decimal notation. An arrow annotated $(a,b)$ denotes the
actual empty rectangle with ordered column sides $a,b$, followed by
the selected inverse matching. In particular the orientation of every
rectangle in the certificate is specified.

We first give an order calculation that makes the certificate exhaustive.
Every matching source with northeast count one has the following data:
\[
\begin{gathered}
 A\in\{0,1\},\quad h\in\{3,4\},\quad L=1-A,\quad s(A)=h,\quad s(L)=t,\\
 \{P,Q,C\}=\{2,3,4\},\quad P<Q,\quad
 s(P)=v<s(Q)=w,\quad s(C)=7-h,\quad
 \{v,w,t\}=\{0,1,2\}.
\end{gathered}
\]
Thus there are exactly $2\cdot2\cdot3\cdot3=36$ such sources.
The native theorem \decl{flatStep_cases} shows that a transition which
preserves northeast count is one of the four flat types below.
A birth can only exchange $(L,t)$ with $(C,7-h)$.
Its target has a decreasing first southwest pair only if either
$C<P$ and $t>v$, or $P<C<Q$ and $t<v$.
In the first case the wrapping-column orientation contains $(P,v)$.
In the second case the nonwrapping-column orientation contains $(P,v)$.
These facts leave the two birth orientations in the table.
A death from a source with northeast count one reaches a retained state,
so it is an exit rather than an internal transition.
Consequently the following six rows include every possible internal
transition from any of the 36 sources.
\[
\begin{array}{c|c|l}
 \text{type}&\text{ordered columns}&\text{necessary and sufficient conditions}\\\hline
 L_+&(L,P)&t>w,\quad A=0,\quad C>P\\
 L_-&(Q,L)&t<v,\quad A=1,\quad C<Q\\
 H_+&(C,P)&C>Q,\quad h=3,\quad v=0\\
 H_-&(Q,C)&C<P,\quad h=4,\quad w=2\\
 B_+&(L,C)&C<P,\quad t>v,\quad (A=0\ \text{or}\ h=4)\\
 B_-&(C,L)&P<C<Q,\quad t<v,\quad (A=1\ \text{or}\ h=3)
\end{array}
\tag{J05.F5.1}\label{J05:eq:size-five-types}
\]
Here is the emptiness check for every row of this table, including its
excluded cases. For $L_+$, the row data must be $(v,w,t)=(0,1,2)$.
The cyclic row interior consists of rows $3,4$.
The northeast point blocks precisely when $A>L$, and the southeast
point blocks precisely when $C<P$. This gives $A=0$ and $C>P$.
For $L_-$, the row data must be $(t,v,w)=(0,1,2)$, and the same two
high rows form the cyclic interior. Its northeast blocker occurs when
$A<L$, and its southeast blocker occurs when $C>Q$.
This gives $A=1$ and $C<Q$.
For $H_+$, the column data are $(P,Q,C)=(2,3,4)$.
The only interior columns are $0,1$. Their two points avoid the
wrapping row interval exactly when $h<7-h$ and $t>v$.
These conditions are $h=3$ and $v=0$.
For $H_-$, the column data are $(C,P,Q)=(2,3,4)$, again with interior
columns $0,1$. The necessary and sufficient inequalities are
$h>7-h$ and $t<w$, namely $h=4$ and $w=2$.
For $B_+$, the column data are $(C,P,Q)=(2,3,4)$.
If $L=1$ there is no interior column. If $L=0$, its only interior
column is $A=1$, whose high row blocks exactly when $h<7-h$.
For $B_-$, the column data are $(P,C,Q)=(2,3,4)$ and $t=0$.
Column $4$ contains the low row $w=2$, which is outside the wrapping
row interval. The only possible further interior point is at column
$A=0$ when $L=1$, and it blocks exactly when $h>7-h$.
These checks prove all six equivalences in
\eqref{J05:eq:size-five-types}. The other orientation of each flat move
is already excluded by \decl{flatStep_cases}.

An entry can reach only a matching source with northeast count one,
by \decl{seamMatchingEntry_northeast_card}.
Two of the 36 sources admit no entry. For $04312$, its possible entry
rectangles are blocked respectively by the points $(2,3)$ and $(0,0)$.
For $32014$, they are blocked respectively by $(1,2)$ and $(4,4)$.
These are the two candidates in \decl{flatEntry_data}.
Let $E_0$ be the 34 other northeast-count-one sources.
Every actual entry source belongs to $E_0$. No assertion that every
member is reached is required for the argument.

Substitution of the three choices of $(v,w,t)$, the three choices of
$(P,Q,C)$, and the two choices each of $A,h$ in
\eqref{J05:eq:size-five-types} gives the following complete list of
nonempty outgoing internal-arrow sets for sources in $E_0$.
The rows not displayed have no internal arrow.
\[
\begin{array}{c|l@{\qquad}c|l}
03142&(3,0):43012,\ (4,0):23041&03412&(4,0):23401\\
04132&(3,0):34012,\ (4,0):24031&13024&(4,2):13402\\
14302&(0,2):34012,\ (4,2):14023&23014&(4,2):23401\\
24301&(0,2):34021&30142&(3,1):34012\\
31024&(4,2):31402&31402&(1,2):34012\\
32041&(1,2):30142&32401&(1,2):34021\\
40312&(4,2):40123&41302&(1,2):43012,\ (4,2):41023\\
42013&(1,2):40123&42031&(1,2):40132\\
42301&(1,2):43021&&
\end{array}
\tag{J05.F5.2}\label{J05:eq:size-five-entry-arrows}
\]
The target in each cell is obtained by swapping the two displayed
columns and then swapping the first two southwest columns.
For example the row $24301$ uses the birth with sides $(0,2)$.
The actual target is $34201$, whose southwest rows at columns $2,3$
are $2,0$. Its selected partner is $34021$, as displayed.
Thus \eqref{J05:eq:size-five-entry-arrows} records state maps as well as
all the emptiness and source tests in \eqref{J05:eq:size-five-types}.

Only four sources with northeast count two occur among these targets:
$34012$, $34021$, $43012$, and $43021$.
Their complete internal-arrow sets are
\[
\begin{array}{c|l}
34012&\varnothing\\
43012&\varnothing\\
34021&(1,2):30412\\
43021&(0,2):03412,\ (1,2):40312.
\end{array}
\tag{J05.F5.3}\label{J05:eq:size-five-high-arrows}
\]
We verify this table without leaving any orientation unchecked.
There is no northwest or southeast point, so there is no birth and
no flat exchange with either of those quadrants.
Swapping the two northeast points leaves the first southwest pair
increasing. Swapping the first two southwest points is the selected
matching and is excluded from an internal arrow.
An interchange of the second and third southwest points also leaves
the first pair increasing. An interchange of the first and third
can yield a decreasing target pair only for southwest row order
$012$. Its nonwrapping rectangle then contains the second southwest
point, and its wrapping rectangle contains a northeast point.
Thus no flat internal exchange remains.

A death must remove column $2$, since removing either of the other
two southwest columns leaves an increasing first pair.
Its remaining pair is decreasing exactly when the southwest order
is $021$. The wrapping-column orientation contains a southwest point.
For the nonwrapping orientation $(A,2)$, $A=1$ has no interior
column and is empty. For $A=0$, column $1$ blocks exactly when its
high row is greater than the high row in column $0$.
This gives precisely the three arrows in
\eqref{J05:eq:size-five-high-arrows} and their displayed targets.

Let $E_j$ be the set of sources reached by $j$ successive internal
arrows starting from $E_0$. The two explicit arrow tables now give
\[
\begin{split}
E_1={}&\{13402,14023,23041,23401,24031,30142,31402,\\
      &\hspace{13mm}34012,34021,40123,40132,41023,43012,43021\},\\
E_2={}&\{03412,30412,34012,40312\},\\
E_3={}&\{23401,40123\},\qquad E_4=\varnothing.
\end{split}
\tag{J05.F5.4}\label{J05:eq:size-five-reachability}
\]
These are finite set substitutions in the displayed tables, rather
than an inference from tests on further sizes.
For completeness neither member of $E_3$ has an exit.
For $23401$, the two exit rectangles in \decl{flatExit_data} contain
respectively $(2,4)$ and $(0,2)$.
For $40123$, they contain respectively $(1,0)$ and $(4,3)$.
Thus a path with exactly three internal transitions cannot exit.
A path with four or more internal transitions would have to pass
through $E_4$, which is empty. This proves the lemma.
\end{proof}

\section{Fixed-second paths have at most two internal transitions}

Use the signed path domain $D$, its nonnegative projections $D_1,D_2$,
and the shared cell $d=D(N-1,n)$ from
Lemma~\ref{J05:lem:seam-positive-projections}.
The next statements concern the first sector, so the retained endpoints
are $p=\decl{seamState}(x,z)$ and $q=\decl{seamState}(y,z)$.

\begin{lemma}[New lemma: the fixed-second projected domain]
\label{J05:lem:seam-first-domain-shape}
Assume $n,m\geq1$ and $x\ne y$. For any complete correction path
from $p$ to $q$, $D_1$ is the indicator of an empty first-factor
rectangle, $d\in\{0,1\}$, and
\begin{equation}\label{J05:eq:seam-second-domain-column}
 D_2(i,j)=d\,\mathbf1_{\{m\}}(i)
 \qquad(0\leq i,j\leq m).
\end{equation}
For every west cell column $c<n$ and every high cell row $r\geq n$,
one also has
\begin{equation}\label{J05:eq:seam-west-high-copy}
 D(c,r)=D_1(c,n).
\end{equation}
\end{lemma}
\begin{proof}
The endpoints of $D_2$ agree. By
Lemma~\ref{J05:lem:seam-domain-index}, write $D_2(i,j)=a_i+b_j$.
Normalize this expression by subtracting $\min_i a_i$ from each
$a_i$ and adding it to each $b_j$. Then all $a_i$ are nonnegative.
At a column where $a_i=0$, positivity of $D_2$ implies $b_j\geq0$
for every $j$. Therefore
\[
 \iota(D_2;z,z)=2\sum_i a_i+2\sum_j b_j
 \geq 2(a_m+b_0)=2d.
\]
The endpoints of $D_1$ differ, so
Lemma~\ref{J05:lem:seam-index-one-rectangle} gives $\iota(D_1)\geq1$.
The index identity \eqref{J05:eq:seam-projected-index} forces equality in
both inequalities. Hence $D_1$ is an empty rectangle and
$d=D_1(n,n)\in\{0,1\}$. Every $a_i$ with $i\ne m$ and every
$b_j$ with $j\ne0$ vanishes. Put $\alpha=a_m$ and $\beta=b_0$.
Then
\[
 D_2(i,j)=\alpha\mathbf1_{\{m\}}(i)
              +\beta\mathbf1_{\{0\}}(j),\qquad
 \alpha+\beta=d.
\]

There are no retained endpoint points in the northeast region.
For $c<n$ and $1\leq j\leq m$, sum
$\partial D=0$ over columns $0,\ldots,c$ and vertex rows
$n+1,\ldots,n+j$. The two telescoping sums give
\[
 D(c,n+j)-D(N-1,n+j)-D(c,n)+D(N-1,n)=0.
\]
By the projection formulas this is
\[
 D(c,n+j)=D_1(c,n)+D_2(m,j)-d=D_1(c,n)-\beta.
\]
Matching rectangles miss these cells. Each signed multiplicity in this
region is therefore nonnegative. It follows that
$D_1(c,n)\geq\beta$ for every $c<n$.
If $\beta=1$, these inequalities and $D_1(n,n)=d=1$ would make the
empty rectangle $D_1$ cover row $n$ in every cell column of the
first-factor torus. A rectangle's cyclic column interval is proper,
which excludes this possibility. Consequently $\beta=0$ and
$\alpha=d$. This proves \eqref{J05:eq:seam-second-domain-column}.
The displayed telescoping equality now proves
\eqref{J05:eq:seam-west-high-copy} for $r>n$. At $r=n$ it is the
definition of $D_1$.
\end{proof}

\begin{lemma}[New lemma: freezing the unchanged state points]
\label{J05:lem:seam-first-frozen-points}
Under the hypotheses of Lemma~\ref{J05:lem:seam-first-domain-shape},
the following original state points of $p$ remain fixed at every
stage of the correction path, including the intermediate matching
targets.
\begin{enumerate}
\item A west point $(c,x(c))$, with $c<n$, whose column is not an
endpoint column of the first-factor rectangle $D_1$.
\item A southeast point $(c,p(c))$ with $n<c<N-1$ and $p(c)>n$.
\end{enumerate}
\end{lemma}
\begin{proof}
The two endpoint states of an empty rectangle agree outside its two
endpoint columns. At such a common point, the four adjacent cell
multiplicities of the rectangle are zero. Indeed the point is not
on an edge, because a state has only one point in each row and each
column, and it cannot be in the rectangle's interior by emptiness.

Let $(c,x(c))$ be a point in item 1. Its four adjacent cells in the
full seam grid have the same multiplicities as the four adjacent
cells at $(c,x(c))$ in $D_1$.
For rows $1,\ldots,n$ and columns $1,\ldots,n-1$, this is the
restriction defining $D_1$. At west column zero, the left adjacent
cell uses global column $N-1$, which is precisely the last-column
restriction defining $D_1$.
At state row zero, the lower adjacent cell uses global row $N-1$.
Equation~\eqref{J05:eq:seam-west-high-copy} identifies its multiplicity
with that at factor row $n$. If both boundary identifications occur,
$D(N-1,N-1)=d=D_1(n,n)$ follows from
\eqref{J05:eq:seam-second-domain-column}.
Thus all four full-grid multiplicities are zero.
Each of these cells is missed by every selected matching rectangle:
a matching rectangle uses only cell columns $n,\ldots,N-2$ and
cell rows $0,\ldots,n-1$.
Consequently the multiplicity at each of these four cells is the
sum of nonnegative contributions of the actual path rectangles.
Every actual rectangle therefore misses each of the four cells.

For a point in item 2, all four adjacent cells lie in the second
projection. Their column indices are different from its last cell
column $m$. Their multiplicities are zero by
\eqref{J05:eq:seam-second-domain-column}. They also lie outside every
selected matching rectangle, because their cell rows are at least
$n$. Every actual rectangle again misses all four adjacent cells.

In either case, suppose the stated point first changed at some stage.
An inverse matching cannot cause that first change. A point which has
not yet changed is either in a west column, or has a high row, whereas
both matching endpoints are southwest points.
Thus the first change would be an actual rectangle using the point
as a corner. Such a rectangle covers one of its four adjacent cells.
This contradicts the vanishing just proved.
\end{proof}

\begin{lemma}[New lemma: deletion of fixed points preserves correction paths]
\label{J05:lem:seam-delete-fixed-points}
Let a complete correction path on a seam of parameters $n,m$ be given.
Delete any collection of original northwest points and original
southeast points which remain fixed throughout the path.
Delete their rows and columns, and relabel the remaining rows and
columns in increasing order.
If $n'$ west columns and $m'$ high rows remain, the result is a
complete correction path, with the same number of internal
transitions, on the native seam of parameters $n',m'$.
\end{lemma}
\begin{proof}
A retained original state has $n$ west low points, one east low point,
and $m$ east high points. Deleting a northwest point removes one west
column and one low row. Deleting a southeast point removes one east
column and one high row. Thus the remaining columns consist of $n'$
west columns and $m'+1$ east columns, while the remaining rows consist
of $n'+1$ low rows and $m'$ high rows. Increasing relabelling identifies
these precisely with the native cuts at $n'$.

All states on the path restrict to bijections on the remaining rows
and columns, since the deleted point values never change.
Every actual rectangle endpoint remains present, since an actual
rectangle changes both of its endpoint points.
The increasing injections of the remaining row and column sets
preserve cyclic order. A surviving point inside a compressed cyclic
rectangle would therefore have been inside the original rectangle.
Every compressed rectangle is consequently empty, has the prescribed
ordered sides, and exchanges its two endpoint values.

A deleted point is never southwest. Thus the first two southwest
columns of every intermediate matching state are preserved by
restriction. Their row order is preserved as well.
The compressed involution is exactly the native partner which swaps
those first two columns, and its sources have increasing rows there.
Two consecutive matching sources cannot become equal after deletion,
because they agreed on all deleted points before deletion already.
The nonmatching condition on internal transitions is therefore
preserved. The original retained endpoints remain retained, and the
entry, internal rectangles, and exit supply the asserted complete
path of the same length.
\end{proof}

\begin{theorem}[New lemma: all fixed-second long paths are absent]
\label{J05:thm:seam-first-long-empty}
Let $n,m,k\in\mathbb N$, let
$x,y\in\decl{GridState}(n+1)$, and let
$z\in\decl{GridState}(m+1)$. If $k\geq3$, then
\[
 \decl{SeamCorrectionPath}\ k\ (x,z)\ (y,z)=\varnothing.
\]
\end{theorem}
\begin{proof}
The northeast count is at most $\min(n,m)$, since the northeast
columns inject both into the $n$ west columns and into the $m$ high
rows. If $\min(n,m)=0$, there is no matching source, so there is no
correction path. If $\min(n,m)=1$, every matching source on a path
has northeast count one. Theorem~\ref{J05:thm:flat} excludes a complete
path with at least two internal transitions.
We may therefore assume $n,m\geq2$.

If $x=y$, the signed path domain has zero mixed difference.
Its index is even by Lemma~\ref{J05:lem:seam-domain-index}, whereas the
signed-domain index of any correction path is one.
This is impossible. Hence assume $x\ne y$.
By Lemma~\ref{J05:lem:seam-first-domain-shape}, the first projection is
an empty rectangle. At most two of its endpoint columns are west
columns. Lemma~\ref{J05:lem:seam-first-frozen-points} freezes every other
original west point, and freezes every original southeast point
except possibly the points in the two boundary columns $n,N-1$.
Keep the at most two west endpoint points, the original east low
point, and those at most two southeast boundary points.
Delete all other original points.

Lemma~\ref{J05:lem:seam-delete-fixed-points} gives a complete correction
path of the same internal length on a seam with $n'\leq2$ and
$m'\leq2$. If one parameter is zero, a matching source cannot exist.
If $\min(n',m')=1$, every matching source is flat and
Theorem~\ref{J05:thm:flat} excludes the path.
The remaining case is $n'=m'=2$, which is excluded by
Lemma~\ref{J05:lem:seam-size-five-truncation} for $k\geq3$.
All possible compressed sizes have been considered.
\end{proof}

\begin{corollary}[New lemma: full J05 acceptance]
\label{J05:cor:seam-first-complete}
Let $n,m\in\mathbb N$ and let $x,y\in\decl{GridState}(n+1)$.
For every $z\in\decl{GridState}(m+1)$, one has
\[
 \decl{seamLongCorrectionCount}((x,z),(y,z))=0
\]
as an ordinary integer, and hence also after coercion to $\mathbb F_2$.
The first conjunct of \decl{SeamPositivePathComparisons}\ $n\ m$
follows.
\end{corollary}
\begin{proof}
Every summand in the native definition of
\decl{seamLongCorrectionCount} is the cardinality of a path type
with at least three internal transitions.
Theorem~\ref{J05:thm:seam-first-long-empty} makes every such type empty.
The finite sum is therefore zero, including an empty index range.
The existing declaration
\decl{seamPositiveCorrectionCount_first_eq_remaining_add_long}
identifies the positive correction count with the number of remaining
first-factor rectangles plus this long count.
Substitution gives exactly the required first comparison.
\end{proof}

\section{Completion of the acceptance statement}\label{J05:sec:completion}

Theorem~\ref{J05:thm:seam-first-long-empty} and
Corollary~\ref{J05:cor:seam-first-complete} discharge M602 and M604.
Corollary~\ref{J05:cor:three} discharges M603.
There is no remaining gap in the stated J05 acceptance statement.
The long count vanishes as an integer, which is stronger than its required
parity vanishing over $\mathbb F_2$.

The exceptional cases are explicit. A zero seam parameter gives an empty
complement. A parameter equal to one gives only flat matching sources.
Equal retained endpoints are ruled out by the even index of a periodic
domain. A first-factor rectangle with a terminal endpoint leaves at most
one west point after compression, so it belongs to the thin case.
If the original east low point is in a boundary column, at most one east
high point is retained in the compressed grid, again giving the thin case.
The only remaining compressed parameters are $n'=m'=2$, covered in full by
Lemma~\ref{J05:lem:seam-size-five-truncation}.
The truncated long-count sum also vanishes when its indexing range is empty.

\section{Lemma ledger}

In the table, ``existing'' refers to native declarations in the supplied
verified Lean material. A declaration used by a supplied proof may reside in
one of its imported modules. The failed three-step module is excluded from
this category. All new mathematical assertions used above have proofs in
this document.

\small
\begin{longtable}{>{\raggedright\arraybackslash}p{0.44\linewidth}
                  >{\raggedright\arraybackslash}p{0.14\linewidth}
                  >{\raggedright\arraybackslash}p{0.34\linewidth}}
Declaration or statement & Status & Use\\\hline\endhead
\decl{SeamCorrectionPath}, \decl{SeamInternalPath}
& Existing & Actual rectangles and the internal length convention.\\
\decl{mem_seamNortheastColumns}, \decl{mem_seamSouthwestColumns}
& Existing & The native quadrant inequalities.\\
\decl{seamSouthwestColumns_card}
& Existing & Exactly two southwest columns in a flat source.\\
\decl{seamComplementFirst_lt_second}, \decl{SeamComplementSource}
& Existing & Ordered columns and increasing southwest rows.\\
\decl{flatEmpty_source}, \decl{isEmpty_iff_forall_notMem_cIoo}
& Existing & Empty-rectangle point tests.\\
\decl{flatStep_cases}, \decl{flatStep_apply}
& Existing & The four transition formulas in Table~\ref{J05:tab:types}.\\
\decl{flatApply_high}
& Existing & Persistence of the northeast point.\\
\decl{flatEntry_data}, \decl{flatExit_data}
& Existing & The two possible entry and exit rectangles.\\
Lemma~\ref{J05:lem:separate}
& New & The northeast point separates low exchanges by column and high exchanges by row.\\
Lemma~\ref{J05:lem:alternate}
& New & The twelve forbidden successive type pairs.\\
Theorem~\ref{J05:thm:flat}
& New & Persistent blockers exclude every completely flat path of length at least two.\\
\decl{seamMatchingEntry_northeast_card},
\decl{seamMatchingExit_northeast_card}
& Existing & Endpoint counts equal one.\\
\decl{seamCorrectionThree_middle_eq_one}
& Existing & A fixed-second three-step path is completely flat.\\
Corollary~\ref{J05:cor:three}
& New & M603 and the fixed-second length-three vanishing.\\
Lemma~\ref{J05:lem:countbound}
& New & Northeast count is at most $\min(n,m)$.\\
\decl{seamNortheastColumns_nonempty_iff}
& Existing & Positive northeast count in the complement.\\
Corollary~\ref{J05:cor:thin}
& New & J05 when one factor has size at most two.\\
\decl{seamLongCorrectionCount},
\decl{seamFirstPositiveComparison_iff_long_paths}
& Existing & Exact long-count definition and the acceptance equivalence.\\
\decl{GridRectangleBetween.target_eq_swapColumns},
\decl{laurentMaslovDegree_add_card_OColumns}
& Existing & Rectangle boundary and native index calculation.\\
\decl{GridDiagram.maslovOℤ_eq_card}
& Existing & Explicit Maslov point-count formula for the projected index identity.\\
\decl{seamState_left_castSucc}, \decl{seamState_right_of_zero},
\decl{seamState_right_of_ne_zero}
& Existing & Boundary projection to the native retained factors.\\
Lemmas~\ref{J05:lem:seam-domain-index} and
\ref{J05:lem:seam-index-one-rectangle}
& New & Index additivity, periodic domains, and positive index-one classification.\\
Lemmas~\ref{J05:lem:seam-signed-domain} and
\ref{J05:lem:seam-positive-projections}
& New & The signed path domain and its nonnegative factor projections.\\
Lemma~\ref{J05:lem:seam-first-domain-shape}
& New & An empty first-factor rectangle and the last-column second-factor annulus.\\
Lemma~\ref{J05:lem:seam-first-frozen-points}
& New & All unchanged west points and all interior southeast points stay fixed.\\
Lemma~\ref{J05:lem:seam-delete-fixed-points}
& New & Deleting fixed points preserves native correction paths and their length.\\
Lemma~\ref{J05:lem:seam-size-five-truncation}
& New & Complete symbolic five-point certificate for arbitrary retained endpoints.\\
Theorem~\ref{J05:thm:seam-first-long-empty}
& New & All fixed-second paths with at least three internal transitions are absent.\\
\decl{seamPositiveCorrectionCount_first_eq_remaining_add_long}
& Existing & Conversion of long-count vanishing to the first positive comparison.\\
Corollary~\ref{J05:cor:seam-first-complete}
& New & Full J05 acceptance, with integer vanishing.\\
\end{longtable}
\normalsize

\section{Source and verification record}

The source is the supplied \decl{combined_01-10.md}, specifically the audit
entry J05, records M602--M604, the mathematical writing rules, and the seam
listings under item 07, together with the native Maslov point-count formula
used in item 08. The definitions and proofs used for flat transitions
belong to \decl{GridConnectedSumSeamFlatTwoStepBarrier.lean}.
The fixed-second three-step classification belongs to
\decl{GridConnectedSumSeamThreeStepSectors.lean}.
The source snapshot is treated as read-only. No Lean code is written or
modified. The supplied material does not include
\decl{snapshot-sources.sha256}, so the audit's repository hash comparison
cannot be performed from these attachments.

No holomorphic Floer theorem is a proof step in this document.
The five-point calculation is a finite written order certificate whose
candidate types, blocker inequalities, state maps, and reachability sets
are all displayed. Its role is justified by the proved support compression
from arbitrary sizes. Diagnostic computer enumerations are not used as
proof steps. All other arguments use finite permutations, integral square
domains, the native rectangle conventions, and the cited existing declarations.


\chapter{J06: Second-sector seam parity}
\label{chapter:J06}



\section{Exact target and status}

The acceptance statement of J06 is the following assertion. For every
$n,m\in\nat$, every $x\in\operatorname{GridState}(n+1)$, and every
$z,w\in\operatorname{GridState}(m+1)$,
\begin{equation}\label{J06:eq:target}
 \begin{aligned}
 &\decl{seamPositiveCorrectionCount}((x,z),(x,w))\\
 &\qquad=\seamcard{\decl{SeamSecondRemainingRectangle}(x,z,w)}
 \quad\text{in }\F.
 \end{aligned}
\end{equation}
The counts in this formula are the native counts in the supplied listings.

This document proves the acceptance statement for all parameters.
In fact the following stronger equality holds over the natural numbers:
\begin{equation}\label{J06:eq:integer}
 \begin{aligned}
 &\decl{seamPositiveCorrectionCount}((x,z),(x,w))\\
 &\qquad=\seamcard{\decl{SeamSecondRemainingRectangle}(x,z,w)}.
 \end{aligned}
\end{equation}
The proof first constructs an explicit bijection between the one-transition
and two-transition correction paths and the remaining second-factor
rectangles. It then proves that every correction with three or more
internal transitions is absent.

The long-path proof uses signed integral domains and their two
nonnegative projections. Its index calculation forces the second
projection to be an empty rectangle. All but at most five state points
must then remain fixed throughout the path. Deleting those fixed points
reduces to four explicitly described small seam grids, whose complete
transition certificates are proved below. This is a uniform reduction
and a finite written proof, with no inference from diagnostic computations.

All chain coefficients and parity statements use $\F$. The integer
cardinalities, domain multiplicities, and gradings retain their native
integer meanings. The proof uses no holomorphic Floer result and no
unproved assertion from the failing seam module.

\section{Native conventions and notation}

Put $N=n+m+1$ and identify $\operatorname{Fin}(N)$ with
$\{0,\ldots,N-1\}$. A state is a permutation from columns to rows.
For distinct $a,b\in\operatorname{Fin}(N)$, write
\[
 I_N(a,b)=
 \begin{cases}
  \{c:a<c<b\},&a<b,\\
  \{c:c>a\text{ or }c<b\},&b<a.
 \end{cases}
\]
This is the open cyclic interval used by \decl{Grid.cIoo}. For a state $s$,
an oriented rectangle with side columns $(a,b)$ has bottom row $s(a)$,
top row $s(b)$, and target $s\circ(a\ b)$. It is empty precisely when
\begin{equation}\label{J06:eq:empty}
 \text{there is no }c\in I_N(a,b)\text{ with }s(c)\in I_N(s(a),s(b)).
\end{equation}
The supplied proofs use these statements as
\decl{bottom_def}, \decl{top_def}, \decl{target_eq_swapColumns}, and
\decl{isEmpty_iff_forall_notMem_cIoo}.
A rectangle is determined by its source and its ordered side columns.

The retained state $p=\decl{seamState}(x,z)$ has coordinates
\begin{equation}\label{J06:eq:seam}
 p(c)=x(c)\quad(0\le c<n),\qquad
 p(n+j)=
 \begin{cases}x(n),&z(j)=0,\\ n+z(j),&z(j)>0.
 \end{cases}
\end{equation}
These are the supplied declarations
\decl{seamState_left_castSucc},
\decl{seamState_right_of_zero}, and
\decl{seamState_right_of_ne_zero}.

A northeast point has column less than $n$ and row greater than $n$.
A southwest point has column at least $n$ and row at most $n$.
For a complementary state $s$, let $P<Q$ be its first two southwest
columns. Its partner, denoted $s^*$, exchanges the values at $P,Q$.
A matching source is a complementary state with $s(P)<s(Q)$.
These are the coordinate properties of the native matching used in the
supplied proofs, including
\decl{seamComplementPartner_apply_first},
\decl{seamComplementPartner_apply_second}, and
\decl{seamComplementPartner_firstBlock_apply}.
The definitions of the foundational sector and matching modules are
imported by the supplied listings rather than reproduced there. The
coordinate properties just listed are also stated in Section 2.3 of the
supplied seam notes and used explicitly throughout the supplied proofs.
No assertion from the failing flat-three-step module is used.

For $k\in\nat$, abbreviate
\[
 \begin{aligned}
 \cP_k(x,z,w)&=
 \decl{SeamCorrectionPath}(k,(x,z),(x,w)),\\
 \cR(x,z,w)&=\decl{SeamSecondRemainingRectangle}(x,z,w).
 \end{aligned}
\]
A path in $\cP_k$ consists of matching sources $s_0,\ldots,s_k$,
an empty entry rectangle $p\longrightarrow s_0^*$, empty internal
rectangles $s_{j-1}\longrightarrow s_j^*$ for $1\le j\le k$, and an
empty exit rectangle $s_k\longrightarrow q$, where
$q=\decl{seamState}(x,w)$. The consecutive matching sources must differ.
This is exactly the definition of \decl{SeamCorrectionPath}, together
with \decl{SeamInternalRectangle} and \decl{SeamInternalPath}.
Thus $k$ counts internal rectangles, and there are $k+2$ rectangles in
all.

\section{The remaining rectangles}

\begin{lemma}[New lemma: coordinate description of the remaining rectangles]
\label{J06:lem:remaining}
Let $n\ge1$, let $x\in\operatorname{GridState}(n+1)$, and let
$z,w\in\operatorname{GridState}(m+1)$. A rectangle from $z$ to $w$
belongs to $\cR(x,z,w)$ if and only if it is empty and its ordered
side columns $\ell,r$ and endpoint rows satisfy
\[
 r<\ell,\qquad 0<z(r)<z(\ell).
\]
For such a rectangle, if $c=z^{-1}(0)$, then $r<c<\ell$.
\end{lemma}

\begin{proof}
By the definition of \decl{SeamSecondRemainingRectangle}, the factor
rectangle is empty, its lifted rectangle is obstructed, and both endpoint
rows are nonzero. The supplied declaration
\decl{seamSecondRemainingRectangle_wraps} gives $r<\ell$.
If $z(\ell)<z(r)$, the rows in the interior of the lifted rectangle are
strictly greater than $n$. No first-block point of \eqref{J06:eq:seam} can
lie there. The second-block points are excluded by emptiness of the
factor rectangle. This contradicts obstruction. Thus
$0<z(r)<z(\ell)$.

Conversely, suppose the displayed inequalities hold. Every first-block
column lies in $I_N(n+\ell,n+r)$, and its row is at most $n$. Such a
row lies in $I_N(n+z(\ell),n+z(r))$. Since $n\ge1$, a first-block
point exists, and the lifted rectangle is obstructed. The nonzero-row
conditions in the definition of $\cR$ hold.

Finally, $0$ lies in $I_{m+1}(z(\ell),z(r))$. Therefore factor emptiness
excludes $c$ from $I_{m+1}(\ell,r)$. It is distinct from both endpoint
columns, so $r<c<\ell$.
\end{proof}

Fix a rectangle as in Lemma~\ref{J06:lem:remaining}. For the rest of the
construction put
\begin{equation}\label{J06:eq:coordinates}
 R=n+r,\quad C=n+c,\quad L=n+\ell,\quad
 K=n+z(r),\quad H=n+z(\ell),\quad t=x(n).
\end{equation}
Then
\begin{equation}\label{J06:eq:orders}
 n\le R<C<L<N,\qquad n<K<H<N,\qquad 0\le t\le n.
\end{equation}
The retained source has $p(R)=K$, $p(C)=t$, and $p(L)=H$.
For every second-block column $D$ with $D<R$ or $D>L$, factor emptiness
also gives
\begin{equation}\label{J06:eq:outside}
 K<p(D)<H.
\end{equation}
Indeed, $D$ is not $C$, so its row is high, and a row outside the open
interval $(K,H)$ would be an interior point of the factor rectangle.
Equality with $K$ or $H$ is excluded by injectivity of the state.

\section{The explicit inverse construction}

\begin{lemma}[New lemma: boundary and interior correction paths]
\label{J06:lem:construction}
Under \eqref{J06:eq:coordinates}--\eqref{J06:eq:orders}, the following rules
associate a native correction path to the chosen remaining rectangle.
If $t=0$ or $t=n$, the path belongs to $\cP_1(x,z,w)$.
If $0<t<n$, it belongs to $\cP_2(x,z,w)$.
\end{lemma}

\begin{proof}
Only the columns in the following tables change. Every omitted column
retains its value from $p$. In each table, the last row is
$q=\decl{seamState}(x,w)$.

Suppose first that $t=0$. Let $B=x^{-1}(n)$. Since $n\ge1$, one has
$B<n$. The values are
\[
\begin{array}{c|rrrr}
 &B&R&C&L\\\hline
 p   &n&K&0&H\\
 s_0 &K&0&n&H\\
 s_1 &K&H&0&n\\
 q   &n&H&0&K
\end{array}
\]
and the ordered rectangle sides are
\begin{equation}\label{J06:eq:zero-path}
 (B,R),\qquad (L,R),\qquad (L,B).
\end{equation}
The first two southwest columns of $s_0$ are $R,C$ with values $0,n$.
Those of $s_1$ are $C,L$ with values $0,n$.
Swapping at the indicated rectangle sides and then applying the matching
inverse gives exactly the next table row. The internal sources differ.

Suppose next that $t=n$. Let $A=x^{-1}(0)$, so $A<n$. The values are
\[
\begin{array}{c|rrrr}
 &A&R&C&L\\\hline
 p   &0&K&n&H\\
 s_0 &H&K&0&n\\
 s_1 &H&0&n&K\\
 q   &0&H&n&K
\end{array}
\]
and the ordered rectangle sides are
\begin{equation}\label{J06:eq:last-path}
 (L,A),\qquad (L,R),\qquad (A,R).
\end{equation}
The matching pairs are $C,L$ in $s_0$ and $R,C$ in $s_1$.
Their values are again $0,n$ in increasing column order.

Finally, suppose $0<t<n$. Let $A=x^{-1}(0)$ and $B=x^{-1}(n)$.
Both columns are less than $n$, and they are distinct. No relative
order between $A$ and $B$ is assumed. The values are
\begin{equation}\label{J06:eq:interior-table}
\begin{array}{c|rrrrr}
 &A&B&R&C&L\\\hline
 p   &0&n&K&t&H\\
 s_0 &H&n&K&0&t\\
 s_1 &H&K&0&n&t\\
 s_2 &0&K&H&t&n\\
 q   &0&n&H&t&K
\end{array}
\end{equation}
and the ordered rectangle sides are
\begin{equation}\label{J06:eq:interior-path}
 (L,A),\qquad (B,R),\qquad (A,R),\qquad (L,B).
\end{equation}
The southwest columns of $s_0$ are $C,L$, with values $0,t$.
Those of $s_1$ are $R,C,L$, with values $0,n,t$.
Those of $s_2$ are $C,L$, with values $t,n$.
In each source the first pair is increasing. The first internal step
increases the northeast count from one to two, and the second decreases
it to one. In particular, the consecutive matching sources differ.
The transpositions in \eqref{J06:eq:interior-path}, followed by the specified
matching inverses, give exactly \eqref{J06:eq:interior-table}.

It remains to verify emptiness of every rectangle. The open row interval
of each rectangle is either $I_N(H,0)$ or $I_N(n,K)$.
The former consists of rows greater than $H$, and the latter consists
of rows strictly between $n$ and $K$.
Every unchanged first-block row is at most $n$, so it belongs to neither
interval. The exceptional high first-block value is $K$ for a rectangle
with row interval $I_N(H,0)$, or $H$ for a rectangle with row interval
$I_N(n,K)$. It also lies outside the relevant open row interval.
Endpoint values do not count as interior points.

For second-block columns, each rectangular column interior is contained
in $\{D:D<R\}\cup\{D:D>L\}$. The changed low columns in the tables
are endpoints or lie strictly between $R$ and $L$, and hence do not
provide an interior point. Every other second-block point in one of
these rectangles has row strictly between $K$ and $H$ by
\eqref{J06:eq:outside}. Such a row lies in neither of the two open row
intervals. Equation~\eqref{J06:eq:empty} proves emptiness in all cases.
The tables therefore define the required native paths.
\end{proof}

\section{Every short path is one of these paths}

\begin{lemma}[New lemma: uniqueness for one transition]
\label{J06:lem:one}
Let $n\ge1$. For fixed $x,z,w$, the one-transition paths in
$\cP_1(x,z,w)$ are in bijection with $\cR(x,z,w)$ when $x(n)$ is
$0$ or $n$. There are no such paths when $0<x(n)<n$.
Under the bijection, the inverse is \eqref{J06:eq:zero-path} or
\eqref{J06:eq:last-path}, respectively.
\end{lemma}

\begin{proof}
If $z=w$, the supplied declaration
\decl{seamTransferPath_self_isEmpty} excludes every correction path,
and a rectangle from $z$ to itself cannot exist. Thus assume $z\ne w$.
The declaration \decl{seamOneStep_all_orientations} assigns one of four
neutral types to a path. The declaration
\decl{seamSecondOneStep_kind_high} excludes the two low types because
the second factor changes. For the high-first type,
\decl{seamSecondOneStep_high_first_boundary} gives $x(n)=0$.
The declarations
\decl{seamSecondOneStep_high_first_row_zero},
\decl{seamSecondOneStep_high_first_row_last}, and
\decl{seamSecondOneStep_high_first_pivot} identify the two southwest
rows as $0,n$ and the first-block pivot as $x^{-1}(n)$.
The corresponding high-second declarations identify $x(n)=n$, the
same extremal southwest rows, and pivot $x^{-1}(0)$.
These declarations give the interior nonexistence assertion, also stated
explicitly as \decl{seamCorrectionOne_second_interior_isEmpty}.

The supplied definitions
\decl{seamSecondOneStepHighFirstRemaining} and
\decl{seamSecondOneStepHighSecondRemaining} project these two cases
to a member of $\cR(x,z,w)$. They include proofs of emptiness,
obstruction, and nonzero endpoint rows.
In the first case, the projected rectangle has left column equal to the
external high column and right column equal to the first southwest
column of the initial matching source. Its zero column is the second
southwest column of that source. The source-value declarations
\decl{seamSecondOneStepHighFirst_source_values} therefore identify
these columns as $L,R,C$ of \eqref{J06:eq:coordinates}.
The orientation declaration gives the three sides
$(B,R),(L,R),(L,B)$, and its state values are the first table in
Lemma~\ref{J06:lem:construction}.
In the second case, the corresponding source-value declaration
\decl{seamSecondOneStepHighSecond_source_values} identifies the
columns $L,R,C$ and gives the sides $(L,A),(L,R),(A,R)$ and the
second table there.

The projected rectangle and $x$ thus determine every side and every
state of the original path. Lemma~\ref{J06:lem:construction} proves that
these recovered data form a path for every remaining rectangle.
The projection and construction are inverse maps.
\end{proof}

\begin{lemma}[New lemma: normal form of a two-transition path]
\label{J06:lem:normal}
Let $p,q$ be retained states, and let a two-transition correction join
them. Write its matching sources as $a_0,a_1,a_2$.
The initial source has a unique northeast column $A<n$, of row $H>n$,
and two southwest columns $P<Q$, of rows $a<b\le n$.
There are columns $C<n$ and $D\ge n$ such that
\[
 C\ne A,\qquad D<P<Q,\qquad a<b<c\le n,
 \qquad a_0(C)=c,\qquad a_0(D)=K>n.
\]
The first internal rectangle has sides $(C,D)$.
The second has sides $(B,D)$, where $B$ is either $A$ or $C$.
The entry has sides $(Q,A)$, and the exit has sides $(Q,E)$, where
$E$ is the remaining northeast column.
If $B=C$, the first retained factors differ.
If the first retained factors are equal, then $B=A$, $H>K$, and the
retained endpoints differ by exchanging the high values in columns
$D,Q$ while fixing every other column.
\end{lemma}

\begin{proof}
The supplied declarations \decl{seamMatchingEntry_northeast_card} and
\decl{seamMatchingExit_northeast_card} give one northeast point at
each end. The declarations \decl{seamMatchingEntry_southwest_card}
and \decl{seamMatchingExit_southwest_card} give two southwest
points there. The supplied theorem
\decl{seamCorrectionTwo_middle_eq_two} gives two at the middle source.
Thus the first step exchanges a low first-block column $C$ with a high
second-block column $D$, and the second exchanges one of the two high
first-block columns $B\in\{A,C\}$ with a low second-block column.
This assertion follows by counting the first-block high rows before
and after a transposition. A transposition inside either block preserves
that count, as formalized by
\decl{seamRectangle_northeast_card_eq_of_same_block}.

Before the first matching inverse, the low second-block columns are
$P,Q,D$, with values $a,b,c$ at those respective columns.
If $Q<D$, the first pair has the increasing values $a,b$. After the
matching inverse this pair would be decreasing, contrary to the
matching-source condition for $a_1$.
If $P<D<Q$, that condition requires $c<a$, and the values of $a_1$
in increasing low-column order are $c,a,b$. Deleting any one of these
three columns leaves an increasing pair. The second matching inverse
would make that pair decreasing. Thus this case admits no death step.
The only possibility is $D<P<Q$. The values of $a_1$ in these columns
are $a,c,b$, with $a<c$. Deleting $P$ or $Q$ again leaves an increasing
pair. Hence the death removes $D$, and the remaining pair must satisfy
$c>b$ before its inverse. This proves $a<b<c$.

The reversed orientation $(D,C)$ for the birth would contain the
point at $P$, whose row $a$ is less than $c$. The reversed orientation
$(D,B)$ for the death would contain the point at $P$, whose row $c$
is between $a$ and the high value at $B$. Emptiness therefore forces
$(C,D)$ and $(B,D)$.

Consider $B=C$ first. If $C<A$, the point at $A$ obstructs the birth
when $H<K$ and obstructs the death when $K<H$. The two values are
distinct, so $A<C$. An entry of sides $(A,P)$ would contain the point
at $C$, whose retained row is $c\in(b,H)$. An exit of sides $(A,P)$
would contain the point at $C$, whose final matching-source row is
$a\in I_N(H,b)$. Hence both outer rectangles have their first-block
endpoint on the right, namely $(Q,A)$ and $(Q,A)$. The retained source
has rows $a,c$ at $A,C$, and the retained target has rows $c,a$ there.
Since $a<c$, their first factors differ.

Now consider $B=A$. If $H<K$ and $C<A$, the point at $A$ obstructs
the birth. If $H<K$ and $A<C$, the point at $C$ obstructs the death.
Thus $K<H$. An entry $(A,P)$ would contain $D$, whose row $K$ lies
strictly between $b$ and $H$. An exit $(C,P)$ would contain $D$, whose
row $H$ lies in the cyclic interval from $K$ to $b$. The outer sides
are therefore $(Q,A)$ and $(Q,C)$.

The values on the five relevant columns are now
\begin{equation}\label{J06:eq:general-table}
\begin{array}{c|rrrrr}
 &A&C&D&P&Q\\\hline
 p   &a&c&K&b&H\\
 a_0 &H&c&K&a&b\\
 a_1 &H&K&a&c&b\\
 a_2 &a&K&H&b&c\\
 q   &a&c&H&b&K .
\end{array}
\end{equation}
Every other column stays fixed. This proves the last assertion and
completes the normal form.
\end{proof}

\begin{lemma}[New lemma: uniqueness for two transitions]
\label{J06:lem:two}
Let $n\ge1$. For fixed $x,z,w$, the two-transition paths in
$\cP_2(x,z,w)$ are in bijection with $\cR(x,z,w)$ when
$0<x(n)<n$. There are no such paths when $x(n)$ is $0$ or $n$.
The inverse bijection is \eqref{J06:eq:interior-path}.
\end{lemma}

\begin{proof}
Apply Lemma~\ref{J06:lem:normal} to a path with first factor $x$ at both
endpoints. Its second case gives \eqref{J06:eq:general-table} and the
sides
\[
 (Q,A),\qquad(C,D),\qquad(A,D),\qquad(Q,C).
\]
The only low second-block value of the retained source is $b$, so
$x(n)=b$.

We first prove $a=0$. If a row less than $a$ existed, its preimage in
$p$ would be a first-block column $V$, distinct from $A,C$.
For $V<A$, the entry $(Q,A)$ would contain that point.
For $A<V$, the death $(A,D)$ would contain that point.
Its value is unchanged at these steps. Both alternatives contradict
emptiness. Since the state contains row $0$, one has $a=0$.
Similarly, if a row greater than $c$ but at most $n$ existed, its
preimage would be a first-block column $V$ distinct from $A,C$.
For $C<V$, the birth $(C,D)$ would contain it.
For $V<C$, the exit $(Q,C)$ would contain it.
Thus $c=n$. In particular, $0<x(n)<n$, $A=x^{-1}(0)$,
and $C=x^{-1}(n)$.

The retained endpoints differ in columns $D,Q$, with rows $K<H$.
Project the rectangle of sides $(Q,D)$ to the second factor.
Its endpoint rows are nonzero and its two intervals wrap.
To prove factor emptiness, consider a second-block column $V<D$.
A row greater than $H$ there would obstruct the death, and a row less
than $K$ would obstruct the birth. For a second-block column $V>Q$,
a row greater than $H$ would obstruct the entry, and a row less than
$K$ would obstruct the exit. The low second-block column $P$ is
strictly between $D$ and $Q$. Thus the projected rectangle is empty.
Lemma~\ref{J06:lem:remaining} makes it a member of $\cR(x,z,w)$.

This projection identifies $D,P,Q$ with $R,C,L$ of
\eqref{J06:eq:coordinates}. Replacing the temporary first-block column
$C$ by $B=x^{-1}(n)$ transforms \eqref{J06:eq:general-table} into
\eqref{J06:eq:interior-table}. It also transforms its four sides into
\eqref{J06:eq:interior-path}. Hence the projection determines every state
and rectangle of the original path. Lemma~\ref{J06:lem:construction}
proves that every remaining rectangle supplies that path. These maps
are inverse.
\end{proof}

\begin{theorem}[New lemma: complete second-sector short-path bijection]
\label{J06:thm:short}
For every $n,m\in\nat$, every $x\in\operatorname{GridState}(n+1)$,
and every $z,w\in\operatorname{GridState}(m+1)$, there is a bijection
\[
 \cP_1(x,z,w)\ \sqcup\ \cP_2(x,z,w)
 \ \longrightarrow\ \cR(x,z,w).
\]
Its inverse uses one transition when $x(n)$ is a boundary row and
two transitions when $0<x(n)<n$.
\end{theorem}

\begin{proof}
For $n\ge1$, Lemmas~\ref{J06:lem:one} and \ref{J06:lem:two} give mutually
exclusive cases according to the value of $x(n)$.
When $n=0$, there are no first-block columns. Every state is retained,
so there are no matching sources or correction paths. The seam state
is $z$, its second-factor rectangle lift is the original rectangle,
and no remaining rectangle is obstructed. All three sets are empty.
\end{proof}

\section{Projected integral domains of a seam correction}

All path coefficients remain in $\mathbb F_2$. Integral domains in this
section record the multiplicities of grid squares before any parity is
taken. Their use does not change the coefficient field of the chain
complex.

For a grid of size $h$, an integral domain is an integer-valued function
$D(c,r)$ on the $h^2$ squares. All cell and vertex indices in the next
two lemmas are cyclic modulo $h$. Define its mixed difference and corner
average by
\[
 \begin{split}
 (\partial D)(c,r)
 &=D(c,r)+D(c-1,r-1)-D(c-1,r)-D(c,r-1),\\
 (\mathcal A D)(c,r)
 &=\tfrac14\bigl(D(c,r)+D(c-1,r-1)+D(c-1,r)+D(c,r-1)\bigr).
 \end{split}
\]
For a state $u$, let $[u]$ be the indicator function of its points.
A domain from $u$ to $v$ satisfies $\partial D=[u]-[v]$. Its index is
\[
 \iota(D;u,v)=\sum_{c=0}^{h-1}
 \bigl((\mathcal A D)(c,u(c))+(\mathcal A D)(c,v(c))\bigr).
\]
This is a definition by finite sums. In particular it makes sense for
domains with negative coefficients.

\begin{lemma}[New lemma: domain index identities]
\label{J06:lem:seam-domain-index}
The following statements hold for integral domains on a toroidal grid.
\begin{enumerate}
\item If $D$ is a domain from $u$ to $v$ and $E$ is a domain from
$v$ to $w$, then
\[
 \iota(D+E;u,w)=\iota(D;u,v)+\iota(E;v,w).
\]
\item A domain with mixed difference zero has the form
$D(c,r)=a_c+b_r$, with integers $a_c,b_r$. Its index between any equal
endpoint states is $2\sum_c a_c+2\sum_r b_r$.
\item An empty native rectangle has index one. For any auxiliary
grid with marking permutation $O$ and any domain from $u$ to $v$,
\begin{equation}\label{J06:eq:seam-domain-maslov}
 \iota(D;u,v)=M_O(u)-M_O(v)+2\sum_c D(c,O(c)).
\end{equation}
Here $M_O$ is the native integer Maslov grading.
\end{enumerate}
\end{lemma}

\begin{proof}
Finite cyclic reindexing gives
\[
 \sum_{c,r}(\partial D)(c,r)(\mathcal A E)(c,r)
 =\sum_{c,r}(\partial E)(c,r)(\mathcal A D)(c,r).
\]
For completeness, after moving the differences from $D$ to $E$, the
coefficient of $D(c,r)$ on the left is
\[
 \tfrac14\bigl(E(c+1,r+1)-E(c+1,r-1)
             -E(c-1,r+1)+E(c-1,r-1)\bigr).
\]
Its coefficients are unchanged when both offsets are negated. This
proves the displayed symmetry. Substituting
$\partial D=[u]-[v]$ and $\partial E=[v]-[w]$ proves the first assertion
by expanding the definition of $\iota$.

If $\partial D=0$, the difference $D(c,r)-D(c-1,r)$ is independent
of $r$. Fixing one row and one column and summing these differences
gives $D(c,r)=a_c+b_r$. Each state occupies every row and every column
once. Substitution into the corner average therefore gives index
$2\sum a_c+2\sum b_r$.

The domain of a rectangle with sides $(a,b)$ and source $u$ is the
indicator of
\[
 [a,b)_{\mathrm{cyc}}\times[u(a),u(b))_{\mathrm{cyc}}.
\]
Its mixed difference is $[u]-[u\circ(a\ b)]$ by the native declaration
\decl{GridRectangleBetween.target_eq_swapColumns}.
If the rectangle is empty, its four endpoint corners each contribute
$1/4$ to the index and every other state point contributes zero.
The native declaration \decl{laurentMaslovDegree_add_card_OColumns}
then proves \eqref{J06:eq:seam-domain-maslov} for an empty rectangle.

To obtain the formula for an arbitrary domain, connect $u$ to $v$ by
adjacent-column transpositions. Such a sequence exists by moving the
desired value into the first incorrect column and repeating on the
remaining columns. Use the rectangle of width one for each transposition.
It is empty because there is no interior column. Their sum $D_0$ has
the required boundary, and the formula for $D_0$ follows by additivity
and telescoping of the native grading differences. The difference
$D-D_0$ has zero mixed difference. For a row or column annulus the
index is two and the sum of the $O$ multiplicities is one.
The second assertion, with arbitrary integer coefficients, now gives
the formula for $D-D_0$ and hence for $D$.
\end{proof}

\begin{lemma}[New lemma: nonnegative domains of index one]
\label{J06:lem:seam-index-one-rectangle}
Let $D$ be a nonnegative integral domain from distinct states $u$ to
$v$. Then $\iota(D;u,v)\ge1$. If its index is one, then $D$ is the
indicator domain of an empty native rectangle from $u$ to $v$.
\end{lemma}

\begin{proof}
Suppose that the states differ in $e$ columns. Then $e\ge2$, and the
mixed difference has a coefficient $1$ or $-1$ at each of their $2e$
distinct changed points. At each such point at least one adjacent cell
has a positive integer coefficient. Its corner average is therefore
at least $1/4$. All remaining contributions are nonnegative, so
$\iota(D;u,v)\ge e/2\ge1$.

If the index is one, this bound gives $e=2$. Write the changed columns
as $A,B$, and put $R=u(A)$ and $S=u(B)$. The target exchanges these
two rows. Let $V$ and $W$ be the nonempty proper cyclic cell intervals
$[A,B)$ and $[R,S)$, and let $D_0$ be the indicator of $V\times W$.
By Lemma~\ref{J06:lem:seam-domain-index},
\[
 D(c,r)=\mathbf1_V(c)\mathbf1_W(r)+a_c+b_r.
\]
Replace $a_c$ by its minimum separately on $V$ and its complement,
and replace $b_r$ by its minimum separately on $W$ and its complement.
The resulting domain $D_b$ is nonnegative, because on each of the four
products of intervals its value is the minimum of $D$ there.
The difference $D-D_b$ is a sum of row and column annuli with
nonnegative integer coefficients. The domain $D_b$ has the same
boundary, so its index is at least one by the preceding paragraph.
Each annulus has index two. Additivity and $\iota(D)=1$ force every
annulus coefficient to be zero. Thus $D=D_b$ is constant on each of
the four products of intervals.

Let the four constants be $a,b,c,d$, with $a$ on $V\times W$ and
$d$ on its opposite product. They are nonnegative integers and satisfy
$a+d-b-c=1$. The four changed points contribute exactly
$a+b+c+d$ to the index. Hence this sum is at most one, and the
only possibilities are $(a,b,c,d)=(1,0,0,0)$ and $(0,0,0,1)$.
These are the two oriented rectangle domains from $u$ to $v$.
The four corners already account for the full index. Every other
source-state corner average is zero, so there is no source point in
the rectangle's interior. This is the native emptiness condition
\decl{isEmpty_iff_forall_notMem_cIoo}.
\end{proof}

Now let $n,m\ge1$ and $N=n+m+1$. For a correction path from retained
state $p=\decl{seamState}(x,z)$ to retained state
$q=\decl{seamState}(y,w)$, define $D$ to be the sum of its entry,
internal, and exit rectangle domains, minus the domains of the selected
matching rectangles at all its matching sources.

\begin{lemma}[New lemma: the signed correction domain]
\label{J06:lem:seam-signed-domain}
The signed domain $D$ of a correction path satisfies
$\partial D=[p]-[q]$ and
\begin{equation}\label{J06:eq:seam-path-index}
 \iota(D;p,q)=1.
\end{equation}
\end{lemma}
\begin{proof}
The matching rectangle at a source with southwest columns $P<Q$ and
rows $a<b\le n$ is $[P,Q)\times[a,b)$. It is empty, because every
column strictly between the first two southwest columns has a high row.
Telescoping the rectangle boundaries gives $\partial D=[p]-[q]$.
There is one more positive rectangle than subtracted matching rectangles.
Each has index one, so additivity in Lemma~\ref{J06:lem:seam-domain-index}
gives \eqref{J06:eq:seam-path-index}.
\end{proof}

Define two restrictions, with the indicated factor indices, by
\begin{equation}\label{J06:eq:seam-domain-projections}
 \begin{aligned}
 D_1(i,j)&=D(i,j) &&(0\le i<n,\ 0\le j\le n),\\
 D_1(n,j)&=D(N-1,j) &&(0\le j\le n),\\
 D_2(i,j)&=D(n+i,n+j) &&(0\le i,j\le m).
 \end{aligned}
\end{equation}
Put $d=D(N-1,n)$.

\begin{lemma}[New lemma: positive projections and the index relation]
\label{J06:lem:seam-positive-projections}
For every correction path as above, $D_1$ is a nonnegative domain from
$x$ to $y$, and $D_2$ is a nonnegative domain from $z$ to $w$.
Their indices satisfy
\begin{equation}\label{J06:eq:seam-projected-index}
 \iota(D_1;x,y)+\iota(D_2;z,w)=1+2d.
\end{equation}
\end{lemma}

\begin{proof}
A matching rectangle uses only columns from $n$ through $N-2$ and
rows from $0$ through $n-1$. It therefore misses both restrictions in
\eqref{J06:eq:seam-domain-projections}. All remaining summands of $D$ are
nonnegative rectangle indicators, proving positivity.

For the boundary assertion, restricting a domain to cyclic lists of
cell columns $c_i$ and cell rows $r_j$ replaces its mixed difference
at $(i,j)$ by
\[
 \sum_{c\in(c_{i-1},c_i]_{\mathrm{cyc}}}
 \sum_{r\in(r_{j-1},r_j]_{\mathrm{cyc}}}(\partial D)(c,r).
\]
This formula follows by telescoping in both coordinates.
For $D_1$, the last column collects the original columns
$n,\ldots,N-1$, and row zero collects the original rows
$n+1,\ldots,N-1,0$. Under these collections the retained state $p$
becomes $[x]+m[(n,0)]$. The extra $m$ copies are the high points in
the second columns. The retained state $q$ gives
$[y]+m[(n,0)]$, so the extra terms cancel.
For $D_2$, column zero collects columns $0,\ldots,n$, and row zero
collects rows $0,\ldots,n$. The retained states become
$[z]+n[(0,0)]$ and $[w]+n[(0,0)]$. Again the extra terms cancel.
These statements use precisely the three supplied coordinate formulas
\decl{seamState_left_castSucc}, \decl{seamState_right_of_zero}, and
\decl{seamState_right_of_ne_zero}.

We prove the index relation by explicit auxiliary markings.
Let $O_1(i)=i$ on the first factor, let
$O_2(j)=j+1\pmod{m+1}$ on the second factor, and let
\[
 O(c)=c\quad(c<n),\qquad
 O(n+j)=n+O_2(j)\quad(0\le j\le m).
\]
Each is a permutation and can be made into a valid auxiliary grid by
choosing $X$ to be its row successor. The Maslov grading depends only on
the specified $O$ markings. The full marking set is the union of the
two projected marking sets with their common square $(N-1,n)$ counted
once. Consequently
\begin{equation}\label{J06:eq:seam-marking-domain-count}
 \sum_c D(c,O(c))=
 \sum_iD_1(i,O_1(i))+\sum_jD_2(j,O_2(j))-d.
\end{equation}

The native Maslov point-count formula gives
\begin{equation}\label{J06:eq:seam-auxiliary-maslov-sum}
 M_O(\decl{seamState}(x,z))=M_{O_1}(x)+M_{O_2}(z).
\end{equation}
Here is a direct verification, so no condition on the auxiliary
$X$ markings is being suppressed. For finite planar point sets $A,B$,
write $I(A,B)$ for the number of ordered pairs that increase strictly
in both coordinates. State points have integral coordinates and marks
have both coordinates shifted by $1/2$. The supplied formula
\(\decl{GridDiagram.maslovO}\mathbb Z\decl{_eq_card}\) is
\[
 M_O(s)=I(s,s)-I(s,O)-I(O,s)+I(O,O)+1.
\]
Splitting each point set into its first $n$ and last $m+1$ columns gives
\[
\begin{array}{rcl}
 I(p,p)&=&I(x,x)+I(z,z)+nm,\\
 I(O,O)&=&I(O_1,O_1)+I(O_2,O_2)+nm,\\
 I(p,O)&=&I(x,O_1)+I(z,O_2)+nm-1,\\
 I(O,p)&=&I(O_1,x)+I(O_2,z)+nm.
\end{array}
\]
To check the first equality, put $t=x(n)$. Removing the last column
from $x$ removes $t$ increasing pairs. The cross pairs with second-column
points contribute $nm+t$. The second equality is the same computation
for the marking permutations.
For the third equality, removing the last marking of $O_1$, whose row
is $n$, removes $n+1$ pairs. All $n(m+1)$ pairs from first-column state
points to second-column markings contribute. The pairs within the second
columns agree with $I(z,O_2)$ because the low state row $t$ is at most
the lowest marking row $n+1/2$.
For the fourth equality, removing the last state point of $x$ removes
$t$ pairs, whereas the cross pairs contribute $nm+t$.
Substituting the four identities into the native formula proves
\eqref{J06:eq:seam-auxiliary-maslov-sum}.

Apply \eqref{J06:eq:seam-domain-maslov} to the three domains. Equations
\eqref{J06:eq:seam-path-index}, \eqref{J06:eq:seam-marking-domain-count}, and
\eqref{J06:eq:seam-auxiliary-maslov-sum} give
\eqref{J06:eq:seam-projected-index}.
\end{proof}

\begin{lemma}[New lemma: the fixed-first-factor domain]
\label{J06:lem:seam-second-domain-shape}
Suppose that $x=y$ and $z\ne w$. Then $D_2$ is the indicator of an
empty second-factor rectangle, $d\in\{0,1\}$, and
\begin{equation}\label{J06:eq:seam-first-domain-row}
 D_1(i,j)=d\,\mathbf1_{\{n\}}(j).
\end{equation}
Furthermore, for every first column $c<n$ and every row $r$,
\begin{equation}\label{J06:eq:seam-first-columns-copy}
 D(c,r)=D(N-1,r).
\end{equation}
\end{lemma}

\begin{proof}
Since $D_1$ has equal endpoints, write $D_1(i,j)=a_i+b_j$.
Subtract the minimum of the $a_i$ from every $a_i$ and add it to
every $b_j$. Then all $a_i$ and all $b_j$ are nonnegative, because
$D_1$ is nonnegative. Lemma~\ref{J06:lem:seam-domain-index} gives
\[
 \iota(D_1)=2\sum_i a_i+2\sum_j b_j
 \ge2(a_n+b_n)=2d.
\]
Lemma~\ref{J06:lem:seam-index-one-rectangle} gives $\iota(D_2)\ge1$.
The index relation forces equality in both inequalities. Thus $D_2$
is an empty rectangle and $d\in\{0,1\}$, and every $a_i$ with $i<n$
and every $b_j$ with $j<n$ vanishes. Write $\alpha=a_n$ and
$\beta=b_n$, so that $\alpha+\beta=d$.

For a first column $c<n$, the retained endpoints agree at that column.
Hence $\partial D(c,r)=0$ for every row $r$. Summing this equality
over columns $0,\ldots,c$ shows that
$D(c,r)-D(N-1,r)$ is independent of $r$.
At row $n$ this difference equals $\beta-(\alpha+\beta)=-\alpha$.
For $r\ge n$, matching rectangles contribute zero, so
$D(c,r)\ge0$. It follows that
\[
 D_2(m,j)=D(N-1,n+j)\ge\alpha\qquad(0\le j\le m).
\]
If $\alpha=1$, the empty rectangle $D_2$ would cover its last column
in every row. Its cyclic row interval is proper, which makes this
impossible. Therefore $\alpha=0$, $\beta=d$, and both conclusions
follow.
\end{proof}

\section{A finite order certificate on the five-point seam grid}

\begin{lemma}[New lemma: the five-point exclusion]
\label{J06:lem:seam-size-five-truncation}
Let $n=m=2$ and let $p,q$ be arbitrary retained indices.
Every native complete correction path from $p$ to $q$ has at most two
internal transitions.
\end{lemma}
\begin{proof}
A state is written as its five rows, in column order. For example
$03142$ denotes the permutation $(0,3,1,4,2)$.
All digits in this proof are single row indices, rather than an integer
written in decimal notation. An arrow annotated $(a,b)$ denotes the
actual empty rectangle with ordered column sides $a,b$, followed by
the selected inverse matching. In particular the orientation of every
rectangle in the certificate is specified.

We first give an order calculation that makes the certificate exhaustive.
Every matching source with northeast count one has the following data:
\[
\begin{gathered}
 A\in\{0,1\},\quad h\in\{3,4\},\quad L=1-A,\quad s(A)=h,\quad s(L)=t,\\
 \{P,Q,C\}=\{2,3,4\},\quad P<Q,\\
 s(P)=v<s(Q)=w,\quad s(C)=7-h,\quad
 \{v,w,t\}=\{0,1,2\}.
\end{gathered}
\]
Thus there are exactly $2\cdot2\cdot3\cdot3=36$ such sources.
The native theorem \decl{flatStep_cases} shows that a transition which
preserves northeast count is one of the four flat types below.
A birth can only exchange $(L,t)$ with $(C,7-h)$.
Its target has a decreasing first southwest pair only if either
$C<P$ and $t>v$, or $P<C<Q$ and $t<v$.
In the first case the wrapping-column orientation contains $(P,v)$.
In the second case the nonwrapping-column orientation contains $(P,v)$.
These facts leave the two birth orientations in the table.
A death from a source with northeast count one reaches a retained state,
so it is an exit rather than an internal transition.
Consequently the following six rows include every possible internal
transition from any of the 36 sources.
\[
\begin{array}{c|c|l}
 \text{type}&\text{ordered columns}&\text{necessary and sufficient conditions}\\\hline
 L_+&(L,P)&t>w,\quad A=0,\quad C>P\\
 L_-&(Q,L)&t<v,\quad A=1,\quad C<Q\\
 H_+&(C,P)&C>Q,\quad h=3,\quad v=0\\
 H_-&(Q,C)&C<P,\quad h=4,\quad w=2\\
 B_+&(L,C)&C<P,\quad t>v,\quad (A=0\ \text{or}\ h=4)\\
 B_-&(C,L)&P<C<Q,\quad t<v,\quad (A=1\ \text{or}\ h=3)
\end{array}
\tag{J06.F5.1}\label{J06:eq:size-five-types}
\]
Here is the emptiness check for every row of this table, including its
excluded cases. For $L_+$, the row data must be $(v,w,t)=(0,1,2)$.
The cyclic row interior consists of rows $3,4$.
The northeast point blocks precisely when $A>L$, and the southeast
point blocks precisely when $C<P$. This gives $A=0$ and $C>P$.
For $L_-$, the row data must be $(t,v,w)=(0,1,2)$, and the same two
high rows form the cyclic interior. Its northeast blocker occurs when
$A<L$, and its southeast blocker occurs when $C>Q$.
This gives $A=1$ and $C<Q$.
For $H_+$, the column data are $(P,Q,C)=(2,3,4)$.
The only interior columns are $0,1$. Their two points avoid the
wrapping row interval exactly when $h<7-h$ and $t>v$.
These conditions are $h=3$ and $v=0$.
For $H_-$, the column data are $(C,P,Q)=(2,3,4)$, again with interior
columns $0,1$. The necessary and sufficient inequalities are
$h>7-h$ and $t<w$, namely $h=4$ and $w=2$.
For $B_+$, the column data are $(C,P,Q)=(2,3,4)$.
If $L=1$ there is no interior column. If $L=0$, its only interior
column is $A=1$, whose high row blocks exactly when $h<7-h$.
For $B_-$, the column data are $(P,C,Q)=(2,3,4)$ and $t=0$.
Column $4$ contains the low row $w=2$, which is outside the wrapping
row interval. The only possible further interior point is at column
$A=0$ when $L=1$, and it blocks exactly when $h>7-h$.
These checks prove all six equivalences in
\eqref{J06:eq:size-five-types}. The other orientation of each flat move
is already excluded by \decl{flatStep_cases}.

An entry can reach only a matching source with northeast count one,
by \decl{seamMatchingEntry_northeast_card}.
Two of the 36 sources admit no entry. For $04312$, its possible entry
rectangles are blocked respectively by the points $(2,3)$ and $(0,0)$.
For $32014$, they are blocked respectively by $(1,2)$ and $(4,4)$.
These are the two candidates in \decl{flatEntry_data}.
Let $E_0$ be the 34 other northeast-count-one sources.
Every actual entry source belongs to $E_0$. No assertion that every
member is reached is required for the argument.

Substitution of the three choices of $(v,w,t)$, the three choices of
$(P,Q,C)$, and the two choices each of $A,h$ in
\eqref{J06:eq:size-five-types} gives the following complete list of
nonempty outgoing internal-arrow sets for sources in $E_0$.
The rows not displayed have no internal arrow.
\[
\begin{array}{c|l@{\qquad}c|l}
03142&(3,0):43012,\ (4,0):23041&03412&(4,0):23401\\
04132&(3,0):34012,\ (4,0):24031&13024&(4,2):13402\\
14302&(0,2):34012,\ (4,2):14023&23014&(4,2):23401\\
24301&(0,2):34021&30142&(3,1):34012\\
31024&(4,2):31402&31402&(1,2):34012\\
32041&(1,2):30142&32401&(1,2):34021\\
40312&(4,2):40123&41302&(1,2):43012,\ (4,2):41023\\
42013&(1,2):40123&42031&(1,2):40132\\
42301&(1,2):43021&&
\end{array}
\tag{J06.F5.2}\label{J06:eq:size-five-entry-arrows}
\]
The target in each cell is obtained by swapping the two displayed
columns and then swapping the first two southwest columns.
For example the row $24301$ uses the birth with sides $(0,2)$.
The actual target is $34201$, whose southwest rows at columns $2,3$
are $2,0$. Its selected partner is $34021$, as displayed.
Thus \eqref{J06:eq:size-five-entry-arrows} records state maps as well as
all the emptiness and source tests in \eqref{J06:eq:size-five-types}.

Only four sources with northeast count two occur among these targets:
$34012$, $34021$, $43012$, and $43021$.
Their complete internal-arrow sets are
\[
\begin{array}{c|l}
34012&\varnothing\\
43012&\varnothing\\
34021&(1,2):30412\\
43021&(0,2):03412,\ (1,2):40312.
\end{array}
\tag{J06.F5.3}\label{J06:eq:size-five-high-arrows}
\]
We verify this table without leaving any orientation unchecked.
There is no northwest or southeast point, so there is no birth and
no flat exchange with either of those quadrants.
Swapping the two northeast points leaves the first southwest pair
increasing. Swapping the first two southwest points is the selected
matching and is excluded from an internal arrow.
An interchange of the second and third southwest points also leaves
the first pair increasing. An interchange of the first and third
can yield a decreasing target pair only for southwest row order
$012$. Its nonwrapping rectangle then contains the second southwest
point, and its wrapping rectangle contains a northeast point.
Thus no flat internal exchange remains.

A death must remove column $2$, since removing either of the other
two southwest columns leaves an increasing first pair.
Its remaining pair is decreasing exactly when the southwest order
is $021$. The wrapping-column orientation contains a southwest point.
For the nonwrapping orientation $(A,2)$, $A=1$ has no interior
column and is empty. For $A=0$, column $1$ blocks exactly when its
high row is greater than the high row in column $0$.
This gives precisely the three arrows in
\eqref{J06:eq:size-five-high-arrows} and their displayed targets.

Let $E_j$ be the set of sources reached by $j$ successive internal
arrows starting from $E_0$. The two explicit arrow tables now give
\[
\begin{split}
E_1={}&\{13402,14023,23041,23401,24031,30142,31402,\\
      &\hspace{13mm}34012,34021,40123,40132,41023,43012,43021\},\\
E_2={}&\{03412,30412,34012,40312\},\\
E_3={}&\{23401,40123\},\qquad E_4=\varnothing.
\end{split}
\tag{J06.F5.4}\label{J06:eq:size-five-reachability}
\]
These are finite set substitutions in the displayed tables, rather
than an inference from tests on further sizes.
For completeness neither member of $E_3$ has an exit.
For $23401$, the two exit rectangles in \decl{flatExit_data} contain
respectively $(2,4)$ and $(0,2)$.
For $40123$, they contain respectively $(1,0)$ and $(4,3)$.
Thus a path with exactly three internal transitions cannot exit.
A path with four or more internal transitions would have to pass
through $E_4$, which is empty. This proves the lemma.
\end{proof}

\section{Uniform reduction of the second sector to five points}

The domains and indices in this section are those of
Lemmas~\ref{J06:lem:seam-domain-index}--\ref{J06:lem:seam-second-domain-shape}.
In particular the cells are indexed by their lower left corners, and
the value at a state point is the average of the four adjacent cells.

\begin{lemma}[New lemma: points forced to remain fixed]
\label{J06:lem:seam-second-fixed-points}
Let $n,m\ge1$, and let a correction path have retained endpoints
$p=\decl{seamState}(x,z)$ and $q=\decl{seamState}(x,w)$ with $z\ne w$.
Put $t=x(n)$, and let $R$ be the empty second-factor rectangle whose
domain is $D_2$. Then the following points of $p$ remain in their original
columns throughout every matching source, matching target, and retained
endpoint of the path.
\begin{enumerate}
\item Every first-column point whose row lies strictly between $0$ and $n$.
\item Every high second-column point whose second-factor point is not
one of the two source corners of $R$.
\end{enumerate}
\end{lemma}

\begin{proof}
Equation~\eqref{J06:eq:seam-first-domain-row} gives
\[
 D(c,r)=0\qquad
 \bigl(c\in\{0,\ldots,n-1,N-1\},\ 0\le r<n\bigr).
\]
A first-column state point of row strictly between $0$ and $n$ has
all four adjacent cells in this set. This includes column zero, whose
left adjacent cell column is $N-1$. Every matching rectangle avoids
these cell columns. Thus each actual entry, internal, or exit rectangle
has coefficient zero at all four cells.

Consider next a high second-column point $(c,r)$, with $c\ge n$ and
$r>n$. If $c>n$, all four adjacent cells lie in the second projection.
If $c=n$, the adjacent cell column $n-1$ is identified with column
$N-1$ by \eqref{J06:eq:seam-first-columns-copy}. In both cases the four
full-domain coefficients equal the four adjacent coefficients of $D_2$
at the corresponding second-factor point.
For an unchanged source point of the empty rectangle $R$, all four
of those coefficients are zero. Indeed the four changed rectangle
corners already contribute its entire index one, so every other source
point has zero corner average, as proved in
Lemma~\ref{J06:lem:seam-index-one-rectangle}. The coefficients are
nonnegative. Matching rectangles avoid every row of these four cells,
because the adjacent cell rows are at least $n$. Hence every actual
rectangle again has coefficient zero at all four cells.

In either case the point cannot change for the first time at an actual
rectangle. Such a change would make it a source corner, and at least
one of its adjacent cells would occur in that rectangle. A matching
inverse does not change a first-column point or a high-row point.
Induction over the finite ordered list of rectangles and matching
inverses proves that each stated point remains fixed throughout the path.
\end{proof}

\begin{lemma}[New lemma: compression of a correction path]
\label{J06:lem:seam-second-compression}
Under the hypotheses of Lemma~\ref{J06:lem:seam-second-fixed-points}, every
correction path of length $k$ induces a correction path of the same
length for seam parameters
\[
 n'\in\{1,2\},\qquad m'\in\{1,2\},
\]
with equal first retained factors at its two endpoints.
\end{lemma}

\begin{proof}
Select every first column whose original row is $0$ or $n$.
There are $a\in\{1,2\}$ such columns. There is only one when
$t=x(n)$ equals $0$ or $n$, and there are two otherwise.
Select the second column $C$ carrying row $t$, together with the
source-corner columns of $R$. Their union has $b\in\{2,3\}$ columns.
It has two when one source-corner row of $R$ is zero, and three otherwise.

Lemma~\ref{J06:lem:seam-second-fixed-points} fixes every original point
outside these $a+b\le5$ selected columns. Since each intermediate
state is a permutation, its restriction to the selected columns is a
bijection onto their original row values. Those values comprise $a+1$
low rows and $b-1$ high rows. The low rows are $0,n$, together with
$t$ if $0<t<n$. The high rows are the nonzero source rows of $R$,
translated by $n$.

Relabel the selected first columns in increasing order by
$0,\ldots,a-1$, and the selected second columns in increasing order
by $a,\ldots,a+b-1$. Relabel the low rows in increasing order by
$0,\ldots,a$, and the high rows in increasing order by
$a+1,\ldots,a+b-1$. Each intermediate state becomes a state of a
seam grid with $n'=a$ and $m'=b-1$.

Every actual rectangle has both endpoint columns among the selected
columns, since otherwise it would change a point that has been proved
fixed. Its two endpoint rows belong to the selected set as well.
An order-preserving cyclic relabeling preserves membership of selected
indices in the cyclic open interval between two selected endpoints.
This statement includes a wrapping interval: its members are exactly
those after the left endpoint or before the right endpoint in the
ordinary order, before and after relabeling.
Therefore the restricted rectangle is empty, because an interior point
after restriction would have been an interior point before restriction.
Its target is the restricted transposed state, by
\decl{GridRectangleBetween.target_eq_swapColumns}.

Every southwest column belongs to the selected set. The remaining
second-column points are high and fixed. Thus the first two southwest
columns are preserved by restriction, and restriction commutes with
matching partners. It preserves the strict inequality between their
low rows. The matching sources remain complementary, because all
northeast points lie in selected first columns. Consecutive restricted
matching sources remain distinct, since every unselected column is fixed.

The restricted retained endpoints are seam states for the new
parameters. Their first-factor permutations agree, because their values
on all selected first columns agree and hence their missing low row
also agrees. The restricted entry, internal rectangles, and exit form
the required correction path of the same length.
\end{proof}

\begin{lemma}[New lemma: the three smaller reduced grids]
\label{J06:lem:seam-thin-small-truncation}
If $(n,m)$ is $(1,1)$, $(1,2)$, or $(2,1)$, then no native internal
path between matching sources has two or more transitions.
\end{lemma}
\begin{proof}
Every complementary state has exactly one northeast point, since its
count is positive and at most $\min(n,m)=1$.
Write that point as $(A,H)$ and the southwest points as $(P,a),(Q,b)$,
where $P<Q$ and $a<b$. The existing declarations
\decl{flatStep_cases} and \decl{flatStep_apply} classify each internal
transition by an additional column $C$ and its row $c$:
\[
\begin{array}{c|c|c|c}
\text{type}&\text{range}&\text{ordered sides}&\text{new southwest data}\\\hline
L_+&C<n,\ c>b&(C,P)&(P,b),(Q,c)\\
L_-&C<n,\ c<a&(Q,C)&(P,c),(Q,a)\\
H_+&C>Q,\ c>n&(C,P)&(Q,a),(C,b)\\
H_-&n\le C<P,\ c>n&(Q,C)&(C,a),(P,b).
\end{array}
\]
The row left at $C$ in the low cases is respectively $a$ and $b$.
The row left at the removed southwest column in the high cases is $c$.
Testing $(A,H)$ against the rectangle interior gives the further
necessary inequalities
\[
 L_+:A<C,\quad L_-:C<A,\quad H_+:H<c,\quad H_-:c<H.
\]
For the two low types the wrapping row interval contains $H$, so the
column of the northeast point must be outside the wrapping or
nonwrapping column interval, respectively. For the two high types the
wrapping column interval contains $A$, so its row must be outside the
specified row interval. These are exact applications of
\decl{isEmpty_iff_forall_notMem_cIoo}.

For $(n,m)=(1,1)$ there is no additional column, so no internal
transition is possible.
For $(n,m)=(1,2)$ the additional column is the unique southeast column,
so only $H_+$ or $H_-$ can occur. After $H_+$, the only southeast
column is the old $P$, with row $c>H$. It lies before the new first
southwest column, so a second $H_+$ is impossible by its column range.
A second $H_-$ would require $c<H$, a contradiction.
After $H_-$ the only southeast column is the old $Q$, with row $c<H$.
It lies after the new second southwest column. A second $H_-$ is
excluded by its column range, and $H_+$ would require $H<c$.

For $(n,m)=(2,1)$ the additional column is the unique northwest column,
so only $L_+$ or $L_-$ can occur. After $L_+$ its row is $a$, and the
southwest rows are $b,c$ with $a<b<c$. A second $L_+$ fails its row
range. A second $L_-$ would require $C<A$, contradicting the first
transition's inequality $A<C$. After $L_-$ the additional row is $b$,
and the southwest rows are $c,a$ with $c<a<b$. A second $L_-$ fails
its row range, while $L_+$ would require $A<C$, contradicting $C<A$.
This exhausts all possibilities for two consecutive transitions.
\end{proof}

\begin{theorem}[New lemma: complete second-sector path truncation]
\label{J06:thm:seam-second-long-empty}
For every $n,m\in\mathbb N$, every first-factor state $x$, every
second-factor states $z,w$, and every $k\ge3$, the native type
\[
 \decl{SeamCorrectionPath}(k,(x,z),(x,w))
\]
has no element.
\end{theorem}

\begin{proof}
If $n=0$, there is no first column and hence no northeast point.
If $m=0$, there is no high row. In either case every state is retained,
so no correction path can have a matching source. Assume $n,m\ge1$.

If $z=w$, the full domain $D$ has equal endpoints. Its mixed difference
is zero, so its index is an even integer by
Lemma~\ref{J06:lem:seam-domain-index}. This contradicts
\eqref{J06:eq:seam-path-index}. Assume $z\ne w$, and let $R=D_2$ be the
empty second-factor rectangle from
Lemma~\ref{J06:lem:seam-second-domain-shape}.

Lemma~\ref{J06:lem:seam-second-compression} produces a correction path
of the same length with $n',m'\in\{1,2\}$.
If $n'=m'=2$, Lemma~\ref{J06:lem:seam-size-five-truncation} excludes
$k\ge3$. Each of the other three parameter pairs is excluded by
Lemma~\ref{J06:lem:seam-thin-small-truncation}. The compression includes
both $t=0,n$ and zero rectangle-endpoint rows, as specified in its proof.
This exhausts every case.
\end{proof}

\begin{corollary}[New lemma: J06 acceptance, with integer equality]
\label{J06:cor:seam-second-complete}
For every $n,m,x,z,w$,
\[
 \decl{seamPositiveCorrectionCount}((x,z),(x,w))
 =\bigl|\decl{SeamSecondRemainingRectangle}(x,z,w)\bigr|
\]
as an equality of natural numbers. In particular this equality holds
in $\mathbb F_2$, which is the J06 acceptance statement and record M605.
\end{corollary}

\begin{proof}
Theorem~\ref{J06:thm:short} identifies
the disjoint union of the one-transition and two-transition path sets
with the remaining second-factor rectangles. Theorem
\ref{J06:thm:seam-second-long-empty} excludes every longer path.
Substitute these facts into the existing declaration
\decl{seamPositiveCorrectionCount_one_two_long}.
\end{proof}

\section{Exceptional cases and explicit coordinate checks}

The case $n=0$ is proved in Theorem~\ref{J06:thm:short} and in fact proves
\eqref{J06:eq:target}, because there are no matching sources.
If $m=0$, every row is at most $n$. Again every state is retained,
there are no matching sources, and the second factor has no rectangle.
Thus \eqref{J06:eq:target} also holds.

If $z=w$, all correction paths are excluded by
\decl{seamTransferPath_self_isEmpty}, and $\cR(x,z,z)$ is empty.
More generally, if $z,w$ have the same permutation sign, the supplied
\decl{seamTransferPath_factor_sign} excludes correction paths with
fixed first factor. The supplied \decl{gridRectangle_sign} excludes
factor rectangles between them. Hence \eqref{J06:eq:target} holds in
this entire case.

The cases $x(n)=0$ and $x(n)=n$ are included explicitly in
Lemma~\ref{J06:lem:seam-second-compression}. They reduce to a seam with
one first column. A zero source-corner row of the projected rectangle
reduces to a seam with only one high row. Lemma
\ref{J06:lem:seam-thin-small-truncation} proves the required exclusions in
both situations. When $n=1$, the interior alternative $0<x(n)<n$ is
empty. When $m\le1$, two distinct nonzero second-factor endpoint rows
do not exist, so $\cR$ is empty. These conclusions agree with the
explicit short-path construction.

For a concrete boundary example, let $n=1$, $m=2$,
$x=(1,0)$, $z=(1,0,2)$, and $w=(2,0,1)$, where tuples list row
values in increasing column order. The only remaining rectangle has
factor sides $(2,0)$. Its seam path has matching sources
\[
 (2,0,1,3),\qquad(2,3,0,1)
\]
and rectangle sides
\[
 (0,1),\qquad(3,1),\qquad(3,0).
\]
These values follow by substitution into \eqref{J06:eq:zero-path}.
The emptiness proof in Lemma~\ref{J06:lem:construction} applies to this
example as a proof, without a numerical inference.

For the first possible interior example, take $n=m=2$,
$x=(0,2,1)$, $z=(1,0,2)$, and $w=(2,0,1)$.
The retained endpoints are $(0,2,3,1,4)$ and $(0,2,4,1,3)$.
The matching sources in \eqref{J06:eq:interior-table} are
\[
 (4,2,3,0,1),\qquad(4,3,0,2,1),\qquad(0,3,4,1,2),
\]
and the four ordered sides are
\[
 (4,0),\qquad(1,2),\qquad(0,2),\qquad(4,1).
\]
This checks explicitly the orientation, the use of the first two
southwest columns at the middle source, and the meaning of internal
length two.

\section{Lemma and declaration table}

All existing declarations below occur in the supplied listings, either
as definitions and proved theorems or as imported declarations explicitly
used in those proofs. No theorem in this document is a Lean declaration.

\small
\begin{longtable}{@{}>{\raggedright\arraybackslash}p{0.42\textwidth}>{\raggedright\arraybackslash}p{0.13\textwidth}>{\raggedright\arraybackslash}p{0.39\textwidth}@{}}
\toprule
Declaration or result & Status & Use\\
\midrule
\endhead
\decl{seamState_left_castSucc},
\decl{seamState_right_of_zero},
\decl{seamState_right_of_ne_zero}
& Existing & Retained coordinates \eqref{J06:eq:seam}.\\
\decl{SeamMatchingSource},
\decl{SeamInternalRectangle},
\decl{SeamInternalPath},
\decl{SeamCorrectionPath}
& Existing & Native path types and their length convention.\\
\decl{bottom_def}, \decl{top_def},
\decl{target_eq_swapColumns},
\decl{isEmpty_iff_forall_notMem_cIoo}
& Existing & Coordinate interpretation and emptiness.\\
\decl{SeamSecondRemainingRectangle},
\decl{seamSecondRemainingRectangle_wraps}
& Existing & Remaining-rectangle definition and wrapping columns.\\
\decl{seamOneStep_all_orientations},
\decl{seamSecondOneStep_kind_high}
& Existing & Exhaustion of second-sector one-transition types.\\
\decl{seamSecondOneStep_high_first_boundary},
\decl{seamSecondOneStep_high_second_boundary}
& Existing & Boundary values of $x(n)$.\\
\decl{seamSecondOneStep_high_first_row_zero},
\decl{seamSecondOneStep_high_first_row_last},
\decl{seamSecondOneStep_high_second_row_zero},
\decl{seamSecondOneStep_high_second_row_last}
& Existing & Extremal southwest rows for one transition.\\
\decl{seamSecondOneStep_high_first_pivot},
\decl{seamSecondOneStep_high_second_pivot}
& Existing & First-block pivot reconstruction.\\
\decl{seamSecondOneStepHighFirstRemaining},
\decl{seamSecondOneStepHighSecondRemaining},
\decl{seamSecondOneStepHighFirst_source_values},
\decl{seamSecondOneStepHighSecond_source_values}
& Existing & One-transition projection and its coordinates.\\
\decl{seamMatchingEntry_northeast_card},
\decl{seamMatchingExit_northeast_card},
\decl{seamMatchingEntry_southwest_card},
\decl{seamMatchingExit_southwest_card},
\decl{seamCorrectionTwo_middle_eq_two},
\decl{seamRectangle_northeast_card_eq_of_same_block}
& Existing & Birth followed by death in every two-transition path.\\
\decl{seamPositiveCorrectionCount_one_two_long},
\decl{seamLongCorrectionCount},
\decl{seamInternalPath_length_lt}
& Existing & Exact splitting of the bounded count.\\
\decl{seamTransferPath_self_isEmpty},
\decl{seamTransferPath_factor_sign},
\decl{gridRectangle_sign}
& Existing & Equal endpoints and equal-sign cases.\\
Lemma~\ref{J06:lem:remaining} & New, proved & Coordinate characterization of all remaining rectangles.\\
Lemma~\ref{J06:lem:construction} & New, proved & Explicit inverse paths and all emptiness checks.\\
Lemma~\ref{J06:lem:one} & New, proved & One-transition inverse and uniqueness.\\
Lemma~\ref{J06:lem:normal} & New, proved & Exhaustive two-transition normal form.\\
Lemma~\ref{J06:lem:two} & New, proved & Interior second-sector projection and uniqueness.\\
Theorem~\ref{J06:thm:short} & New, proved & Complete short-path bijection.\\
\decl{flatStep_cases}, \decl{flatStep_apply},
\decl{flatEntry_data}, \decl{flatExit_data}
& Existing & Exhaustive flat transitions and finite-grid entry and exit tests.\\
\decl{laurentMaslovDegree_add_card_OColumns},
\(\decl{GridDiagram.maslovO}\mathbb Z\decl{_eq_card}\)
& Existing & Native rectangle index and Maslov point-count formulas.\\
Lemma~\ref{J06:lem:seam-domain-index}
& New, proved & Additivity of the corner-average index and periodic-domain formula.\\
Lemma~\ref{J06:lem:seam-index-one-rectangle}
& New, proved & Every nonnegative domain of index one is an empty rectangle.\\
Lemma~\ref{J06:lem:seam-signed-domain}
& New, proved & The signed domain of every correction has index one.\\
Lemma~\ref{J06:lem:seam-positive-projections}
& New, proved & Positive factor projections and the exact index sum.\\
Lemma~\ref{J06:lem:seam-second-domain-shape}
& New, proved & Fixed-first-factor projection and the boundary-row annulus.\\
Lemma~\ref{J06:lem:seam-size-five-truncation}
& New, proved & Complete five-point transition certificate.\\
Lemmas~\ref{J06:lem:seam-second-fixed-points} and
\ref{J06:lem:seam-second-compression}
& New, proved & Fixed points and a path-preserving reduction to at most five points.\\
Lemma~\ref{J06:lem:seam-thin-small-truncation}
& New, proved & The three smaller reduced grids.\\
Theorem~\ref{J06:thm:seam-second-long-empty} and
Corollary~\ref{J06:cor:seam-second-complete}
& New, proved & Absence of all longer paths and the full J06 acceptance.\\
\bottomrule
\end{longtable}
\normalsize

\section{Remaining gap}

There is no remaining mathematical gap in the J06 acceptance statement
under the supplied native coordinate conventions. Theorem
\ref{J06:thm:seam-second-long-empty} excludes every individual correction
with at least three internal transitions. Theorem~\ref{J06:thm:short}
identifies every shorter positive correction with exactly one remaining
second-factor rectangle. Corollary~\ref{J06:cor:seam-second-complete}
therefore proves the stronger integer equality \eqref{J06:eq:integer}
and removes M605. This is a written proof and does not assert that the
new lemmas have been formalized in Lean.



\chapter{J07: Mixed-sector seam parity}
\label{chapter:J07}



\section{Exact statement and conventions}

All chain coefficients in this document are in $\F$.
Integral multiplicities will be used only to record unions and differences
of rectangle domains before any path cardinality is reduced modulo two.
No change of the coefficient field is involved.

\begin{theorem}[New lemma: mixed-sector exclusion at every length]
\label{J07:thm:j07-complete}
For every $n,m,k\in\mathbb N$ and every two retained indices
$p=(p_1,p_2)$ and $q=(q_1,q_2)$, if
$p_1\ne q_1$ and $p_2\ne q_2$, then the native type
\[
 \Paths_k(p,q)=\decl{SeamCorrectionPath}\ k\ p\ q
\]
is empty.
Consequently the exact acceptance statement of J07 holds:
\begin{equation}\label{J07:eq:j07-acceptance}
 \bigl(\decl{seamPositiveCorrectionCount}(p,q):\F\bigr)=0.
\end{equation}
\end{theorem}

The theorem proves catalogue record M606, including all long-path
contributions in the mixed sector. It does not assert the universal
truncation statement M607 when a factor is fixed.
The proof applies to each chosen path, so it proves that the integer
number of mixed correction paths is zero before reduction to $\F$.
It does not require an involution pairing paths or a bound on their length.

Put $N=n+m+1$ and identify the native index set
$\operatorname{Fin}(N)$ with $\{0,\ldots,N-1\}$.
A state is a permutation from columns to rows.
The native declarations
\decl{seamState_left_castSucc},
\decl{seamState_right_of_zero}, and
\decl{seamState_right_of_ne_zero} give the following retained state:
\begin{equation}\label{J07:eq:j07-retained}
 s_p(c)=
 \begin{cases}
 p_1(c),&0\le c<n,\\
 p_1(n),&c=n+d\text{ and }p_2(d)=0,\\
 n+p_2(d),&c=n+d\text{ and }p_2(d)\ne0.
 \end{cases}
\end{equation}
Thus every retained point in a column $c<n$ has row at most $n$.

For a complementary state $u$, its southwest columns are the columns
$c\ge n$ with $u(c)\le n$.
If its first two southwest columns are $P<Q$, the native matching
partner $\tau u$ exchanges their rows. This is the operation
\decl{seamComplementPartner}.
A matching source has $u(P)=a<b=u(Q)$.
Its matching rectangle has cell domain
$[P,Q)\times[a,b)$.
These are the coordinate meanings of
\decl{SeamMatchingSource},
\decl{seamComplementFirst},
\decl{seamComplementSecond}, and
\decl{seamComplementRectangle} used explicitly in the supplied seam
listings and described in Section~2.3 of the supplied I5 notes.
The two defining modules \decl{GridConnectedSumSeamSector} and
\decl{GridConnectedSumSeamMatching} are referenced but not reproduced
in the combined attachment. The proof uses the coordinate descriptions
just stated, which also appear in the supplied proofs involving those
objects.

For distinct indices $a,b$, the notation $[a,b)_{\rm cyc}$ means the
cyclic interval containing $a$ and ending immediately before $b$.
Its open version is denoted by $(a,b)_{\rm cyc}$.
An ordered rectangle from a state $u$ with column sides $(a,b)$ has
cell domain
\[
 [a,b)_{\rm cyc}\times[u(a),u(b))_{\rm cyc}
\]
and target obtained by exchanging the rows at $a,b$.
It is empty when no other state point lies in the product of the two
open intervals. These are the supplied conventions of
\decl{GridRectangleBetween.target_eq_swapColumns},
\decl{GridRectangleBetween.bottom_def},
\decl{GridRectangleBetween.top_def}, and
\decl{isEmpty_iff_forall_notMem_cIoo}.
The path definitions impose no restriction on marking squares.

By the native definitions \decl{SeamInternalRectangle},
\decl{SeamInternalPath}, and \decl{SeamCorrectionPath}, a correction
with $k$ internal transitions consists of matching sources
$u_0,\ldots,u_k$, an entry rectangle
$s_p\longrightarrow\tau u_0$, internal rectangles
$u_j\longrightarrow\tau u_{j+1}$, and an exit rectangle
$u_k\longrightarrow s_q$.
All these rectangles are empty, and successive internal sources are
unequal. The argument below uses the selected rectangles themselves.

\section{Projected integral domains of a seam correction}

All path coefficients remain in $\mathbb F_2$. Integral domains in this
section record the multiplicities of grid squares before any parity is
taken. Their use does not change the coefficient field of the chain
complex.

For a grid of size $h\ge2$, an integral domain is an integer-valued function
$D(c,r)$ on the $h^2$ squares. All cell and vertex indices in the next
two lemmas are cyclic modulo $h$. Define its mixed difference and corner
average by
\[
 \begin{split}
 (\partial D)(c,r)
 &=D(c,r)+D(c-1,r-1)-D(c-1,r)-D(c,r-1),\\
 (\mathcal A D)(c,r)
 &=\tfrac14\bigl(D(c,r)+D(c-1,r-1)+D(c-1,r)+D(c,r-1)\bigr).
 \end{split}
\]
For a state $u$, let $[u]$ be the indicator function of its points.
A domain from $u$ to $v$ satisfies $\partial D=[u]-[v]$. Its index is
\[
 \iota(D;u,v)=\sum_{c=0}^{h-1}
 \bigl((\mathcal A D)(c,u(c))+(\mathcal A D)(c,v(c))\bigr).
\]
This is a definition by finite sums. In particular it makes sense for
domains with negative coefficients.

\begin{lemma}[New lemma: domain index identities]
\label{J07:lem:seam-domain-index}
The following statements hold for integral domains on a toroidal grid.
\begin{enumerate}
\item If $D$ is a domain from $u$ to $v$ and $E$ is a domain from
$v$ to $w$, then
\[
 \iota(D+E;u,w)=\iota(D;u,v)+\iota(E;v,w).
\]
\item A domain with mixed difference zero has the form
$D(c,r)=a_c+b_r$, with integers $a_c,b_r$. Its index between any equal
endpoint states is $2\sum_c a_c+2\sum_r b_r$.
\item An empty native rectangle has index one. For any auxiliary
grid with marking permutation $O$ and any domain from $u$ to $v$,
\begin{equation}\label{J07:eq:seam-domain-maslov}
 \iota(D;u,v)=M_O(u)-M_O(v)+2\sum_c D(c,O(c)).
\end{equation}
Here $M_O$ is the native integer Maslov grading.
\end{enumerate}
\end{lemma}

\begin{proof}
Finite cyclic reindexing gives
\[
 \sum_{c,r}(\partial D)(c,r)(\mathcal A E)(c,r)
 =\sum_{c,r}(\partial E)(c,r)(\mathcal A D)(c,r).
\]
For completeness, after moving the differences from $D$ to $E$, the
coefficient of $D(c,r)$ on the left is
\[
 \tfrac14\bigl(E(c+1,r+1)-E(c+1,r-1)
             -E(c-1,r+1)+E(c-1,r-1)\bigr).
\]
Its coefficients are unchanged when both offsets are negated. This
proves the displayed symmetry. Substituting
$\partial D=[u]-[v]$ and $\partial E=[v]-[w]$ proves the first assertion
by expanding the definition of $\iota$.

If $\partial D=0$, the difference $D(c,r)-D(c-1,r)$ is independent
of $r$. Fixing one row and one column and summing these differences
gives $D(c,r)=a_c+b_r$. Each state occupies every row and every column
once. Substitution into the corner average therefore gives index
$2\sum a_c+2\sum b_r$.

The domain of a rectangle with sides $(a,b)$ and source $u$ is the
indicator of
\[
 [a,b)_{\mathrm{cyc}}\times[u(a),u(b))_{\mathrm{cyc}}.
\]
Its mixed difference is $[u]-[u\circ(a\ b)]$ by the native declaration
\decl{GridRectangleBetween.target_eq_swapColumns}.
If the rectangle is empty, its four endpoint corners each contribute
$1/4$ to the index and every other state point contributes zero.
The native declaration \decl{laurentMaslovDegree_add_card_OColumns}
then proves \eqref{J07:eq:seam-domain-maslov} for an empty rectangle.

To obtain the formula for an arbitrary domain, connect $u$ to $v$ by
adjacent-column transpositions. Such a sequence exists by moving the
desired value into the first incorrect column and repeating on the
remaining columns. Use the rectangle of width one for each transposition.
It is empty because there is no interior column. Their sum $D_0$ has
the required boundary, and the formula for $D_0$ follows by additivity
and telescoping of the native grading differences. The difference
$D-D_0$ has zero mixed difference. For a row or column annulus the
index is two and the sum of the $O$ multiplicities is one.
The second assertion, with arbitrary integer coefficients, now gives
the formula for $D-D_0$ and hence for $D$.
\end{proof}

\begin{lemma}[New lemma: a positive domain joining distinct states has index at least one]
\label{J07:lem:seam-positive-index-lower-bound}
Let $h\ge2$, let $u,v$ be distinct states of size $h$, and let $D$
be a nonnegative integral domain from $u$ to $v$.
Then $\iota(D;u,v)\ge1$.
\end{lemma}

\begin{proof}
Suppose that $u$ and $v$ differ in $e$ columns.
They cannot differ in exactly one column, because both are permutations,
so $e\ge2$. Their $2e$ changed points are distinct vertices.
At each such vertex the mixed difference of $D$ equals $1$ or $-1$.
Since the four adjacent cell multiplicities are nonnegative integers,
at least one of them is positive. The corner average at that vertex
is consequently at least $1/4$.
Each changed vertex occurs once in the defining sum for
$\iota(D;u,v)$, and every remaining summand is nonnegative.
Therefore $\iota(D;u,v)\ge 2e/4\ge1$.
\end{proof}

Now let $n,m\ge1$ and $N=n+m+1$. For a correction path from retained
state $p=\decl{seamState}(x,z)$ to retained state
$q=\decl{seamState}(y,w)$, define $D$ to be the sum of its entry,
internal, and exit rectangle domains, minus the domains of the selected
matching rectangles at all its matching sources. The matching rectangle
at a source with southwest columns $P<Q$ and rows $a<b\le n$ is
$[P,Q)\times[a,b)$. It is empty, because every column strictly between the first two southwest columns
has a high row and therefore lies outside the open row interval.
Telescoping the rectangle boundaries gives $\partial D=[p]-[q]$.
There is one more positive rectangle than subtracted matching rectangles,
so Lemma~\ref{J07:lem:seam-domain-index} gives
\begin{equation}\label{J07:eq:seam-path-index}
 \iota(D;p,q)=1.
\end{equation}

Define two restrictions, with the indicated factor indices, by
\begin{equation}\label{J07:eq:seam-domain-projections}
 \begin{aligned}
 D_1(i,j)&=D(i,j) &&(0\le i<n,\ 0\le j\le n),\\
 D_1(n,j)&=D(N-1,j) &&(0\le j\le n),\\
 D_2(i,j)&=D(n+i,n+j) &&(0\le i,j\le m).
 \end{aligned}
\end{equation}
Put $d=D(N-1,n)$.

\begin{lemma}[New lemma: positive projections and the index relation]
\label{J07:lem:seam-positive-projections}
For every correction path as above, $D_1$ is a nonnegative domain from
$x$ to $y$, and $D_2$ is a nonnegative domain from $z$ to $w$.
Their indices satisfy
\begin{equation}\label{J07:eq:seam-projected-index}
 \iota(D_1;x,y)+\iota(D_2;z,w)=1+2d.
\end{equation}
\end{lemma}

\begin{proof}
A matching rectangle uses only columns from $n$ through $N-2$ and
rows from $0$ through $n-1$. It therefore misses both restrictions in
\eqref{J07:eq:seam-domain-projections}. All remaining summands of $D$ are
nonnegative rectangle indicators, proving positivity.

For the boundary assertion, restricting a domain to cyclic lists of
cell columns $c_i$ and cell rows $r_j$ replaces its mixed difference
at $(i,j)$ by
\[
 \sum_{c\in(c_{i-1},c_i]_{\mathrm{cyc}}}
 \sum_{r\in(r_{j-1},r_j]_{\mathrm{cyc}}}(\partial D)(c,r).
\]
This formula follows by telescoping in both coordinates.
For $D_1$, the last column collects the original columns
$n,\ldots,N-1$, and row zero collects the original rows
$n+1,\ldots,N-1,0$. Under these collections the retained state $p$
becomes $[x]+m[(n,0)]$. The extra $m$ copies are the high points in
the second columns. The retained state $q$ gives
$[y]+m[(n,0)]$, so the extra terms cancel.
For $D_2$, column zero collects columns $0,\ldots,n$, and row zero
collects rows $0,\ldots,n$. The retained states become
$[z]+n[(0,0)]$ and $[w]+n[(0,0)]$. Again the extra terms cancel.
These statements use precisely the three supplied coordinate formulas
\decl{seamState_left_castSucc}, \decl{seamState_right_of_zero}, and
\decl{seamState_right_of_ne_zero}.

We prove the index relation by explicit auxiliary markings.
Let $O_1(i)=i$ on the first factor, let
$O_2(j)=j+1\pmod{m+1}$ on the second factor, and let
\[
 O(c)=c\quad(c<n),\qquad
 O(n+j)=n+O_2(j)\quad(0\le j\le m).
\]
Each is a permutation and can be made into a valid auxiliary grid by
choosing $X$ to be its row successor. The Maslov grading depends only on
the specified $O$ markings. The full marking set is the union of the
two projected marking sets with their common square $(N-1,n)$ counted
once. Consequently
\begin{equation}\label{J07:eq:seam-marking-domain-count}
 \sum_c D(c,O(c))=
 \sum_iD_1(i,O_1(i))+\sum_jD_2(j,O_2(j))-d.
\end{equation}

The native Maslov point-count formula gives
\begin{equation}\label{J07:eq:seam-auxiliary-maslov-sum}
 M_O(\decl{seamState}(x,z))=M_{O_1}(x)+M_{O_2}(z).
\end{equation}
Here is a direct verification, so no condition on the auxiliary
$X$ markings is being suppressed. For finite planar point sets $A,B$,
write $I(A,B)$ for the number of ordered pairs that increase strictly
in both coordinates. State points have integral coordinates and marks
have both coordinates shifted by $1/2$. The supplied formula
% Exact native identifier: GridDiagram.maslovOℤ_eq_card.
\decl{GridDiagram.maslovO}\ensuremath{\mathbb Z}\decl{_eq_card} is
\[
 M_O(s)=I(s,s)-I(s,O)-I(O,s)+I(O,O)+1.
\]
Splitting each point set into its first $n$ and last $m+1$ columns gives
\[
\begin{array}{rcl}
 I(p,p)&=&I(x,x)+I(z,z)+nm,\\
 I(O,O)&=&I(O_1,O_1)+I(O_2,O_2)+nm,\\
 I(p,O)&=&I(x,O_1)+I(z,O_2)+nm-1,\\
 I(O,p)&=&I(O_1,x)+I(O_2,z)+nm.
\end{array}
\]
To check the first equality, put $t=x(n)$. Removing the last column
from $x$ removes $t$ increasing pairs. The cross pairs with second-column
points contribute $nm+t$. The second equality is the same computation
for the marking permutations.
For the third equality, removing the last marking of $O_1$, whose row
is $n$, removes $n+1$ pairs. All $n(m+1)$ pairs from first-column state
points to second-column markings contribute. The pairs within the second
columns agree with $I(z,O_2)$ because the low state row $t$ is at most
the lowest marking row $n+1/2$.
For the fourth equality, removing the last state point of $x$ removes
$t$ pairs, whereas the cross pairs contribute $nm+t$.
Substituting the four identities into the native formula proves
\eqref{J07:eq:seam-auxiliary-maslov-sum}.

Apply \eqref{J07:eq:seam-domain-maslov} to the three domains. Equations
\eqref{J07:eq:seam-path-index}, \eqref{J07:eq:seam-marking-domain-count}, and
\eqref{J07:eq:seam-auxiliary-maslov-sum} give
\eqref{J07:eq:seam-projected-index}.
\end{proof}


\section{Removing two annuli}

\begin{lemma}[New lemma: the unoccupied northeast region]
\label{J07:lem:j07-northeast-domain}
Let $n,m\ge1$, let $D$ be the signed domain of a native correction
between retained states, and let $D_1,D_2,d$ be defined by
\eqref{J07:eq:seam-domain-projections}.
For every $0\le c<n$ and every $0\le j\le m$,
\begin{equation}\label{J07:eq:j07-northeast-formula}
 D(c,n+j)=D_1(c,n)+D_2(m,j)-d\ge0.
\end{equation}
If
\[
 a=\min_{0\le i\le n}D_1(i,n),\qquad
 b=\min_{0\le j\le m}D_2(m,j),
\]
then $a,b$ are nonnegative integers and $a+b\ge d$.
\end{lemma}

\begin{proof}
At every vertex $(c,r)$ with $c<n$ and $r>n$, both retained endpoint
indicator functions vanish by \eqref{J07:eq:j07-retained}.
Hence $\partial D(c,r)=0$ there.
Fix $c<n$. Summing over columns $0,\ldots,c$ gives, for $n<r<N$,
\[
 D(c,r)-D(N-1,r)=D(c,r-1)-D(N-1,r-1).
\]
Telescoping from row $n$ to row $n+j$ proves
\eqref{J07:eq:j07-northeast-formula}, because
$D(c,n)=D_1(c,n)$ and $D(N-1,n+j)=D_2(m,j)$.
All the subtracted matching rectangles have rows less than $n$ and
columns at least $n$. Therefore $D(c,n+j)$ is a sum of nonnegative
rectangle indicators. This proves the inequality in that equation.

The two minima exist because their finite index sets are nonempty.
They are nonnegative integers by
Lemma~\ref{J07:lem:seam-positive-projections}.
If a minimum for $a$ occurs at $i=n$, then
$a=D_1(n,n)=d$, so $a+b\ge d$.
Otherwise choose $c<n$ attaining $a$, and choose $j$ attaining $b$.
Equation~\eqref{J07:eq:j07-northeast-formula} gives $a+b-d\ge0$.
These alternatives also include the case in which the first minimum
is attained both at $n$ and at another column.
\end{proof}

\begin{lemma}[New lemma: normalization of the factor domains]
\label{J07:lem:j07-annulus-normalization}
Under the hypotheses and notation of
Lemma~\ref{J07:lem:j07-northeast-domain}, define
\[
 E_1(i,j)=D_1(i,j)-a\,\mathbf1_{\{n\}}(j),\qquad
 E_2(i,j)=D_2(i,j)-b\,\mathbf1_{\{m\}}(i).
\]
Then $E_1$ and $E_2$ are nonnegative integral domains with the same
respective endpoints as $D_1$ and $D_2$, and
\begin{equation}\label{J07:eq:j07-index-budget}
 \iota(E_1;p_1,q_1)+\iota(E_2;p_2,q_2)
 =1+2d-2a-2b\le1.
\end{equation}
\end{lemma}

\begin{proof}
Only the last row is altered in $D_1$, and $a$ is the minimum
multiplicity on that row. Only the last column is altered in $D_2$,
and $b$ is the minimum on that column. Thus every cell of $E_1,E_2$
is nonnegative.
A complete cell row or cell column has mixed difference zero,
so the endpoints are unchanged.
Each such annulus has index two by
Lemma~\ref{J07:lem:seam-domain-index}.
Additivity of that lemma and
\eqref{J07:eq:seam-projected-index} give the equality in
\eqref{J07:eq:j07-index-budget}. The inequality follows from $a+b\ge d$,
proved in Lemma~\ref{J07:lem:j07-northeast-domain}.
\end{proof}

\begin{proof}[Proof of Theorem~\ref{J07:thm:j07-complete}]
If $n=0$, the first-factor grid state is the unique permutation of a
singleton. The hypothesis $p_1\ne q_1$ is impossible.
If $m=0$, the same argument applies to $p_2\ne q_2$.
It remains to consider $n,m\ge1$.

Suppose that a chosen correction path of some length $k\ge0$ exists.
Form its signed domain and the nonnegative normalized domains
$E_1,E_2$ in Lemma~\ref{J07:lem:j07-annulus-normalization}.
Their endpoint pairs are distinct in both factors by hypothesis.
Lemma~\ref{J07:lem:seam-positive-index-lower-bound} therefore gives
\[
 \iota(E_1;p_1,q_1)\ge1,\qquad
 \iota(E_2;p_2,q_2)\ge1.
\]
Their sum is at least two, contradicting
\eqref{J07:eq:j07-index-budget}.
This proves that every $\Paths_k(p,q)$ is empty.

The native definition \decl{seamPositiveCorrectionCount} is the finite
sum of the cardinalities of $\Paths_{j+1}(p,q)$ for
$0\le j<\#\decl{SeamMatchingSource}(n,m)-1$, with natural-number
subtraction in the upper bound.
Every summand is zero by the first assertion.
The sum is therefore the integer zero, and its image in $\F$ is zero.
This proves \eqref{J07:eq:j07-acceptance}, including M606 exactly as stated.
\end{proof}

\section{Small sizes, orientations, and dependency ledger}

The cases $n=0$ and $m=0$ were treated in the proof and require no
auxiliary grid of size one.
For $n=m=1$, the first projection uses columns $0,2$ and rows $0,1$
of the size-three seam grid. The second uses columns $1,2$ and rows
$1,2$. Their common square is $(2,1)$.
Every matching rectangle is contained in the remaining square $(1,0)$,
so the positivity argument and the shared-square subtraction have
exactly the stated signs in this smallest nonvacuous case.
The same projection formulas apply when just one of $n,m$ equals one.

On a size-two factor, the two cyclic orientations of a rectangle can
both occur. The signed-domain argument treats each chosen rectangle
separately and uses its half-open cell interval. The four endpoint
corner contributions still total one.
A row or column annulus still has index two.
No orientation or multiplicity is discarded in the passage from a
path to its signed domain.
The argument also treats $k=0$ and arbitrary larger $k$ without a
separate classification of short paths.

All exact declaration names in the table refer to the supplied Lean
listings. The main seam declarations lie in
\decl{OnLeConjecture.KnotFloer.GridFull}.
The rectangle and point-count declarations are cited with the
namespaces displayed in those listings.

\small
\begin{longtable}{p{0.49\textwidth}p{0.43\textwidth}}
\toprule
Lemma or declaration & Status and use \\
\midrule
\endhead
\decl{seamState_left_castSucc},
\decl{seamState_right_of_zero},
\decl{seamState_right_of_ne_zero}
& Existing Lean declarations. They give the retained coordinates and
both factor boundary projections. \\
\decl{SeamInternalRectangle}, \decl{SeamInternalPath},
\decl{SeamCorrectionPath}, \decl{seamPositiveCorrectionCount}
& Existing Lean declarations. They specify the exact path choices,
the count of internal transitions, and the finite acceptance sum. \\
\decl{seamComplementRectangle_isEmpty},
\decl{seamComplementRectangle_left},
\decl{seamComplementRectangle_right}
& Existing Lean declarations used in the supplied matching proofs.
They identify the selected empty southwest rectangle and its column
interval. The increasing source condition supplies its row interval. \\
\decl{GridRectangleBetween.target_eq_swapColumns},
\decl{GridRectangleBetween.bottom_def},
\decl{GridRectangleBetween.top_def},
\decl{isEmpty_iff_forall_notMem_cIoo}
& Existing Lean declarations. They fix the corner orientation,
rectangle boundary, and native emptiness convention. \\
\decl{laurentMaslovDegree_add_card_OColumns}
& Existing Lean declaration. It gives the Maslov identity for an empty
rectangle in each explicitly constructed auxiliary grid. \\
\decl{GridDiagram.maslovO}\ensuremath{\mathbb Z}\decl{_eq_card}
& Existing Lean declaration. Its finite point-count formula proves the Maslov sum for the
chosen auxiliary markings. \\
Lemma~\ref{J07:lem:seam-domain-index}
& New lemma, proved here. It proves index additivity, the structure of
zero-boundary domains, and the arbitrary-domain Maslov formula. \\
Lemma~\ref{J07:lem:seam-positive-index-lower-bound}
& New lemma, proved here. It bounds the index of a nonnegative domain
between distinct states using its changed corners. \\
Lemma~\ref{J07:lem:seam-positive-projections}
& New lemma, proved here. It proves both positive factor projections
and their exact index relation. \\
Lemma~\ref{J07:lem:j07-northeast-domain}
& New lemma, proved here. It derives the northeast interpolation
formula and the inequality for the two annulus multiplicities. \\
Lemma~\ref{J07:lem:j07-annulus-normalization}
& New lemma, proved here. It produces two nonnegative domains with
total index at most one. \\
Theorem~\ref{J07:thm:j07-complete}
& New lemma, proved here. It excludes every mixed correction path
and proves the exact J07 acceptance statement. \\
\bottomrule
\end{longtable}
\normalsize

\section{Remaining gap}

There is no remaining gap in the written proof of J07 or M606 under
the native coordinate definitions recorded above.
The statements in this document are written proofs and are not claims
that their new lemmas have already been formalized in Lean.
The proof uses no holomorphic Heegaard Floer result and no computation
as a proof step.
It makes no assertion here about the fixed-factor long-path counts in
J05 or J06, or about universal truncation M607.




\section{The combined seam conclusion: M601 and M607}

\begin{theorem}[New lemma: universal truncation at internal length two]
\label{seam:universal-truncation}
For every $n,m\in\mathbb N$, every retained indices $p,q$, and every
$k\geq3$, the native type $\decl{SeamCorrectionPath}\ k\ p\ q$ is empty.
\end{theorem}
\begin{proof}
If $p_2=q_2$, apply Theorem~\ref{J05:thm:seam-first-long-empty}.
If $p_2\ne q_2$ and $p_1=q_1$, apply
Theorem~\ref{J06:thm:seam-second-long-empty}.
If both factors differ, apply Theorem~\ref{J07:thm:j07-complete}.
These cases exhaust all retained endpoint pairs, including equal pairs.
Each cited theorem includes the parameters zero and the smallest
positive sizes. This is precisely M607.
\end{proof}

\begin{theorem}[New lemma: the exact seam positive-path comparisons]
\label{seam:positive-comparisons}
For every $n,m\in\mathbb N$, the native proposition
\decl{SeamPositivePathComparisons}\ $n\ m$ holds.
In fact its three cardinality comparisons hold before reduction to
$\mathbb F_2$:
\begin{align*}
 \decl{seamPositiveCorrectionCount}((x,z),(y,z))
 &=\#\decl{SeamFirstRemainingRectangle}(x,y,z),\\
 \decl{seamPositiveCorrectionCount}((x,z),(x,w))
 &=\#\decl{SeamSecondRemainingRectangle}(x,z,w),\\
 \decl{seamPositiveCorrectionCount}(p,q)&=0
 \quad\text{if }p_1\ne q_1\text{ and }p_2\ne q_2.
\end{align*}
All variables in each displayed line range over the corresponding
native factor state sets.
\end{theorem}
\begin{proof}
The first equality is Corollary~\ref{J05:cor:seam-first-complete},
using the existing integer identity
\decl{seamPositiveCorrectionCount_first_eq_remaining_add_long}.
The second equality is Corollary~\ref{J06:cor:seam-second-complete},
which uses the explicitly proved short-path bijection and the absence
of every longer path. The third equality follows because every
summand is the cardinality of an empty type by
Theorem~\ref{J07:thm:j07-complete}. Reducing these equalities modulo
two gives exactly the three conjuncts of the supplied definition of
\decl{SeamPositivePathComparisons}. This is M601.
\end{proof}

The universal theorem also implies the arbitrary-endpoint M603
statement, since it excludes every three-transition path outright.
It implies the sufficient emptiness form of M604 for all lengths at
least four. Thus all records M601--M607 named in the seam jobs are
settled by the displayed written proofs.

\begin{longtable}{@{}>{\raggedright\arraybackslash}p{0.58\linewidth}>{\raggedright\arraybackslash}p{0.34\linewidth}@{}}
\textbf{Lemma or declaration}&\textbf{Status and role}\\\hline\endhead
Theorems~\ref{J05:thm:seam-first-long-empty},
\ref{J06:thm:seam-second-long-empty}, and
\ref{J07:thm:j07-complete}&New, proved in the preceding chapters.
They cover every retained endpoint pair.\\
\decl{seamPositiveCorrectionCount_first_eq_remaining_add_long}
&Existing Lean declaration. It identifies the first short count.\\
\decl{SeamPositivePathComparisons}
&Existing Lean definition. It is the exact three-sector target.\\
Theorems~\ref{seam:universal-truncation} and
\ref{seam:positive-comparisons}&New, proved here. They assemble
M607 and M601 without additional assumptions.\\\hline
\end{longtable}

There is no remaining mathematical gap in these combined seam
statements. They are not assertions that the new written lemmas have
already been entered into the Lean development.


\chapter{J21: An explicit second torus representative}
\label{chapter:J21}



\section*{Chapter overview}
For every $n\geq2$, we construct a finite Laurent correction of the native
first torus cycle over $\mathbb F_2$. We compute its complete boundary by
empty-rectangle cancellation. The resulting cycle represents exactly
\decl{reflectedTorusAnchorClass} and lies in the bifiltration rectangle
\[
             \left(3,\frac{(n-2)(n-1)}2\right).
\]
This proves the strict-range acceptance statement of J21 and the explicit
correction and displacement assertions M902 and M903. Every new combinatorial
identity used in the construction is proved below. The argument uses the
native full differential, including rectangles containing $X$ markings.
It uses no identification with a holomorphic complex.




\section{The exact statement and the native conventions}
\label{J21:sec:torus-target}

Fix a natural number $n\geq2$, and put
\[
             p=n+1,\qquad N=2p+1=2n+3,\qquad
             g=\frac{p(p-1)}2=\frac{n(n+1)}2.
\]
All chains in this document have coefficients in $\F[U,U^{-1}]$.
Columns and rows are elements of $\Z/N\Z$, with positive cyclic order
induced by $0,1,\ldots,N-1$. Whenever ordinary inequalities or absolute
values are used, the corresponding residues are represented in this range.
For columns $a,b$, put
\[
 d(a,b)=(b-a)\bmod N\in\{0,1,\ldots,N-1\}.
\]
For distinct endpoints, $[a,b)$ and $(a,b)$ denote the positive half-open
and open cyclic intervals of columns. An indicator $[P]$ is one when
the proposition $P$ holds and is zero otherwise. Indicators in coefficient
identities belong to $\F$. A bar over an integer denotes its reduction
modulo two. Counts, gradings, and displacement weights are ordinary
integers or rational numbers, before any explicitly indicated reduction.

The native grid $G=\decl{reflectedTorusGrid}\,n$ has markings
\begin{equation}
              O(c)=-c-1,\qquad X(c)=p-c-1\pmod N.
\label{J21:eq:torus-native-markings}
\end{equation}
The marking with index $c$ occupies the square with center
$(c+\tfrac12,O(c)+\tfrac12)$ or $(c+\tfrac12,X(c)+\tfrac12)$.
For a permutation $\sigma$ of the columns, define the state
\[
                         x_\sigma(c)=-\sigma(c)\pmod N.
\]
Every state has exactly one such relative permutation, since negation is
a bijection of the rows. The anchor is $x_0=x_{\mathrm{id}}$.
These identifications are the existing declarations
\decl{reflectedTorusGrid_O_apply}, \decl{reflectedTorusGrid_X_apply},
\decl{reflectedTorusAnchor_val}, and \decl{reflectedTorusRelativeColumns}.
A cycle word $(a_0\,a_1\,\cdots\,a_{q-1})$ denotes the permutation
sending $a_i$ to $a_{i+1}$ with indices modulo $q$ and fixing every
unlisted column. Products of disjoint cycles have their usual meaning.
A word such as $01342$ denotes a cycle on ordered labels, not a one-line
permutation.

Write $M(x)$ for the native $O$-Maslov grading and $A(x)$ for the native
Alexander grading. The monomial $U^k x$ has Maslov degree $M(x)-2k$
and bifiltration coordinates $(-k,A(x)-k)$.
The grading convention is \decl{laurentMaslovDegree}. The support criterion
for these two filtration coordinates is
\decl{mem_laurentMaslovBifiltrationLevel_iff}.
The native full Laurent differential, denoted by $D$, is
\begin{equation}
 D x=\sum_y\sum_{R\in\operatorname{Rect}^{\circ}(x,y)}U^{o(R)}y
     =\sum_{j\geq0}U^j d_jx.
\label{J21:eq:torus-native-differential}
\end{equation}
Here $\operatorname{Rect}^{\circ}(x,y)$ is the set of directed rectangles
whose open interior contains no source-state point, and $o(R)$ is the
number of covered $O$ squares. Rectangles containing $X$ markings are
included. The operators $d_j$ act on the finite $\F$-span of the states.
Equation~\eqref{J21:eq:torus-native-differential} is the existing
\decl{laurentCoefficient_eq}. Its Laurent linear extension is
\decl{laurentDifferential_monomialGenerator}.
For every state-empty rectangle from $x$ to $y$, the existing
\decl{laurentMaslovDegree_add_card_OColumns} gives
\begin{equation}
                         M(y)=M(x)-1+2o(R).
\label{J21:eq:torus-rectangle-degree}
\end{equation}

We will use the following exact native displacement identity.
For $0\leq d<N$, define
\[
 W(d)=pd-N\max\{0,d-p-1\}.
\]
This is \decl{reflectedTorusDisplacementWeight}. The declarations
\decl{alexanderTwo_reflectedTorus_eq_sum_displacement},
\decl{alexander_reflectedTorusAnchor}, and the native conversion
\decl{GridDiagram.two_mul_alexander_eq_intCast} give
\begin{equation}
 A(x_\sigma)=g-\frac1N\sum_{c\in\Z/N\Z}W(d(c,\sigma(c))).
\label{J21:eq:torus-native-alexander}
\end{equation}
Indeed, the first declaration gives the equality after multiplying both
sides by $2N$. The conversion replaces twice the rational Alexander
grading with its integer version, and $N>0$ permits cancellation.
The imported conversion is used explicitly in the supplied proof of
\decl{alexander_reflectedTorus_le_of_displacement_weight}.
In particular $W(0)=0$, and for $1\leq d\leq p$ one has
\begin{equation}
                       W(d)=pd,\qquad W(N-d)=(p+1)d.
\label{J21:eq:torus-weight-two-slopes}
\end{equation}
Substitution in the defining formula for $W$ proves these two equalities.

Let $z_1=\decl{firstTorusCycle}\,n\,\F$ and let
$\xi=\decl{reflectedTorusAnchorClass}\,n\,\F$.
For a chain $z$, the notation $[U^k x]z$ denotes its coefficient at that
Laurent monomial. For a chain with Laurent exponent zero, $[x]z$ denotes
the coefficient of the state $x$.

\begin{theorem}[New theorem: the exact J21 acceptance statement]
\label{J21:thm:j21-main}
For every natural number $n\geq2$, there is an explicitly specified degree-one
Laurent chain $w_2$ such that $z_2=z_1+Dw_2$ is a degree-zero cycle,
$[z_2]=\xi$, and every monomial $U^k y$ with $[U^k y]z_2=1$ satisfies
\begin{equation}
                 -k\leq3,\qquad A(y)-k\leq g-2p+3.
\label{J21:eq:j21-strong-corner}
\end{equation}
For every integer $m>0$ with $2m<p$, the same monomials satisfy
\begin{equation}
                  A(y)-mk\leq g+3m-2p
\label{J21:eq:j21-weighted-target}
\end{equation}
and the exact displacement inequality of M903,
\begin{equation}
 N\bigl(2p-m(3+k)\bigr)
       \leq\sum_{c\in\Z/N\Z}W(d(c,\sigma_y(c))),
 \qquad y=x_{\sigma_y}.
\label{J21:eq:j21-M903-target}
\end{equation}
Consequently $\xi$ is represented below
$\decl{Staircase.consecutiveWeight}(n,2/m,2)$ in the weighted filtration
at $t=2/m$.
\end{theorem}

The formulas for the correction chains and the complete coefficient
calculation occupy the intervening sections. The proof of
Theorem~\ref{J21:thm:j21-main} is completed in
Section~\ref{J21:sec:torus-final-bridge}. Equation~\eqref{J21:eq:j21-strong-corner}
is the parameter-free strengthening recorded in M901. The target of J21
keeps precisely its stated hypotheses.

\paragraph{The smallest native convention check.}
For $n=2$, one has $p=3$, $N=7$, and the three native coordinate lists are
\[
\begin{array}{c|rrrrrrr}
 c&0&1&2&3&4&5&6\\\hline
 O(c)&6&5&4&3&2&1&0\\
 X(c)&2&1&0&6&5&4&3\\
 x_0(c)&0&6&5&4&3&2&1
\end{array}
\]
These lists follow by substitution in~\eqref{J21:eq:torus-native-markings}.
The existing \decl{maslovO_reflectedTorusAnchor} gives $M(x_0)=0$,
and \decl{alexander_reflectedTorusAnchor} gives $A(x_0)=3$.
The proofs below keep the half-open square intervals and open state
interiors that these coordinates require.



\section{A relative-permutation formula for the Maslov grading}
\label{J21:sec:relative-permutation-maslov}

The following calculation uses the native state coordinates and the native
half-integer marking coordinates. It does not use a Floer-theoretic
identification. Let $N\geq 1$, let $O(c)=N-1-c$, and let $\sigma$ be a
permutation of $\{0,\ldots,N-1\}$. Define
\[
 x_\sigma(c)=(-\sigma(c))\bmod N,
 \qquad
 \operatorname{fr}(\sigma)=\sum_{c=0}^{N-1}|\sigma(c)-c|.
\]
Let $\operatorname{inv}(\sigma)$ denote the number of pairs $c<d$ for which
$\sigma(c)>\sigma(d)$.

For finite sets of points $P,Q$ in the coordinate square, define
$I(P,Q)=\#\{(s,t)\in P\times Q:s_1<t_1\text{ and }s_2<t_2\}$.
Thus the inequalities in this pair count are strict.

\begin{lemma}[New lemma: relative-permutation Maslov formula]
\label{J21:lem:torus-permutation-maslov}
For every such $N$ and $\sigma$, the native $O$-Maslov grading satisfies
\[
 M_O(x_\sigma)=\operatorname{inv}(\sigma)-\operatorname{fr}(\sigma).
\]
\end{lemma}
\begin{proof}
Put $a=\sigma^{-1}(0)$. For $c\neq a$ the row of the state point is
$N-\sigma(c)$, whereas the row at column $a$ is zero. Therefore
\[
 I(x_\sigma,x_\sigma)
 =\operatorname{inv}(\sigma)+N-1-2a.
\]
Indeed, deleting the value zero removes exactly $a$ inversions, and the
state point in row zero contributes exactly $N-1-a$ increasing pairs.
The marking points have coordinates
$(d+\tfrac12,N-d-\tfrac12)$. For a state point in column $c\neq a$,
its contribution to $I(x_\sigma,O)$ is
$\max\{\sigma(c)-c,0\}$. The state point in column $a$ contributes
$N-a$. Thus
\[
 I(x_\sigma,O)=
 \sum_c\max\{\sigma(c)-c,0\}+N-a.
\]
Likewise the contribution of the state point in column $c\neq a$ to
$I(O,x_\sigma)$ is $\max\{c-\sigma(c),0\}$, and the contribution of
column $a$ is zero. Hence
\[
 I(O,x_\sigma)=\sum_c\max\{c-\sigma(c),0\}-a.
\]
The markings form a decreasing sequence, so $I(O,O)=0$. Substitution in the existing finite-count formula
\(\decl{GridDiagram.maslovO}\mathbb Z\decl{_eq_card}\)

\[
 M_O(x)=I(x,x)-I(x,O)-I(O,x)+I(O,O)+1
\]
proves the identity. Every equality above includes the case $N=1$.
\end{proof}
\begin{lemma}[New lemma: deletion of fixed columns]
\label{J21:lem:torus-fixed-column-deletion}
Suppose $t$ is a fixed point of $\sigma$. Delete $t$ from both the domain
and the range and reindex the remaining ordered sets increasingly. The
resulting permutation $\bar\sigma$ satisfies
\[
 \operatorname{inv}(\bar\sigma)-\operatorname{fr}(\bar\sigma)
 =\operatorname{inv}(\sigma)-\operatorname{fr}(\sigma).
\]
\end{lemma}
\begin{proof}
Let $r$ be the number of columns $c<t$ with $\sigma(c)>t$. The number
of columns $c>t$ with $\sigma(c)<t$ also equals $r$, since $\sigma$
is a bijection fixing $t$. Deleting $t$ removes exactly $2r$
inversions. It also reduces the absolute displacement of each of these
$2r$ crossing arrows by one and leaves every other absolute displacement
unchanged. Thus the inversion number and the footrule sum decrease by the same
quantity.
\end{proof}
\begin{lemma}[New lemma: circular rotations]
\label{J21:lem:torus-rotation-maslov}
Let $S=\{s_0<\cdots<s_{q-1}\}$ be a subset of columns, with $q\geq2$.
For $1\leq j<q$, let $\sigma$ fix the complement of $S$ and satisfy
$\sigma(s_i)=s_{(i+j)\bmod q}$. Then
\[
 M_O(x_\sigma)=-j(q-j).
\]
\end{lemma}
\begin{proof}
Delete every fixed column using the preceding lemma. The resulting
permutation is the cyclic shift by $j$ on $q$ elements. Its inversion
number is $j(q-j)$, since its two increasing blocks have respective
lengths $q-j$ and $j$. Its absolute displacement is
$(q-j)j+j(q-j)=2j(q-j)$. The first lemma gives the result.
\end{proof}

\begin{corollary}[New lemma: small permutation types]
\label{J21:cor:torus-small-types}
Every transposition state has Maslov value $-1$. Every three-cycle
state has value $-2$, and every four-cycle state has value $-3$.
On four columns $a<b<c<d$ in circular order, the two-transposition
states $(a\,b)(c\,d)$ and $(a\,d)(b\,c)$ have value $-2$, while
$(a\,c)(b\,d)$ has value $-4$.
\end{corollary}
\begin{proof}
Delete fixed columns by Lemma~\ref{J21:lem:torus-fixed-column-deletion}.
The transposition and the two three-cycles are the rotations covered by
Lemma~\ref{J21:lem:torus-rotation-maslov}. For the remaining permutations,
the following table evaluates the two counts in
Lemma~\ref{J21:lem:torus-permutation-maslov} on the four ordered labels.
\[
\begin{array}{c|r|r}
 \text{cycle word or pairing}&\operatorname{fr}&\operatorname{inv}\\\hline
 (0123),(0132),(0231),(0321)&6&3\\
 (0213),(0312)&8&5\\
 (01)(23)&4&2\\
 (03)(12)&8&6\\
 (02)(13)&8&4
\end{array}
\]
These six words are every four-cycle, and the last three entries are
every pairing. Subtracting the second column from the third proves all
the assertions.
\end{proof}


\section{The native rectangle coordinates}
\begin{lemma}[New lemma: rectangle coordinates]
\label{J21:lem:torus-relative-rectangle}
For a rectangle with ordered column sides $a,b$ and source $x_\sigma$, its
number of $O$ markings is
\[
 o(a,b)=\bigl|[a,b)\cap[\sigma(b),\sigma(a))\bigr|.
\]
It is empty if and only if there is no $c\in(a,b)$ such that
$\sigma(c)\in(\sigma(b),\sigma(a))$.
\end{lemma}
\begin{proof}
The native vertical interval runs from $-\sigma(a)$ to $-\sigma(b)$.
Negating reverses its orientation. An interior state point in column $c$
therefore lies in the rectangle precisely when
$c\in(a,b)$ and $\sigma(c)\in(\sigma(b),\sigma(a))$.
An $O$ marking in column $c$ has row $-c-1$.
Its membership in the half-open vertical interval is equivalent to
$c\in[\sigma(b),\sigma(a))$. Intersecting this condition with the half-open
column interval gives the asserted count. The convention that the lower
and left edges are included is the convention used by
\decl{OColumns} and \decl{laurentCoefficient_eq}.
\end{proof}



\section{Cyclic relabeling and the exact insertion index}
\label{J21:sec:torus-insertion-geometry}

For the native grid with $N=2p+1$ and $p\geq3$, retain the relative
states $x_\sigma(c)=-\sigma(c)$ and put
\[
 \ell(\sigma)=-M_O(x_\sigma)
       =\sum_{c=0}^{N-1}|\sigma(c)-c|-\operatorname{inv}(\sigma).
 \tag{IG.1}
\]
The equality is the proved new Maslov formula in
Section~\ref{J21:sec:relative-permutation-maslov}.
The following lemmas concern the native toroidal rectangle conventions.
They do not use any coefficient assigned to a correction chain.

\begin{lemma}[New lemma: cyclic conjugation]
\label{J21:lem:torus-cyclic-conjugation}
Let $\rho(c)=c+1\bmod N$. For every permutation $\sigma$,
\[
                     \ell(\rho\sigma\rho^{-1})=\ell(\sigma).
\]
\end{lemma}
\begin{proof}
If $\sigma(N-1)=N-1$, the last fixed point moves to the first fixed
point. Deleting these fixed points gives the same increasingly
reindexed permutation. The fixed-column deletion lemma in
Section~\ref{J21:sec:relative-permutation-maslov} proves the assertion.

Suppose $\sigma(N-1)\neq N-1$. Put $v=\sigma(N-1)$ and
$b=\sigma^{-1}(N-1)$, so $v,b<N-1$.
Moving the last position to the first changes the inversion number
by $2v-(N-1)$. Then moving the largest value to the smallest, in its
new position $b+1$, changes that number by
$2(b+1)-(N-1)$. Their total is $2(v+b-N+2)$.
All absolute displacements except the two incident to the last
column or last value are unchanged. The changes of these two
absolute displacements are $2v-N+2$ and $2b-N+2$.
Their sum is again $2(v+b-N+2)$. Subtracting the two changes in
(IG.1) proves the equality.
\end{proof}

\begin{lemma}[New lemma: cycle-count lower bound]
\label{J21:lem:torus-cycle-lower-bound}
For every permutation $\sigma$, one has
\[
                \ell(\sigma)\geq N-\#\{\text{cycles of }\sigma\},
\]
where fixed points are included among the cycles.
In particular a single $k$-cycle has $\ell\geq k-1$.
\end{lemma}
\begin{proof}
If the permutation is not the identity, choose its least nonfixed
column $a$. Put $c=\sigma(a)$ and $b=\sigma^{-1}(a)$.
All columns smaller than $a$ are fixed, so $b,c>a$.
Swapping the values in columns $a,b$ fixes $a$ and deletes it from
its old permutation cycle. Thus the number of cycles increases by one.
The footrule sum decreases by $2\min\{b-a,c-a\}$.
The inversion number decreases by $2t+1$, where
\[
                  t=\#\{i:a<i<b\text{ and }a<\sigma(i)<c\}.
\]
Both $b-a-1$ eligible columns and $c-a-1$ eligible values bound $t$.
Hence $t\leq\min\{b-a,c-a\}-1$, and $\ell$ decreases by at least
one. Iterating this operation reaches the identity, whose $\ell$
is zero, after exactly $N-\#\{\text{cycles of }\sigma\}$ operations.
The assertion follows by adding these inequalities.
\end{proof}

\begin{lemma}[New lemma: exact insertion index]
\label{J21:lem:torus-exact-insertion-index}
Let $a,b,q$ be pairwise distinct columns. Suppose
$\sigma(q)=q$ and $\sigma(a)=b$, and put
$\tau=\sigma\circ(a\ q)$. Thus $\tau$ replaces the arrow
$a\mapsto b$ by $a\mapsto q\mapsto b$.
Exactly one of the two candidate rectangles from $x_\sigma$ to
$x_\tau$ avoids all $O$ markings. If its number of interior source
points is $r$, then
\[
                         \ell(\tau)-\ell(\sigma)=1+2r.
 \tag{IG.2}
\]
In particular this rectangle is state-empty if and only if the
left side equals one.
\end{lemma}
\begin{proof}
Use Lemma~\ref{J21:lem:torus-cyclic-conjugation} to make $a=0$.
There are two possible orders of $b,q$.
If $0<q<b$, swapping the two values changes the footrule sum by
zero and decreases the inversion number by
\[
             1+2\#\{i:0<i<q\text{ and }q<\sigma(i)<b\}.
\]
The rectangle from column zero to column $q$ has covered column
interval $[0,q)$ and possible $O$-column interval $[q,b)$.
These are disjoint, so it avoids $O$.
Its source interior is precisely the displayed set.
The complementary rectangle has $N-b>0$ $O$ markings.
Equation (IG.2) follows in this case.

If $0<b<q$, the footrule sum increases by $2(q-b)$ and the
inversion number increases by
\[
             1+2\#\{i:0<i<q\text{ and }b<\sigma(i)<q\}.
\]
There are $q-b-1$ values strictly between $b$ and $q$.
None is the image of column zero or column $q$.
Subtracting the displayed count from $q-b-1$ therefore gives
\[
                     \#\{i:q<i<N\text{ and }b<\sigma(i)<q\}.
\]
This is the interior count for the rectangle from column $q$ to
column zero. Its column interval $[q,0)$ is disjoint from its
possible marking interval $[b,q)$, so it avoids $O$.
The complementary rectangle has $b>0$ $O$ markings.
Subtracting the inversion change from the footrule change proves
(IG.2) in the second case.

These assertions about the intervals are exactly the native
half-open-square and open-state-interior conventions used in
\decl{GridRectangle.mem_coveredSquares} and
\decl{isEmpty_iff_forall_notMem_cIoo}. The two cases exhaust all
orders of the three distinct columns.
\end{proof}

\begin{lemma}[New lemma: deletion from a minimal cycle]
\label{J21:lem:torus-minimal-cycle-deletion}
Let $3\leq k\leq N$, and suppose $\tau$ is a single $k$-cycle
with $\ell(\tau)=k-1$. For every vertex $q$ of its support, delete
$q$ from its cyclic word and then fix it. Let $\sigma=\tau\setminus q$
be the resulting permutation. Then
\[
                              \ell(\sigma)=k-2,
\]
and exactly one zero-$O$ rectangle goes from $x_\sigma$ to
$x_\tau$. That rectangle is state-empty.
\end{lemma}
\begin{proof}
The permutation $\sigma$ is a $(k-1)$-cycle, so
Lemma~\ref{J21:lem:torus-cycle-lower-bound} gives $\ell(\sigma)\geq k-2$.
Since $k\geq3$, the two neighbors $a,b$ of $q$ in the deleted
cycle are distinct, and $a,b,q$ meet the hypotheses of
Lemma~\ref{J21:lem:torus-exact-insertion-index}.
That lemma gives
$\ell(\tau)-\ell(\sigma)=1+2r$ for a nonnegative integer $r$.
Together with $\ell(\tau)=k-1$, the two inequalities force
$\ell(\sigma)=k-2$ and $r=0$. The same lemma gives uniqueness.
\end{proof}

\begin{lemma}[New lemma: the minimal-cycle boundary coefficient]
\label{J21:lem:torus-minimal-cycle-coefficient}
Let $3\leq k\leq N$, and let $H$ be any chain supported on
single $(k-1)$-cycle states. Write $h(\sigma)=[x_\sigma]H$.
For every single $k$-cycle $\tau$ with $\ell(\tau)=k-1$,
\[
             [x_\tau]d_0H
                  =\sum_{q\in\operatorname{supp}\tau}h(\tau\setminus q).
 \tag{IG.3}
\]
\end{lemma}
\begin{proof}
A swap taking a single $(k-1)$-cycle to a single $k$-cycle must
insert a previously fixed vertex into the cycle.
Swapping two moved columns splits the old cycle, and swapping two
fixed columns creates an additional transposition, so neither can
produce the target.
Thus all possible sources are exactly the vertex deletions in
(IG.3). Each contributes one empty zero-$O$ rectangle by
Lemma~\ref{J21:lem:torus-minimal-cycle-deletion}.
The native coefficient formula \decl{laurentCoefficient_eq}
now gives (IG.3).
\end{proof}

\paragraph{Small cases.}
The exact insertion lemma requires $a,b,q$ distinct.
Accordingly the deletion and coefficient lemmas start at $k=3$.
For $k=2$, the source obtained by deleting a vertex is the anchor,
and its two complementary zero-$O$ rectangles cancel.
No assertion of uniqueness is made in that case.
The applications below use $k=5$, and the converse insertion
applications use source type $(3,2)$, so all their neighbors are
distinct on the smallest permitted grid $N=7$.


\section{Permutation types occurring in the correction}
\label{J21:sec:torus-types}

The standing assumptions are $p\geq3$, $N=2p+1$, and coefficients in
$\mathbb F _2$. For a permutation $\sigma$ of $\mathbb Z/N\mathbb Z$,
write $x_\sigma(c)=-\sigma(c)$ and $\ell(\sigma)=-M_O(x_\sigma)$.
The relative-permutation grading formula gives
\[
 \ell(\sigma)=\sum_i|\sigma(i)-i|-\operatorname{inv}(\sigma).
\]
We use the previously proved invariance of this expression under deletion
of fixed columns and simultaneous cyclic translation of all labels.


For a permutation $\sigma$, define
$r(\sigma)=N-\#\{\text{cycles of }\sigma\}$, including fixed points
in the count. Lemma~\ref{J21:lem:torus-cycle-lower-bound} states
$\ell(\sigma)\geq r(\sigma)$.

\begin{lemma}[New lemma: the change in cycle count under a swap]
\label{J21:lem:torus-transposition-cycle-count}
Multiplying a permutation by a transposition changes $r$ by exactly
one. In particular, a target of a single swap from a four-cycle or
from a permutation of type $(3,2)$ has $r\geq2$.
\end{lemma}
\begin{proof}
If the swapped columns belong to distinct permutation cycles, swapping
their outgoing arrows joins those two cycles into one cycle. If they
belong to the same cycle, write its cyclic word as $(a,A,b,B)$, where
$A,B$ are the lists of entries between the two columns. Swapping the
arrows from $a,b$ separates this word into the two cycles $(a,B)$ and
$(b,A)$. An empty list gives a fixed point. Thus the number of cycles
changes by one in either case. A four-cycle and a permutation of type
$(3,2)$ both have $r=3$, which gives the last assertion.
\end{proof}

\begin{lemma}[New lemma: insertion of a fixed column]
\label{J21:lem:torus-insertion-increases}
Suppose that $q$ is fixed by $\sigma$, that $a\neq q$, and that
$b=\sigma(a)\neq q$. Let $\tau$ be the permutation obtained by replacing
the arrow $a\mapsto b$ with $a\mapsto q\mapsto b$. Then
$\ell(\tau)-\ell(\sigma)$ is a positive odd integer.
Every state-empty rectangle from $x_\sigma$ to $x_\tau$ avoids all
$O$ markings.
\end{lemma}
\begin{proof}
Translate the labels cyclically so that $q=0$. Then $a,b>0$.
The change in the sum of absolute displacements is
\[
 a+b-|a-b|=2\min\{a,b\}.
\]
Put
\[
 t=\#\{i:0<i<a\text{ and }0<\sigma(i)<b\}.
\]
Swapping the values $0,b$ in columns $0,a$ increases the inversion
number by $2t+1$. There are at most $a-1$ eligible columns and at most
$b-1$ eligible values, so $t\leq\min\{a,b\}-1$. Consequently
\[
 \ell(\tau)-\ell(\sigma)=2\min\{a,b\}-2t-1\geq1.
\]
This difference is odd. For an empty rectangle with $o$ markings, the
existing native rectangle grading identity gives
\[
 \ell(\tau)-\ell(\sigma)=1-2o.
\]
Since $o\geq0$, the preceding positive lower bound forces $o=0$.
\end{proof}

\begin{lemma}[New lemma: swapping two fixed columns]
\label{J21:lem:torus-fixed-columns-cancel}
For every permutation $\sigma$ and every distinct fixed columns $a,b$,
the coefficient of $x_{\sigma\circ(a\ b)}$ in $D x_\sigma$ is zero.
\end{lemma}
\begin{proof}
The two candidate rectangles have complementary half-open column and
marking intervals in the relative rectangle formula, so both have
$O$ count zero. Let $I=(a,b)$ and $J=(b,a)$ be the two open cyclic
intervals. Their emptiness conditions are respectively that no arrow
of $\sigma$ goes from $I$ to $J$ and that no arrow goes from $J$ to $I$.
The bijection $\sigma$ preserves $I\cup J$, since it fixes $a,b$.
The two numbers of crossing arrows are therefore equal. Thus both
rectangles are empty or neither is empty. In the former case their
identical terms cancel in $\mathbb F _2$.
\end{proof}

\begin{lemma}[New lemma: inversion of relative permutations]
\label{J21:lem:torus-inverse-commutes}
The linear map $\iota(x_\sigma)=x_{\sigma^{-1}}$ commutes with the
full native Laurent differential and preserves $O$ counts of individual
rectangles.
\end{lemma}
\begin{proof}
On the grid torus, use the coordinate map
\[
 (c,r)\longmapsto(-r,-c).
\]
It sends the point $(c,-\sigma(c))$ to $(\sigma(c),-c)$, which belongs
to $x_{\sigma^{-1}}$. Every $O$ square has center
$(c+\tfrac12,-c-\tfrac12)$, which this map fixes.
The map interchanges the southwest and northeast source corners of
a directed rectangle, and likewise interchanges its two target corners.
It therefore preserves the direction of the rectangle. It is a bijection
on covered squares and on state interiors. Thus it pairs the native
empty rectangles in the two differentials and preserves their weights.
\end{proof}

\begin{lemma}[New lemma: the relevant three-cycle and transposition states]
\label{J21:lem:torus-three-two-type}
Let $\sigma$ be a permutation whose nontrivial cycles consist of one
three-cycle and one transposition. Then $M_O(x_\sigma)$ is either $-3$
or $-5$. It equals $-5$ precisely when the five moved columns have a
unique positive circular ordering
\[
 u,a,v,b,c
\]
in which the transposition is $(u\ v)$ and the three-cycle is either
$(a\ b\ c)$ or $(a\ c\ b)$.
\end{lemma}
\begin{proof}
The two endpoints of the transposition split the other three moved
columns into two circular intervals. If all three belong to the same
interval, the endpoints are adjacent in the circular order of the five
moved columns. Delete the fixed columns and rotate the labels so that
the pair is $(0\ 1)$. The other cycle is $(2\ 3\ 4)$ or its inverse.
For both permutations the footrule sum is $6$ and the inversion number
is $3$. Hence $\ell=3$.

Otherwise one interval contains one triangle vertex and the other
contains two. Orient the pair uniquely so that its positive interval
contains the single vertex, and call that vertex $a$. Call the two
vertices in the remaining interval $b,c$ in positive order. After
deleting fixed columns, the two one-line permutations are
\[
 (2,3,0,4,1),\qquad(2,4,0,1,3).
\]
Both have footrule sum $10$ and inversion number $5$. Hence $\ell=5$.
These alternatives exhaust the distributions of three vertices between
the two intervals, and the construction proves uniqueness of the
intrinsic ordering in the second case.
\end{proof}

\begin{lemma}[New lemma: five-cycle Maslov values]
\label{J21:lem:torus-five-cycle-table}
On any five cyclically ordered columns, the two permutations that
rotate these columns by two or three positions have Maslov value $-6$.
Every other five-cycle on those columns has Maslov value $-4$.
\end{lemma}
\begin{proof}
Delete fixed columns and label the remaining columns $0,1,2,3,4$ in
positive circular order. The following table gives representatives of
the orbits under simultaneous cyclic translation of these five labels.
The last two columns are the footrule sum and the inversion number.
\[
\begin{array}{c|c|r|r|r}
 \text{cycle word}&\text{one-line permutation}&\text{orbit size}
       &\text{footrule}&\operatorname{inv}\\\hline
 (0,1,2,3,4)&(1,2,3,4,0)&1&8&4\\
 (0,1,2,4,3)&(1,2,4,0,3)&5&8&4\\
 (0,1,3,4,2)&(1,3,0,4,2)&5&8&4\\
 (0,1,4,3,2)&(1,4,0,2,3)&5&8&4\\
 (0,2,1,4,3)&(2,4,1,0,3)&5&10&6\\
 (0,2,4,1,3)&(2,3,4,0,1)&1&12&6\\
 (0,3,1,4,2)&(3,4,0,1,2)&1&12&6\\
 (0,4,3,2,1)&(4,0,1,2,3)&1&8&4
\end{array}
\]
The four orbits of size one are the four nonidentity circular rotations.
For each other row, conjugating by the five label rotations gives five
distinct permutations. Comparing the displayed cycle words up to adding
a constant to every entry shows that these four orbits are distinct.
Their sizes, together with the four singleton orbits, sum to $24$,
which is the number of five-cycles on five labels. Thus the table
exhausts the permutations under consideration. Subtracting its last
column from its preceding column and using cyclic-label invariance of
$\ell$ proves the assertion.
\end{proof}

\begin{lemma}[New lemma: the possible differentials of these states]
\label{J21:lem:torus-three-two-boundaries}
Let $H$ be a chain supported on the states of
Lemma~\ref{J21:lem:torus-three-two-type} with Maslov value $-5$, and suppose
that $\iota H=H$. Write $D=\sum_{j\geq0}U^j d_j$.
Then $d_jH=0$ for $j\geq3$. The chain $d_2H$ is supported on
three-cycle states. The chain $d_1H$ is supported on single five-cycle
states of Maslov value $-4$. The chain $d_0H$ is supported on the
following permutation types of Maslov value $-6$: a single five-cycle,
a four-cycle and a transposition, or two three-cycles.
\end{lemma}
\begin{proof}
A rectangle swaps two columns. If both are fixed, its contribution
cancels by Lemma~\ref{J21:lem:torus-fixed-columns-cancel}. If exactly one
is fixed, Lemma~\ref{J21:lem:torus-insertion-increases} gives $O$ count zero.
The new vertex extends either the three-cycle to a four-cycle or the
transposition to a three-cycle. These are the two listed six-column
types in $d_0H$.

If both columns are moved, a swap within the transposition leaves the
three-cycle, a swap within the three-cycle leaves two transpositions,
and a swap between the two cycles joins them to one five-cycle.
The possible Maslov values are $-2$ for a three-cycle, $-2$ or $-4$
for two transpositions, and $-4$ or $-6$ for a five-cycle by
Lemma~\ref{J21:lem:torus-five-cycle-table}.

The rectangle grading identity gives target Maslov value $-6+2j$ for
$d_j$. Thus $j\leq2$, and the listed types have exactly the indicated
degrees. A two-transposition target is fixed by $\iota$. Its possible
sources in $H$ occur in distinct inverse pairs, since a three-cycle is
not equal to its inverse. Lemma~\ref{J21:lem:torus-inverse-commutes} and
$\iota H=H$ make their contributions equal. All such terms cancel.
This proves the assertions.
\end{proof}


\section{The first boundary correction and four-cycle deletion}
\label{J21:sec:torus-first-identities}
\begin{lemma}[New lemma: the transposition differential]
\label{J21:lem:torus-transposition-differential}
Let $a\ne b$, put $T_{ab}=x_{(a\ b)}$, and let
$C_{abc}^{+}$ and $C_{abc}^{-}$ denote the two three-cycle states on
$\{a,b,c\}$. The plus sign denotes the positive ambient
circular orientation, and the minus sign denotes its inverse. Then
\[
 D T_{ab}
 =U\,[d(a,b)=1\ \text{or}\ d(b,a)=1]x_{\mathrm{id}}
   +\sum_{c\notin\{a,b\}}(C_{abc}^{+}+C_{abc}^{-}).
\]
\end{lemma}
\begin{proof}
Every target is obtained by one column transposition, by the definition of
\decl{GridRectangleBetween.emptyRectangles} appearing in
\decl{laurentCoefficient_eq}.
If that transposition is $(a\ b)$, the two possible rectangles have
$O$ counts $d(a,b)$ and $d(b,a)$, respectively. Every strictly intermediate
column on either cyclic interval is fixed by $(a\ b)$ and contributes an
interior state point. Hence the corresponding rectangle is empty exactly
when that interval has length one. This proves the anchor coefficient.

Suppose that the transposition uses one of $a,b$ and a third column $c$.
To exhibit the rectangles without assuming a linear cut, first arrange
$a,c,b$ in this cyclic order. For the swap $(a\ c)$, the rectangle with
ordered sides $(a,c)$ has disjoint half-open intervals $[a,c)$ and $[c,b)$
in Lemma~\ref{J21:lem:torus-relative-rectangle}. It is empty and has $O$ count
zero. For the swap $(b\ c)$, the rectangle with ordered sides $(c,b)$ has
disjoint half-open intervals $[c,b)$ and $[a,c)$. It is empty and has $O$
count zero. These two swaps produce the two cyclic orientations on the
three columns. The alternative rectangle in each case has positive $O$
count. It cannot be empty, since the existing declaration
\decl{maslovO_swap_reflectedTorusAnchor} gives source Maslov grading
$-1$, Corollary~\ref{J21:cor:torus-small-types} gives target Maslov grading $-2$,
and \decl{laurentMaslovDegree_add_card_OColumns} would then force its
$O$ count to be zero. Interchanging $a$ and $b$ covers the other cyclic
order of the three columns.

Finally suppose that the swapped columns $c,d$ are disjoint from
$\{a,b\}$. The two half-open intervals for the rectangle with sides
$(c,d)$ are complementary, so its $O$ count is zero. Every fixed state
point misses its interior. The rectangle is obstructed precisely when
one of $a,b$ belongs to $(c,d)$ and the other belongs to $(d,c)$.
The other rectangle is obstructed under the same condition. Thus the two
rectangles are either both empty or both nonempty. Their contributions
cancel over $\mathbb F_2$. This treats all possible target transpositions.
\end{proof}

\begin{lemma}[New lemma: a balanced-triangle representative]
\label{J21:lem:torus-balanced-first}
For an unordered pair let
$e(a,b)=\min\{d(a,b),d(b,a)\}\bmod2$, and define
\[
 H_e=\sum_{\{a,b\}} e(a,b)T_{ab}.
\]
Then
\[
 x_{\mathrm{id}}+U^{-1}D H_e
 =U^{-1}\sum_{\substack{\{a,b,c\}\text{ with cyclic gaps }\alpha,\beta,\gamma\\
                       \max\{\alpha,\beta,\gamma\}\leq p}}
       (C_{abc}^{+}+C_{abc}^{-}).
\]
The right-hand side is a degree-zero cycle in the exact anchor class.
Writing $g=p(p-1)/2$, every positive cyclic orientation in this sum has
Alexander grading $g-p$, and every negative cyclic orientation has
Alexander grading $g-p-1$.
\end{lemma}
\begin{proof}
There are $N$ adjacent unordered pairs. Their coefficients in $H_e$ are
one, and $N$ is odd. Lemma~\ref{J21:lem:torus-transposition-differential}
therefore makes the anchor coefficient of $D H_e$ equal to $U$.
The coefficient of either orientation on a triple is
$e(a,b)+e(b,c)+e(c,a)$.
If all three cyclic gaps are at most $p$, this is the parity of their
sum $N$, which is one. Otherwise exactly one gap exceeds $p$.
Replacing that gap by the sum of the other two makes the sum of the
three shorter pair distances twice the sum of those two gaps. Its
parity is zero. This proves the chain identity.

The existing declaration \decl{laurentDifferential_reflectedTorusAnchor}
and the square-zero declaration
\decl{laurentDifferential_sq_eq_zero} prove the cycle assertion.
The displayed boundary proves equality with the exact anchor homology
class. Corollary~\ref{J21:cor:torus-small-types} proves degree zero term by term.

Use \decl{alexanderTwo_reflectedTorus_eq_sum_displacement}.
On the positive orientation the three nonzero displacements are
$\alpha,\beta,\gamma\leq p$. Their displacement weights are
$p\alpha,p\beta,p\gamma$, with total $pN$.
On the inverse orientation the displacements are
$N-\alpha,N-\beta,N-\gamma\geq p+1$. Their weights are
$(p+1)\alpha,(p+1)\beta,(p+1)\gamma$, with total $(p+1)N$.
The stated Alexander gradings follow from the native anchor
normalization \decl{alexander_reflectedTorusAnchor}.
\end{proof}

\begin{lemma}[New lemma: deleting a vertex of a four-cycle]
\label{J21:lem:torus-fourcycle-deletion}
Let $\sigma$ be a four-cycle on four distinct columns. Let $v$ be one of
its columns, let $a=\sigma^{-1}(v)$ and $b=\sigma(v)$, and let $\tau$ be
the three-cycle obtained by deleting $v$ from this cyclic word.
The coefficient of $U x_\tau$ in $D x_\sigma$ is
\[
 [d(v,a)=1\text{ or }d(a,v)=1]
 +[d(v,b)=1\text{ or }d(b,v)=1].
\]
Consequently the three-cycle part with $O$ count one in the differential
of a four-cycle is the sum, over its ambient-adjacent cycle edges, of
the two three-cycles obtained by deleting either endpoint.
\end{lemma}
\begin{proof}
The target is obtained by swapping columns $v,a$.
Lemma~\ref{J21:lem:torus-relative-rectangle} gives the two possible $O$ counts
\[
 \min\{d(v,a),d(v,b)\},\qquad
 \min\{d(a,v),d(b,v)\}.
\]
The first count is one precisely when $a$ or $b$ is the next column after
$v$. The second is one precisely when $a$ or $b$ is the previous column.
In either case the corresponding rectangle has column width one or row
height one, so its open interior contains no state point.
The source and target Maslov gradings are $-3$ and $-2$ by
Lemma~\ref{J21:lem:torus-permutation-maslov}. The existing rectangle grading
identity forces every empty rectangle between them to have $O$ count
one. These are therefore all contributions. Adding the two indicators
over $\mathbb F_2$ proves the coefficient formula.
\end{proof}



\section{An explicit coefficient for the three-cycle--transposition layer}
\label{J21:sec:h32-definition}

All coefficient sums in this section take values in $\mathbb F _2$.
Indicators used in coordinate formulas are the ordinary integers zero and one. Fix an integer $p\geq3$ and put $N=2p+1$.
We identify the columns with $\mathbb Z/N\mathbb Z$ and use the native
markings $O(c)=N-1-c$. For distinct columns $r,s$, write
\[
 d(r,s)=(s-r)\bmod N\in\{1,\ldots,N-1\}.
\]
For a permutation $\sigma$ of the columns, put
$x_\sigma(c)=-\sigma(c)\bmod N$.
The state $x_{\mathrm{id}}$ is the native anchor.
These formulas agree with the attached declarations
\decl{reflectedTorusGrid_O_apply} and
\decl{reflectedTorusAnchor_val}.


\subsection{The intrinsic order of the five moved columns}

By Lemma~\ref{J21:lem:torus-three-two-type}, every permutation of type
$(3,2)$ and Maslov value $-5$ has a unique positive circular ordering
\[
                         (u,a,v,b,c)
 \tag{H32.2}
\]
in which its transposition is $(u\,v)$ and its three-cycle is either
$(a\,b\,c)$ or $(a\,c\,b)$. This is the intrinsic ordering used in
the coefficient formula below.

\subsection{The sixteen half-turn chambers}

Take a tuple (H32.2), and use $u$ as the origin.
For $z\in\{a,v,b,c\}$, define
\[
 r_z=d(u,z),\qquad
 \chi_z=[r_z>p]\in\{0,1\},\qquad
 \theta_z=r_z-\chi_z\left(p+\frac12\right).
 \tag{H32.3}
\]
In (H32.3), $\chi_z$ denotes the integer representative zero or one
of the indicator, so the formula for $\theta_z$ is an equality of real
numbers. The four numbers $\theta_z$ belong to $(0,p+\tfrac12)$ and are
pairwise distinct. Distinct numbers from the same half-circle remain
distinct after the translation in (H32.3).
Numbers from opposite half-circles cannot coincide because one
translation differs from the other by a half-integer.
Arrange these four numbers in increasing order and record their
$\chi$ values. This gives a word
\[
                  w=w_1w_2w_3w_4\in\{0,1\}^4.
 \tag{H32.4}
\]
For each of the four vertices let
$t_z=\overline{r_z}\in\mathbb F _2$ be the parity of its distance
from $u$.

\paragraph{New definition.}
Define $h_{32}(u,a,v,b,c)\in\mathbb F _2$ by the following complete table.
\[
\begin{array}{c|c@{\qquad}c|c}
 w&h_{32}&w&h_{32}\\ \hline
 0000&0&1000&1+t_v\\
 0001&1+t_c&1001&t_v+t_c\\
 0010&1+t_c&1010&t_a+t_v\\
 0011&t_b+t_c&1011&0\\
 0100&1+t_a&1100&0\\
 0101&t_a+t_c&1101&1+t_b\\
 0110&0&1110&t_b+t_c\\
 0111&0&1111&0
\end{array}
 \tag{H32.5}
\]
Every entry in this table is an explicit affine expression in four
parities. No choice of a solution of a linear system enters the
definition.

\paragraph{New definition.}
The raw correction chain is
\[
 H_{32}=
 \sum_{\substack{(u,a,v,b,c)\\
                  \text{in positive circular order}}}
 h_{32}(u,a,v,b,c)
 \left(
 x_{(u\,v)(a\,b\,c)}+
 x_{(u\,v)(a\,c\,b)}
 \right).
 \tag{H32.6}
\]
In this sum all five columns are distinct.
The distinguished first entry $u$ is part of the indexing tuple.
By the intrinsic-order lemma, each permutation of cycle type
$(3,2)$ and Maslov value minus five occurs exactly once.
Thus $U^{-3}H_{32}$ has Maslov degree one by the existing declaration
\decl{laurentMaslovDegree}.

\begin{lemma}[New lemma: coefficient symmetries]
\label{J21:lem:h32-invariance}
The coefficient in (H32.6) is preserved by simultaneous cyclic
translation of all columns and by replacing the relative permutation
with its inverse.
\end{lemma}
\begin{proof}
Every quantity in (H32.3)--(H32.5) uses a cyclic distance from the
distinguished column $u$. Simultaneous cyclic translation preserves
these quantities. Inversion changes $(a\,b\,c)$ into $(a\,c\,b)$
and preserves the transposition and its intrinsic order.
The same coefficient is assigned to the two summands in (H32.6).
\end{proof}

\subsection{Inversion and self-inverse targets}

\begin{corollary}[New lemma: inverse pairing of the sources]
\label{J21:cor:h32-selfinverse}
Every self-inverse relative permutation has coefficient zero in
$D H_{32}$.
\end{corollary}
\begin{proof}
Pair every source in (H32.6) with its inverse. The pair has equal
coefficients by the definition, and its two permutations are distinct
because their three-cycles are inverses of each other. For a self-inverse
target, the differential coefficients of the two sources are equal by
Lemma~\ref{J21:lem:torus-inverse-commutes}. They therefore cancel in $\F$.
\end{proof}

\subsection{The complete two-\texorpdfstring{$O$}{O} rectangle coefficient}

Write
\[
                         D=\sum_{j\geq0}U^j d_j,
 \tag{H32.8}
\]
where $d_j$ counts native state-empty rectangles containing exactly
$j$ $O$ markings. This is the coefficient expansion of the attached
\decl{laurentCoefficient_eq}.

\begin{lemma}[New lemma: the two-O rectangle list]
\label{J21:lem:h32-two-O-neighbours}
Let $\tau$ be a three-cycle with support $T$.
The coefficient of $x_\tau$ in $d_2H_{32}$ is
\[
 \sum_{\substack{a\in T\\a-1\notin T,\ a+1\notin T}}
 h_{32}(a-1,a,a+1,b,c),
 \tag{H32.9}
\]
where $b,c$ are the other two elements of $T$ in positive circular
order after $a+1$. This formula is valid for both orientations of
the three-cycle.
\end{lemma}
\begin{proof}
A rectangle swaps two columns. In a source with one three-cycle and
one transposition, swapping columns from different cycles joins them
to a five-cycle. Swapping two columns of the three-cycle gives two
transpositions. The only swap producing a single three-cycle is
therefore the swap that deletes the transposition.
Every possible source has the form $\tau(u\,v)$ with $u,v\notin T$.

For this source, the rectangle restoring the fixed columns $u,v$
has the same positive cyclic interval $[u,v)$ for its column interval
and for the columns of its covered $O$ markings.
Its $O$ count is $d(u,v)$.
The complementary rectangle has $O$ count $d(v,u)$.
Since $N\geq7$, at most one of these counts can equal two.
For a contribution to $d_2$, one of them does equal two.
Relabel its endpoints as $u=a-1$ and $v=a+1$.

This rectangle has width two and height two.
Its only interior column is $a$, and its only interior grid row is
the anchor row $-a$.
It is state-empty exactly when the source does not contain the anchor
point $(a,-a)$.
The transposition fixes the column $a$, so this condition is exactly
$a\in T$.
The requirement that the transposition be disjoint from $\tau$ is
$a-1,a+1\notin T$.
Thus each summand in (H32.9) contributes exactly one rectangle, and
there are no other contributing rectangles.
Its intrinsic order is $(a-1,a,a+1,b,c)$.
The two orientations of the source three-cycle have the same
coefficient by (H32.6), which proves the last assertion.
\end{proof}

\paragraph{New definition.}
For positive integers $r,s$ with $r+s\leq2p$, define
$L_p(r,s)\in\mathbb F _2$ as follows.
Set $L_p(r,s)=0$ if $r=1$ or $s=1$.
For $r,s\geq2$, use the complete table
\[
\begin{array}{c|c}
 \text{Condition}&L_p(r,s)\\ \hline
 s\geq p+2&0\\
 s=p+1&1\\
 s=p&[r=p]\\
 s\leq p-1,\ r\leq p-1&1+\overline{s}\\
 s\leq p-1,\ r=p&\overline{s}\\
 s\leq p-1,\ r=p+1&1+\overline{s}\\
 s\leq p-1,\ r\geq p+2&1+\overline{r}+\overline{s}
\end{array}
 \tag{H32.10}
\]
Here bars denote integer parity.
The conditions in the table are disjoint and exhaustive.
Some rows have no instances for the smallest values of $p$.

\begin{lemma}[New lemma: the complete two-O gap coefficient]
\label{J21:lem:h32-two-O-gap}
Let $a,b,c$ be in positive circular order and put
\[
 \alpha=d(a,b),\qquad
 \beta=d(b,c),\qquad
 \gamma=d(c,a).
\]
The coefficient of either $x_{(a\,b\,c)}$ or $x_{(a\,c\,b)}$
in $d_2H_{32}$ is
\[
 L_p(\alpha,\gamma)+L_p(\beta,\alpha)+L_p(\gamma,\beta).
 \tag{H32.11}
\]
\end{lemma}
\begin{proof}
Consider the contribution centered at $a$ in (H32.9).
It is absent if $\alpha=1$ or $\gamma=1$, which agrees with
the first clause of (H32.10).
Suppose $\alpha,\gamma\geq2$.
Translate the column $a-1$ to zero.
The five intrinsic coordinates become
\[
       (u,a,v,b,c)=(0,1,2,B,C),\qquad
       B=1+\alpha,\qquad C=2p+2-\gamma.
 \tag{H32.12}
\]
They satisfy $3\leq B<C\leq2p$.
In particular $t_a=1$, $t_v=0$,
$t_b=1+\overline{\alpha}$, and
$t_c=\overline{\gamma}$.

The following table obtains the half-turn word directly from
(H32.12). It lists every possible location of $B,C$.
\[
\begin{array}{c|c|c}
 \text{Condition}&w&h_{32}\\ \hline
 C\leq p&0000&0\\
 C=p+1&1000&1\\
 C=p+2,\ B\leq p&0100&0\\
 C=p+2,\ B=p+1&1010&1\\
 C\geq p+3,\ B\leq p&0010\text{ or }0001&1+\overline{\gamma}\\
 C\geq p+3,\ B=p+1&1001&\overline{\gamma}\\
 C\geq p+3,\ B=p+2&0101&1+\overline{\gamma}\\
 C\geq p+3,\ B\geq p+3&0011&
                         1+\overline{\alpha}+\overline{\gamma}
\end{array}
 \tag{H32.13}
\]
For example, a point at coordinate $p+1$ projects to $\tfrac12$,
a point at coordinate $p+2$ projects to $\tfrac32$, and every point
at coordinate at least $p+3$ projects beyond $2$.
The two fixed projected coordinates from $a,v$ are $1,2$.
These facts determine the displayed order in every row.
In the fifth row, the projected coordinate from $C$ lies either before
or after the projected coordinate from $B$.
The two words have the same entry in (H32.5).
Substitution of the four parities into (H32.5) gives the last column.

The equalities in (H32.12) transform (H32.13) exactly into
$L_p(\alpha,\gamma)$ in (H32.10).
Applying the same argument with the vertices cyclically relabelled
gives $L_p(\beta,\alpha)$ and $L_p(\gamma,\beta)$.
Their sum is (H32.9), proving (H32.11).
\end{proof}

\begin{corollary}[New lemma: no further positive-$O$ layers]
\label{J21:cor:h32-positive-types}
The chain $d_2H_{32}$ is supported on three-cycle states, and
$d_jH_{32}=0$ for every $j\geq3$.
\end{corollary}
\begin{proof}
The chain (H32.6) satisfies every hypothesis of
Lemma~\ref{J21:lem:torus-three-two-boundaries}: each source has Maslov value
$-5$, and its coefficient agrees with that of its inverse. That lemma
gives both assertions.
\end{proof}


\section{The four-cycle correction defined by adjacent pairs}
\label{J21:sec:h4-density}

Fix $p\geq3$ and put $N=2p+1$. All coefficients in this section
belong to $\mathbb F _2$. Columns are elements of
$\mathbb Z/N\mathbb Z$. For distinct columns $u,v$, the symbol
$d(u,v)$ denotes the representative of $v-u$ in
$\{1,\ldots,N-1\}$. A bar denotes integer parity.
For a relative permutation $\sigma$, write $x_\sigma(c)=-\sigma(c)$,
as in the native reflected-grid conventions established above.

\subsection{The explicit coefficient and chain}

\paragraph{New definition.}
Define $h_p:\{2,\ldots,2p\}\longrightarrow\mathbb F _2$ by
\[
 h_p(d)=
 \begin{cases}
  \overline d,&2\leq d\leq p,\\
  1+\overline p,&d=p+1,\\
  1,&p+2\leq d\leq2p.
 \end{cases}
 \tag{H4D.1}
\]
For integers $2\leq A<B\leq2p$, put
\[
                    C_p^0(A,B)=h_p(A)+h_p(B).
 \tag{H4D.2}
\]
Define $T_p^0(A,B)$ by the following disjoint and exhaustive table.
\[
\begin{array}{c|c}
 \text{condition}&T_p^0(A,B)\\\hline
 B\leq p&1+\overline B\\
 A\leq p, B\geq p+2, B-A\geq p+1&1+\overline B\\
 A\leq p, B\geq p+2, B-A\leq p&0\\
 B=p+1&1+\overline A+\overline p\\
 A=p+1&1\\
 A\geq p+2&1+\overline A
\end{array}
 \tag{H4D.3}
\]
To verify exhaustiveness, first separate $B\leq p$, $B=p+1$, and
$B\geq p+2$. In the last case, separate $A\leq p$, $A=p+1$, and
$A\geq p+2$, and split $A\leq p$ according to whether $B-A$ is at
most $p$. Thus the definition covers all stated arguments, including
$p=3$.

\paragraph{New definition.}
Let $S$ be a four-element subset of the columns and let $\sigma$ be
a four-cycle with support $S$. For every $u\in S$ with $u+1\in S$,
there are unique integers $2\leq A<B\leq2p$ such that
\[
                         S=\{u,u+1,u+A,u+B\}
 \tag{H4D.4}
\]
with the displayed columns in positive circular order.
Label these four columns by $0,1,2,3$
in that order. The contribution $w_p(\sigma,u)$ is given by the
following complete table. A word in the first column is the cycle
word beginning with label zero.
\[
\begin{array}{c|c@{\qquad}c|c}
 \text{cycle word}&w_p&\text{cycle word}&w_p\\\hline
 0123&C_p^0(A,B)&0321&0\\
 0132&C_p^0(A,B)&0231&0\\
 0213&T_p^0(A,B)&0312&T_p^0(A,B)+C_p^0(A,B)
\end{array}
 \tag{H4D.5}
\]
Define
\[
 f_4(\sigma)=\sum_{\substack{u\in S\\u+1\in S}}w_p(\sigma,u),
 \qquad
 H_4=\sum_{\substack{S\subset\mathbb Z/N\mathbb Z\\|S|=4}}
              \ \sum_{\substack{\sigma\text{ a four-cycle}\\
                                 \operatorname{supp}(\sigma)=S}}
                         f_4(\sigma)x_\sigma.
 \tag{H4D.6}
\]
All indexing sets are finite. For each summand, the adjacent pair
determines its normalization (H4D.4) uniquely, so these formulas
involve no auxiliary choice.

Equivalently, the complete raw correction chain is the explicit sum
\begin{align*}
 H_4=\sum_{u\in\mathbb Z/N\mathbb Z}\ \sum_{2\leq A<B\leq2p}
 \bigl[&C_p^0(A,B)
   \bigl(x_{(u\,u+1\,u+A\,u+B)}
        +x_{(u\,u+1\,u+B\,u+A)}
        +x_{(u\,u+B\,u+1\,u+A)}\bigr)\\
       &+T_p^0(A,B)
   \bigl(x_{(u\,u+A\,u+1\,u+B)}
        +x_{(u\,u+B\,u+1\,u+A)}\bigr)\bigr].
 \tag{H4D.7}
\end{align*}
Repeated states in this expression are added in $\mathbb F _2$.
Equation (H4D.7) follows by distributing the six entries of (H4D.5)
and then interchanging the two finite sums in (H4D.6).

\begin{lemma}[New lemma: grading of the four-cycle correction]
\label{J21:lem:h4-density-degree}
Every state in $H_4$ has Maslov value $-3$. Thus $U^{-2}H_4$ has
Maslov degree one.
\end{lemma}
\begin{proof}
Delete fixed columns and apply the proved new relative-permutation
Maslov formula. The six cycle words in (H4D.5) have the following
footrule sums and inversion numbers.
\[
\begin{array}{c|c|c}
 \text{cycle words}&\sum_i|\sigma(i)-i|&\operatorname{inv}(\sigma)\\\hline
 0123, 0321, 0132, 0231&6&3\\
 0213, 0312&8&5
\end{array}
\]
For each word, subtracting the footrule sum from the inversion number
gives $-3$. Fixed-column deletion and simultaneous cyclic translation
preserve that difference by the new grading lemmas already proved.
Finally the existing declaration \decl{laurentMaslovDegree} gives
$-3-2(-2)=1$ for every term of $U^{-2}H_4$.
\end{proof}

\subsection{The coefficient identities used in the rectangle proof}

\begin{lemma}[New lemma: support on a unit gap]
\label{J21:lem:h4-unit-support}
If a four-element support contains no pair $u,u+1$, then $f_4$ is
zero on every four-cycle with that support.
\end{lemma}
\begin{proof}
The indexing set in the coefficient sum (H4D.6) is empty.
\end{proof}

\begin{lemma}[New lemma: common inverse sum]
\label{J21:lem:h4-common-inverse-sum}
Let $S$ be any four-element subset of the columns. Let $\sigma_+$
be its positive circular four-cycle and let $\sigma_- =\sigma_+^{-1}$.
For every four-cycle $\sigma$ on $S$, one has
\[
 f_4(\sigma_-)=0,\qquad
 f_4(\sigma)+f_4(\sigma^{-1})=f_4(\sigma_+).
 \tag{H4D.8}
\]
\end{lemma}
\begin{proof}
Fix an adjacent pair $u,u+1$ in $S$ and normalize as in (H4D.4).
The inverse pairs of cycle words are precisely the pairs of entries
in the three rows of (H4D.5). In each row their sum is $C_p^0(A,B)$.
The positive circular cycle has word $0123$ in every such
normalization, so its contribution is also $C_p^0(A,B)$.
The negative circular cycle has word $0321$, whose contribution is
zero. Sum these three identities over every adjacent pair of $S$.
Equation (H4D.6) gives exactly (H4D.8).
\end{proof}

\paragraph{New definition.}
For a three-element subset $T$ of the columns, define
\[
 P_3(T)=\sum_{\substack{u\in T\\u+1\in T}}
          h_p(d(u,z)),\qquad
          \{z\}=T\setminus\{u,u+1\}.
 \tag{H4D.9}
\]
The third vertex $z$ is distinct from the adjacent endpoints, so
$2\leq d(u,z)\leq2p$ and every summand belongs to the domain of
(H4D.1).

\begin{lemma}[New lemma: a primitive for the convex coefficient]
\label{J21:lem:h4-convex-primitive}
For every four-element subset $S$, let $C(S)$ denote $f_4$ of its
positive circular cycle. Then
\[
                     C(S)=\sum_{q\in S}P_3(S\setminus\{q\}).
 \tag{H4D.10}
\]
For every five-element subset $T$, one consequently has
\[
                        \sum_{q\in T}C(T\setminus\{q\})=0.
 \tag{H4D.11}
\]
\end{lemma}
\begin{proof}
Expand the right side of (H4D.10) by (H4D.9) and group terms by
their adjacent pair $u,u+1$. If the other two columns of $S$ are
$u+A,u+B$ in the normalization (H4D.4), that pair occurs in exactly
two of the three-element subsets. Its total contribution is
$h_p(A)+h_p(B)=C_p^0(A,B)$.
This is its contribution to the positive circular four-cycle by
(H4D.5). Summing over all adjacent pairs proves (H4D.10).

Expand the left side of (H4D.11) using (H4D.10). Each three-element
subset of $T$ occurs twice, once for each order of deleting its two
complementary vertices. Its coefficient is therefore zero in
$\mathbb F _2$. This proves (H4D.11).
\end{proof}

\paragraph{Exceptional cases.}
The construction applies for every $p\geq3$. It includes adjacent
pairs crossing the cyclic cut because all normalizations use the
cyclic distance. When a four-element support contains several
adjacent pairs, (H4D.6) includes each pair once, including the case
of four consecutive columns. The definition and the proofs above
require no nonempty interval between the distinguished values
$2,p,p+1,p+2,2p$, so they apply at the smallest allowed value $p=3$.


\section{The complete lower boundary of the four-cycle correction}
\label{J21:sec:h4-lower-boundary}

Fix an integer $p\geq3$ and put $N=2p+1$.
All coefficients and all displayed polynomial identities in this section
are in $\mathbb F _2$. A bar denotes reduction of an integer modulo two.
The function $f_4$ and the chain $H_4$ are the explicit four-cycle
coefficient and chain defined in (H4D.5)--(H4D.7).
The chain $H_{32}$ is defined in (H32.6).
We use the decomposition
\[
                         D=\sum_{j\geq0}U^j d_j
\]
of the native Laurent differential, where $d_j$ counts state-empty
rectangles with exactly $j$ covered $O$ markings.
This is the existing declaration
\decl{laurentCoefficient_eq} with coefficients grouped by their
nonnegative $O$ counts.

\subsection{The coefficient at a prescribed unit gap}

For positive integers $a,b,c$ with $a+b+c=2p$, define
\[
 K_p(a,b,c)=
 f_4\bigl((0\,1\,1+a+b\,1+a)\bigr).
 \tag{H4L.1}
\]
The four columns in positive circular order are
$0,1,1+a,1+a+b$.
For every integer $u$ with $1\leq u\leq2p-2$, define
\[
 F_p(u)=
 \begin{cases}
  1+[u=p-1],&1\leq u\leq p,\\
  \overline u,&p+1\leq u\leq2p-2.
 \end{cases}
 \tag{H4L.2}
\]
Also define
\[
 B_p(a,b,c)=
 \begin{cases}
  [a\leq p]\,\overline a,&c<p,\\
  \overline{a+p},&c=p,\\
  \overline b,&c>p.
 \end{cases}
 \tag{H4L.3}
\]

\begin{lemma}[New lemma: the unit-gap coefficient]
\label{J21:lem:h4-unit-coefficient}
For all positive $a,b,c$ with $a+b+c=2p$,
\[
 K_p(a,b,c)=B_p(a,b,c)+[a=1]F_p(b)+[c=1]F_p(a).
 \tag{H4L.4}
\]
\end{lemma}
\begin{proof}
Put $A=1+a$ and $B=1+a+b=N-c$.
The support in (H4L.1) is $(0,1,A,B)$ in positive circular order.
Apply the explicit adjacent-pair density (H4D.5) to the cycle
$(0\,1\,B\,A)$.
The pair $(0,1)$ is a forward cycle edge, with contribution
$C_p^0(A,B)$.
The pair $(1,A)$ occurs precisely when $A=2$, or equivalently $a=1$.
After normalizing this pair to $(0,1)$, the cycle word is $0213$,
so its contribution is $T_p^0(B-1,2p)$.
The pair $(A,B)$, when it occurs, is a reverse cycle edge and has
zero contribution.
The pair $(B,0)$ occurs precisely when $B=2p$, or equivalently $c=1$.
Its normalized cycle word is $0312$, so its contribution is
$T_p^0(2,A+1)+C_p^0(2,A+1)$.
These are all possible ambient adjacent pairs in the support.
Consequently
\begin{align*}
 K_p(a,b,c)={}&C_p^0(A,B)
       +[a=1]T_p^0(B-1,2p)\\
       &+[c=1]\bigl(T_p^0(2,A+1)+C_p^0(2,A+1)\bigr).
 \tag{H4L.5}
\end{align*}
This formula also holds when more than one of these unit pairs occurs,
because the coefficient is defined as the sum of their separate
contributions.

The definition of $h_p$ in (H4D.1) gives
\[
 C_p^0(1+a,N-c)=B_p(a,b,c).
 \tag{H4L.6}
\]
Indeed, if $c<p$, then $h_p(N-c)=1$.
Adding $h_p(1+a)$ gives $\overline a$ for $a\leq p$ and zero for
$a>p$.
If $c=p$, then $a+b=p$ and the sum is
$\overline{a+1}+1+\overline p=\overline{a+p}$.
If $c>p$, both $1+a$ and $N-c$ are at most $p$.
The sum is $\overline{1+a+N-c}=\overline{a+c}=\overline b$.
These are the three rows of (H4L.3).

For $1\leq u\leq2p-2$, the table (H4D.3) also gives the two identities
\[
 T_p^0(u+1,2p)=F_p(u),\qquad
 T_p^0(2,u+2)+h_p(u+2)=F_p(u).
 \tag{H4L.7}
\]
Here are all their cases.
\[
\begin{array}{c|c|c}
 \text{range of }u&T_p^0(u+1,2p)&T_p^0(2,u+2)+h_p(u+2)\\ \hline
 1\leq u\leq p-2&1&1\\
 u=p-1&0&0\\
 u=p&1&1\\
 p+1\leq u\leq2p-2&\overline u&\overline u
\end{array}
 \tag{H4L.8}
\]
In the first column of values, the first range has
$2p-(u+1)\geq p+1$, the next range has $u+1=p$, and the third has
$u+1=p+1$.
The final range has $u+1\geq p+2$.
These respectively use the second, third, fifth, and sixth rows of
(H4D.3).
For the second column of values, the subrange $u\leq p-2$ has
$u+2\leq p$, so its terms are
$1+\overline{u+2}$ and $\overline{u+2}$.
At $u=p-1$ they are $1+\overline p$ and $1+\overline p$.
At $u=p$ they are zero and one.
In the final range they are $1+\overline{u+2}$ and one.
This proves the entire table and (H4L.7).

If $a=1$, then $B-1=b+1$, so the second summand of (H4L.5) is
$[a=1]F_p(b)$.
If $c=1$, then $A+1=a+2$ and $h_p(2)=0$, so the last summand is
$[c=1]F_p(a)$.
Using (H4L.6) proves (H4L.4).
\end{proof}

\subsection{One-\texorpdfstring{$O$}{O} rectangles ending at a triangle}

\begin{lemma}[New lemma: the complete triangle coefficient]
\label{J21:lem:h4-triangle-coefficient}
Let $A,B,C$ be distinct columns in positive circular order, and put
\[
 a=d(A,B),\qquad b=d(B,C),\qquad c=d(C,A).
\]
Thus $a,b,c>0$ and $a+b+c=2p+1$.
The coefficient of either triangle orientation on this support in
$d_1H_4$ is
\[
 J_p(a,b,c)=
 \sum_{\mathrm{cyc}\,(a,b,c)}
 [a\geq2]\bigl(K_p(a-1,b,c)+K_p(b,c,a-1)\bigr).
 \tag{H4L.10}
\]
\end{lemma}
\begin{proof}
We use the proved new rectangle deletion formula,
Lemma~\ref{J21:lem:torus-fourcycle-deletion}.
If a four-cycle has consecutive cycle entries $(r,q,s)$ and deleting
$q$ produces the target triangle, its one-$O$ rectangle coefficient
is
\[
               [d(r,q)=1\text{ or }d(q,r)=1]
              +[d(s,q)=1\text{ or }d(q,s)=1].
 \tag{H4L.11}
\]
This formula follows from the native empty-rectangle criterion and
the $O$ interval formula. In particular every contributing removed
column is an ambient neighbor of one of the three target columns.

Group the terms of (H4L.11) by that adjacent target column.
In the open gap from $A$ to $B$, there are no such terms if $a=1$.
If $a\geq2$, the neighbor after $A$ is $q=A+1$.
For the positive triangle orientation, the two terms belonging to
this neighbor are the four-cycles
\[
                       (A\,q\,B\,C),\qquad
                       (A\,B\,C\,q).
\]
Their sum of coefficients equals the coefficient of
$(A\,q\,C\,B)$.
Indeed, Lemma~\ref{J21:lem:h4-common-inverse-sum} says that the sum of
the coefficients of a four-cycle and its inverse is the forward
convex coefficient, and the inverse of the forward convex cycle has
coefficient zero by (H4D.5).
After translating $A$ to zero, the resulting coefficient is
$K_p(a-1,b,c)$.
For the negative triangle orientation, the same two source cycles
are replaced by their inverses.
One is the reverse convex cycle with coefficient zero, and the other
has coefficient $K_p(a-1,b,c)$.
Thus both target orientations give the same result.

The neighbor before $B$ is $q'=B-1$.
Starting the positive circular order at $q'$ gives
$(q',B,C,A)$ with gaps $(1,b,c,a-1)$.
The identical argument gives $K_p(b,c,a-1)$ for either target
orientation.
If $a=2$, then $q=q'$.
Counting these as two neighbor incidences is required by
(H4L.11). A source adjacent to both $A$ and $B$ is then counted twice
and has zero coefficient, as (H4L.11) prescribes.
Applying the same grouping to the other two gaps proves (H4L.10).
\end{proof}

\begin{lemma}[New lemma: cancellation of the entire triangle layer]
\label{J21:lem:h4-h32-triangle-cancellation}
Let $T_{\mathrm{bal}}$ be the sum of both orientations of all
three-cycles whose three positive circular gaps are at most $p$.
The triangle coefficients satisfy
\[
                   d_1H_4+d_2H_{32}=T_{\mathrm{bal}}.
 \tag{H4L.12}
\]
Here the equality is restricted to three-cycle targets. The complete
chain equality will follow below after the other targets are excluded.
\end{lemma}
\begin{proof}
For $a\geq2$, set
\[
 T_p(a,b,c)=K_p(a-1,b,c)+K_p(b,c,a-1).
\]
Formula (H4L.4) gives
\begin{align*}
 T_p(a,b,c)={}&[c<p]\,\overline c+[b>p]\,\overline b
       +[a=p+1]\,\overline p\\
     &\quad+[c=1]F_p(a-1)+[b=1]F_p(c).
 \tag{H4L.13}
\end{align*}
For completeness, put $x=a-1$, $y=b$, and $z=c$, so
$x+y+z=2p$.
The two terms involving $[x=1]F_p(y)$ cancel.
The sum $B_p(x,y,z)+B_p(y,z,x)$ is zero if $z\geq p$.
For $z<p$, it is respectively
\[
 \overline x\quad\text{if }y>p,\qquad
 \overline y\quad\text{if }x=p,\qquad
 \overline z\quad\text{otherwise}.
\]
These disjoint cases follow from the three rows of (H4L.3).
Since $\overline{x+y+z}=0$, their common indicator expression is
the first line of (H4L.13).
This proves (H4L.13) without any restriction on the sizes of the gaps.

Suppose first that all three gaps $a,b,c$ are at least two.
The second line of (H4L.13) vanishes.
Summing the first line cyclically yields
\[
 J_p(a,b,c)=1+
 \overline p\left(\sum_{v\in\{a,b,c\}}[v=p]
                    +\sum_{v\in\{a,b,c\}}[v=p+1]\right).
 \tag{H4L.14}
\]
The initial one is $\overline{a+b+c}$.
If exactly one gap is one, rotate the triple to $(1,b,c)$, where
$b,c\geq2$ and $b+c=2p$.
The two nonzero summands in (H4L.10) give
\[
 J_p(1,b,c)=1+\overline p\bigl([b=p+1]+[c=p+1]\bigr)
                 +F_p(b-1)+F_p(b).
 \tag{H4L.15}
\]
The terms $[c>p]\overline c$ and $[b<p]\overline b$ cancel
because their indicators are equal and $\overline b=\overline c$.
If two gaps are one, rotate them to $(1,1,2p-1)$.
Equation (H4L.13), with (H4L.2), gives $J_p=0$.

The coefficient of $d_2H_{32}$ is the proved expression
\[
                L_p(a,c)+L_p(b,a)+L_p(c,b)
 \tag{H4L.16}
\]
from (H32.10)--(H32.11).
The following table displays all cases of (H4L.14)--(H4L.16).
Write $\eta=\overline p$.
\[
\begin{array}{c|c|c|c}
 \text{Gap condition}&J_p&d_2H_{32}&[\max(a,b,c)\leq p]\\ \hline
 a,b,c\geq2,\ \max(a,b,c)<p&1&0&1\\
 a,b,c\geq2,\ \max(a,b,c)=p&1+\eta&\eta&1\\
 a,b,c\geq2,\ c=p+1&1+\eta&1+\eta&0\\
 a,b,c\geq2,\ c\geq p+2&1&1&0\\ \hline
 (a,b,c)=(1,b,2p-b),\ 2\leq b\leq p-2&1&1&0\\
 (a,b,c)=(1,p-1,p+1)&\eta&\eta&0\\
 (a,b,c)=(1,p,p)&0&1&1\\
 (a,b,c)=(1,p+1,p-1)&1&1&0\\
 (a,b,c)=(1,b,2p-b),\ p+2\leq b\leq2p-2&0&0&0\\ \hline
 (a,b,c)=(1,1,2p-1)&0&0&0
\end{array}
 \tag{H4L.17}
\]
The third and fourth rows permit cyclic rotation to put the unique
gap greater than $p$ in the last position.
All other rows are also understood up to the cyclic rotation already
specified. No permutation of the gap order is used.

Here is a direct verification of the third column.
If all gaps are less than $p$, its three terms are
$1+\overline c$, $1+\overline a$, and $1+\overline b$,
whose sum is zero.
If all gaps are at least two and one equals $p$, exactly one does.
Rotating it to $c=p$ gives the three terms
$0$, $1+\overline a$, and $\overline b$, whose sum is $\eta$.
For $c=p+1$, the three terms are
$1$, $1+\overline a$, and $1+\overline b$.
Since $a+b=p$, their sum is $1+\eta$.
For $c\geq p+2$, the terms are
$0$, $1+\overline a$, and
$1+\overline c+\overline b$, whose sum is one.
When $a=1$, the first two terms of (H4L.16) vanish.
Substitution of the five possible ranges of $b$ into $L_p(c,b)$
gives respectively $1,\eta,1,1,0$.
When two gaps are one, each term of (H4L.16) contains an argument
equal to one and is zero.
The second column follows from (H4L.14), (H4L.15), and the two rows
of (H4L.2).
Thus every row is an exact algebraic case, rather than an enumeration
at finitely many grid sizes.

In each row, the sum of the second and third columns is the fourth
column. Lemma~\ref{J21:lem:h4-triangle-coefficient} and (H32.11) hold
for both orientations of a triangle, so this proves (H4L.12) for
every three-cycle target.
\end{proof}

\subsection{Cancellation at targets with two transpositions}

\begin{lemma}[New lemma: the two-transposition coefficients]
\label{J21:lem:h4-double-transposition-cancellation}
Every relative permutation consisting of two disjoint transpositions
has coefficient zero in $D H_4$.
\end{lemma}
\begin{proof}
Let the four moved columns be $A<B<C<D$ after a cyclic cut.
The three possible pairings are
\[
 (A\,B)(C\,D),\qquad
 (A\,D)(B\,C),\qquad
 (A\,C)(B\,D).
\]
Their Maslov values, from the proved new inversion-minus-footrule
formula, are respectively $-2,-2,-4$.
A rectangle from a four-cycle can reach such a target only by
cutting that four-cycle into its two transpositions.
Conversely, joining two columns that belong to different target
transpositions produces every possible four-cycle source.
There are four sources, partitioned into two inverse pairs.

For the first pairing, these inverse pairs are
\[
 \{(A\,B\,C\,D),(A\,D\,C\,B)\},\qquad
 \{(A\,B\,D\,C),(A\,C\,D\,B)\}.
\]
For either inverse pair, each source has one-$O$ differential
coefficient
\[
                  e=[C-B=1]+[N-D+A=1].
 \tag{H4L.18}
\]
To verify (H4L.18), the source $(A\,B\,C\,D)$ swaps the columns
$B,D$.
Its two possible rectangle $O$ counts are $C-B$ and $N-D+A$.
The source $(A\,B\,D\,C)$ swaps $B,C$, with the same two
$O$ counts.
When the relevant count is one, the open overlap of the column
interval and the corresponding anchor-row interval contains no fixed
state point. The two remaining moved points lie outside the
rectangle interior. Thus that rectangle is state-empty.
When the count is not one, it does not contribute to $d_1$.
This proves the assertion for these two sources.
The proved new inversion symmetry of the rectangle differential
gives the same assertion for their inverses.
By Lemma~\ref{J21:lem:h4-common-inverse-sum}, the coefficient sum
within each source pair is the same forward convex coefficient $C_0$.
The total coefficient of the target is therefore $eC_0+eC_0=0$.
The native rectangle Maslov relation excludes $O$ counts other than
one for this target.

For the second pairing, take the positive circular order
$(B,C,D,A)$.
It becomes the first pairing in that order, so the same proof applies.
Equivalently, its common rectangle coefficient is
$[D-C=1]+[B-A=1]$, and it again multiplies $C_0+C_0$.

For the crossing pairing $(A\,C)(B\,D)$, the four sources again
form two inverse pairs. Each source has exactly one state-empty
rectangle with no $O$ marking.
For one inverse pair, choose the source $(A\,D\,B\,C)$ and
swap $A,B$.
The rectangle with column interval $[A,B)$ and anchor-row interval
$[C,D)$ has disjoint intervals. It contains no $O$ marking and no
state point in its interior.
The complementary rectangle has positive $O$ count.
For the other inverse pair, choose $(A\,B\,D\,C)$ and swap
$A,D$.
The rectangle with column interval $[D,A)$ and anchor-row interval
$[B,C)$ likewise has disjoint intervals and is state-empty.
Its complement has positive $O$ count.
Inversion supplies the same rectangle count for the other two
sources.
The common inverse-sum lemma therefore gives total coefficient
$C_0+C_0=0$.
The Maslov relation requires $O$ count zero for this target.
All three pairings have been treated, proving the lemma.
\end{proof}

\begin{corollary}[New lemma: the full lower-layer identity]
\label{J21:cor:h4-h32-lower-identity}
The following equalities of complete chains hold.
\[
                    d_1H_4+d_2H_{32}=T_{\mathrm{bal}},
             \qquad d_jH_4=0\quad(j\geq2).
 \tag{H4L.19}
\]
Moreover $d_0H_4$ has no target consisting of two transpositions.
\end{corollary}
\begin{proof}
Every source in $H_4$ is a four-cycle of Maslov value $-3$.
Swapping two columns of its support cuts the cycle into a three-cycle
with a fixed column or into two transpositions.
Swapping a moved column with a fixed column inserts that column into
the cycle and therefore produces a five-cycle.
The proved new insertion inequality implies that an empty rectangle
inserting a fixed column has $O$ count zero.
Swapping two fixed columns contributes zero by the proved new
two-fixed-column rectangle pairing.
These alternatives exhaust all swaps.
Consequently the only possible targets in $d_1H_4$ are three-cycles
and two transpositions.
The latter vanish by
Lemma~\ref{J21:lem:h4-double-transposition-cancellation}.
The proved new inversion pairing for $H_{32}$ and (H32.11) show
that $d_2H_{32}$ is supported entirely on three-cycles.
Lemma~\ref{J21:lem:h4-h32-triangle-cancellation} now proves the first
complete chain equality in (H4L.19).

For the second equality, a four-cycle has $r=3$.
By Lemma~\ref{J21:lem:torus-transposition-cycle-count}, multiplication by
one transposition leaves $r\geq2$. The proved new inequality
$-M(x_\sigma)\geq r(\sigma)$ therefore gives target Maslov value
at most $-2$.
The native empty-rectangle Maslov relation gives target Maslov value
$-4+2j$ for an $O$ count $j$.
It follows that $j\leq1$, as required.
The last assertion is part of
Lemma~\ref{J21:lem:h4-double-transposition-cancellation}.
\end{proof}

\paragraph{Smallest sizes and the scope of the argument.}
All formulas above hold at $p=3$, corresponding to the smallest
allowed $n=2$.
Intervals such as $2\leq b\leq p-2$ are then empty and contribute
no case. The values $F_p(p-1)=0$ and $F_p(p)=1$ are explicitly
retained, including the extreme balanced triple $(1,p,p)$.
The proof counts both incidences when a triangle gap equals two.
It therefore also covers the coincidence of the two inserted
neighbor columns.
This section proves all lower-layer cancellations.
The remaining Maslov-minus-four five-cycle coefficients are computed
in Sections~\ref{J21:sec:h4-m4-A} and~\ref{J21:sec:m4-B}.


\section{The complete Maslov minus six boundary of the explicit \texorpdfstring{$(3,2)$}{(3,2)} correction}
\label{J21:sec:h32-m6}

This section proves a support statement for the explicit chain $H_{32}$
defined in Section~\ref{J21:sec:h32-definition}. Every coefficient below belongs
to $\mathbb F _2$. We retain the integers $p\geq3$, $N=2p+1$, and
$g=p(p-1)/2$, the native markings $O(c)=-c-1$ and $X(c)=p-c-1$, and
$x_\sigma(c)=-\sigma(c)$, with all coordinates reduced modulo $N$.
We write
\[
 S(a,b)=[d(a,b)\leq p]\qquad(a\neq b).
\]
Thus $S(a,b)+S(b,a)=1$. This identity follows because the two positive
cyclic distances sum to $N=2p+1$.
The operators $d_j$ are defined by $D=\sum_{j\geq0}U^j d_j$ as in
(H32.8). The raw chain $d_0H_{32}$ has Maslov value minus six by the
existing declaration
\decl{laurentMaslovDegree_add_card_OColumns}.

For three distinct vertices $a,b,c$ in positive circular order, define
\[
 B(a,b,c)=S(a,b)S(b,c)S(c,a).
 \tag{M6.1}
\]
We call their triangle balanced if $B(a,b,c)=1$.
For five distinct vertices $c_0,\ldots,c_4$ in positive circular order,
we call their pentagon balanced if
\[
             S(c_i,c_{i+2})=1\quad\text{for every }i\pmod5.
 \tag{M6.2}
\]
These are definitions by inequalities in the native column coordinates.
They impose no condition on the gaps being equal.

\subsection{Two elementary primitives for the coefficient table}

A half-turn word for a positive circularly ordered tuple is defined by
(H32.3)--(H32.4), with the first vertex as origin.
The same definition applies to tuples of three or four vertices.
For a four-tuple $(u,a,v,z)$ put $t_x=\overline{d(u,x)}$ and define
$F(u,a,v,z)$ by
\[
\begin{array}{c|c@{\qquad}c|c}
 w&F&w&F\\ \hline
 000&0&100&1+t_v\\
 001&1+t_z&101&0\\
 010&1+t_a&110&1+t_z\\
 011&0&111&0
\end{array}
 \tag{M6.3}
\]
For a positive triple $(u,a,v)$ define $K(u,a,v)$ by
\[
\begin{array}{c|cccc}
 w&00&01&10&11\\ \hline
 K&0&0&t_v&t_a+t_v.
\end{array}
 \tag{M6.4}
\]
The quantities in (M6.3)--(M6.4) are new explicitly defined functions.

\begin{lemma}[New lemma: primitive identities]\label{J21:lem:m6-primitives}
For every positive circularly ordered five-tuple $(u,a,v,b,c)$,
\[
 h_{32}(u,a,v,b,c)=F(u,a,v,b)+F(u,a,v,c).
 \tag{M6.5}
\]
For every positive circularly ordered four-tuple $(u,a,v,b)$,
\[
 F(u,a,v,b)+F(v,b,u,a)=K(u,a,v)+K(v,b,u).
 \tag{M6.6}
\]
Moreover the function
\[
 G(u,a,v,b)=F(u,a,v,b)+F(a,v,b,u)
\]
satisfies
\[
\begin{split}
 G(u,a,v,b)={}&S(u,a)S(a,v)\,\overline{d(u,a)}
       +S(u,v)(1+S(a,b))\\
       &+(1+S(v,b))\,\overline{d(a,b)}.
\end{split}
 \tag{M6.7}
\]
\end{lemma}
\begin{proof}
We give all finite algebraic cases used in this lemma.
Write a tuple as $(c_0,\ldots,c_{q-1})$ in positive circular order,
with origin $c_0$, and put $t_i=\overline{d(c_0,c_i)}$, so $t_0=0$.
If the origin changes to $c_i$, then
\[
 \overline{d(c_i,c_j)}=t_i+t_j+[j<i].
 \tag{M6.8}
\]
The last summand is present because a wrap adds the odd integer $N$.
Consequently all the identities are identities in formal parity
variables, with no assumption on their numerical values.

For (M6.5), use the indexing $(u,a,v,b,c)=(c_0,c_1,c_2,c_3,c_4)$.
The next table lists the two four-point values obtained by deleting
$c$ or $b$. The sum of its last two columns is precisely the
corresponding entry of (H32.5).
\[
\begin{array}{c|c|c@{\qquad}c|c|c}
 w&F(u,a,v,b)&F(u,a,v,c)&w&F(u,a,v,b)&F(u,a,v,c)\\ \hline
 0000&0&0&1000&0&1+t_2\\
 0001&0&1+t_4&1001&1+t_2&1+t_4\\
 0010&0&1+t_4&1010&1+t_2&1+t_1\\
 0011&1+t_3&1+t_4&1011&0&0\\
 0100&0&1+t_1&1100&1+t_2&1+t_2\\
 0101&1+t_1&1+t_4&1101&1+t_3&0\\
 0110&1+t_1&1+t_1&1110&1+t_3&1+t_4\\
 0111&0&0&1111&0&0
\end{array}
 \tag{M6.9}
\]
Deleting a point deletes its projected position from the ordered
half-turn list. This describes directly how the entries in (M6.9)
follow from (M6.3).

For the remaining identities use
$(u,a,v,b)=(c_0,c_1,c_2,c_3)$.
The following table substitutes (M6.3)--(M6.4), changing origins by
(M6.8) where necessary.
\[
\begin{array}{c|c|c|c|c}
 w&F(u,a,v,b)+F(v,b,u,a)&K(u,a,v)&K(v,b,u)&G(u,a,v,b)\\ \hline
 000&0&0&0&t_1\\
 001&1+t_3&0&1+t_3&1+t_3\\
 010&1+t_2&0&1+t_2&1+t_1\\
 011&0&0&0&0\\
 100&1+t_2&0&1+t_2&t_1\\
 101&t_2&t_2&0&t_1\\
 110&t_2&t_2&0&t_1\\
 111&t_1+t_2&t_1+t_2&0&0
\end{array}
 \tag{M6.10}
\]
This proves (M6.6). To verify (M6.7) without any unlisted geometric
cases, the six indicators
\[
 (S(u,a),S(u,v),S(u,b),S(a,v),S(a,b),S(v,b))
\]
have respective values
\[
\begin{array}{c|c@{\qquad}c|c}
 w&(S_{01},S_{02},S_{03},S_{12},S_{13},S_{23})
 &w&(S_{01},S_{02},S_{03},S_{12},S_{13},S_{23})\\ \hline
 000&111111&100&110111\\
 001&110100&101&100101\\
 010&110101&110&100111\\
 011&100001&111&000111
\end{array}
\]
Substituting these eight strings in the right side of (M6.7), with
$\overline{d(a,b)}=t_1+t_3$, gives its last column in (M6.10).
The eight words exhaust the four-point half-turn orders. This proves
all three identities.
\end{proof}

\subsection{Which rectangles can contribute}

The next arguments use the relative-permutation Maslov formula and
fixed-column deletion lemma in Section~\ref{J21:sec:relative-permutation-maslov},
and the new exact insertion-index identity (IG.2) in
Lemma~\ref{J21:lem:torus-exact-insertion-index}. In the notation used there,
\[
 \ell(\sigma)=-M(x_\sigma)
       =\sum_c|\sigma(c)-c|-\operatorname{inv}(\sigma).
 \tag{M6.11}
\]
An empty zero-$O$ rectangle increases $\ell$ by one by
\decl{laurentMaslovDegree_add_card_OColumns}.


\begin{corollary}[New lemma: the Maslov minus six types]
\label{J21:lem:m6-types}
Every nonzero term of $d_0H_{32}$ has relative cycle type $(5)$,
$(4,2)$, or $(3,3)$ and Maslov value $-6$. For type $(5)$, the
only possibilities are the two star permutations $c_i\mapsto c_{i+2}$
and $c_i\mapsto c_{i+3}$ on five circularly ordered columns.
\end{corollary}
\begin{proof}
Apply Lemma~\ref{J21:lem:torus-three-two-boundaries} to the chain (H32.6).
Its inverse symmetry follows from Lemma~\ref{J21:lem:h32-invariance}.
The five-cycle assertion is Lemma~\ref{J21:lem:torus-five-cycle-table}.
\end{proof}

For later use we record why deleting a vertex gives every contribution
to the targets with two nontrivial cycles.
If a target of type $(4,2)$ or $(3,3)$ has $\ell=6$, delete a vertex from its four-cycle in the $(4,2)$ case, or from
either three-cycle in the $(3,3)$ case. This gives a possible
source of type $(3,2)$.
It occurs in $H_{32}$ exactly when its two transposition endpoints
belong to different complementary arcs of its triangle.
This is the intrinsic-order lemma of
Section~\ref{J21:sec:h32-definition}. Such a source has $\ell=5$.
Lemma~\ref{J21:lem:torus-exact-insertion-index} then states that its unique zero-$O$ reinsertion
rectangle has no interior state points, because
$\ell(\text{target})-\ell(\text{source})=1$.
Thus each eligible deletion contributes its coefficient exactly once.
No other source of type $(3,2)$ can have that target, since the swap
must remove precisely the newly fixed vertex.

\subsection{Cancellation of every four-cycle--transposition target}

\begin{lemma}[New lemma: the $(4,2)$ coefficient]\label{J21:lem:m6-four-two}
Every relative permutation of cycle type $(4,2)$ has coefficient zero
in $d_0H_{32}$.
\end{lemma}
\begin{proof}
It suffices to consider a target with $\ell=6$.
Let its transposition be $(u\ v)$.
If all four other vertices belong to the same complementary arc of
$u,v$, deleting any one of them leaves the transposition endpoints
in the same arc of the remaining triangle. There are no eligible
sources, so the coefficient is zero.

Suppose the two arcs contain one and three of the four-cycle vertices.
Write their positive order as $(u,a,v,b_1,b_2,b_3)$.
Deleting $a$ is ineligible. Deleting one of the three $b$ vertices
leaves the three eligible sources, and their total coefficient is
\[
 h_{32}(u,a,v,b_1,b_2)+h_{32}(u,a,v,b_1,b_3)
                         +h_{32}(u,a,v,b_2,b_3)=0.
\]
Indeed, (M6.5) writes this sum as a sum in which each of
$F(u,a,v,b_1),F(u,a,v,b_2),F(u,a,v,b_3)$ occurs twice.
The orientation of the remaining three-cycle does not affect its
coefficient by (H32.6).

Finally suppose the two arcs contain two vertices each, with order
$(u,a_1,a_2,v,b_1,b_2)$. All four deletions are eligible. Their sum is
\[
\begin{split}
 &h_{32}(u,a_1,v,b_1,b_2)+h_{32}(u,a_2,v,b_1,b_2)\\
 &\qquad+h_{32}(v,b_1,u,a_1,a_2)+h_{32}(v,b_2,u,a_1,a_2)\\
 &=\sum_{i,j=1}^2\bigl(F(u,a_i,v,b_j)+F(v,b_j,u,a_i)\bigr)\\
 &=\sum_{i,j=1}^2\bigl(K(u,a_i,v)+K(v,b_j,u)\bigr)=0.
\end{split}
\]
The first identity is (M6.5), the second is (M6.6), and every term
in the last sum occurs twice. The three distributions exhaust the
positions of the four-cycle vertices.
\end{proof}

\subsection{The star coefficient}

\begin{lemma}[New lemma: the pentagon coefficient]\label{J21:lem:m6-pentagon}
Let $c_0,\ldots,c_4$ be five vertices in positive circular order.
The coefficient in $d_0H_{32}$ of each of their two star states is
one if the pentagon is balanced and is zero otherwise.
\end{lemma}
\begin{proof}
Write the indices modulo five. The five possible sources are the
permutations with transposition $(c_i\ c_{i+2})$ and a three-cycle
on the other three vertices. Their intrinsic orders are
$(c_i,c_{i+1},c_{i+2},c_{i+3},c_{i+4})$.
For the forward star, the source at $i=0$ is
\[
 (c_0\ c_2)(c_1\ c_3\ c_4).
\]
Its rectangle from column $c_2$ to column $c_3$ is empty and avoids
$O$. There are no moved columns in that open column interval.
The columns of its covered $O$ markings would have to belong to
both $[c_2,c_3)$ and $[c_4,c_0)$, which are disjoint.
An interior fixed state point would force its adjacent $O$ center
inside the rectangle, so no fixed point is in its interior either.
Its complementary rectangle contains $O$.
For the inverse star the source is
$(c_0\ c_2)(c_1\ c_4\ c_3)$, and the same argument applies to
the rectangle from $c_4$ to $c_0$, whose possible marking columns
belong to both $[c_4,c_0)$ and $[c_2,c_3)$.
Cyclic reindexing proves the same statement for all five sources.
A transposition splitting a five-cycle into lengths three and two
has exactly five choices, so this list is exhaustive.
Thus either coefficient is
\[
 \sum_{i=0}^4h_{32}(c_i,c_{i+1},c_{i+2},c_{i+3},c_{i+4}).
 \tag{M6.12}
\]

Changing the origin cyclically permutes the sixteen half-turn
chambers in the following four orbits.
The words in each row are listed in the order of successive origins.
\[
\begin{array}{c|l}
 \text{Representative}&\text{Successive words}\\ \hline
 0000&0000,0001,0011,0111,1111\\
 0010&0010,0101,1011,1110,1000\\
 0100&0100,1001,0110,1101,1100\\
 1010&1010,1010,1010,1010,1010
\end{array}
\]
Here is the explicit rule producing the table.
Adjoin the initial bit $b_0=0$ to a word $b_1b_2b_3b_4$.
If there is another zero bit, let $j$ be the least positive index
with $b_j=0$. Otherwise put $j=1$.
The vertex at projected position $j$ is the next vertex in the
original circular order. After changing the origin to this vertex,
the new word is
\[
 (b_{j+1}+b_j,\ldots,b_4+b_j,
       b_0+b_j+1,\ldots,b_{j-1}+b_j+1).
\]
The additions here are modulo two. Projected positions before the
new origin wrap once through a half-circle, which accounts for the
extra one. This rule gives the five words in each displayed row.
The table lists all sixteen words, with no omissions or repetitions
between rows.
Using (M6.8), the four corresponding sums in (M6.12) are respectively
\[
\begin{aligned}
 &0+t_1+t_1+0+0=0,\\
 &(1+t_4)+(1+t_2)+0+(t_1+t_2)+(t_1+t_4)=0,\\
 &(1+t_1)+(1+t_3)+0+(t_1+t_3)+0=0,\\
 &(t_1+t_2)+(t_2+t_3)+(t_3+t_4)+(1+t_4)+t_1=1.
\end{aligned}
\]
These are direct substitutions of all entries of (H32.5), so they
prove that the coefficient is the indicator of chamber $1010$.

It remains to identify this chamber geometrically.
Put $L=N/2$ and use $c_0=0$ as origin.
If (M6.2) holds, then
$0<c_1<c_2<L<c_3<c_4<N$.
The other two-step inequalities give
\[
       0<c_3-L<c_1<c_4-L<c_2<L.
\]
Conversely these inequalities imply that all five positive two-step
distances are less than $L$, and hence at most $p$ because they are
integers. This projected order has half-circle bits $1010$.
Therefore the indicator just obtained is exactly (M6.2).
\end{proof}

\subsection{The two-triangle coefficient}

\begin{lemma}[New lemma: the $(3,3)$ coefficient]\label{J21:lem:m6-three-three}
Let $\tau$ consist of two disjoint three-cycles and have
$M(x_\tau)=-6$. Let $T_1,T_2$ be their two unordered vertex triples.
Then
\[
                    [x_\tau]d_0H_{32}=B(T_1)B(T_2).
 \tag{M6.13}
\]
Here $B(T)$ means (M6.1) after placing the triple in positive circular
order. The formula holds for every orientation of the two cycles
having the stated Maslov value.
\end{lemma}
\begin{proof}
Write the six vertices in positive circular order as $c_0,\ldots,c_5$.
For brevity, write $S(i,j)=S(c_i,c_j)$, and similarly suppress the
$c$ in $F,G,h_{32}$.
Up to a cyclic change of origin and interchange of the two triples,
the underlying partition has one of three forms:
\[
 \{0,1,2\}\mid\{3,4,5\},\qquad
 \{0,1,3\}\mid\{2,4,5\},\qquad
 \{0,2,4\}\mid\{1,3,5\}.
 \tag{M6.14}
\]
For a direct enumeration, the ten unordered partitions into triples
consist of the three cyclic placements of the first type, the six
placements of the second type, and the last type.
Eligibility of a deletion depends only on this partition, by the
intrinsic-order criterion and the insertion-index discussion above.
It does not depend on the orientations of the three-cycles.

For the first partition in (M6.14), every deletion leaves the two
transposition endpoints in the same complementary arc of the other
triangle. Therefore there are no sources and the coefficient is zero.
Both triangles cannot be balanced: balance would require both
$d(c_0,c_2)>p$ and $d(c_3,c_5)>p$, but these disjoint arcs cannot
have total length greater than $N$. Thus (M6.13) holds in this case.

For the second partition, the four eligible intrinsic source orders
are
\[
 (0,2,3,4,5),\quad (1,2,3,4,5),\quad
 (2,3,5,0,1),\quad (2,3,4,0,1).
\]
Expanding their coefficients by (M6.5) and pairing each four-tuple
with its one-step cyclic rotation gives the rectangular sum
\[
                    \sum_{i\in\{0,1\}}\sum_{j\in\{4,5\}}G(i,2,3,j).
\]
In (M6.7) the first term is independent of $j$, and the third term
is independent of $i$. Each of those terms occurs twice. The sum is
therefore
\[
                 (S(0,3)+S(1,3))(S(2,4)+S(2,5)).
 \tag{M6.15}
\]
Because $0,1,2,3,4,5$ are circularly ordered, the first factor equals
one exactly when $S(0,3)=0$ and $S(1,3)=1$. The second equals one
exactly when $S(2,4)=1$ and $S(2,5)=0$.
Under these four conditions $S(0,1)=1$ and $S(4,5)=1$.
Indeed, if $d(c_0,c_1)>p$, it and the disjoint arc from $c_2$ to
$c_5$ would have total length at least $2p+2>N$.
If $d(c_4,c_5)>p$, the same contradiction follows with the arc
from $c_0$ to $c_3$.
Consequently (M6.15) is one exactly when both triangles
$(0,1,3)$ and $(2,4,5)$ are balanced. Conversely their balance
implies the four displayed conditions immediately from (M6.1).
This proves (M6.13) for the second partition.

For the last partition in (M6.14), every deletion is eligible.
The six intrinsic orders are
\[
              (i,i+1,i+2,i+3,i+5),\qquad i\pmod6.
\]
By (M6.5), their sum is
\[
                       \sum_{i=0}^5G(i,i+1,i+2,i+3).
 \tag{M6.16}
\]
To see this equality, the first four-point term is the consecutive
four-tuple starting at $i$, and the second is the one-step rotation
of the consecutive four-tuple starting at $i+5$.

Let $\delta_i=d(c_i,c_{i+1})$, with indices modulo six, and put
\[
 A_i=[\delta_i\leq p],\qquad B_i=[\delta_i+\delta_{i+1}\leq p].
\]
Formula (M6.7) gives
\[
 G(i,i+1,i+2,i+3)
 =A_iA_{i+1}\overline{\delta_i}
   +B_i(1+B_{i+1})
   +(1+A_{i+2})(\overline{\delta_{i+1}}+
                              \overline{\delta_{i+2}}).
\]
In the sum (M6.16), the coefficient of $\overline{\delta_i}$ is
$A_iA_{i+1}+A_i+A_{i+1}$.
Two successive positive gaps cannot both exceed $p$, since their
sum would exceed $N$. Thus this coefficient is one.
Since $\sum_i\delta_i=N$ is odd, (M6.16) becomes
\[
                         1+\sum_iB_i+\sum_iB_iB_{i+1}.
 \tag{M6.17}
\]
Two disjoint two-gap arcs cannot both have length greater than $p$.
Their total would again be at least $2p+2>N$.
Therefore the set of indices with $B_i=0$ is empty, a singleton,
or two adjacent indices. A cycle on six indices has no three
pairwise adjacent distinct indices, so these are all possibilities.
If that set is empty, (M6.17) is $1+6+6=1$ in $\mathbb F _2$.
For one zero it is $1+5+4=0$.
For two adjacent zeros it is $1+4+3=0$.
Hence (M6.17) equals $\prod_iB_i$.
The latter product is
$B(c_0,c_2,c_4)B(c_1,c_3,c_5)$, since the six positive two-step
gaps are the successive gaps in those two triangles.
This proves (M6.13) in the final case.
\end{proof}

\subsection{The resulting native Alexander bound}

\begin{theorem}[New theorem: exact Maslov minus six support]
\label{J21:thm:h32-m6-support}
The chain $d_0H_{32}$ is the sum of the following distinct states,
each with coefficient one:
\begin{enumerate}
 \item The two star states on every balanced pentagon.
 \item Every state consisting of two balanced three-cycles and having
       native Maslov value minus six.
\end{enumerate}
Every one of these states satisfies
\[
                              A(x)\leq g-2p.
 \tag{M6.18}
\]
Consequently, for every integer $m>0$, every term of
$U^{-3}d_0H_{32}$ satisfies the exact second-vertex bound
\[
                 A(x)-m(-3)\leq g+3m-2p.
 \tag{M6.19}
\]
\end{theorem}
\begin{proof}
The exact coefficient statement follows from
Lemmas~\ref{J21:lem:m6-types}, \ref{J21:lem:m6-four-two},
\ref{J21:lem:m6-pentagon}, and \ref{J21:lem:m6-three-three}.
It remains to compute the native Alexander values, including the
orientation signs.
The attached declarations
\decl{alexanderTwo_reflectedTorus_eq_sum_displacement}
and
\decl{alexander_reflectedTorusAnchor}
give
\[
 A(x_\sigma)=g-\frac1N\sum_c W(d(c,\sigma(c))),
 \qquad
 W(d)=pd-N\max\{0,d-p-1\},
 \tag{M6.20}
\]
where fixed columns contribute $W(0)=0$.
In particular, if $1\leq d\leq p$, then
\[
                    W(d)=pd,\qquad W(N-d)=(p+1)d.
 \tag{M6.21}
\]

For a balanced triangle, its three positive successive gaps are at
most $p$ and sum to $N$.
The positive cycle therefore has total displacement weight $pN$,
and its inverse has total weight $(p+1)N$ by (M6.21).
Their Alexander drops from $g$ are respectively $p$ and $p+1$.
For two disjoint balanced triangles the weights add.
Their total drop is $2p$, $2p+1$, or $2p+2$, according to their
orientations. This proves (M6.18) for the second family.

For a balanced pentagon, its five positive two-step distances are
at most $p$ and sum to $2N$.
The star $c_i\mapsto c_{i+2}$ has total weight $2pN$, and the
inverse star $c_i\mapsto c_{i+3}$ has total weight $2(p+1)N$.
Their Alexander drops are exactly $2p$ and $2p+2$, respectively.
This proves (M6.18) for the first family.
Finally (M6.19) is obtained by adding $3m$ to (M6.18).
No restriction stronger than $p\geq3$ and $m>0$ has been used.
\end{proof}

\paragraph{Exceptional cases and scope.}
The smallest allowed grid is $N=7$, corresponding to $p=3$.
All tuple lengths in this section are at most six, so no proof
requires a nonexistent column in this case.
For example, the pentagon $(0,1,3,4,5)$ has successive gaps
$(1,2,1,1,2)$ and is balanced. Its forward and inverse stars have
Maslov value minus six and Alexander values $-3$ and $-5$,
respectively, since $g=3$.
Thus the forward term $U^{-3}x$ attains the bound in (M6.19).
This checks the exponent and chirality conventions on the smallest
native grid. The proof above applies to every larger grid without
changing the coefficient tables.
Theorem~\ref{J21:thm:h32-m6-support} proves the complete $d_0H_{32}$
part of the support calculation. The $d_2H_{32}$ cancellation was proved in
Corollary~\ref{J21:cor:h4-h32-lower-identity}. The $d_1H_{32}$ cancellation
is proved in the next two sections.



\section{Two diagonal identities and one five-cycle cancellation}
\label{J21:sec:h4-m4-A}

Fix $p\geq3$ and put $N=2p+1$.
All coefficients in this section belong to $\mathbb F _2$.
We use the functions $h_p$, $C_p^0$, and $T_p^0$ from
(H4D.1)--(H4D.3), the four-cycle coefficient $f_4$ from
(H4D.5)--(H4D.6), and the four-point function $F$ from (M6.3).
The primitive identity
\[
 h_{32}(u,a,v,b,c)=F(u,a,v,b)+F(u,a,v,c)
 \tag{M4A.1}
\]
is the proved new Lemma~\ref{J21:lem:m6-primitives}.

\subsection{The two diagonal identities}

\begin{lemma}[New lemma: first diagonal identity]
\label{J21:lem:h4-diagonal-primitive-first}
For every pair of integers $2\leq A<B\leq2p$,
\[
 T_p^0(A,B)+h_p(A)+h_p(B)+F(0,1,A,B)
       =1+\overline A+[A=p+1].
 \tag{M4A.2}
\]
\end{lemma}
\begin{proof}
Put $\eta=\overline p$.
The following six disjoint cases are the cases defining $T_p^0$.
The second column is the half-turn word for $(0,1,A,B)$.
The next two columns follow from the explicit definitions of $F$
and $h_p$.
\[
\begin{array}{c|c|c|c}
 \text{condition}&w&F(0,1,A,B)&h_p(A)+h_p(B)\\ \hline
 B\leq p&000&0&\overline A+\overline B\\
 A\leq p,\ B\geq p+2,\ B-A\geq p+1
       &001&1+\overline B&1+\overline A\\
 A\leq p,\ B\geq p+2,\ B-A\leq p
       &010&0&1+\overline A\\
 B=p+1&100&1+\overline A&1+\overline A+\eta\\
 A=p+1&101&0&\eta\\
 A\geq p+2&011&0&0
\end{array}
 \tag{M4A.3}
\]
To obtain the half-turn words, the projected position of the column
$1$ is $1$, the projected position of $p+1$ is $\tfrac12$, and the
projected position of every column at least $p+2$ is greater than
$1$.
When $A\leq p$ and $B\geq p+2$, the projected position of $B$ is
greater than $A$ precisely when $B-A\geq p+1$.
These comparisons give every word in the table, including the
boundary cases.

Add the value of $T_p^0(A,B)$ in (H4D.3) to the last two columns.
The result is $1+\overline A$ in the first, second, third, fourth,
and sixth rows.
In the fifth row it is $1+\eta$.
Here $A=p+1$, so
$1+\overline A+[A=p+1]=1+(\eta+1)+1=1+\eta$.
This proves (M4A.2) in every case.
\end{proof}

\begin{lemma}[New lemma: second diagonal identity]
\label{J21:lem:h4-diagonal-primitive-second}
Let $r,z$ be integers with
\[
                       1\leq r,\qquad r+2\leq z\leq2p.
\]
Then
\[
 T_p^0(z-r,N-r)+F(0,r,r+1,z)
                   =\overline r+\overline p\,[r=p].
 \tag{M4A.4}
\]
\end{lemma}
\begin{proof}
The arguments of $T_p^0$ belong to its domain, because
$2\leq z-r<N-r\leq2p$.
The following table gives every case of the half-turn word and the
two summands. Every parity is measured from the first column zero.
\[
\begin{array}{c|c|c|c}
 \text{condition}&w&F(0,r,r+1,z)&T_p^0(z-r,N-r)\\ \hline
 r\leq p-1,\ z\leq p
       &000&0&\overline r\\
 r\leq p-1,\ p+1\leq z\leq p+r
       &100&\overline r&0\\
 r\leq p-1,\ z=p+r+1
       &010&1+\overline r&1\\
 r\leq p-1,\ z\geq p+r+2
       &001&1+\overline z&1+\overline z+\overline r\\
 r=p&110&1+\overline z&1+\overline z\\
 r\geq p+1&111&0&\overline r
\end{array}
 \tag{M4A.5}
\]
In the first four rows, the columns $r,r+1$ lie in the first
half-circle.
For $z\geq p+1$, its projected position lies before $r$,
strictly between $r$ and $r+1$, or after $r+1$ according as
$z\leq p+r$, $z=p+r+1$, or $z\geq p+r+2$.
This gives the indicated words.
For $r=p$, the projected order is $r+1,z,r$, which gives $110$.
For $r\geq p+1$, all three columns are in the second half-circle,
and their projected order is unchanged.
The entries for $F$ now follow from (M6.3).

For the entries in the last column, put $A=z-r$ and $B=N-r$.
In the first row, $A\leq p$ and $B-A=N-z\geq p+1$, so
$T_p^0=1+\overline B=\overline r$.
In the second row, $A\leq p$, $B\geq p+2$, and $B-A\leq p$.
In the third row, $A=p+1$.
In the fourth row, $A\geq p+2$.
In the fifth row, $B=p+1$, so
$T_p^0=1+\overline{z-p}+\overline p=1+\overline z$.
In the last row, $B\leq p$, so again
$T_p^0=1+\overline B=\overline r$.
These deductions verify the entire last column.

Adding the two summands gives $\overline r$ in every row except
$r=p$, where it gives zero.
This is exactly the right side of (M4A.4).
\end{proof}

\subsection{A useful expression for a nonconvex four-cycle}

For a positive circularly ordered four-tuple $(a,b,d,e)$, put
\[
                   G'(a,b,d,e)=f_4((a\,d\,e\,b)).
\]
Write $E(x,y)=[d(x,y)=1]$ for a positive ambient adjacent pair.
Define the two auxiliary functions on positive distances by
\[
 \psi_1(s)=1+\overline s+[s=p+1],\qquad
 \psi_2(s)=\overline s+\overline p\,[s=p].
 \tag{M4A.6}
\]

\begin{lemma}[New lemma: the nonconvex four-cycle formula]
\label{J21:lem:h4-A-four-coefficient}
For every positive circularly ordered $(a,b,d,e)$,
\begin{align*}
 G'(a,b,d,e)={}&
 E(b,d)\bigl(F(b,d,e,a)+\psi_1(d(b,e))\bigr)\\
 &+E(d,e)\bigl(h_p(d(d,a))+h_p(d(d,b))\bigr)\\
 &+E(e,a)\bigl(F(d,e,a,b)+\psi_2(d(d,e))\bigr).
 \tag{M4A.7}
\end{align*}
Each expression multiplied by an adjacent-pair indicator is used
only when that indicator is one.
Its arguments then belong to the stated domains of the corresponding
functions.
\end{lemma}
\begin{proof}
There are four possible positive ambient adjacent pairs in this
support, namely $(a,b)$, $(b,d)$, $(d,e)$, and $(e,a)$.
Normalize each pair in turn to labels $(0,1)$ and read its cycle word
from $(a\,d\,e\,b)$.
The four resulting words are respectively
\[
                        0231,\qquad0312,\qquad0132,\qquad0213.
\]
By (H4D.5), their densities are respectively zero,
$T_p^0+C_p^0$, $C_p^0$, and $T_p^0$.

If $E(b,d)=1$, translate $b$ to zero.
The four columns become $(0,1,A,B)$ with
$A=d(b,e)$ and $B=d(b,a)$.
Lemma~\ref{J21:lem:h4-diagonal-primitive-first} expresses its density as
$F(b,d,e,a)+\psi_1(d(b,e))$.
If $E(d,e)=1$, its two remaining coordinates relative to $d$
are $d(d,a),d(d,b)$, so its density is the second line of (M4A.7).
If $E(e,a)=1$, use $d$ as origin and put
$r=d(d,e)$ and $z=d(d,b)$.
The four columns are $(0,r,r+1,z)$.
Relative to the adjacent pair $(e,a)$, the other two coordinates
are $z-r$ and $N-r$.
Lemma~\ref{J21:lem:h4-diagonal-primitive-second} gives its density as
$F(d,e,a,b)+\psi_2(d(d,e))$.
Summing the densities over all adjacent pairs gives (M4A.7) by the
definition (H4D.6).
\end{proof}

\subsection{The five-cycle of circular type 01243}

\begin{lemma}[New lemma: cancellation for the first nonconvex type]
\label{J21:lem:h4-h32-A-cancellation}
Let $A,B,C,D,E$ be distinct columns in positive circular order and
put
\[
                         \tau=(A\,B\,C\,E\,D).
\]
Then the coefficient of $x_\tau$ in $d_0H_4+d_1H_{32}$ is zero.
\end{lemma}
\begin{proof}
The new five-cycle grading table,
Lemma~\ref{J21:lem:torus-five-cycle-table}, gives $M(x_\tau)=-4$.
The proved new minimal-cycle deletion formula,
Lemma~\ref{J21:lem:torus-minimal-cycle-coefficient}, says that every
four-cycle obtained by deleting one vertex of $\tau$ contributes
through exactly one empty zero-$O$ reinsertion rectangle.
Thus the coefficient of $x_\tau$ in $d_0H_4$ is
\begin{align*}
 &f_4((B\,C\,E\,D))+f_4((A\,C\,E\,D))
       +f_4((A\,B\,E\,D))\\
 &\hspace{20mm}
       +C(\{A,B,C,E\})+C(\{A,B,C,D\}),
 \tag{M4A.8}
\end{align*}
where $C(S)$ is the forward convex coefficient on a four-element
support $S$.
The common inverse-sum lemma replaces each of the first three terms
by the corresponding $G'$ term plus its convex coefficient.
The five convex coefficients then sum to zero by
Lemma~\ref{J21:lem:h4-convex-primitive}.
Consequently (M4A.8) equals
\[
                  G'(B,C,D,E)+G'(A,C,D,E)+G'(A,B,D,E).
 \tag{M4A.9}
\]

Expand (M4A.9) using (M4A.7).
The indicator $E(B,D)$ is zero because $C$ lies strictly between
$B$ and $D$.
Likewise $E(E,B)=0$ because $A$ lies strictly between $E$ and $B$.
The two terms with indicator $E(C,D)$ have the same
$\psi_1(d(C,E))$, which cancels in their sum.
Their remaining sum is
\[
 E(C,D)\bigl(F(C,D,E,B)+F(C,D,E,A)\bigr)
                 =E(C,D)h_{32}(C,D,E,A,B)
\]
by (M4A.1).
The two terms with indicator $E(E,A)$ have the same
$\psi_2(d(D,E))$, which also cancels.
They leave
\[
 E(E,A)\bigl(F(D,E,A,C)+F(D,E,A,B)\bigr)
                 =E(E,A)h_{32}(D,E,A,B,C).
\]
Finally, each of $h_p(d(D,A))$, $h_p(d(D,B))$, and
$h_p(d(D,C))$ occurs exactly twice with indicator $E(D,E)$.
Those six terms cancel.
We have proved
\[
 [x_\tau](d_0H_4)
   =E(C,D)h_{32}(C,D,E,A,B)
       +E(E,A)h_{32}(D,E,A,B,C).
 \tag{M4A.10}
\]

We now compute the one-$O$ rectangles from $H_{32}$ to this target.
Cutting the five-cycle $\tau$ into a three-cycle and a transposition
has exactly five possibilities.
The following table lists the source permutation and its Maslov
value, obtained from the proved new $(3,2)$ interlacing criterion.
\[
\begin{array}{c|c|c}
 \text{swapped columns}&\text{source}&M\\ \hline
 A,C&(A\,E\,D)(B\,C)&-3\\
 A,E&(A\,D)(B\,C\,E)&-5\\
 B,D&(A\,B)(C\,E\,D)&-3\\
 B,E&(A\,B\,D)(C\,E)&-5\\
 C,D&(A\,B\,C)(D\,E)&-3
\end{array}
 \tag{M4A.11}
\]
Only the two rows of Maslov value minus five belong to $H_{32}$.
For the source $(A\,D)(B\,C\,E)$, the swap of columns $A,E$
has $O$ counts $d(B,D)$ and $d(E,A)$.
The first is at least two because $C$ lies strictly between $B,D$.
The second is one exactly when $E(E,A)=1$.
In that event its column interval has width one, so the rectangle
is state-empty.
The intrinsic order of this source is $(D,E,A,B,C)$.
Its contribution is therefore the second term of (M4A.10).

For the source $(A\,B\,D)(C\,E)$, the swap of columns $B,E$
has $O$ counts $d(C,D)$ and $d(E,B)$.
The second is at least two because $A$ lies strictly between $E,B$.
The first is one exactly when $E(C,D)=1$.
Its anchor-row interval then has height one, so the rectangle is
state-empty.
The intrinsic order of this source is $(C,D,E,A,B)$.
Its contribution is the first term of (M4A.10).
The native relative-rectangle interval formula gives these counts
directly, and the five cuts in (M4A.11) exhaust every source of
cycle type $(3,2)$ reaching the target.
Thus (M4A.10) is also the coefficient of $x_\tau$ in $d_1H_{32}$.
The two equal coefficients cancel in $\mathbb F _2$, proving the
lemma.
\end{proof}

\paragraph{Smallest cases.}
The diagonal identities include $A=p+1$, $B=p+1$, and $r=p$
as explicit separate cases.
They hold at $p=3$ with any empty open ranges omitted.
The four-cycle and five-cycle arguments permit every collection of
adjacent pairs, including consecutive runs of three or more support
columns. Such pairs are counted separately by (H4D.6), and the
cancellations above are equalities of these finite sums.


\section{The remaining five-cycle coefficients}
\label{J21:sec:m4-B}

Throughout this section $p\geq3$, $N=2p+1$, and all sums of coefficients
are taken in $\mathbb F _2$. Put
\[
                    R_4=d_0H_4+d_1H_{32}.
\]
The preceding boundary-type lemmas and the cancellation of the
two-transposition targets show that $R_4$ is supported on single
five-cycles of Maslov value $-4$. We determine these coefficients exactly.

\subsection{Inversion and the circular five-cycles}

\begin{lemma}[New lemma: inversion invariance of the residual]
\label{J21:lem:m4-residual-inverse}
For every five-cycle $\tau$ of Maslov value $-4$, its coefficient in
$R_4$ equals that of $\tau^{-1}$. The positive and negative circular
five-cycles both have coefficient zero.
\end{lemma}
\begin{proof}
By the five-cycle deletion formula, the coefficient in $d_0H_4$ is the
sum of the five four-cycle coefficients obtained by deleting a vertex.
The common inverse-sum lemma therefore makes the sum of the coefficients
of $\tau,\tau^{-1}$ equal to
\[
                     \sum_{q\in T} C(T\setminus\{q\}),
\]
where $T$ is their common support. This sum vanishes by
Lemma~\ref{J21:lem:h4-convex-primitive}. The chain $H_{32}$ is invariant
under inversion, which commutes with every $d_j$ by the native rectangle
bijection. Thus its $d_1$ coefficients are also invariant under inversion.
This proves the first assertion.

For the positive circular five-cycle, the $d_0H_4$ coefficient is the
same vanishing sum of $C$ coefficients. A transposition splitting this
cycle into cycles of lengths three and two produces a pair whose two
endpoints are adjacent in the circular order of the five moved columns.
Such a source has Maslov value $-3$, by the three-cycle--transposition
classification, and is absent from $H_{32}$. Hence its $d_1H_{32}$
coefficient is zero. Inversion gives the assertion for the negative
circular cycle.
\end{proof}

\subsection{An exact sum over the adjacent pairs}

For this subsection, fix five positively ordered columns and assume that
one selected adjacent pair is $0,1$. Write the remaining three columns
as $A,B,C$, so
\[
                       2\leq A<B<C\leq2p.
\]
These coordinates are obtained by cyclic translation from any selected
pair $u,u+1$. Abbreviate
\[
 h(d)=h_p(d),\qquad T(A,B)=T_p^0(A,B),\qquad
 C^0(A,B)=h(A)+h(B).
\]
A local contribution means the summands indexed by this one adjacent
pair in the definition of $H_4$, together with the one-$O$ rectangles
from $H_{32}$ whose unique covered $O$ square is in this gap.
The total coefficient is the sum of these local contributions over all
adjacent pairs in the five-element support.

\begin{lemma}[New lemma: two identities for the local contribution]
\label{J21:lem:m4-local-two-identities}
For every $2\leq A<B<C\leq2p$ one has
\begin{align}
 h_{32}(C,0,1,A,B)&=T(A,C)+T(B,C),
                                  \label{J21:eq:m4-local-first}\\
 h_{32}(0,1,A,B,C)&=T(A,B)+T(A,C)+h(B)+h(C).
                                  \label{J21:eq:m4-local-second}
\end{align}
\end{lemma}
\begin{proof}
The first diagonal-primitive identity proved above states
\[
             T(A,Z)+h(A)+h(Z)+F(0,1,A,Z)
                   =1+\overline A+[A=p+1].
\]
Use this identity with $Z=B$ and $Z=C$, and add the two equations.
The terms depending only on $A$ cancel. The primitive formula for
$h_{32}$ then gives \eqref{J21:eq:m4-local-second}.

For \eqref{J21:eq:m4-local-first}, translate $C$ to zero in the second
diagonal-primitive identity. Its parameters are
$r=N-C$ and $z=r+A$. They satisfy $1\leq r$, $r+1<z<N$.
That identity gives
\[
 T(A,C)+F(C,0,1,A)
             =\overline{N-C}+\overline p\,[N-C=p],
\]
whose right side is independent of $A$. Add the corresponding equation
with $A$ replaced by $B$ and apply the primitive formula for
$h_{32}(C,0,1,A,B)$. This proves \eqref{J21:eq:m4-local-first}.
\end{proof}

\begin{lemma}[New lemma: the five local coefficients]
\label{J21:lem:m4-B-unit-table}
Let the five columns $0,1,A,B,C$ be labelled $0,1,2,3,4$ in that
order. For the following five cycle words, the local contribution of
the selected adjacent pair is as stated.
\[
\begin{array}{c|c}
 \text{cycle word}&\text{local contribution}\\\hline
 01342&0\\
 02314&0\\
 03412&0\\
 01423&0\\
 03124&[\{0,1,A,B,C\}\text{ is a balanced pentagon}]
\end{array}
\]
Here a pentagon is balanced when each sum of two successive positive
circular gaps is at most $p$.
\end{lemma}
\begin{proof}
The following table records the complete contributing terms. Its
middle column lists the three adjacent-pair densities obtained by
deleting one of labels $2,3,4$. Its last column gives the intrinsic
orders of the sources in $H_{32}$ whose restoring rectangle has its
one $O$ marking in the gap $0,1$. An empty entry is a zero sum.
\[
\begin{array}{c|c|l}
 \text{word}&\text{four-cycle densities}&
             \text{three-cycle--transposition coefficients}\\\hline
 01342&C^0(B,C)+C^0(A,C)+C^0(A,B)&0\\
 02314&T(B,C)+T(A,C)&h_{32}(C,0,1,A,B)\\
 03412&T(A,C)+T(A,B)+h(B)+h(C)&h_{32}(0,1,A,B,C)\\
 01423&C^0(B,C)+C^0(A,C)+C^0(A,B)&0\\
 03124&T(B,C)+T(A,B)+h(B)+h(C)&
       \begin{aligned}
       &h_{32}(1,A,B,C,0)\\
       &{}+h_{32}(A,B,C,0,1)\\
       &{}+h_{32}(B,C,0,1,A)
       \end{aligned}
\end{array}
\tag{M4B.1}
\]
We justify that this table is complete from the native rectangles.
Deleting either endpoint of the selected pair gives no density for
that pair. Deleting any of the other three columns gives exactly the
middle-column words in the six-word density definition.

For the last column, a source splitting a five-cycle into lengths three
and two is obtained by one of the five cuts separated by two positions
in the cycle word. The relative rectangle formula says that its marking
columns are the intersection of two half-open intervals whose endpoints
are moved columns. If this intersection contains just the one marking
in the gap $0,1$, it must be exactly $[0,1)$. Any interval containing
two moved-column gaps has length at least two. Applying this criterion
to the five cuts gives precisely the following list. The source cycles
and the ordered rectangle columns are written in the labels $0,\ldots,4$.
\[
\begin{array}{c|c|c|c}
 \text{target word}&\text{source cycles}&
 \text{ordered columns}&[\sigma(j),\sigma(i))\\\hline
 02314&(0\ 2\ 3)(1\ 4)&(4,3)&[0,1)\\
 03412&(0\ 2)(1\ 3\ 4)&(0,1)&[3,2)\\
 03124&(0\ 2\ 4)(1\ 3)&(0,1)&[3,2)\\
 03124&(0\ 3\ 1)(2\ 4)&(4,1)&[0,2)\\
 03124&(0\ 3)(1\ 2\ 4)&(4,3)&[0,1)
\end{array}
\tag{M4B.2}
\]
In the first and fifth rows the row interval has height one. In the
second and third rows the column interval has width one. In the fourth
row, the intersection $[4,1)\cap[0,2)$ is $[0,1)$, and the only
moved point in its column interior is in column zero. Its image is
column three, outside the required row interval. Fixed points cannot
lie in its interior because the open intersection contains no integer
column. Thus all five listed rectangles are empty and have one $O$
marking. Their source permutations all have the interlaced $(3,2)$
type of Maslov value $-5$. The other cuts either have Maslov value
$-3$ or have marking intersection different from $[0,1)$, so none
contributes to this selected pair. This proves (M4B.1) without any
restriction on the other circular gaps.


The following table also records all excluded cuts of Maslov value $-5$.
Let $c_0,c_1,c_2,c_3,c_4$ be the five support columns in positive order,
and let $I_s=[c_s,c_{s+1})$, with indices modulo five.
For a cut $i<j$, the last two entries list the sets of indices $s$ whose
intervals $I_s$ form the marking-column intersections for the ordered
rectangles $(i,j)$ and $(j,i)$ respectively. Thus an entry $\{0\}$
is precisely the selected marking square. An entry containing two
indices contains at least two squares, because all five support gaps
have positive integer length.
\[
\begin{array}{c|c|c|c|c}
 \text{target}&(i,j)&\text{source}&(i,j)\text{ gaps}&(j,i)\text{ gaps}\\\hline
 01342&(0,3)&(0\ 4\ 2)(1\ 3)&\{1,2\}&\{4\}\\
 01342&(0,4)&(0\ 2)(1\ 3\ 4)&\{1\}&\{4\}\\
 01342&(2,3)&(0\ 1\ 3)(2\ 4)&\{2\}&\{4\}\\
 02314&(1,2)&(0\ 2\ 4)(1\ 3)&\{1\}&\{3\}\\
 02314&(2,4)&(0\ 2)(1\ 4\ 3)&\{3\}&\{0,1\}\\
 02314&(3,4)&(0\ 2\ 3)(1\ 4)&\{3\}&\{0\}\\
 03412&(0,1)&(0\ 2)(1\ 3\ 4)&\{0\}&\{2\}\\
 03412&(1,3)&(0\ 3\ 2)(1\ 4)&\{2\}&\{0,4\}\\
 03412&(2,3)&(0\ 3)(1\ 2\ 4)&\{2\}&\{4\}\\
 01423&(0,2)&(0\ 3)(1\ 4\ 2)&\{1\}&\{3,4\}\\
 01423&(0,4)&(0\ 2\ 3)(1\ 4)&\{1\}&\{4\}\\
 01423&(1,2)&(0\ 1\ 3)(2\ 4)&\{1\}&\{3\}\\
 03124&(0,1)&(0\ 2\ 4)(1\ 3)&\{0\}&\{2\}\\
 03124&(1,4)&(0\ 3\ 1)(2\ 4)&\{2,3\}&\{0\}\\
 03124&(3,4)&(0\ 3)(1\ 2\ 4)&\{3\}&\{0\}
\end{array}
\tag{M4B.2a}
\]
Each row is obtained by substituting its displayed permutation in
$[i,j)\cap[\sigma(j),\sigma(i))$ and in the complementary ordered
rectangle formula. These are all three interlaced cuts of each target.
The other two cuts have Maslov value $-3$ by
Lemma~\ref{J21:lem:torus-three-two-type}. Thus the five entries $\{0\}$
are exactly the five rectangles in (M4B.2).

The first and fourth rows of (M4B.1) vanish because each $h$ value
appears twice. The second and third rows vanish by
Lemma~\ref{J21:lem:m4-local-two-identities}.
For the final row, those same identities give
\begin{align*}
 T(B,C)+T(A,B)+h(B)+h(C)
   ={}&h_{32}(C,0,1,A,B)\\
      &{}+h_{32}(0,1,A,B,C).
\end{align*}
Its total local contribution is therefore the sum of $h_{32}$ over
all five cyclic rotations of $(0,1,A,B,C)$.
The pentagon coefficient lemma, Lemma~\ref{J21:lem:m6-pentagon}, proves
that this sum is exactly the balanced-pentagon indicator. This proves
the displayed five-row conclusion.
\end{proof}

\subsection{The exact residual and its Alexander grading}

\begin{theorem}[New theorem: the Maslov minus four residual]
\label{J21:thm:m4-exact}
Define, for $u\in\mathbb Z/N\mathbb Z$ and integers
$2\leq b\leq p$, $p+2\leq c\leq p+b$, the five-cycle
\[
             \tau_{u,b,c}=(u\,u+p+1\,u+1\,u+b\,u+c).
\]
All its five entries are distinct. Then
\[
 R_4=\sum_{u\in\mathbb Z/N\mathbb Z}
          \ \sum_{b=2}^{p}\ \sum_{c=p+2}^{p+b}
                   \bigl(x_{\tau_{u,b,c}}
                         +x_{\tau_{u,b,c}^{-1}}\bigr).
 \tag{M4B.3}
\]
Every displayed state has Maslov value $-4$ and Alexander value $g-2p$.
\end{theorem}
\begin{proof}
The five-cycle orbit table shows that the cycles of Maslov value $-4$
are the two circular cycles and the circular relabellings of
$01243$, $01342$, and their inverses. The circular-cycle coefficients
vanish by Lemma~\ref{J21:lem:m4-residual-inverse}. Lemma~\ref{J21:lem:h4-h32-A-cancellation} shows that all its coefficients vanish. Inversion
gives the same conclusion for its inverse type.

Consider the $01342$ type and its five circular relabellings. For each
adjacent pair in its support, normalize that pair to $0,1$. The five
possible relabelled words are exactly those in
Lemma~\ref{J21:lem:m4-B-unit-table}. Thus only word $03124$ can contribute,
and it contributes one precisely when the support is balanced.
For a fixed cycle of this type, at most one adjacent pair gives word
$03124$. The circular-label orbit has size five, so two different
choices of first vertex cannot yield that same word.

In the normalization $0,1,A,B,C$, balance forces $B=p+1$.
Indeed the two-step inequalities from zero and from one give
$A\leq p$ and $B-1\leq p$. The two-step inequality along the
complementary arc from $B$ to zero gives $N-B\leq p$.
Together these last two inequalities imply $B=p+1$.
The remaining balance inequalities are exactly
\[
                 2\leq A\leq p,\qquad
                 p+2\leq C\leq p+A.
\]
To check the converse, these inequalities give the five two-step
distances
$A$, $p$, $C-A$, $p$, and $N-C+1$, all at most $p$.
Thus they are also sufficient for balance.
The contributing cycle word $03124$ is therefore
$(0,p+1,1,A,C)$. Translating back and using inversion invariance gives
exactly (M4B.3).

The distinguished adjacent pair also proves that no state is repeated
within the sum in (M4B.3). It determines $u$, and the increasing support
then determines $b,c$. The five inverse words of the words in
Lemma~\ref{J21:lem:m4-B-unit-table} are respectively
$02431$, $04132$, $02143$, $03241$, and $04213$.
None belongs to the original five-word list. Hence the inverse states
are distinct from the forward states as well. Every state in (M4B.3)
therefore has coefficient one.

The distinctness of its five entries follows from
$0<1<b<p+1<c<N$. Its compressed cycle word is $03124$, which has
Maslov value $-4$ by the five-cycle orbit table.
The five positive displacements of $\tau_{u,b,c}$ in cycle order are
\[
                 p+1, p+1, b-1, c-b, N-c.
\]
The last three belong to $\{1,\ldots,p\}$ and sum to $2p$.
All five are at most $p+1$, and their total is $2N$.
The native displacement formula therefore gives Alexander drop $2p$.
For the inverse cycle, the two first displacements become $p$, and
the remaining three become $N-d$, with $d$ ranging over those last
three positive displacements. The weight of each such inverse
displacement is $(p+1)d$. Hence its total displacement weight is
\[
                    2p^2+(p+1)(2p)=2pN.
\]
Its drop is also $2p$. The anchor normalization $A(x_0)=g$ proves the
stated Alexander values.
\end{proof}


\section{The native second representative and the three acceptance statements}
\label{J21:sec:torus-final-bridge}

All coefficients in this section lie in $\mathbb F _2=\mathbb Z/2$.
Let $n\geq2$, put
\[
                  p=n+1,\qquad N=2n+3=2p+1,\qquad
                  g=\frac{n(n+1)}2=\frac{p(p-1)}2,
\]
and let $G=\decl{reflectedTorusGrid}\,n$.
We write $U^k$ for the native Laurent monomial
\decl{LaurentPolynomial.T} at exponent $k\in\mathbb Z$.
Thus multiplication by $U^k$ changes the Maslov degree by $-2k$,
as specified by \decl{laurentMaslovDegree}.
Let $x_0=\decl{reflectedTorusAnchor}\,n$ and let
$D=\decl{laurentDifferential}\,G\,(\mathbb Z/2)$.
The markings and the anchor are
\[
       O(c)=N-1-c,\qquad
       X(c)=N-1-((c+p+1)\bmod N),\qquad
       x_0(c)=-c\bmod N.
 \tag{FB.1}
\]
These are exactly \decl{reflectedTorusGrid_O_apply},
\decl{reflectedTorusGrid_X_apply}, and
\decl{reflectedTorusAnchor_val}.
The existing declarations \decl{maslovO_reflectedTorusAnchor} and
\decl{alexander_reflectedTorusAnchor} give $M(x_0)=0$ and $A(x_0)=g$.

\subsection{The correction relative to the native first cycle}

The native triangular chain has the following formula in the present
notation:
\[
 H_{\triangle}
   =\sum_{a=0}^{p-2}\ \sum_{b=p-1}^{p-1+a}x_{(a\,b)}.
 \tag{FB.2}
\]
Here $a,b$ are ordinary representatives in $\{0,\ldots,N-1\}$,
and $x_{(a\,b)}$ is the anchor with columns $a,b$ interchanged.
Indeed, \decl{firstTorusTrianglePairs} uses precisely the inequalities
$a<n$, $n\leq b$, and $b\leq n+a$.
Substitution of $n=p-1$ gives (FB.2), and
\decl{firstTorusTriangle} assigns coefficient one to every listed
transposition.
The lower bound $p-1>a$ shows that the two columns are distinct.
The upper bound $b\leq2p-3<N$ verifies that every pair is in the
native column range.
Thus
\[
 \decl{firstTorusCycle}\,n\,(\mathbb Z/2)
                       =x_0+U^{-1}D H_{\triangle}
 \tag{FB.3}
\]
is an equality with the native definition, with no change of cyclic cut.

Let $H_e$ be the explicit odd-shorter-distance transposition sum of
Lemma~\ref{J21:lem:torus-balanced-first}.
Let $H_4$ be the explicit adjacent-pair sum (H4D.7), and let
$H_{32}$ be the intrinsic five-column sum (H32.6).
Define
\[
 \boxed{\displaystyle
 w_2=U^{-1}(H_e+H_{\triangle})+U^{-2}H_4+U^{-3}H_{32}}
 \tag{FB.4}
\]
and
\[
 z_2=\decl{firstTorusCycle}\,n\,(\mathbb Z/2)+D w_2.
 \tag{FB.5}
\]
Equations (FB.2), (H4D.7), (H32.6), and the displayed definition of
$H_e$ specify every term and coefficient of $w_2$ by finite sums.

\begin{theorem}[New theorem: the explicit correction and exact class]
\label{J21:thm:torus-second-native-cycle}
For every $n\geq2$, the chain $w_2$ in (FB.4) has native Maslov
degree one. The chain $z_2$ in (FB.5) has native Maslov degree zero,
is a cycle, and represents exactly
\decl{reflectedTorusAnchorClass}\,$n\,(\mathbb Z/2)$.
Its full coefficient formula is
\[
                         z_2=U^{-2}R_4+U^{-3}R_6,
 \tag{FB.6}
\]
where $R_4$ is the explicit sum in (M4B.3) and $R_6$ is the explicit
sum in Theorem~\ref{J21:thm:h32-m6-support}.
In particular, the only Laurent exponents in its support are
$-2$ and $-3$.
\end{theorem}
\begin{proof}
Every state of $H_e$ and $H_{\triangle}$ has Maslov value $-1$.
For $H_{\triangle}$ this is also the existing declaration
\decl{firstTorusTriangle_mem_maslovPiece}, together with
\decl{maslovO_reflectedTorusAnchor}.
Every state of $H_4$ has Maslov value $-3$ by
Lemma~\ref{J21:lem:h4-density-degree}, and every state of $H_{32}$ has
Maslov value $-5$ by the intrinsic-order calculation in
Section~\ref{J21:sec:h32-definition}.
The shifts in (FB.4) therefore give respectively
\[
                       -1+2=1,\qquad -3+4=1,\qquad -5+6=1.
\]
The exact native membership follows from
\decl{mem_laurentMaslovPiece} and
\decl{laurentMonomial_smul_mem_maslovPiece}.

Since $D$ is Laurent-linear and the coefficients have characteristic
two, (FB.3)--(FB.5) give
\[
 z_2=x_0+D\bigl(U^{-1}H_e+U^{-2}H_4+U^{-3}H_{32}\bigr).
 \tag{FB.7}
\]
The existing declaration \decl{laurentDifferential_mem_maslovPiece}
states that $D$ lowers degree by one.
Together with $M(x_0)=0$, this proves degree zero for $z_2$.
The existing declarations
\decl{laurentDifferential_reflectedTorusAnchor} and
\decl{laurentDifferential_sq_eq_zero} then give
\[
                         D z_2=D x_0+D^2(\cdots)=0.
\]

We next verify the exact native homology class.
The existing map \decl{nativeCornerCycleOfChain} takes this
homogeneous cycle to a cycle in degree zero of
\decl{laurentMaslovComplex}\,$G\,(\mathbb Z/2)$.
The degree-one chain (FB.4) determines a native degree-one element
using the inverse of \decl{laurentMaslovChainEquiv}.
By \decl{laurentMaslovChainEquiv_d}, its native boundary is exactly
$D w_2$ on underlying Laurent chains.
Equation (FB.5) therefore identifies the native cycle of $z_2$
with \decl{firstTorusNativeCycle}\,$n\,(\mathbb Z/2)$ plus that
specific boundary.
Ordinary homology kills this boundary.
The existing declaration \decl{firstTorusNativeCycle_homology}
identifies the remaining class with
\decl{reflectedTorusAnchorClass}\,$n\,(\mathbb Z/2)$.
This proves equality with the prescribed class itself.
Equivalently, (FB.7) supplies a degree-one boundary witness directly
from the anchor cycle used to define that same class.

It remains to compute the support of $z_2$.
Write $D=\sum_{j\geq0}U^j d_j$.
Lemma~\ref{J21:lem:torus-balanced-first} gives
$x_0+U^{-1}D H_e=U^{-1}T_{\mathrm{bal}}$.
Corollary~\ref{J21:cor:h4-h32-lower-identity} gives
\[
 d_1H_4+d_2H_{32}=T_{\mathrm{bal}},\qquad d_jH_4=0\quad(j\geq2).
\]
The two-$O$ calculation in Section~\ref{J21:sec:h32-definition} gives
$d_jH_{32}=0$ for $j\geq3$.
Substitute these identities into (FB.7).
The coefficient of $U^{-1}$ is
$T_{\mathrm{bal}}+d_1H_4+d_2H_{32}=0$.
The remaining terms are
\[
             U^{-2}(d_0H_4+d_1H_{32})+U^{-3}d_0H_{32}.
\]
Theorem~\ref{J21:thm:m4-exact} and
Theorem~\ref{J21:thm:h32-m6-support} identify these two raw chains
with $R_4$ and $R_6$ respectively.
This proves (FB.6), including every coefficient, and proves M902.
\end{proof}

\subsection{The native bifiltration and weighted level}

\begin{theorem}[New theorem: the full second-vertex bound]
\label{J21:thm:torus-second-bifiltration}
For every $n\geq2$, the chain $z_2$ satisfies the native coefficient
bounds
\[
                -k\leq3,\qquad
                A(y)-k\leq g-2p+3
 \tag{FB.8}
\]
for every monomial $U^k y$ with nonzero coefficient.
Thus it belongs to
\[
 \decl{laurentBifiltrationLevel}\,G\,(\mathbb Z/2)\,
       3\,\left(\frac{(n-2)(n-1)}2\right).
 \tag{FB.9}
\]
Its native degree-zero cycle belongs to the corresponding level of
\decl{laurentMaslovBifiltration}.
\end{theorem}
\begin{proof}
By Theorem~\ref{J21:thm:m4-exact}, every term of $R_4$ has
$A(y)=g-2p$.
Its term in (FB.6) has $k=-2$, so its two coordinates are
\[
                           2,\qquad g-2p+2.
\]
By Theorem~\ref{J21:thm:h32-m6-support}, every term of $R_6$ has
$A(y)\leq g-2p$.
Its term in (FB.6) has $k=-3$, so its two coordinates are bounded
by
\[
                           3,\qquad g-2p+3.
\]
This proves (FB.8).
The identity
$g-2p+3=(n-2)(n-1)/2$ is obtained by substituting
$p=n+1$ and $g=n(n+1)/2$.
The existing declaration
\decl{mem_laurentMaslovBifiltrationLevel_iff}
states exactly that the native homogeneous bifiltration is given
by these two coefficient bounds.
Its defining pullback along \decl{laurentMaslovChainEquiv} gives
the final assertion.
\end{proof}

\begin{theorem}[New theorem: exact acceptance M901]
\label{J21:thm:torus-second-M901}
Let $n,m$ be natural numbers with $2\leq n$, $0<m$, and
$2m<n+1$.
In degree zero of
\decl{laurentMaslovComplex}\,$(\decl{reflectedTorusGrid}\,n)\,
(\mathbb Z/2)$, the class
\decl{reflectedTorusAnchorClass}\,$n\,(\mathbb Z/2)$ is
\decl{RealFiltration.RepresentedBelow} the level
\[
       \decl{Staircase.consecutiveWeight}\,
                   n\,\left(\frac2m\right)\,\langle2\rangle
                       =\frac{g+3m-2p}{m}
 \tag{FB.10}
\]
for the weighted native bifiltration at $t=2/m$.
More generally, the same cycle represents the class at the second
vertex for every $t\in[0,2]$.
\end{theorem}
\begin{proof}
The index $\langle2\rangle$ belongs to $\mathrm{Fin}(n+1)$ because
$n\geq2$.
The exact native definitions
\decl{Staircase.consecutiveFirst} and
\decl{Staircase.consecutiveSecond} give its coordinates
\[
                     3,\qquad \frac{(n-2)(n-1)}2=g-2p+3.
\]
The definition \decl{Bifiltration.weighted} includes every element
of this bifiltration level in weighted level
\[
                \left(1-\frac t2\right)3
                    +\frac t2(g-2p+3)
                =3+\frac t2(g-2p).
 \tag{FB.11}
\]
Indeed, this bifiltered element is one of the generators in the
span defining that weighted level.
Theorem~\ref{J21:thm:torus-second-bifiltration} supplies the required
membership, and Theorem~\ref{J21:thm:torus-second-native-cycle} supplies
both a cycle and its exact homology class.
These are precisely the two conditions in the existing definition
\decl{RealFiltration.RepresentedBelow}.

For $m>0$, put $t=2/m$.
Equation (FB.11) becomes $(g+3m-2p)/m$, which proves (FB.10) and
the acceptance statement.
Equivalently, the weight of an individual monomial is
\[
       \left(1-\frac1m\right)(-k)+\frac1m(A(y)-k)
                         =\frac{A(y)}m-k,
\]
and (FB.6) gives directly
\[
                        A(y)-mk\leq g+3m-2p.
 \tag{FB.12}
\]
The proof supplies (FB.12) for every positive integer $m$.
The additional strict inequality is therefore included without
requiring a further construction.
\end{proof}

\subsection{The exact integer inequality of M903}

For a native state $y$, abbreviate the integer appearing on the
right side of M903 by
\[
 \mathcal W_n(y)=\sum_{c\in\mathrm{Fin}(N)}
       W_p\bigl(d(c,\sigma_y(c))\bigr),\qquad
 W_p(d)=p\,d-N\max\{d-p-1,0\},
 \tag{FB.13}
\]
where $d(c,c)=0$ and
$\sigma_y=\decl{reflectedTorusRelativeColumns}\,n\,y$.
For representatives in $\{0,\ldots,N-1\}$, the function $d$ is
exactly \decl{reflectedTorusCyclicDisplacement}\,$n$.
The integer weight $W_p$ is exactly
\decl{reflectedTorusDisplacementWeight}\,$n$, because $n+2=p+1$
and the subtraction inside the native natural-number expression
is truncated at zero.
The existing identities
\decl{alexanderTwo_reflectedTorus_eq_sum_displacement} and
\decl{alexander_reflectedTorusAnchor} give, after embedding the
integer quantities in $\mathbb Q$,
\[
                            \mathcal W_n(y)=N(g-A(y)).
 \tag{FB.14}
\]

\begin{theorem}[New theorem: exact acceptance M903]
\label{J21:thm:torus-second-M903}
Under the hypotheses of Theorem~\ref{J21:thm:torus-second-M901}, every
nonzero monomial $U^k y$ of $z_2$ satisfies
\[
                  N\bigl(2(n+1)-m(3+k)\bigr)
                                      \leq \mathcal W_n(y).
 \tag{FB.15}
\]
\end{theorem}
\begin{proof}
For $k=-2$, Theorem~\ref{J21:thm:m4-exact} gives
$\mathcal W_n(y)=2pN$.
The left side of (FB.15) is $N(2p-m)$, which is at most $2pN$
because $m>0$.
For $k=-3$, Theorem~\ref{J21:thm:h32-m6-support} gives
$\mathcal W_n(y)\geq2pN$, and the left side is exactly $2pN$.
These two cases exhaust the support by (FB.6).
Equation (FB.15) is therefore the native M903 inequality with all
three of its named functions unchanged.
Applying the existing declaration
\decl{alexander_reflectedTorus_le_of_displacement_weight} with
integer parameter $2p-3m-mk$ recovers (FB.12).
\end{proof}

\begin{proof}[Completion of the proof of Theorem~\ref{J21:thm:j21-main}]
Take the finite chain $w_2$ in (FB.4). Theorem~\ref{J21:thm:torus-second-native-cycle}
proves its degree, the cycle condition, and the exact anchor-class identity.
Theorem~\ref{J21:thm:torus-second-bifiltration} proves
\eqref{J21:eq:j21-strong-corner}. Theorems~\ref{J21:thm:torus-second-M901}
and~\ref{J21:thm:torus-second-M903} prove \eqref{J21:eq:j21-weighted-target},
\eqref{J21:eq:j21-M903-target}, and the required represented-below statement.
These statements establish every assertion of Theorem~\ref{J21:thm:j21-main}.
\end{proof}

\subsection{The smallest parameters}

The following checks are direct specializations of the proved formulas.
They are not used to infer any assertion for larger sizes.

For $n=2$, one has $p=3$, $N=7$, and $g=3$.
The second vertex is $(3,0)$.
The only positive integer in the strict range $2m<3$ is $m=1$,
and the requested weighted parameter is $t=2$.
The native marking rows, in increasing column order, are
\[
 O=(6,5,4,3,2,1,0),\qquad
 X=(2,1,0,6,5,4,3),\qquad
 x_0=(0,6,5,4,3,2,1).
\]
Formula (FB.2) gives
$H_{\triangle}=x_{(0\,2)}+x_{(1\,2)}+x_{(1\,3)}$.
Every term of $U^{-2}R_4$ has coordinates $(2,-1)$, and every
term of $U^{-3}R_6$ has coordinates bounded by $(3,0)$.
For example, the positive circular order $(0,1,3,4,5)$ is a balanced
pentagon. Its positive star is $(0,3,5,1,4)$.
Its five displacements are $3,2,3,3,3$, of total $2N=14$.
It consequently has Alexander value $g-2p=-3$ and occurs in
$U^{-3}R_6$ with weighted level zero at $t=2$.
This meets the asserted bound at the smallest allowed grid.

If $m\geq2$, the first parameters satisfying the strict inequality
are $p=5$, $n=4$, and $m=2$.
Here $N=11$, $g=10$, the vertex is $(3,3)$, and $t=1$.
The level (FB.10) is $3$, or equivalently the monomial inequality is
$A(y)-2k\leq6$.
The terms of $U^{-2}R_4$ have $A=0$ and satisfy $A-2k=4$.
The terms of $U^{-3}R_6$ have $A\leq0$ and satisfy $A-2k\leq6$.
The pentagon $(0,2,4,6,8)$ is balanced because its successive
two-gap distances are $4,4,4,5,5$.
Its positive star $(0,4,8,2,6)$ has displacements $4,4,5,4,5$,
of total $2N=22$, and Alexander value zero.
The corresponding term of $U^{-3}R_6$ has weighted level exactly
$3$.
The native triangle in this case is
\[
 H_{\triangle}
  =x_{(0\,4)}
   +x_{(1\,4)}+x_{(1\,5)}
   +x_{(2\,4)}+x_{(2\,5)}+x_{(2\,6)}
   +x_{(3\,4)}+x_{(3\,5)}+x_{(3\,6)}+x_{(3\,7)},
\]
which checks the native index shift independently of the cyclic
formulas used for the other correction terms.

\paragraph{Acceptance and remaining gap.}
Formula (FB.4) and Theorem~\ref{J21:thm:torus-second-native-cycle}
prove M902. Theorem~\ref{J21:thm:torus-second-bifiltration} proves the
stronger bifiltration statement accompanying M901, and
Theorem~\ref{J21:thm:torus-second-M901} proves its exact weighted
acceptance statement. Theorem~\ref{J21:thm:torus-second-M903} proves
the stated native integer inequality.
There is no remaining mathematical gap in J21.


\section{Catalogue of the proof inputs}
\label{J21:sec:j21-lemma-catalogue}

All coefficient calculations in this proof take place in
$\mathbb F _2$. The following tables distinguish existing native
declarations from the new lemmas proved in this document.
No holomorphic Heegaard Floer theorem is an input.

Unqualified names in rows supplied in full belong to
\decl{OnLeConjecture.KnotFloer.GridFull}, unless a different namespace
is displayed. The prefixes \decl{Staircase.}, \decl{Bifiltration.},
and \decl{RealFiltration.} are relative to
\decl{OnLeConjecture.KnotFloer}.
Source filenames are relative to \decl{OnLeConjecture/KnotFloer/}.
They match the supplied file headings, except for the explicitly marked
location supplied only in the research notes.
The rows marked imported explicitly distinguish imported
primitive names from declarations whose defining statements are
included in the attachment.

\providecommand{\jtwentyonefile}[2]{\texttt{#1}\allowbreak\texttt{#2}}

\subsection{Existing Lean declarations}

\begingroup
\small
\renewcommand{\arraystretch}{1.14}
\begin{longtable}{@{}>{\raggedright\arraybackslash}p{0.43\linewidth}
>{\raggedright\arraybackslash}p{0.26\linewidth}
>{\raggedright\arraybackslash}p{0.23\linewidth}@{}}
\textbf{Existing declaration or definition}&
\textbf{Verified supplied source}&\textbf{Use in this proof}\\\hline
\endfirsthead
\textbf{Existing declaration or definition}&
\textbf{Verified supplied source}&\textbf{Use in this proof}\\\hline
\endhead

\decl{consecutiveTorusGridSize},
\decl{reflectedTorusGrid},
\decl{reflectedTorusAnchor},
\decl{reflectedTorusGrid_O_apply},
\decl{reflectedTorusGrid_X_apply}
&\jtwentyonefile{GridReflectedTorus}{Basic.lean}
&The exact size and native $O,X$ marking conventions.\\[3pt]

\decl{reflectedTorusAnchor_val},
\decl{maslovO_reflectedTorusAnchor}
&\jtwentyonefile{GridReflectedTorus}{Anchor.lean}
&The state $x_0(c)=-c$ and its Maslov value zero.\\[3pt]

\decl{LaurentChain},
\decl{laurentDifferential},
\decl{laurentCoefficient_eq},
\decl{laurentDifferential_monomialGenerator},
\decl{laurentDifferential_sq_eq_zero}
&\jtwentyonefile{GridFull}{Laurent.lean}
&The native chain module, the sum over state-empty rectangles with
weight $U^{\#O}$, and $D^2=0$.\\[3pt]

\decl{laurentMaslovDegree},
\decl{mem_laurentMaslovPiece},
\decl{laurentMaslovComplex},
\decl{laurentMaslovDegree_add_card_OColumns},
\decl{laurentMonomial_smul_mem_maslovPiece},
\decl{laurentDifferential_mem_maslovPiece}
&\jtwentyonefile{GridFullLaurent}{Maslov.lean}
&The degree $M_O(x)-2k$, the coefficient criterion for homogeneity,
and the differential's degree minus one.\\[3pt]

\decl{maslovO_swap_reflectedTorusAnchor}
&\jtwentyonefile{GridReflectedTorus}{FirstCycle.lean}
&The Maslov value minus one of the transposition states used in $H_e$.\\[3pt]

\decl{laurentDifferential_reflectedTorusAnchor}
&\jtwentyonefile{GridReflectedTorus}{Basic.lean}
&The anchor is a cycle of the full Laurent differential.\\[3pt]

\decl{alexander_reflectedTorusAnchor}
&\jtwentyonefile{GridReflectedTorus}{AnchorAlexander.lean}
&The absolute normalization $A(x_0)=n(n+1)/2=g$.\\[3pt]

\decl{reflectedTorusCyclicDisplacement},
\decl{reflectedTorusRelativeColumns},
\decl{reflectedTorusDisplacementWeight},
\decl{alexanderTwo_reflectedTorus_eq_sum_displacement},
\decl{alexander_reflectedTorus_le_of_displacement_weight}
&\jtwentyonefile{GridReflectedTorus}{SecondAlexander.lean}
&The exact displacement formula and the native implication from
the integer inequality M903 to the Alexander bound.\\[3pt]

\decl{laurentMaslovChainEquiv},
\decl{laurentMaslovChainEquiv_d}
&\jtwentyonefile{GridFullLaurent}{Spectral.lean}
&Identification of the native graded complex and its differential
with homogeneous Laurent chains and $D$.\\[3pt]

\decl{nativeCornerCycleOfChain},
\decl{nativeCornerCycleOfChain_apply},
\decl{reflectedTorusAnchorCycle},
\decl{reflectedTorusAnchorClass}
&\jtwentyonefile{GridReflectedTorus}{FirstClass.lean}
&Passage from the displayed degree-zero cycle to the native cycle
module and the exact anchor class.\\[3pt]

\decl{mem_laurentMaslovBifiltrationLevel_iff},
\decl{laurentMaslovBifiltration}
&\jtwentyonefile{GridFullLaurent}{Spectral.lean}
&The coefficient bounds $-k\leq a$ and $A(x)-k\leq b$ are precisely
the native bifiltration bounds.\\[3pt]

\decl{Bifiltration.weighted},
\decl{RealFiltration.RepresentedBelow}
&\jtwentyonefile{Spectral}{Threshold.lean}
&The weighted level and the statement that the specified homology
class has a representative at that level.\\[3pt]

\decl{Staircase.consecutiveFirst},
\decl{Staircase.consecutiveSecond},
\decl{Staircase.consecutiveWeight}
&\jtwentyonefile{Consecutive}{Staircase.lean}
&The exact vertex-two coordinates and its weighted value.\\[3pt]

\decl{firstTorusTrianglePairs},
\decl{firstTorusTriangle},
\decl{firstTorusCycle},
\decl{firstTorusTriangle_mem_maslovPiece}
&\jtwentyonefile{GridReflectedTorus}{FirstCycle.lean}
&The exact triangular index set, its degree, and the definition of
the existing first representative used in the explicit correction.\\[3pt]

\decl{firstTorusNativeCycle},
\decl{firstTorusNativeCycle_homology}
&\jtwentyonefile{GridReflectedTorus}{FirstClass.lean}
&Identification of the native first cycle with the exact anchor
homology class in the boundary comparison.\\[3pt]

Imported native filtration definition
\decl{laurentBifiltrationLevel}
&\jtwentyonefile{GridFullLaurent}{Bifiltration.lean},
identified in the supplied notes and imported by
\jtwentyonefile{GridFullLaurent}{Spectral.lean}
&The ungraded native bifiltration named in the stronger conclusion.
Its defining listing is omitted. The fully supplied homogeneous
coefficient criterion above supplies the needed support test.\\[3pt]

Imported rectangle and monomial primitives
\decl{LaurentPolynomial.T},
\decl{OColumns},
\decl{GridRectangleBetween.emptyRectangles},
\decl{GridState.swapColumns},
\decl{GridRectangle.mem_coveredSquares},
\decl{isEmpty_iff_forall_notMem_cIoo}
&Use sites in
\jtwentyonefile{GridFull}{Laurent.lean},
\jtwentyonefile{GridReflectedTorus}{FirstCycle.lean}
&The Laurent monomial notation, native half-open marking intervals,
open state interiors, and column-swap convention.
Their defining files are not separately
listed in the attachment.\\[3pt]

Imported grading primitives
\(\decl{GridDiagram.maslovO}\mathbb Z\decl{_eq_card}\),
\decl{GridDiagram.two_mul_alexander_eq_intCast}
&Use sites in
\jtwentyonefile{GridReflectedTorus}{Anchor.lean},
\jtwentyonefile{GridReflectedTorus}{AnchorAlexander.lean},
\jtwentyonefile{GridReflectedTorus}{SecondAlexander.lean}
&The four finite counts defining integer Maslov grading and the
identity $2A=\operatorname{alexanderTwo}$. The supplied uses verify
these names, but do not identify their own defining filenames.\\\hline
\end{longtable}
\endgroup

The imported grading identifier containing $\mathbb Z$ uses the native
Unicode character \texttt{U+2124} in its Lean spelling.
The table does not assign an unverified defining filename to an
imported primitive. In particular, the supplied criterion
\decl{mem_laurentMaslovBifiltrationLevel_iff} in
\decl{GridFullLaurentSpectral.lean} is used for the filtration
conversion, so the omitted
\decl{GridFullLaurentBifiltration.lean} listing is not needed as a
separate citation for that conversion.

\subsection{New lemmas and explicit constructions}

Every entry in the following table is new and is proved in this
document. The references point to the corresponding mathematical
statements rather than to diagnostic calculations.

\begingroup
\small
\renewcommand{\arraystretch}{1.14}
\begin{longtable}{@{}>{\raggedright\arraybackslash}p{0.35\linewidth}
>{\raggedright\arraybackslash}p{0.57\linewidth}@{}}
\textbf{New lemma group and reference}&\textbf{Conclusion used}\\\hline
\endfirsthead
\textbf{New lemma group and reference}&\textbf{Conclusion used}\\\hline
\endhead

Rectangle coordinates and finite-count Maslov formula,
Lemmas~\ref{J21:lem:torus-relative-rectangle},
\ref{J21:lem:torus-permutation-maslov},
\ref{J21:lem:torus-fixed-column-deletion},
\ref{J21:lem:torus-rotation-maslov}, and
Corollary~\ref{J21:cor:torus-small-types}
&The exact relative interval test and
$M(x_\sigma)=\operatorname{inv}(\sigma)-\sum_c|\sigma(c)-c|$,
including deletion of fixed columns, cyclic-support rotation
gradings, and the small permutation types.\\[3pt]

Cyclic relabelling, cycle-count bound, and exact insertion index,
Lemmas~\ref{J21:lem:torus-cyclic-conjugation},
\ref{J21:lem:torus-cycle-lower-bound},
\ref{J21:lem:torus-exact-insertion-index},
\ref{J21:lem:torus-minimal-cycle-deletion}, and
\ref{J21:lem:torus-minimal-cycle-coefficient}
&The unique zero-$O$ insertion rectangle has interior count
$(\ell(\tau)-\ell(\sigma)-1)/2$. Deleting vertices computes the
boundary coefficients of cycles with $\ell=k-1$.\\[3pt]

Transposition differential and balanced first correction,
Lemmas~\ref{J21:lem:torus-transposition-differential} and
\ref{J21:lem:torus-balanced-first}
&The explicit chain $H_e$ satisfies $DH_e=Ux_0+T_{\mathrm{bal}}$.\\[3pt]

Structural rectangle identities and cycle types,
Lemmas~\ref{J21:lem:torus-transposition-cycle-count},
\ref{J21:lem:torus-insertion-increases},
\ref{J21:lem:torus-fixed-columns-cancel},
\ref{J21:lem:torus-inverse-commutes},
\ref{J21:lem:torus-three-two-type},
\ref{J21:lem:torus-five-cycle-table}, and
\ref{J21:lem:torus-three-two-boundaries}
&Two-fixed-column terms cancel, inversion preserves rectangle
coefficients, and the possible boundary types and Maslov values
are completely classified.\\[3pt]

The explicit $H_{32}$ coefficient and its two-$O$ boundary,
Section~\ref{J21:sec:h32-definition},
Lemmas~\ref{J21:lem:h32-invariance},
\ref{J21:lem:h32-two-O-neighbours},
\ref{J21:lem:h32-two-O-gap}, and
Corollaries~\ref{J21:cor:h32-selfinverse} and
\ref{J21:cor:h32-positive-types}
&The intrinsic five-column indexing, the complete coefficient
table, inversion pairing, and the exact triangle coefficient
$L_p(\alpha,\gamma)+L_p(\beta,\alpha)+L_p(\gamma,\beta)$,
together with the exclusion of higher $O$ counts.\\[3pt]

The explicit adjacent-pair correction $H_4$,
Section~\ref{J21:sec:h4-density} and
Lemma~\ref{J21:lem:h4-density-degree}
&A finite formula for every correction coefficient and the degree
one of $U^{-2}H_4$.\\[3pt]

Adjacent-pair coefficient identities,
Lemmas~\ref{J21:lem:h4-unit-support},
\ref{J21:lem:h4-common-inverse-sum}, and
\ref{J21:lem:h4-convex-primitive}
&Support on adjacent pairs, a common coefficient sum for inverse
cycles, and the three-point primitive that cancels the circular
five-cycle boundary.\\[3pt]

Four-cycle deletion and the full lower cancellation,
Lemmas~\ref{J21:lem:torus-fourcycle-deletion},
\ref{J21:lem:h4-unit-coefficient},
\ref{J21:lem:h4-triangle-coefficient},
\ref{J21:lem:h4-h32-triangle-cancellation},
\ref{J21:lem:h4-double-transposition-cancellation}, and
Corollary~\ref{J21:cor:h4-h32-lower-identity}
&The complete equality $d_1H_4+d_2H_{32}=T_{\mathrm{bal}}$,
the cancellation of every two-transposition target, and
$d_jH_4=0$ for $j\geq2$.\\[3pt]

The four-point primitive identities,
Lemma~\ref{J21:lem:m6-primitives}
&The functions $F,K,G$ express the coefficient of $H_{32}$ in
forms that telescope in the boundary calculation.\\[3pt]

The Maslov-minus-six coefficient and support calculation,
Lemmas~\ref{J21:lem:m6-types}, \ref{J21:lem:m6-four-two},
\ref{J21:lem:m6-pentagon}, \ref{J21:lem:m6-three-three}, and
Theorem~\ref{J21:thm:h32-m6-support}
&The exact chain $R_6=d_0H_{32}$ consists of the stated balanced
pentagon and two-triangle states. Every term has Alexander drop
at least $2p$.\\[3pt]

Diagonal identities and the first nonconvex five-cycle type,
Lemmas~\ref{J21:lem:h4-diagonal-primitive-first},
\ref{J21:lem:h4-diagonal-primitive-second},
\ref{J21:lem:h4-A-four-coefficient}, and
\ref{J21:lem:h4-h32-A-cancellation}
&The explicit $T_p^0$ identities and the cancellation of the
five-cycle type $01243$.\\[3pt]

The remaining five-cycle coefficients,
Lemmas~\ref{J21:lem:m4-residual-inverse},
\ref{J21:lem:m4-local-two-identities},
\ref{J21:lem:m4-B-unit-table}, and
Theorem~\ref{J21:thm:m4-exact}
&The exact residual $R_4=d_0H_4+d_1H_{32}$ and the Alexander
value $g-2p$ of each surviving term.\\[3pt]

The explicit native cycle and acceptance statements,
Theorems~\ref{J21:thm:j21-main},
\ref{J21:thm:torus-second-native-cycle},
\ref{J21:thm:torus-second-bifiltration},
\ref{J21:thm:torus-second-M901}, and
\ref{J21:thm:torus-second-M903}
&The degree-one correction $w_2$, its complete coefficient formula,
the exact anchor class, the second-vertex bifiltration bound,
the required weighted representation, and the precise native
integer displacement inequality.\\\hline
\end{longtable}
\endgroup

The proof of the requested representative uses these identities in
the single chain equation
\[
 x_0+D\bigl(U^{-1}H_e+U^{-2}H_4+U^{-3}H_{32}\bigr)
                  =U^{-2}R_4+U^{-3}R_6.
\]
The displayed boundary gives the exact anchor homology class.
The existing square-zero identity gives the cycle condition.
The new residual calculations give the support bound.




\chapter{J10: One X switch for a deleted braid letter}
\label{chapter:J10}



\section{Exact statement and source status}

The J10 acceptance statement is the following assertion. Let $m\ge2$,
let $u,v$ be positive adjacent words on $m$ strands, and let
$i\in\{0,\ldots,m-2\}$. Suppose that every adjacent generator occurs
in $u v$. Put $w=u\sigma_i v$, $w'=uv$, and $j=|u|$.
Here the subscripts are the native zero-based indices, so $\sigma_0$
denotes the generator usually written $\sigma_1$.
There exists a grid $H$ of size $3|w|$ such that
\begin{equation}\label{J10:eq:target}
 \operatorname{XSwitchSequence}(\operatorname{wordGrid}(w),H,1)
 \text{ is nonempty},\qquad
 H\Mov\operatorname{wordGrid}(w').
\end{equation}
The switch in the proof is in the consecutive columns $3j+1,3j+2$.
The grid $H$ has exactly the same $O$ marking as the initial grid.

We prove an explicit statement for every pair of word grids satisfying
the marking and row-order equations of the supplied declaration
\decl{exists_positiveRectilinearWordGrid}. Thus the proof is independent
of every choice of distinct row perturbations. The supplied audit defines
$\operatorname{wordGrid}(a)$ as the grid field of
\decl{positiveRoundedWordSelectedData} $a$. Instantiating the
proved construction at that field requires its incidence equation
and its strict height-order property. The defining structure and
selector are not reproduced in the attachment, and the audit's
field inventory does not state those two equations. They are isolated
precisely in Section~\ref{J10:sec:native-identifications}.

There is a source limitation. The combined attachment does not reproduce
\nolinkurl{PositiveBraidGridSelectedRepresentative.lean}, where the selected
data are defined. It also imports rather than reproduces the defining
modules for \decl{RowsNoninterleaving}, \decl{ColumnsNoninterleaving},
and \decl{IsStabilization}. The supplied move module describes its
generators as the grid moves of Chapter~3 of the attached OSS book.
The proof specifies every local pattern and every noninterleaving
condition. A review of the omitted selected-data fields and primitive
move definitions is still needed to certify the literal native
instantiation. No equality of smooth link types is used as a substitute
for a grid-move sequence.

This is a combinatorial proof. It uses no coefficient field other than
$\F$ if one subsequently applies the supplied chain-map results.

\section{The row-separated word grids}

Fix a positive word $a$ of length $L$. Its columns are the pairs
$(h,c)$ with $0\le h<L$ and $c\in\{0,1,2\}$, ordered by
\begin{equation}\label{J10:eq:columns}
 (h,c)\longmapsto 3h+c.
\end{equation}
This is the order proved by
\decl{positiveRectilinearWordColumn_index_lt}.
For the letter in position $h$, let $r$ be its generator index.
Its three outgoing horizontal arcs have unperturbed heights
\[
 r+\tfrac12,\qquad r,\qquad r+1
\]
at columns $(h,0),(h,1),(h,2)$, respectively.
They are the $O$ corners. The corresponding $X$ corners have
unperturbed heights
\[
 r,\qquad r+1,\qquad r+\tfrac12.
\]
These formulas follow from
\decl{positiveRectilinearOutgoingArc_zero},
\decl{positiveRectilinearOutgoingArc_one},
\decl{positiveRectilinearOutgoingArc_two}, their incoming versions,
and \decl{positiveRectilinearHorizontalArcRow}.

Each horizontal arc has its own row after separation.
If its unperturbed height is $h(\alpha)$, write $R(\alpha)$ for its
finite grid row. Distinct arcs have distinct rows, and
\begin{equation}\label{J10:eq:roworder}
 h(\alpha)<h(\beta)\quad\Longrightarrow\quad
 R(\alpha)<R(\beta).
\end{equation}
The row of the $X$ corner at a column is the row of the incoming arc
matched to that column. The marking equations in
\decl{exists_positiveRectilinearWordGrid_of_injective} are exactly
\begin{equation}\label{J10:eq:markings}
 O(c)=R(\text{outgoing arc at }c),\qquad
 X(c)=R(\text{incoming arc at }c).
\end{equation}
The strict order in~\eqref{J10:eq:roworder} follows from
\decl{exists_positiveRectilinearWordGrid}.

For an integer height $r$, the occurrences are the letters involving
rank $r$. At each occurrence the incoming corner precedes the outgoing
corner in the column order. The outgoing arc at one occurrence goes
to the incoming corner of the next occurrence, in cyclic word order.
This is precisely
\decl{positiveRectilinearWordMatching_occurrence} and
\decl{positiveRectilinearSortedSuccessor_next}.
At a half-integer height, the arc belonging to position $h$ goes from
$(h,0)$ to $(h,2)$, by
\decl{positiveRectilinearWordMatching_middle}.

\begin{definition}
An admissible grid of a word $a$ is a grid with the column order
\eqref{J10:eq:columns}, the incidence equations~\eqref{J10:eq:markings}, and
a total order of arc rows satisfying~\eqref{J10:eq:roworder}.
\end{definition}

The definition records the part of the supplied word-grid construction
used below. The actual magnitudes of the real row perturbations do not
enter a finite grid diagram.

\section{Elementary move lemmas}

We use the following conventions for the elementary moves. Two adjacent
rows or columns commute when their endpoint pairs do not alternate
cyclically and have no common endpoint. Equivalently, after choosing
a cut of the circle, the two closed intervals are disjoint or one is
strictly inside the other. All commutations below have distinct
endpoints. A destabilization replaces three marked corners of an
adjacent $2\times2$ block by the single marking of the duplicated type
at the fourth corner, and deletes the row and column through the
corner of the other type. The fourth corner must be unmarked.
These are the conventions of
\cite[Definitions~3.1.6 and 3.1.7, Section~3.2]{OSSbook}.
They are also the move types
named in the supplied \decl{GridDiagram.IsMove}.

\begin{lemma}[New lemma: commutations of a short row or column]
\label{J10:lem:short}
Suppose the two markings of a row lie in adjacent columns $a,b$.
That row can be commuted with an adjacent row having neither endpoint
in columns $a,b$. The analogous statement holds for a column whose
two marking rows are adjacent.
\end{lemma}

\begin{proof}
No other column lies in the cyclic open interval from $a$ to $b$ in
their adjacent direction. Thus the endpoints of the other row are on
the same side of that interval. Its two endpoints cannot alternate
with $a,b$. They are distinct from $a,b$ by hypothesis, so the
commutation criterion holds. Interchanging rows and columns proves
the second assertion. Each commutation gives \decl{MovesTo} by
\decl{movesTo_swapRows} or \decl{movesTo_swapColumns}.
\end{proof}

\begin{lemma}[New lemma: same-height rows do not interleave]
\label{J10:lem:bands}
In an admissible word grid, any two different horizontal arcs with
the same unperturbed height have distinct, nonalternating endpoints.
Consequently, any two admissible grids of the same word are related
by row commutations.
\end{lemma}

\begin{proof}
At a half-integer height, the two arcs belong to different letters.
Each runs from column $3h$ to column $3h+2$ of its letter.
These intervals are disjoint for different letters, and their
endpoints are distinct.

At an integer height, list the occurrences in cyclic order.
The incoming and outgoing endpoints of each occurrence are distinct,
with incoming first. An arc joins an outgoing endpoint to the next
incoming endpoint. The interiors of the resulting arcs are mutually
disjoint on the column circle. No incoming endpoint equals any
outgoing endpoint, because the two columns at an occurrence are
distinct and different letters occupy disjoint column triples.
Thus all endpoint pairs are nonalternating and disjoint. This also
follows from the explicit endpoint exclusion in
\decl{positiveRectilinearIntegerRow_arc_excludes_occurrence}.

Condition~\eqref{J10:eq:roworder} makes the rows of each unperturbed height
a consecutive group. Two admissible row orders differ only in the
orders inside these groups. A finite permutation of a linearly ordered
set is a product of adjacent exchanges. One proof is to move the first
desired element left to its desired place, then repeat on the shorter
remaining list. Every exchange used here is between same-height
rows, so the preceding endpoint argument verifies its commutation
condition. Their composition proves the final assertion by
\decl{MovesTo.trans}.
\end{proof}

\begin{lemma}[New lemma: deleting a continuation column]
\label{J10:lem:continuation}
Consider a grid with an integer-height group of horizontal arcs indexed
in cyclic order. Assume the rows of this height group occupy a consecutive
interval in the row order. Suppose one column $c$ joins the end of an incoming
arc $\alpha_-$ to the start of an outgoing arc $\alpha_+$ at that
height, with
\[
 X(c)=R(\alpha_-),\qquad O(c)=R(\alpha_+),\qquad
 \alpha_-\ne\alpha_+.
\]
Let $d$ be the $X$ endpoint of $\alpha_+$ and $e$ the $O$ endpoint of
$\alpha_-$. Assume $c,d,e$ are distinct. Suppose the other arcs in
this height group have endpoint pairs disjoint from and nonalternating
with the pairs of $\alpha_-$ and $\alpha_+$.
Suppose also that, in the chosen cyclic direction of $\alpha_+$ from
$c$ to $d$, the column $e$ is not passed.

Then grid moves delete column $c$ and row $R(\alpha_+)$, and replace
the two horizontal arcs by a single arc from $e$ to $d$ on the row
of $\alpha_-$. The moves use row commutations within this height
group, column commutations involving only $c$, cyclic column
permutations if needed, and one destabilization. The cyclic order
of every surviving column is preserved.
\end{lemma}

\begin{proof}
Move the row of $\alpha_+$ toward the row of $\alpha_-$ by adjacent
row commutations, stopping when they are adjacent. All intervening
rows belong to the same height group. The hypothesis verifies each
noninterleaving test, and the two rows sharing column $c$ are never
commuted with one another.

The column $c$ now has adjacent endpoint rows. Move it along the
cyclic column order toward $d$, stopping immediately before $d$.
A cyclic column permutation can first place this portion of the
circle in a linear interval. No crossed column has a marking in either
of the two endpoint rows, because their other marked columns are
$d,e$, and the route crosses neither. Lemma~\ref{J10:lem:short} makes
each column commutation legal.

The two adjacent rows and the adjacent columns $c,d$ now contain
\[
\begin{array}{c|cc}
 &c&d\\ \hline
 R(\alpha_-)&X&\varnothing\\
 R(\alpha_+)&O&X.
\end{array}
\]
The fourth corner is empty because the other marking of the row of
$\alpha_-$ is in column $e\ne d$. This is an $X$ destabilization.
It deletes the column through the displayed $O$ and the row through
that $O$, and places the remaining $X$ in column $d$ on the row of
$\alpha_-$. Every deleted or moved column was $c$, so the cyclic
order of the surviving columns is unchanged. The row and incidence
claims follow from the displayed markings.
\end{proof}

\section{The single switch}

For the remainder of the proof let $w=u\sigma_i v$ and $w'=uv$ be as
in the acceptance statement. Let $G$ be any admissible grid of $w$.
Write
\[
 a=3j,\qquad b=3j+1,\qquad c=3j+2
\]
for the columns of the deleted letter.
Denote its middle arc by $M$, its outgoing lower arc by $L_+$,
and its outgoing upper arc by $U_+$.
Let $L_-$ and $U_-$ be the arcs incoming to the lower and upper corners
of this occurrence. Since $\sigma_i$ occurs in $w'$, each of the
two integer ranks $i,i+1$ has another occurrence. Its cyclic
predecessor is not this deleted occurrence. Hence
\begin{equation}\label{J10:eq:distinct}
 L_-\ne L_+,\qquad U_-\ne U_+.
\end{equation}
The outgoing and incoming arc descriptions give the following complete
local marking table:
\begin{equation}\label{J10:eq:initial}
\begin{array}{c|ccc}
 &a&b&c\\ \hline
 O&R(M)&R(L_+)&R(U_+)\\
 X&R(L_-)&R(U_-)&R(M).
\end{array}
\end{equation}
All lower rows are below $R(M)$, and all upper rows are above $R(M)$,
by~\eqref{J10:eq:roworder}.

\begin{lemma}[New lemma: the candidate switch is a grid]
\label{J10:lem:switch}
Define $H.O=G.O$ and let $H.X$ be $G.X$ with its values in columns
$b,c$ exchanged. Then $H$ is a grid diagram of size $3|w|$, and there
is a one-step native $X$-switch sequence from $G$ to $H$.
\end{lemma}

\begin{proof}
The two marking maps remain permutations. The only disjointness
conditions that change are those in columns $b,c$.
In column $b$ the new $X$ row is $R(M)$ and the $O$ row is $R(L_+)$.
Their unperturbed heights differ, so they are distinct.
In column $c$ the new $X$ row is $R(U_-)$ and the $O$ row is $R(U_+)$.
These are distinct by~\eqref{J10:eq:distinct} and injectivity of the row
assignment. All other column conditions hold because $G$ is a grid.

The indices $b,c$ are distinct and consecutive, with
$c=\decl{finRotate}(3|w|)(b)$. Indeed, $j<|w|$ and
$b=3j+1<3|w|-1$. The constructor \decl{XSwitchSequence.cons},
with the empty tail \decl{XSwitchSequence.nil} on $H$, gives the
required sequence using the two defining marking equalities.
\end{proof}

After the switch the local table is
\begin{equation}\label{J10:eq:switched}
\begin{array}{c|ccc}
 &a&b&c\\ \hline
 O&R(M)&R(L_+)&R(U_+)\\
 X&R(L_-)&R(M)&R(U_-).
\end{array}
\end{equation}
The proof now removes exactly these three columns by grid moves.

\section{The explicit sequence of grid moves}

\begin{theorem}[New lemma: deletion for every admissible word grid]
\label{J10:thm:deletion}
Let $m\ge2$, let $u,v$ be positive words on $m$ strands, and let
$i\in\{0,\ldots,m-2\}$. Suppose every generator occurs in $uv$.
For every admissible grid $G$ of $u\sigma_i v$ and every admissible
grid $G'$ of $uv$, the grid $H$ of Lemma~\ref{J10:lem:switch} satisfies
\[
 H\Mov G'.
\]
The sequence from $H$ uses exactly three destabilizations, together
with row and column commutations and cyclic column permutations.
\end{theorem}

\begin{proof}
\textit{Removal of the middle arc.}
The row of $M$ has its $O$ and $X$ in the adjacent columns $a,b$.
Both $R(L_-)$ and $R(L_+)$ lie below it.
Let $T$ be whichever of these two rows is higher.
Move the row of $M$ downward until it is immediately above $T$.
The only other rows with a marking in columns $a,b$ are precisely
$R(L_-)$ and $R(L_+)$, by~\eqref{J10:eq:switched}.
Neither is crossed, since the process stops at the higher of them.
Lemma~\ref{J10:lem:short} therefore verifies every row commutation used.

If $T=R(L_+)$, the adjacent block has markings
\[
\begin{array}{c|cc}
 &a&b\\ \hline
 R(M)&O&X\\
 R(L_+)&\varnothing&O.
\end{array}
\]
Its fourth corner is empty because $L_-\ne L_+$.
An $O$ destabilization deletes row $R(M)$ and column $b$, and
places $O$ at $(a,R(L_+))$.
If $T=R(L_-)$, the block is instead
\[
\begin{array}{c|cc}
 &a&b\\ \hline
 R(M)&O&X\\
 R(L_-)&X&\varnothing.
\end{array}
\]
An $X$ destabilization deletes row $R(M)$ and column $a$, and
places $X$ at $(b,R(L_-))$.
In either case one surviving column, denoted $d_0$, has
\begin{equation}\label{J10:eq:lowercont}
 X(d_0)=R(L_-),\qquad O(d_0)=R(L_+).
\end{equation}
All markings outside the deleted letter's columns are unchanged.
The upper continuation is still column $c$, with rows $R(U_-),R(U_+)$.

\textit{Removal of the lower continuation.}
In cyclic occurrence order at rank $i$, the preceding occurrence,
the deleted occurrence, and the next occurrence appear consecutively,
with the preceding and next allowed to be the same other occurrence.
Let $e_-$ be the outgoing corner of the preceding occurrence and
$e_+$ the incoming corner of the next occurrence.
The arcs $L_-$ and $L_+$ now run from $e_-$ to $d_0$ and from
$d_0$ to $e_+$, respectively.
The three columns are distinct. In particular, if the preceding and
next occurrences coincide, their incoming and outgoing corners are
distinct by \decl{positiveRectilinearIntegerRowIn_lt_out}.
The directed arc from $d_0$ to $e_+$ does not pass $e_-$.
Every other rank-$i$ arc has disjoint endpoints and disjoint interior
from these two arcs, by the cyclic successor description.

The rows of the integer-height group remain consecutive.
Moving and deleting $M$ has not changed the relative order of any
surviving rows. Lemma~\ref{J10:lem:continuation} therefore applies to
\eqref{J10:eq:lowercont}. Its moves delete $d_0$ and the outgoing row
$R(L_+)$, keeping $R(L_-)$ and joining $e_-$ directly to $e_+$.
This is exactly the change to the rank-$i$ cyclic successor when one
occurrence is deleted.

\textit{Removal of the upper continuation.}
The same argument applies to rank $i+1$, with incoming and outgoing
arcs $U_-,U_+$. The previous moves have not changed the cyclic order
of its surviving columns or the relative order of its rows.
Lemma~\ref{J10:lem:continuation} deletes the remaining local column $c$
and the outgoing row $R(U_+)$, and bypasses the deleted occurrence
in the upper cyclic successor.

Three columns and three arc rows have now been removed. They are
exactly the columns of the deleted letter and its outgoing arcs
$M,L_+,U_+$. The cyclic order of all other columns is the original
order with this triple removed. A cyclic column permutation restores
the conventional first column if needed.
Every half-integer arc belonging to another letter is unchanged.
At integer ranks other than $i,i+1$, every successor is unchanged.
At ranks $i,i+1$, the preceding and next occurrences have been joined,
as the definition of the successor for $uv$ requires.
Thus the resulting incidence is exactly that of the word $uv$.

The relative order of surviving rows has only been changed by
commutations inside an integer-height group. The temporarily moved
row $M$ was deleted. Hence the remaining row order still satisfies
\eqref{J10:eq:roworder} for $uv$. The resulting grid is admissible for
$uv$. Lemma~\ref{J10:lem:bands} relates it to the prescribed admissible
grid $G'$ by row commutations. Composition of these elementary moves
using \decl{MovesTo.trans} proves the theorem.
\end{proof}

The construction is an explicit finite algorithm. Each row is moved
one position at a time toward its specified target row. Each short
column is moved one position at a time along the specified arc until
it is adjacent to its target column. The number of intervening rows
or columns decreases at each such step, so the algorithm terminates.
The proof records the relevant noninterleaving test before each loop
and the three marked corners before each destabilization.

\section{Smallest cases and the role of the hypothesis}

The smallest allowed strand number is $m=2$. The shorter word must
contain its unique generator, so the shortest example is
$w=\sigma_0\sigma_0$, $w'=\sigma_0$.
Choose increasing occurrence order within each equal-height row group.
The initial six-column grid has
\[
 O=(2,0,4,3,1,5),\qquad X=(1,5,2,0,4,3).
\]
Deleting the first letter uses the switch in columns $1,2$.
The complete movie is recorded below. A tuple lists row values in
increasing column order, and all subsequent indices are renumbered
after each deletion.
\[
\begin{array}{c|c|c}
\text{operation}&O&X\\ \hline
\text{initial grid}&(2,0,4,3,1,5)&(1,5,2,0,4,3)\\
\text{switch }1,2&(2,0,4,3,1,5)&(1,2,5,0,4,3)\\
\text{delete row }2\text{ and column }0&(0,3,2,1,4)&(1,4,0,3,2)\\
\text{commute columns }0,1&(3,0,2,1,4)&(4,1,0,3,2)\\
\text{delete row }0\text{ and column }1&(2,1,0,3)&(3,0,2,1)\\
\text{commute columns }0,1&(1,2,0,3)&(0,3,2,1)\\
\text{delete row }2\text{ and column }1&(1,0,2)&(0,2,1).
\end{array}
\]
At the two commutations the vertical endpoint intervals are disjoint.
The three deletions are the $X$ destabilizations specified in the
general proof. The final grid is the three-column grid of
$\sigma_0$. Every row and every column in each displayed grid has one
marking of each type at distinct positions, which can also be checked
directly from the tuples. This is a finite convention check, not the
proof for arbitrary words.

The assumption on the shorter word matters to the validity of the
candidate switch. If one instead deletes the sole letter of
$w=\sigma_0$, then $O=(1,0,2)$ and $X=(0,2,1)$.
The candidate switch would give $X=(0,1,2)$, and column $2$ would
have coincident $O$ and $X$ markings. This excluded example explains
why the proof verifies $U_-\ne U_+$ before constructing $H$.

If the deleted letter is the first or last letter, the cyclic predecessor
or successor can wrap across the word boundary. The continuation lemma
uses the cyclic order and permits a cyclic column permutation, so the
same proof applies. If the shorter word has only one occurrence of
one of the two involved ranks, the previous and next occurrences
coincide, but their incoming and outgoing columns are distinct.
The empty fourth corner in the final destabilization remains empty.
No assumption on the number of link components is used.

\section{Exact missing native identifications}
\label{J10:sec:native-identifications}

The proved deletion movie requires only the following selected-grid and
primitive-move identifications. Their source declarations are absent from
the combined attachment. They are stated as unverified implications,
not as additional proved lemmas or silently assumed hypotheses.

\subsection{The precise selected-grid implication that is absent}

Fix a word $w$ of length $L$ and write
\[
 \begin{gathered}
 G_w=(\decl{positiveRoundedWordSelectedData}(w)).\mathrm{grid},\\
 C=\operatorname{Fin}(L)\times\operatorname{Fin}(3).
 \end{gathered}
\]
Let $\jmath:C\to\operatorname{Fin}(3L)$ be
\decl{finProdFinEquiv}, let $M_w$ be
\decl{positiveRectilinearWordMatching} $w$, and let
$h_w:C\to\mathbb R$ be
\decl{positiveRectilinearWordArcRow} $w$.
The exact data required to instantiate the proved deletion algorithm are
\begin{align}
 G_w.X(\jmath a)&=G_w.O(\jmath M_w^{-1}(a))
       &&\text{for every }a\in C,\label{J10:eq:j10-required-selected-incidence}\\
 h_w(a)<h_w(b)&\Longrightarrow
 G_w.O(\jmath a)<G_w.O(\jmath b)
       &&\text{for every }a,b\in C.\label{J10:eq:j10-required-selected-order}
\end{align}
No separate row-bijection hypothesis is required.
Indeed, $R(a)=G_w.O(\jmath a)$ is automatically a bijection because
$G_w.O$ and $\jmath$ are bijections. The two displayed properties say
exactly that this $R$ has the incoming marking equation and preserves
the strict unperturbed height order.

The reproduced theorem
\decl{exists_positiveRectilinearWordGrid} proves the existence of
some grid satisfying these properties. Its witness is a quantified
variable $G$, not the selected field $G_w$.
The supplied audit record L411 gives the type
\decl{positiveRoundedWordSelectedData} $w$ and states that its
structure contains a grid, rows, row-offset data, and profiles.
It does not reproduce a selector equation, a soundness theorem for
that grid field, or the propositions stored in the structure.
Thus L410 does not identify its existential witness with $G_w$.
The same-word independence of L412 concerns smooth representatives,
and L414 recovers a matching only under the marking-equation hypotheses.
Neither supplies \eqref{J10:eq:j10-required-selected-incidence} or
\eqref{J10:eq:j10-required-selected-order} for $G_w$.

The irreducible missing source is therefore the portion of
\decl{PositiveBraidGridSelectedRepresentative.lean} defining
\decl{PositiveRoundedWordRepresentativeData} and
\decl{positiveRoundedWordSelectedData}, or existing consequences
of those definitions that imply the two displayed properties.
The smallness of the perturbation, the smooth offset, the profiles,
and the one-cycle hypothesis of the smooth representative are not
needed by J10 once those two properties are known.

\subsection{The precise native move implications that are absent}

The supplied module
\decl{OnLeConjecture/External/TauCeti/KnotTheory/Grid/Move/Basic.lean}
imports
\decl{Grid.Commutation.Move} and
\decl{Grid.Stabilization.Basic}.
It proves that a hypothesis
\decl{RowsNoninterleaving} or
\decl{ColumnsNoninterleaving} yields a row or column move, and
that a hypothesis \decl{IsDestabilization} yields a move.
The combined source does not reproduce definitions of these hypotheses
or lemmas recognizing them from local marking inequalities.
The exact implications required by the written proof are the following.

\begin{enumerate}
\item For two cyclically adjacent rows, if their four marked columns
are distinct and their endpoint pairs do not alternate around the
column circle, then the native predicate
\decl{RowsNoninterleaving} holds. The analogous implication is
required for columns.

\item Suppose two adjacent rows and two adjacent columns have three
marked corners, with equal types at opposite corners and the other
type at a third corner. Suppose the fourth corner is empty.
Deleting the row and column through the unique-type corner and placing
the duplicated marking type at the fourth corner satisfies the native
predicate \decl{IsDestabilization}.
The proof uses an $O$ duplicated pattern in one possible first step,
and $X$ duplicated patterns in the other first step and the two
continuation steps. Both orders of the adjacent rows are allowed in
the continuation lemma.
\end{enumerate}

The move-module docstring explicitly says that its presentation follows
OSS Chapter~3. This is substantial evidence for the intended conventions.
It does not, as a matter of deduction from the supplied declarations,
provide the constructors or recognition implications above.
No equality of grid-link classes or smooth isotopy can replace those
missing implications in the literal acceptance statement.

\section{Lemma table and remaining verification}

\small
\begin{longtable}{p{0.57\textwidth}p{0.35\textwidth}}
\toprule
Declaration or lemma & Status and role\\
\midrule
\endhead
\decl{PositiveAdjacentWord}, \decl{PositiveAdjacentWord.toSigned}
& Existing Lean definitions. They fix the word and generator convention.\\
\decl{positiveRectilinearWordColumn_index_lt}
& Existing Lean declaration. It fixes the column order $3h+c$.\\
\decl{positiveRectilinearOutgoingArc},
\decl{positiveRectilinearIncomingArc},
\decl{positiveRectilinearHorizontalArcRow}
& Existing Lean definitions. They fix corner types and unperturbed heights.\\
\decl{positiveRectilinearWordMatching_middle},
\decl{positiveRectilinearWordMatching_occurrence},
\decl{positiveRectilinearSortedSuccessor_next}
& Existing Lean declarations. They identify the literal horizontal incidence.\\
\decl{positiveRectilinearIntegerRowIn_lt_out},
\decl{positiveRectilinearIntegerRow_arc_excludes_occurrence}
& Existing Lean declarations. They give endpoint separation at integer heights.\\
\decl{exists_positiveRectilinearWordGrid_of_injective},
\decl{exists_positiveRectilinearWordGrid}
& Existing Lean declarations. They give the admissible marking equations and row order.\\
\decl{XSwitchSequence.cons}, \decl{XSwitchSequence.nil}
& Existing Lean constructors. They give the sequence of exactly one switch.\\
\decl{GridDiagram.movesTo_swapRows},
\decl{GridDiagram.movesTo_swapColumns},
\decl{GridDiagram.cyclicColumns_movesTo},
\decl{GridDiagram.IsDestabilization.movesTo},
\decl{GridDiagram.MovesTo.trans}
& Existing Lean declarations. They assemble the elementary move sequence.\\
Lemmas~\ref{J10:lem:short}, \ref{J10:lem:bands}, and \ref{J10:lem:continuation}
& New lemmas. Their proofs verify every commutation and deletion used.\\
Lemma~\ref{J10:lem:switch}
& New lemma. It proves validity of the specified X switch with unchanged O markings.\\
Theorem~\ref{J10:thm:deletion}
& New theorem. It gives the full finite deletion algorithm for all admissible word grids.\\
\decl{positiveRoundedWordSelectedData}
& Existing object identified in the supplied audit. Its defining file is absent from the combined source attachment.\\
\bottomrule
\end{longtable}
\normalsize

The written combinatorial argument has no remaining parity or geometric
isotopy step. It proves the stated move theorem for the explicit word-grid
marking equations supplied in L410. The unresolved source check is the
literal instantiation in~\eqref{J10:eq:target}: the omitted definition of
\decl{positiveRoundedWordSelectedData} must supply those equations and
the row-order property, and the omitted primitive native move definitions
must agree with the elementary patterns used above. The audit records the selected object, and the move-module documentation
names the intended grid moves, but neither supplies the two selected-grid
equations or the primitive recognition implications isolated above.
The literal native instantiation therefore remains unverified from the
provided sources.
This document does not claim a Lean proof or completion of M702's
separate chain-map and canonical-class transport obligations.





\chapter{J15: Braid length, discriminants, and polar intersection}
\label{chapter:J15}



\section{Exact statements and scope}

The analytic rings in this document are the native rings of complex analytic
germs. The braid words are the native adjacent-generator words. No Floer
homology result or fibre-surface argument is used. The exponent of a signed
word is an integer, as required by the native type
\decl{SignedAdjacentWord}, whose signs belong to $\Z^\times=\{1,-1\}$.
This arithmetic is not a change of coefficient field in a chain complex.

\begin{theorem}[New lemma: M1155]
\label{J15:thm:discriminant}
Let $m\ge1$, and let
\[
 P(\xi,\zeta)=\zeta^m+a_1(\xi)\zeta^{m-1}+\cdots+a_m(\xi)
\]
have coefficients holomorphic near $0$, with $a_j(0)=0$ for every $j$.
Assume its discriminant $D(\xi)=\Disc_\zeta P(\xi,\zeta)$ is not
identically zero. For every sufficiently small $r>0$, the roots along the
positively oriented circle $\xi=r e^{\ii\theta}$ define a closed braid
of $m$ distinct sheets. For every signed event word representing this braid,
with the native positive-crossing convention,
\begin{equation}\label{J15:eq:main-discriminant}
 \Exp(w)=\ord_0 D=i_0(P,P_\zeta).
\end{equation}
Here $i_0$ is the native complex dimension of the analytic local pair quotient.
The word may be obtained by resolving simultaneous projected events.
\end{theorem}

The conditions $a_j(0)=0$ are the distinguished condition in the term
``Weierstrass polynomial.'' They are necessary for the last equality to be
an intersection number at the single point $(0,0)$. A counterexample without
that condition is given in Section~\ref{J15:sec:exceptions}.

\begin{theorem}[New lemma: reduced-germ length identity]
\label{J15:thm:reduced}
Let $f\in\C\{\xi,\zeta\}$ be reduced, with $f(0,0)=0$, and suppose
$f(0,\zeta)$ is not the zero germ. Put $m=i_0(f,\xi)$ and
\[
 \mu(f)=\dim_\C\C\{\xi,\zeta\}/(f_\xi,f_\zeta).
\]
For every signed word $w$ representing the root braid of the projection
$\xi$ on a sufficiently small positively oriented base circle,
\begin{equation}\label{J15:eq:exp-polar}
 \Exp(w)+1=\mu(f)+m.
\end{equation}
If $w$ is positive, then $\length(w)+1=\mu(f)+m$.
If $f$ has multiplicity $S$ and the projection is transverse, then $m=S$.
\end{theorem}

For a branch, the theorem has the following consequence in the exact form
needed for the existing analytic event words.

\begin{theorem}[New lemma: branch word and J15(a)]
\label{J15:thm:branch}
Let $f$ be an irreducible plane curve germ, let
$\beta:(\C,0)\longrightarrow(\C^2,0)$ have an injective analytic
representative with pullback kernel $(f)$, and suppose one coordinate of
$\beta$ has finite positive order $a$. Choose the positive event word $w$
and the normalized equation-link representative supplied by
\decl{PointedMapGerm.exists_equation_link_canonical_positive_braid}.
The choices can be made on a sufficiently small base circle so that
\begin{equation}\label{J15:eq:branch-length}
 \length(w)+1=2\delta(f)+a.
\end{equation}
The same word has one permutation cycle, uses every adjacent generator,
and has closure ambient isotopic to the normalized equation link with the
orientation in that existing declaration.

In particular, for a resolution level $L$ and one of its singular points $j$,
the existing data in
\decl{SingularSectionPointData.exists_branch_component_order} provide
$a$ with
\[
 2\le a\le (g.\decl{criticalJacobianFactorization}(hi)).
                 \decl{maxMultiplicity}.
\]
For that $a$ there is a native positive word satisfying all the word-level
conjuncts of the acceptance statement
\decl{ResolutionSectionLevel.exists_branch_word}, including
\eqref{J15:eq:branch-length} for the local section equation at $j$ and the link
identification for the same word.
\end{theorem}

For a branch of multiplicity $a$ with a transverse coordinate, the $a$ in
Theorem~\ref{J15:thm:branch} is that multiplicity. The theorem also allows a
finite projection of larger degree, which is useful because the catalogued
branch-component-order declaration specifies a coordinate order.

\begin{corollary}[New lemma: J15(b) for the J14 word]
\label{J15:cor:fulltwist}
Let $f$ be a reduced germ of multiplicity $S\ge1$ with $b$ irreducible
branches, and use a transverse projection. Suppose the word of J14 is
$d\,v$, where $d=\decl{fullTwistWord}(S)$ and $v$ is positive, and the
word represents the same root braid through closed $S$-braids in the
solid torus. Then
\begin{equation}\label{J15:eq:full-length}
 \length(d\,v)+1=\mu(f)+S,
\end{equation}
and
\begin{equation}\label{J15:eq:tail-length}
 \length(v)+S^2+b=2\delta(f)+2S.
\end{equation}
\end{corollary}

This corollary first proves the length of the actual analytic word.
Section~\ref{J15:sec:fulltwist-invariant} then proves the broader proposed
API in M407 for every positive full-twist word with an ambiently isotopic
closure, by an additional calculation over $\mathbb F_2$.

\section{Conventions, continuous arguments, and crossing signs}

For a signed native word
$w=((i_1,e_1),\ldots,(i_q,e_q))$, define
$\Exp(w)=\sum_{\nu=1}^q e_\nu$, where each $e_\nu$ is $1$ or $-1$.
The generator indices are $0,\ldots,m-2$, as in
\decl{SignedAdjacentWord} and \decl{adjacentSwap}.
A positive native word becomes a signed word by giving every letter sign
$1$. Therefore its exponent is its list length.

Fix a labelled lift $y_1(\theta),\ldots,y_m(\theta)$ of a braid over
$[b,b+2\pi]$. Labels may be permuted at the end. For a fixed nonzero
complex scalar $c$, projection means ordering the real parts of $c y_j$.
At a projected crossing of two sheets, put
$d=c(y_j-y_k)$. Then $\Re d=0$ and $\Im d\ne0$.
The positive crossing convention is
\begin{equation}\label{J15:eq:cross-sign}
 -\Im(d)\frac{d}{d\theta}\Re(d)>0.
\end{equation}
This is the positive angular-rotation convention in
\decl{exists_radius_planeBranchAngularDifference_projected_crossings_pos},
recorded in L4.16 of the supplied notes. Its analytic input, reproduced in
the supplied source, is
\decl{exists_radius_planeBranchSheetDifference_re_pos}.
It is consistent with the positive event words in
\decl{PointedMapGerm.exists_equation_link_canonical_positive_braid}.

\begin{lemma}[New lemma: continuous argument and its derivative]
\label{J15:lem:argument}
Let $d:[u,v]\longrightarrow\C\setminus\{0\}$ be continuous.
There is a continuous real function $\alpha$ with
$d/|d|=e^{\ii\alpha}$. If $d$ is continuously differentiable, then
$\alpha$ is continuously differentiable and
\begin{equation}\label{J15:eq:argument-derivative}
 \alpha'=\Im(d'/d).
\end{equation}
At a projected crossing, the sign in \eqref{J15:eq:cross-sign} equals the sign
of $\alpha'$.
\end{lemma}

\begin{proof}
The path has a positive minimum modulus. Uniform continuity permits a
finite subdivision into intervals on which the quotient by the value at
the left endpoint lies in $\{z:|z-1|<1\}$. On each such interval,
use the argument with values in $(-\pi/2,\pi/2)$ for that quotient,
and add the argument already chosen at the left endpoint. Patching these
finite choices gives a continuous argument on $[u,v]$.

If $d$ is continuously differentiable, choose an initial argument
$\alpha(u)$ and put
\[
 \alpha(\theta)=\alpha(u)+\int_u^\theta\Im(d'(s)/d(s))\,ds.
\]
Differentiating the unit complex number
$e^{-\ii\alpha(\theta)}d(\theta)/|d(\theta)|$ gives zero, because
its logarithmic derivative is
$d'/d-\Re(d'/d)-\ii\Im(d'/d)=0$.
Its initial value is $1$. This proves the asserted argument formula.
At a crossing $d=\ii h$ with $h\ne0$, one has
\[
 \alpha'=-\frac{\Re d'}h
 =\frac{-h\Re d'}{h^2}.
\]
The denominator is positive, which proves the sign assertion.
\end{proof}

The elementary two-strand positive half-turn can be written
$y_1(s)=-e^{\ii\pi s}$, $y_2(s)=e^{\ii\pi s}$ for $0\le s\le1$.
The argument of the difference increases by $\pi$, its squared difference
winds once positively, and \eqref{J15:eq:cross-sign} is positive at $s=1/2$.
Thus one positive letter contributes $+1$, rather than $-1$ or $2$, to
the winding of the squared-difference product.

For a closed nonzero complex loop $q$, define
$\wind(q)=(\alpha(v)-\alpha(u))/(2\pi)$ using Lemma~\ref{J15:lem:argument}.
This is an integer because $q(v)=q(u)$. The value is independent of the
initial argument. It is additive under multiplication and concatenation,
by addition of arguments.

\begin{lemma}[New lemma: homotopy invariance of the winding]
\label{J15:lem:homotopy}
The winding of a nonzero closed complex loop is constant under a continuous
homotopy through nonzero closed loops.
\end{lemma}

\begin{proof}
If two loops $q_0,q_1$ are sufficiently uniformly close relative to the
positive minimum of $|q_0|$, their pointwise quotient lies in the disc
$\{z:|z-1|<1\}$. That disc admits a single-valued continuous argument,
for example the argument in $(-\pi/2,\pi/2)$. This argument returns to
its initial value on the closed quotient loop. Thus the two windings agree.
For a homotopy on a compact parameter square, nonvanishing gives a positive
minimum modulus and uniform continuity. Partition the homotopy parameter
into finitely many intervals for which this closeness condition holds.
Equality on consecutive intervals proves the assertion.
\end{proof}

\section{The boundary terms of all unordered pairs cancel}

\begin{lemma}[New lemma: pairwise crossing count]
\label{J15:lem:pair-count}
Suppose $y_1,\ldots,y_m$ are continuously differentiable, pairwise distinct
on $[b,b+2\pi]$, and
\[
 y_j(b+2\pi)=y_{\tau(j)}(b)
\]
for a permutation $\tau$. Choose $c\ne0$ so that no projected coincidence
occurs at either endpoint and every pairwise projected crossing is
transverse. Count simultaneous events separately for each unordered pair.
The signed crossing count is
\begin{equation}\label{J15:eq:pair-winding}
 \sum_{j<k}\ \sum_{\substack{\theta\in(b,b+2\pi)\\
                 \Re(c(y_j-y_k))=0}}
  \operatorname{sgn}\!\left(-\Im(c(y_j-y_k))
       \frac d{d\theta}\Re(c(y_j-y_k))\right)
 =\wind\!\left(\prod_{j<k}(y_j-y_k)^2\right).
\end{equation}
\end{lemma}

\begin{proof}
Choose a continuous argument $\alpha_{jk}$ of $c(y_j-y_k)$ for each
$j<k$, using Lemma~\ref{J15:lem:argument}. Its crossing levels are
$\pi/2+\pi\Z$. At an upward crossing the integer function
\[
 F_{jk}(\theta)=
   \left\lfloor\frac{\alpha_{jk}(\theta)}\pi-\frac12\right\rfloor
\]
increases by one, and at a downward crossing it decreases by one.
It has no other jumps. The endpoints are not crossing levels, and the
crossings are finite and transverse. Telescoping its jumps gives the signed
count for this pair as
\[
 F_{jk}(b+2\pi)-F_{jk}(b).
\]
Let $\{u\}=u-\lfloor u\rfloor$ denote fractional part.
Consequently the count for this pair is
\begin{align*}
 \frac{\alpha_{jk}(b+2\pi)-\alpha_{jk}(b)}\pi
 &-\left\{\frac{\alpha_{jk}(b+2\pi)}\pi-\frac12\right\}\\
 &+\left\{\frac{\alpha_{jk}(b)}\pi-\frac12\right\}.
\end{align*}
The endpoint permutation permutes the unordered pairs. It may reverse a
chosen ordering of a pair, but reversal changes an argument by an integer
multiple of $\pi$. Therefore the multiset of the endpoint fractional parts
in this display is the same as the multiset of the initial fractional parts.
They cancel after summation over all unordered pairs. The total count is
\[
 \frac1\pi\sum_{j<k}
       \bigl(\alpha_{jk}(b+2\pi)-\alpha_{jk}(b)\bigr).
\]
A continuous argument of the squared-difference product is
$2\sum_{j<k}\alpha_{jk}$, up to a constant argument for
$c^{-m(m-1)}$. Dividing its total change by $2\pi$ proves
\eqref{J15:eq:pair-winding}. This also proves that the total is independent
of the starting labels and the nonzero constant $c$.
\end{proof}

\begin{lemma}[New lemma: generic projection and multiple events]
\label{J15:lem:events}
For a real-analytic collision-free braid on a compact period, a constant
projection direction can be chosen so that all pairwise projected crossings
are transverse and the endpoints are regular. Simultaneous crossings may
be separated by an arbitrarily small collision-free perturbation supported
away from the endpoints. The exponent of a word for the perturbed braid
is the pairwise signed count in Lemma~\ref{J15:lem:pair-count}.
If every pairwise crossing is positive, its positive event word has length
equal to the sum of the numbers of coincident unordered pairs at all events.
\end{lemma}

\begin{proof}
For each pair, the derivative of its argument is real analytic. Either it
vanishes identically or its zero set on a period is finite. In the first
case the argument is constant, and one excludes projection directions
making that pair project to the same real coordinate. In the second case,
one excludes the finitely many directions that make a critical argument
a crossing level. Also exclude the finitely many directions producing an
endpoint coincidence. There are only finitely many pairs, so a direction
outside this finite set exists. Every crossing is then simple, and there
are finitely many crossings on the compact period.

Take disjoint small intervals around the finitely many event times.
At an event, add small, separately chosen real offsets to the projected
coordinates of the labelled sheets, constant near that event and zero near
the interval endpoints. A simple zero of the difference of two projected
coordinates persists with its sign. The first variation of its time is
\[
 -\frac{h_j-h_k}{(\Re(c(y_j-y_k)))'(\theta_0)},
\]
where $h_j,h_k$ are the chosen offsets. For two distinct unordered pairs,
the corresponding linear forms in the offsets are different. Their
coincidence excludes a proper hyperplane of offset choices. Choose one direction $(h_1,\ldots,h_m)$ outside the finitely many
such hyperplanes and use offsets $\tau h_j$. The first derivatives
at $\tau=0$ of the competing crossing times are then different.
For every sufficiently small $\tau>0$, those times are distinct,
so all crossings in an event interval are separated. The positive minimum
separation of the original complex sheets ensures that these perturbations
and their straight-line interpolation create no complex collision.
They produce one ordinary adjacent crossing for each original coincident
pair, with the same sign.

Reading the real order between successive simple events gives a signed
adjacent-generator word. Its exponent is the sum of these signs by
definition. If all signs are positive, its length is the number of pairs
counted with multiplicity over the events. The native branch theorem
already supplies exactly this pair-count length formula for its chosen
event words, so its event resolution does not need to be replaced.
\end{proof}

\begin{corollary}[New lemma: exponent is preserved by braid homotopy]
\label{J15:cor:exp-homotopy}
The exponent in Lemma~\ref{J15:lem:events} is unchanged under a homotopy of
closed braids through configurations of distinct complex points.
\end{corollary}

\begin{proof}
The squared-difference product is symmetric in the sheets, so it is a
well-defined nonzero loop for an unordered configuration braid. A
collision-free homotopy induces a nonzero homotopy of these products.
Apply Lemmas~\ref{J15:lem:pair-count} and \ref{J15:lem:homotopy}.
\end{proof}

\section{The local argument principle and distinguished preparation}

\begin{lemma}[New lemma: local argument principle]
\label{J15:lem:local-principle}
Let $h$ be a holomorphic germ at $0\in\C$ which is not zero, and put
$d=\ord_0h$. For every sufficiently small $r>0$,
\[
 \wind\bigl(\theta\longmapsto h(r e^{\ii\theta})\bigr)=d.
\]
More generally, if a holomorphic function has no zero on the boundary of a
closed disc, the winding of its boundary values is the sum of the orders
of its interior zeros.
\end{lemma}

\begin{proof}
The Taylor series factors as $h(\xi)=\xi^d u(\xi)$, with $u(0)\ne0$.
Shrink the disc until $u$ is nonzero there. The holomorphic function $u'/u$
has a primitive on this disc, obtained by integrating its power series.
If $L'=u'/u$, then differentiation gives $(u e^{-L})'=0$.
Thus $u$ is a nonzero constant times $e^L$, and its argument has total
change zero on every concentric circle. The factor $\xi^d$ contributes
$2\pi d$, which proves the first assertion.

For the second assertion, there are finitely many zeros in the disc,
since they are isolated and there are none on its boundary. Factor their
linear factors with their orders from the function. The remaining
holomorphic function is nonzero on the disc and has a logarithm by the
same primitive argument. Each factor $\xi-a$ with $a$ in the interior
has winding one. This follows by moving $a$ to the centre through the
disc and applying Lemma~\ref{J15:lem:homotopy}. Multiplication adds the
windings, proving the stated sum.
\end{proof}

\begin{lemma}[New lemma: preparation in the projection coordinate]
\label{J15:lem:preparation}
Let $f(\xi,\zeta)$ be holomorphic near $(0,0)$, with $f(0,0)=0$
and $f(0,\zeta)$ nonzero of order $m\ge1$ at $\zeta=0$.
There are holomorphic germs $a_1,\ldots,a_m$, all vanishing at zero,
and a holomorphic unit $u$ such that
\begin{equation}\label{J15:eq:preparation}
 f(\xi,\zeta)=u(\xi,\zeta)
       \bigl(\zeta^m+a_1(\xi)\zeta^{m-1}+\cdots+a_m(\xi)\bigr).
\end{equation}
Moreover $i_0(f,\xi)=m$.
\end{lemma}

\begin{proof}
Choose $R>0$ so that $f(0,\zeta)$ has no zero other than $0$ in
$|\zeta|\le R$ and has no zero on $|\zeta|=R$.
Uniform continuity on that circle permits a base disc in which
$f(\xi,\zeta)$ has no boundary zero. Its boundary winding is constant
as $\xi$ varies in that disc, by Lemma~\ref{J15:lem:homotopy}, and equals
$m$ at $\xi=0$. Lemma~\ref{J15:lem:local-principle} gives exactly $m$
interior zeros counted with multiplicity for each such $\xi$.

Their power sums are holomorphic functions of $\xi$. Indeed, for
$q=1,\ldots,m$, the power sum is
\[
 p_q(\xi)=\frac1{2\pi\ii}
    \int_{|v|=R}v^q\frac{f_\zeta(\xi,v)}{f(\xi,v)}\,dv.
\]
The contour identities needed here follow directly from power series.
If $g$ is holomorphic on a neighbourhood of $|v|\le R$, expand $g$
on a slightly larger disc and integrate the series term by term on
$|v|=R$. Expanding $(v-z)^{-1}=\sum_{k\ge0}z^k v^{-k-1}$ for
$|z|<R$ gives
$\frac1{2\pi\ii}\int g(v)/(v-z)\,dv=g(z)$.
Termwise integration also gives zero for the integral of a holomorphic
function without this denominator. These are all the contour formulas
used below.
For a fixed $\xi$, factor the interior zeros as in the proof of
Lemma~\ref{J15:lem:local-principle}. Each logarithmic derivative term
$\nu/(v-v_0)$ contributes $\nu v_0^q$ by the power-series form of
Cauchy's formula. The nonzero factor contributes zero.
The integrand is holomorphic in $\xi$ on a neighbourhood of the contour,
so differentiating under this compact contour integral proves holomorphy
of $p_q$.

Newton's identities recursively define the elementary symmetric functions
$e_q(\xi)$ by
\[
 q e_q=\sum_{j=1}^q(-1)^{j-1}e_{q-j}p_j,\qquad e_0=1.
\]
They are holomorphic. The monic polynomial with coefficients
$a_q=(-1)^q e_q$ has precisely these interior zeros, with their
multiplicities. At $\xi=0$ all of them are zero, so $a_q(0)=0$.
Call this polynomial $P$.

For fixed $\xi$, the quotient $f/P$ has removable singularities at
all these roots and is holomorphic and nonzero on $|\zeta|<R$.
Its jointly holomorphic representative is the contour integral
\[
 u(\xi,\zeta)=\frac1{2\pi\ii}
  \int_{|v|=R}\frac{f(\xi,v)}{P(\xi,v)(v-\zeta)}\,dv.
\]
Cauchy's formula identifies this with the quotient for every fixed
$\xi$. Differentiation under the integral proves joint holomorphy.
The quotient has no zero, so $u(0,0)\ne0$, and \eqref{J15:eq:preparation}
follows.
Finally, the quotient by $(f,\xi)$ is isomorphic to
$\C\{\zeta\}/(\zeta^m)$, with basis
$1,\zeta,\ldots,\zeta^{m-1}$. Its dimension is $m$.
\end{proof}

This proof of the required preparation statement is included because the
supplied file imports the project's preparation modules without reproducing
their exact declarations. It fixes the distinguished hypothesis explicitly.

\section{Proof of the discriminant and polar equalities}

\begin{proof}[Proof of Theorem~\ref{J15:thm:discriminant}]
Factor $D(\xi)=\xi^d v(\xi)$ with $v(0)\ne0$.
For every sufficiently small $r>0$, the polynomial has distinct roots for
all $|\xi|=r$, since its discriminant does not vanish there.
A simple polynomial root varies holomorphically locally in $\xi$ by the
analytic implicit function theorem. Applying this fact successively along
the compact circle gives labelled roots on $[b,b+2\pi]$. They cannot
escape to infinity because they are roots of a monic polynomial with
bounded coefficients. Their endpoints are permuted because the polynomial
at the two endpoints is the same.

For completeness, the simple-root fact used here follows by writing a
root equation near $(\xi_0,\zeta_0)$ as
$\zeta-\zeta_0=T(\xi,\zeta-\zeta_0)$ with
$\partial_\zeta T(\xi_0,0)=0$. On a sufficiently small polydisc this
is a contraction in the second variable. Iteration from zero converges
uniformly on a smaller polydisc to its unique holomorphic fixed point.
This proves the local analytic root assertion needed above.

For a monic polynomial, the algebraic discriminant identity is
\begin{equation}\label{J15:eq:disc-product}
 D(\xi)=\prod_{j<k}(y_j(\theta)-y_k(\theta))^2,
 \qquad \xi=r e^{\ii\theta}.
\end{equation}
One can verify this identity by differentiating
$P(\xi,\zeta)=\prod_j(\zeta-y_j)$ and taking the product of
$P_\zeta(\xi,y_j)$ over all roots. Each unordered pair contributes
the negative of its squared difference. Thus
\begin{equation}\label{J15:eq:resultant}
 \Res_\zeta(P,P_\zeta)=(-1)^{m(m-1)/2}D.
\end{equation}
This is also the monic specialization of the existing algebraic declaration
\decl{Polynomial.resultant_deriv}, cited in M407 and M1155.
The standard definition of the discriminant fixes the sign in
\eqref{J15:eq:disc-product}.

Lemmas~\ref{J15:lem:pair-count} and \ref{J15:lem:events} identify the exponent
with the winding of the right side of \eqref{J15:eq:disc-product}.
Lemma~\ref{J15:lem:local-principle} makes this winding $d$.
Corollary~\ref{J15:cor:exp-homotopy} permits any event resolution representing
the same braid. This proves the first equality in
\eqref{J15:eq:main-discriminant}, including its sign and its factor of two.

For the second equality, take the polynomial degree parameter $d_0=m-1$
in the supplied declarations from \decl{IntersectionResultant.lean}.
The polynomial $P$ is monic of degree $d_0+1$, and every lower coefficient
belongs to the maximal ideal, by the hypotheses. Equation~\eqref{J15:eq:resultant}
and $D\ne0$ make the resultant nonzero. Therefore
\decl{Germ.hasFiniteIntersection_of_resultant_ne_zero} supplies finite
intersection, and
\decl{Germ.intersectionMultiplicity_eq_resultant_order_of_ne_zero}
gives
\[
 i_0(P,P_\zeta)=\ord_0\Res_\zeta(P,P_\zeta)=\ord_0 D.
\]
The last equality uses the nonzero constant in \eqref{J15:eq:resultant}.
When applying these native declarations, the polynomial variable is native
coordinate $0$ and the coefficient variable is coordinate $1$.
If the named coordinates are in the other order, exchange the coordinates,
which identifies the two pair quotients and preserves their dimensions.
\end{proof}

\begin{lemma}[New lemma: the unit factor preserves the polar pair ideal]
\label{J15:lem:unit}
If $f=uP$ with $u$ a holomorphic unit, then
\[
 (f,f_\zeta)=(P,P_\zeta)
\]
in the local analytic ring. In particular their pair quotient dimensions
agree whenever either one is finite.
\end{lemma}

\begin{proof}
The product rule gives $f_\zeta=u_\zeta P+uP_\zeta$.
Because $u$ is invertible, $(f)=(P)$. Modulo this common principal ideal,
$f_\zeta$ and $P_\zeta$ differ by multiplication by $u$.
Each of the two displayed ideals therefore contains the generators of the
other, proving equality.
\end{proof}

\begin{proof}[Proof of Theorem~\ref{J15:thm:reduced}]
Use Lemma~\ref{J15:lem:preparation} to write $f=uP$ with $P$ distinguished
monic of degree $m=i_0(f,\xi)$. The curve is reduced and pointed, and
$(f,\xi)$ has finite colength by that lemma. These are exactly the
hypotheses of the supplied theorem
\decl{PlaneCurve.exists_polar_intersection_identity} with parameter
$\xi$. Its polar is $f_\zeta$ or $-f_\zeta$, according to the coordinate
order. This sign is a unit and leaves the pair ideal unchanged. That existing
theorem supplies finite polar intersection and the identity
\begin{equation}\label{J15:eq:teissier}
 i_0(f,f_\zeta)+1=\mu(f)+i_0(f,\xi)=\mu(f)+m.
\end{equation}
Lemma~\ref{J15:lem:unit} gives finite intersection for $(P,P_\zeta)$.
The supplied \decl{Germ.resultant_ne_zero_of_hasFiniteIntersection}
then makes its resultant nonzero. Equation~\eqref{J15:eq:resultant} makes
$D$ nonzero, so Theorem~\ref{J15:thm:discriminant} applies.
The roots of $P$ are the sheets of $f=0$, since $u$ does not vanish.
Consequently
\[
 \Exp(w)=\ord_0 D=i_0(P,P_\zeta)=i_0(f,f_\zeta).
\]
Combining with \eqref{J15:eq:teissier} proves \eqref{J15:eq:exp-polar}.
For a positive word its exponent equals its length.

Finally, let $S$ be the order of $f$, and let $f_S$ be its first nonzero
homogeneous part. Transversality of the projection means
$f_S(0,1)\ne0$ in the named coordinates $(\xi,\zeta)$.
Thus $f(0,\zeta)$ has exact order $S$, and
Lemma~\ref{J15:lem:preparation} gives $m=S$.
\end{proof}

\section{The branch word and the exact resolution-level conclusion}

\begin{lemma}[New lemma: the normalized branch has exactly the required degree]
\label{J15:lem:branch-degree}
Let $f$ be irreducible, let an injective analytic representative $F$ have
pullback kernel $(f)$, and suppose its projection coordinate is
$\xi(F(t))=t^a$ with $a\ge1$. Then $i_0(f,\xi)=a$.
For every sufficiently small nonzero $\xi$, the roots of its distinguished
polynomial are exactly the complementary coordinates of
$F(t)$ with $t^a=\xi$, each occurring once.
\end{lemma}

\begin{proof}
The germ $f(0,\zeta)$ cannot be identically zero. Otherwise the Taylor
series would make $\xi$ divide $f$. Irreducibility would make $f$
associated to $\xi$, and its pullback would force $t^a=0$, a contradiction.
Prepare $f=uP$ by Lemma~\ref{J15:lem:preparation}, and let the degree of $P$
be $M=i_0(f,\xi)$.

Write the complementary coordinate of $F(t)$ as $Y(t)$. The polynomial
\[
 Q(t^a,\zeta)=\prod_{\omega^a=1}(\zeta-Y(\omega t))
\]
has coefficients invariant under $t\mapsto\omega t$.
In a convergent power series, this invariance makes every nonzero
coefficient have exponent divisible by $a$. Thus the coefficients are
convergent power series in $\xi=t^a$. They vanish at $\xi=0$.
The polynomial $Q(\xi,\zeta)$ is distinguished monic of degree $a$.
Its pullback along $F$ is zero, so the kernel hypothesis gives $Q=fh$
for an analytic germ $h$.

For small nonzero $\xi$, the $a$ points $F(t)$ with $t^a=\xi$ are
distinct by injectivity. Their complementary coordinates are distinct
because their projection coordinates agree. They are roots of $P$, so
$M\ge a$.
On the other hand, the identity $Q=uPh$ shows that every zero of $P$
in a sufficiently small fibre disc is a zero of $Q$ with at least the
same multiplicity. All $M$ roots of $P$ and all $a$ roots of $Q$ lie
in this disc after shrinking the base. Hence $M\le a$.
It follows that $M=a$, and the $a$ distinct displayed roots exhaust
those of $P$. This also proves the intersection assertion.
\end{proof}

\begin{proof}[Proof of Theorem~\ref{J15:thm:branch}]
Apply the supplied normalization and positive-braid construction to
$\beta$ and its chosen coordinate. The resulting representative $F$
satisfies $\xi(F(t))=t^a$, is injective, and has the same equation
kernel after the source reparametrization. Lemma~\ref{J15:lem:branch-degree}
identifies its sheet braid with the roots of the distinguished polynomial
of $f$, and gives $i_0(f,\xi)=a$. The product polynomial $Q$ constructed
in that lemma is holomorphic for every $|\xi|<R^a$ on which the chosen
representative is defined and injective. It has distinct roots for
$0<|\xi|<R^a$. Thus its discriminant has no zero in that punctured
disc. The local-order winding formula applies to every circle there,
including the radius selected by the existing link construction.

The positivity agrees with the convention used above. If $h(t)$ is a
difference of two parameter sheets, the supplied declaration
\decl{planeBranchSheetDifference_exists_positive_order} and its
radius consequence give
$\Re(t h'(t)/h(t))>0$ for every sufficiently small nonzero $t$.
Along the base angle one has $t=r_0e^{\ii\theta/a}$, up to a root of
unity, and hence
\[
 \frac d{d\theta}\arg h(t)
   =\frac1a\Re\left(\frac{t h'(t)}{h(t)}\right)>0.
\]
Multiplication by the constant scale used for the native sheets does not
alter this derivative. Lemma~\ref{J15:lem:argument} therefore makes every
projected pair event positive. A sufficiently small common radius exists
because only finitely many roots of unity and sheet pairs occur.

Retain the regular base angle and positive event words in
\decl{PointedMapGerm.exists_equation_link_canonical_positive_braid}.
Denote their flattened word by $w$. The positive projected-crossing
input for these same event data is the existing
\decl{PointedMapGerm.exists_equation_link_finite_braid_projection},
quoted as L404 in the supplied audit. The angular calculation above
checks its sign convention. No different word or radius is substituted
for the chosen native word.
The canonical-word declaration states exactly that its length is the sum, over
all event times, of the numbers of coincident unordered pairs. The use of
final-position labels in that formula does not change the number of
unordered pairs, since each sorting is a permutation.
Lemma~\ref{J15:lem:pair-count} applies, and the positivity just established
gives
\[
 \length(w)=\ord_0 D=i_0(f,f_\zeta)=\mu(f)+a-1.
\]
Equivalently, the same conclusion follows from Theorem~\ref{J15:thm:reduced}.

The supplied Milnor identity
\decl{PlaneCurve.jacobianColength_add_irreducibleBranchCount}
applies because $f$ is pointed and irreducible, hence reduced.
An irreducible equation has one branch, so it gives
$\mu(f)=2\delta(f)$. Substitution proves \eqref{J15:eq:branch-length}.
The endpoint permutation equations are unchanged because the word and
the event data have not been replaced. Apply
\decl{braid_word_cycleCount_eq_one_of_endpoint_rotation} and
\decl{braid_word_generatorIndices_eq_univ_of_endpoint_rotation}.
They give one cycle and all adjacent generators for this same $w$.
The normalized link and its diffeotopy to the canonical word model are
also the ones in the existing construction. No new link identification
is inferred from the numerical equality.

For the resolution-level assertion, fix $j$. The supplied notes record
that \decl{ResolutionSectionLevel.singular} supplies
\decl{SingularSectionPointData} at $L.\decl{point}(j)$ with upper bound
$T=(g.\decl{criticalJacobianFactorization}(hi)).\decl{maxMultiplicity}$.
The quoted existing declaration
\decl{SingularSectionPointData.exists_branch_component_order}
provides the injective parameter germ, the local equation kernel, the
coordinate index, and $2\le a\le T$. Use those data in the result just
proved. It gives a word on these same $a$ strands with the required length,
cycle count, generator set, and normalized link. The inequalities for $a$
are preserved from that declaration. The defining modules of the resolution
level and singular-point data are catalogued in the attachment rather than
reproduced there. This last assembly uses their quoted statements and does
not assert a new unverified size bound.
\end{proof}

\section{The full-twist word}

\begin{proof}[Proof of Corollary~\ref{J15:cor:fulltwist}]
The multiplicity and transversality hypotheses give $m=S$ in
Theorem~\ref{J15:thm:reduced}. The J14 word represents the same root braid,
and its letters are all positive. Consequently \eqref{J15:eq:full-length}
follows from \eqref{J15:eq:exp-polar} and
Corollary~\ref{J15:cor:exp-homotopy}.

For the full rotation of a configuration of $S$ points, each difference
is multiplied by $e^{\ii\theta}$ over one turn. Its squared-difference
product is multiplied by $e^{\ii S(S-1)\theta}$. Thus the full twist
has exponent $S(S-1)$ by Lemma~\ref{J15:lem:pair-count}.
The positive word $d$ representing that rotation in J14 therefore has
length $S(S-1)$. Subtracting this length from \eqref{J15:eq:full-length}
gives
\[
 \length(v)=\mu(f)-S^2+2S-1.
\]
The supplied Milnor identity gives $\mu(f)+b=2\delta(f)+1$.
Substitution yields \eqref{J15:eq:tail-length}.
\end{proof}

There is also a direct check on the strict-transform word used in J14.
If $u_j=y_j/\xi$, then on the base circle
\[
 \prod_{j<k}(u_j-u_k)^2
   =\xi^{-S(S-1)}\prod_{j<k}(y_j-y_k)^2.
\]
Hence the strict-transform braid has exponent
$\ord_0D-S(S-1)$. This equals the tail length because J14 makes that
word positive. Distinct tangent clusters cause no difficulty in this
identity. The product is over all unordered pairs, including pairs from
different clusters, and the endpoint permutation still permutes those
pairs.

\section{The broader link-isotopy quantifier in M407}
\label{J15:sec:fulltwist-invariant}

The discriminant determines the exponent of the analytic root braid.
The proposed M407 API also permits a positive word known only through
an ambient isotopy of its closure. The following algebraic argument
supplies the additional comparison. Every coefficient in its modules
and matrices belongs to $\mathbb F_2$. The indeterminate $T$ below is
an auxiliary variable for an infinite cyclic covering. It is unrelated
to the Floer variable $U$.

\begin{lemma}[New lemma: the elementary braid matrices]
\label{J15:lem:burau-elementary}
Let $s\geq2$, put $R=\mathbb F_2[T,T^{-1}]$, and let $B_i$ be the
$s$ by $s$ identity matrix with its rows and columns $i,i+1$ replaced
by the block
\[
 \begin{pmatrix}1+T&T\\1&0\end{pmatrix},
 \qquad 0\leq i\leq s-2.
\]
Write $v=(1,\ldots,1)^{\mathsf t}$ and
$q=(1,T,\ldots,T^{s-1})$.
The matrices satisfy the positive braid relations, fix $v$, preserve
the row vector $q$, and induce matrices $\overline B_i$ on the free
quotient $R^s/Rv$ of rank $s-1$. They have the following properties.
\begin{enumerate}
\item $\det\overline B_i=T$.
\item Every minor of every $\overline B_i$ has degree at most one
as a polynomial in $T$.
\item If $c_s=\sigma_0\cdots\sigma_{s-2}$ and $d_s=c_s^s$,
then $\overline B_{d_s}=T^sI$.
\item For a positive word $w$ of length $l$, every minor of
$\overline B_w$ has degree at most $l$, and
$\det\overline B_w=T^l$.
\end{enumerate}
\end{lemma}
\begin{proof}
The determinant of the displayed block is $T$, which is invertible
in $R$. Blocks on disjoint pairs commute. On three successive
coordinates, the two products in the adjacent braid relation both
equal
\[
 \begin{pmatrix}
  1+T&T+T^2&T^2\\
  1+T&T&0\\
  1&0&0
 \end{pmatrix}.
\]
This verifies all relations by matrix multiplication.
The identities $B_iv=v$ and $qB_i=q$ follow on the same two
coordinates from
\[
 \begin{pmatrix}1+T&T\\1&0\end{pmatrix}
 \binom11=\binom11,
 \qquad
 (1,T)\begin{pmatrix}1+T&T\\1&0\end{pmatrix}=(1,T).
\]
The quotient is free, with basis the classes of
$e_0,\ldots,e_{s-2}$, since the last basis vector equals their sum
there. The invariant line has eigenvalue one, so the determinant on
the quotient equals the determinant on the whole module, namely $T$.

On $R^s$, the difference $B_i-I$ is a rank-one matrix
$u_i(T)\lambda_i$, where $u_i(T)$ has degree at most one and
$\lambda_i$ is a constant row vector supported in coordinates
$i,i+1$. Since $\lambda_i(v)=0$, it descends to the quotient with
the same property. In the multilinear expansion of any minor of
$I+u_i(T)\lambda_i$, a term using the rank-one summand in two
columns vanishes because those columns are proportional.
Every surviving term uses it at most once. Thus every minor has
degree at most one. The Cauchy--Binet formula for a product expresses
each minor as a sum of products of minors of the factors.
Induction over the letters of $w$ gives the degree bound $l$.
Multiplying determinants gives $T^l$.

The product $C=B_0\cdots B_{s-2}$ has first row
\[
 (1+T,\ T(1+T),\ldots,T^{s-2}(1+T),\ T^{s-1})
\]
and has ones immediately below the diagonal, with all its other
entries zero. Expansion of $\det(\lambda I-C)$ along its last
column, followed by induction on $s$, gives
\[
 \det(\lambda I-C)
  =(\lambda+1)
     (\lambda^{s-1}+T\lambda^{s-2}+\cdots+T^{s-1}).
\]
The invariant line contributes $\lambda+1$, so the second factor
is the characteristic polynomial of $\overline C$.
The Cayley--Hamilton identity and multiplication by
$\overline C+TI$ give
\[
 \overline C^{,s}+T^sI=0.
\]
Over $\mathbb F_2$ this says
$\overline B_{d_s}=\overline C^{,s}=T^sI$.
This argument also covers even $s$, when the characteristic
polynomial can have repeated roots.
\end{proof}

\begin{lemma}[New lemma: a closed-braid determinant is a link invariant]
\label{J15:lem:alexander-determinant}
Let $\beta$ be an $s$-strand braid, with $s\geq2$, and suppose
\[
 P_\beta(T)=\det(I+\overline B_\beta(T))\ne0.
\]
Then
\begin{equation}\label{J15:eq:alexander-quotient}
 A_{\widehat\beta}(T)
   =\frac{P_\beta(T)}{1+T+\cdots+T^{s-1}}
\end{equation}
is a nonzero Laurent polynomial, defined up to multiplication by a
power of $T$ as an invariant of the oriented closure.
Its breadth, meaning its highest exponent minus its lowest exponent,
is unchanged by an orientation-preserving ambient isotopy.
It is also unchanged if every component orientation is reversed.
\end{lemma}
\begin{proof}
We prove the formula from a cellular presentation, so no property of
a fibre surface is used. Let $X$ be the complement of the closed
braid and let $\widetilde X$ be the infinite cyclic covering
specified by sending each oriented meridian to $1\in\mathbb Z$.
The homology $H_1(\widetilde X,\mathbb F_2)$ is an $R$-module,
where $T$ is the deck transformation. An oriented link isotopy
identifies these coverings and their modules. Simultaneously reversing
all component orientations replaces $T$ by $T^{-1}$, which preserves
the breadth of every Laurent polynomial.

Here is a direct presentation of the module. Choose successive
meridians $x_0,\ldots,x_{s-1}$ in a transverse punctured disc.
Transport across a positive adjacent crossing replaces
\[
 x_i\longmapsto x_ix_{i+1}x_i^{-1},\qquad
 x_{i+1}\longmapsto x_i.
\]
This is obtained by carrying the two small meridian loops through
the half twist while keeping their paths to the boundary base point.
One loop acquires the path around $x_i$ and its inverse, and the
other acquires the original path of $x_i$.
For the auxiliary module calculation, fix this transport convention.
There is no sign ambiguity in the length conclusion. If the native
positive half twist is drawn with the inverse meridian transport,
all its matrices are $B_i(T)^{-1}$. Let $J$ reverse the order of the
$s$ basis vectors. Direct inversion of the two by two block gives
\[
 J B_i(T)^{-1}J=B_{s-2-i}(T^{-1}).
\]
Thus reversing the meridian order and the choice of deck generator
reduces that convention to the displayed polynomial matrices as well.
These changes preserve the module and its breadth. They replace the
Coxeter product by its reversed product, which has the same full-twist
scalar as checked below. The analytic exponent sign in the preceding
discriminant proof is still the fixed native positive sign.

If $\phi_\beta$ is the resulting automorphism of the free group,
identifying the top and bottom meridians in the closed diagram gives
the presentation
\[
 \pi_1(X)=\left\langle x_0,\ldots,x_{s-1}\ \middle|\
       \phi_\beta(x_i)x_i^{-1}=1\ (0\leq i<s)\right\rangle.
\]
To justify the presentation directly, apply the meridian presentation
of each crossing in succession. Its relation eliminates the outgoing
meridian in terms of the incoming ones by the displayed conjugation.
After all crossings have been processed, closure identifies each
outgoing meridian with the corresponding incoming meridian.
There are no other relations in the diagram complement, by gluing
the crossing neighborhoods and their arc neighborhoods along discs.

The boundary coefficients for the cyclic covering of this
presentation have an elementary word description. A positive letter
$x_j$ after a prefix of exponent $a$ contributes $T^a$ in the
$j$th edge coordinate. A negative letter $x_j^{-1}$ after that
prefix contributes $T^{a-1}$ in that coordinate, since signs are
equal in $\mathbb F_2$. These are the lifted directed edges of the
word path. For a concatenation, the coefficients of the second word
are multiplied by the exponent of the first. This rule is the
abelianized derivative rule
\[
 \partial_j(uv)=\partial_j u+T^{\exp u}\partial_j v.
\]
Applied to the two transported meridians above, it gives exactly the
block $B_i$. Applied successively to a braid word, it gives
$B_\beta$, with matrix multiplication in the order of this transport
convention. If word composition is recorded in the opposite order,
use the reversed matrix product throughout. The minor degree bound
is valid in either order. The reversed Coxeter product has the same
characteristic polynomial: with
$E=\operatorname{diag}(1,T,\ldots,T^{s-1})$, every generator satisfies
$B_i^{\mathsf t}=EB_iE^{-1}$, so that product is conjugate to the
transpose of the original Coxeter product. It also fixes $v$. The
same quotient characteristic polynomial and Cayley--Hamilton argument
therefore give the same scalar $T^sI$ for its full twist. No comparison
of arbitrary word orders is required.

Each relation $\phi_\beta(x_i)x_i^{-1}$ therefore has lifted
boundary equal to row $i$ of $B_\beta-I$.
The boundary of each one-cell is $T-1$, so the one-cycles form
\[
 K=\{(a_0,\ldots,a_{s-1})\in R^s:\ \sum_i a_i=0\}.
\]
Here $T-1\ne0$ and $R$ has no zero divisors.
Deleting the last coordinate identifies $K$ with $R^{s-1}$.
Thus the module is presented by the rows of
$A=B_\beta-I$ with its last column deleted. This computation of
$H_1$ depends only on the fundamental group and its homomorphism to
$\mathbb Z$: the covering fundamental group is the kernel, and
its abelianization with $\mathbb F_2$ coefficients is $H_1$.
Consequently using the displayed presentation instead of a full
cell decomposition of $X$ does not change the module.

Work temporarily over the field $\mathbb F_2(T)$.
The hypothesis $P_\beta\ne0$ says that the induced operator
$I-\overline B_\beta$ on $R^s/Rv$ is invertible over this field.
Hence $I-B_\beta$ has rank $s-1$, right kernel spanned by $v$,
and left kernel spanned by $q$. The last statement follows from
$qB_\beta=q$ and the rank. Its adjugate therefore has the form
\[
 \operatorname{adj}(I-B_\beta)=a(T)vq.
\]
Since the first entry of $q$ and every entry of $v$ equal one,
the scalar $a(T)$ is one of the adjugate entries and belongs to
$R$. The maximal minors in the module presentation just described
are its last-row adjugate entries
$a(T),a(T)T,\ldots,a(T)T^{s-1}$.
Their greatest common divisor in the principal ideal domain $R$
is $a(T)$, up to a unit. It is nonzero.

For completeness, this greatest common divisor is an isomorphism
invariant of the presented torsion module. Euclidean row and column
operations diagonalize a full-rank presentation over $R$.
The diagonal entries present cyclic summands, and the product of
the nonunit diagonal entries generates the ideal of maximal minors.
For each irreducible polynomial, the exponents in those summands
are recovered from the dimensions of the successive quotients by
its powers. The product is therefore independent of the chosen
presentation. The only units of $R$ are powers of $T$.

It remains to compute $a(T)$. The characteristic polynomial satisfies
\[
 \det(\lambda I-B_\beta)
   =(\lambda-1)\det(\lambda I-\overline B_\beta).
\]
Differentiate at $\lambda=1$. Expansion of a determinant by its
columns gives the derivative as the trace of its adjugate, including
in characteristic two. Hence
\[
 P_\beta(T)=\operatorname{tr}\operatorname{adj}(I-B_\beta)
    =a(T)(1+T+\cdots+T^{s-1}).
\]
This proves \eqref{J15:eq:alexander-quotient} and all its asserted
invariance properties.
\end{proof}

\begin{lemma}[New lemma: length detected by a full-twist prefix]
\label{J15:lem:fulltwist-breadth}
Let $s\geq2$ and let $w$ be a native positive adjacent word on
$s$ strands. Put $l=\length(w)$ and $\beta=d_sw$.
Then the invariant in Lemma~\ref{J15:lem:alexander-determinant} is
nonzero and has breadth
\begin{equation}\label{J15:eq:fulltwist-breadth}
 \operatorname{br}A_{\widehat\beta}
    =s(s-1)+l-s+1=\length(\beta)-s+1.
\end{equation}
\end{lemma}
\begin{proof}
Lemma~\ref{J15:lem:burau-elementary} gives
\[
 P_\beta(T)=\det(I+T^s\overline B_w).
\]
Its constant coefficient is one, since the entries of
$\overline B_w$ are polynomials. Expanding by principal minors,
the term using all $s-1$ rows and columns is
\[
 T^{s(s-1)}\det\overline B_w=T^{s(s-1)+l}.
\]
Every term using $k\leq s-2$ rows and columns has degree at most
$sk+l\leq s(s-2)+l$, by the minor bound.
Thus the displayed highest term cannot cancel.
The lowest degree of $P_\beta$ is zero and its highest degree is
exactly $s(s-1)+l$. Both extreme coefficients are one in
$\mathbb F_2$.

The denominator in \eqref{J15:eq:alexander-quotient} has lowest degree
zero and highest degree $s-1$. Since the quotient belongs to $R$,
additivity of the lowest and highest exponents under multiplication
gives lowest degree zero and highest degree
$s(s-1)+l-s+1$ for that quotient. This proves the formula.
\end{proof}

\begin{corollary}[New lemma: isotopic positive full-twist words have equal length]
\label{J15:cor:isotopic-fulltwist-length}
Let $s\geq1$ and let $v,w$ be native positive words on $s$ strands.
If the oriented closures of $d_sv$ and $d_sw$ are ambiently isotopic,
then $\length(v)=\length(w)$.
The same conclusion holds if the isotopy is given without component
orientations.
\end{corollary}
\begin{proof}
For $s=1$ there is no adjacent generator, so both words are empty.
For $s\geq2$, the oriented assertion follows by applying
Lemma~\ref{J15:lem:fulltwist-breadth} to the two equal invariant breadths.

For the assertion without specified orientations, give each closure
its braid orientation. If there is only one component, an isotopy
either preserves or reverses it, and both cases preserve breadth.
If there are several components, the linking number of every pair
is strictly positive. Indeed the full twist has two positive
crossings for each pair of its strands, as proved by the rotation
calculation in J14. Two different closure components are unions of
disjoint strand sets. Their crossings in the prefix alone therefore
contribute a positive linking number. The remaining positive letters
add only nonnegative contributions. The integer linking number is
half the signed number of crossings between the two components.

Reversing exactly one member of a component pair reverses this
integer, whereas an ambient isotopy preserves it. Thus an isotopy
between the two unoriented links must either preserve the orientations
of all components or reverse all of them, after its permutation of
components. Both possibilities preserve the invariant breadth by
Lemma~\ref{J15:lem:alexander-determinant}. The same length equality
follows. Linking numbers here are ordinary integer topological
counts. Every homology and matrix coefficient remains in
$\mathbb F_2$.
\end{proof}

\begin{theorem}[New lemma: the full positive-word quantifier in M407]
\label{J15:thm:all-fulltwist-words}
Let $f$ be a reduced plane curve germ of multiplicity $S\geq1$,
let $b$ be its irreducible branch count, and choose a transverse
projection. Let $w$ be any positive adjacent word on $S$ strands
whose full-twist closure $\widehat{d_Sw}$ is ambiently isotopic to
the small link of $f$. Assume the geometric identification of that
link with its transverse root braid supplied by M403(b), whenever
the native sphere formulation is intended. Then
\begin{equation}\label{J15:eq:all-fulltwist-length}
 \length(w)+S^2+b=2\delta(f)+2S.
\end{equation}
The theorem does not require a braid isotopy between the specified
word and the analytic root braid.
\end{theorem}
\begin{proof}
Theorem~\ref{J14:thm:transverse-germ} produces a positive
word $v$ such that $d_Sv$ is the same analytic root braid.
Its geometric link is the specified small link under the stated
M403(b) identification. The two closures $\widehat{d_Sv}$ and
$\widehat{d_Sw}$ are therefore ambiently isotopic.
Corollary~\ref{J15:cor:isotopic-fulltwist-length} gives
$\length(v)=\length(w)$.
The discriminant, resultant, and polar-intersection argument already
proved in Corollary~\ref{J15:cor:fulltwist} gives
\[
 \length(v)+S^2+b=2\delta(f)+2S.
\]
Substitution proves \eqref{J15:eq:all-fulltwist-length}.

For the proposed resolution-assembly API, set
\[
 f=A.\decl{imageEquation},\qquad S=A.\decl{squareSum}.
\]
The field names here denote the native assembly fields, rather than
additional analytic choices.
The fields \decl{imageEquation_mem} and
\decl{imageEquation_not_mem}, recorded in L431 of the supplied
audit, say that $S$ is the exact order of this equation.
The proposed closure-isotopy predicate in M407 supplies precisely
the closure comparison used above. Thus
\eqref{J15:eq:all-fulltwist-length} is its full displayed conclusion,
including the quantifier over every such positive word.
The native definition of that proposed predicate still belongs to
the geometric assembly, rather than to this word-length proof.
\end{proof}

This comparison uses only a cyclic covering, finite matrix algebra,
and the discriminant calculation for the analytic word. It uses
neither Stallings' theorem nor Rudolph's fibre-surface argument.
The classical determinant formula is also stated in
\cite[Introduction]{MortonBraid}. The complete $\mathbb F_2$
derivation required here is Lemma~\ref{J15:lem:alexander-determinant},
so that published formula is not an unproved step of this argument.

The smallest matrix cases agree with the analytic checks below.
For $s=2$ and $w$ empty, the numerator is $1+T^2$ and the denominator
is $1+T$, giving $A=1+T$ of breadth one. For $s=2$ and
$w=\sigma_0$, they give $A=1+T+T^2$ of breadth two. These are
applications of the displayed matrices and quotient formula, without
an appeal to a knot table. For $s=1$ the adjacent word is empty,
as handled separately in Corollary~\ref{J15:cor:isotopic-fulltwist-length}.


\section{Exceptional cases and convention checks}\label{J15:sec:exceptions}

If $m=1$, the discriminant is $1$, the signed root braid has exponent zero,
and the positive word is empty. For a transverse smooth branch, $a=1$,
$\mu=\delta=0$, and \eqref{J15:eq:branch-length} reads $0+1=0+1$.
Its generator set is the empty set of indices, which is the full set
$\operatorname{Fin}(0)$, and its permutation has one fixed-point cycle.
The lower bound $a\ge2$ in the resolution-level acceptance comes from
singular-point data and is not imposed on arbitrary smooth branches.

For the cusp $f(\xi,\zeta)=\zeta^2-\xi^3$, the discriminant is
$4\xi^3$. The two sheet differences increase in argument by $3\pi$
over one base turn. Its Jacobian quotient has basis $1,\xi$, so
$\mu=2$. The supplied Milnor identity and irreducibility give
$\delta=1$. There are three positive pair crossings and
$\length(w)=3=2\delta+2-1$.
For the transverse ordinary double point
$f=(\zeta-\xi)(\zeta+\xi)$, the discriminant is $4\xi^2$.
There are two positive pair crossings, giving the two-letter positive
full twist. Here $S=b=2$. The Jacobian ideal is $(\xi,\zeta)$, so $\mu=1$,
and the supplied Milnor identity gives $\delta=1$. The tail is empty.
These computations are exact substitutions in the proved identities,
not evidence for an unproved general statement.

For an ordinary $S$-fold point with distinct tangent slopes, every pair
has a difference of order one in $\xi$. Thus $D$ has order $S(S-1)$,
which is the full-twist length, and \eqref{J15:eq:tail-length} gives an empty
tail. This agrees with the separated strict-transform clusters in J14.

For a simultaneous event involving $q\ge3$ sheets with equal projected
real part, the pair count contains exactly $\binom q2$ terms from that
cluster. Lemma~\ref{J15:lem:events} resolves it into that many simple pair
crossings. The endpoint cancellation in Lemma~\ref{J15:lem:pair-count}
already counts it pairwise, so no division by $q$ or $2$ occurs.

The distinguished condition cannot be omitted from the local-intersection
assertion. For
\[
 P(\xi,\zeta)=\zeta\bigl((\zeta-1)^2-\xi\bigr),
\]
one has
$\Disc P=4\xi(1-\xi)^2$, of order one at zero.
However $P_\zeta(0,0)=1$, so the local pair quotient at $(0,0)$ is zero
and $i_0(P,P_\zeta)=0$. The other two roots coalesce at $(0,1)$.
Thus the discriminant order of a general monic polynomial measures all
central collisions, not only the collision at the origin.

Transversality is necessary when the strand number is identified with the
multiplicity of the germ. The smooth germ $f=\zeta^2-\xi$ has multiplicity
one and $\mu=0$, but projection to $\xi$ has degree two. Its discriminant
has order one and its root braid has one positive letter.
Formula~\eqref{J15:eq:exp-polar} gives $1=0+2-1$, as required.
Replacing the projection degree by the multiplicity would give a false
equality. This is also the obstruction to asserting a positive full-twist
prefix on two strands from regularity of projection degree two alone.

\section{Lemma and declaration table}

\small
\begin{longtable}{@{}>{\raggedright\arraybackslash}p{0.43\textwidth}>{\raggedright\arraybackslash}p{0.15\textwidth}>{\raggedright\arraybackslash}p{0.36\textwidth}@{}}
\toprule
Declaration or result & Status & Exact use\\
\midrule\endhead
\decl{SignedAdjacentWord}, \decl{adjacentSwap}
& Existing Lean & Signs $\Z^\times$ and zero-based generator indices.\\
\decl{PointedMapGerm.exists_equation_link_canonical_positive_braid}
& Existing Lean & Chosen branch word, pair-count length, endpoint equations, and unchanged link identification.\\
\decl{planeBranchSheetDifference_exists_positive_order},
\decl{exists_radius_planeBranchSheetDifference_re_pos}
& Existing Lean & Positive angular derivative for branch pairs on a common radius.\\
\decl{exists_radius_planeBranchAngularDifference_projected_crossings_pos},
\decl{PointedMapGerm.exists_equation_link_finite_braid_projection}
& Existing Lean, catalogued & Native positive-crossing convention and positive determinants for the chosen event data. Their modules are cited in L4.16 and L404 rather than reproduced.\\
\decl{braid_word_cycleCount_eq_one_of_endpoint_rotation},
\decl{braid_word_generatorIndices_eq_univ_of_endpoint_rotation}
& Existing Lean & One cycle and every generator for the same branch word.\\
\decl{SingularSectionPointData.exists_branch_component_order},
\decl{ResolutionSectionLevel.singular}
& Existing Lean, catalogued & Exact branch data and inherited bounds $2\le a\le T$.\\
\decl{Polynomial.resultant_deriv}
& Existing Lean, catalogued & Discriminant--resultant sign relation.\\
\decl{Germ.hasFiniteIntersection_of_resultant_ne_zero},
\decl{Germ.resultant_ne_zero_of_hasFiniteIntersection},
\decl{Germ.intersectionMultiplicity_eq_resultant_order_of_ne_zero}
& Existing Lean & Distinguished local resultant formula, with all hypotheses checked.\\
\decl{PlaneCurve.exists_polar_intersection_identity}
& Existing Lean & Teissier identity applied directly to $f$ and parameter $\xi$.\\
\decl{PlaneCurve.jacobianColength_add_irreducibleBranchCount}
& Existing Lean & Milnor identity for the branch and reduced cases.\\
Lemmas~\ref{J15:lem:argument}--\ref{J15:lem:events}
& New, proved & Continuous arguments, sign, pair cancellation, event resolution, and winding invariance.\\
Lemma~\ref{J15:lem:local-principle}
& New, proved & Winding equals analytic vanishing order.\\
Lemma~\ref{J15:lem:preparation}
& New, proved & Distinguished preparation and projection intersection number.\\
Lemma~\ref{J15:lem:unit}
& New, proved & Equality of the polar pair ideals after a unit factor.\\
Lemma~\ref{J15:lem:branch-degree}
& New, proved & The normalized $a$-sheet parameterization gives degree $a$.\\
Theorems~\ref{J15:thm:discriminant}, \ref{J15:thm:reduced}, \ref{J15:thm:branch}
& New, proved & M1155, polar length identity, and branch acceptance.\\
Corollary~\ref{J15:cor:fulltwist}
& New, proved conditional on the specified J14 word & Reduced full-twist length and tail identity.\\
J14 transverse full-twist construction
& New, in companion document & Existence and positivity of the word representing the same root braid.\\
\bottomrule
Lemmas~\ref{J15:lem:burau-elementary}--\ref{J15:lem:fulltwist-breadth}
& New & Finite braid matrices, the cyclic-covering module, and the exact
breadth formula over $\mathbb F_2$.\\
Corollary~\ref{J15:cor:isotopic-fulltwist-length}
& New & Equal length for ambiently isotopic positive full-twist closures.\\
Theorem~\ref{J15:thm:all-fulltwist-words}
& New & The full positive-word quantifier in M407, with the audit's
explicit reduced-link realization input.\\
\decl{ResolutionSectionAssembly.imageEquation_mem},
\decl{ResolutionSectionAssembly.imageEquation_not_mem}
& Existing, audit L431 & The assembly's square sum is the exact order
of its image equation.\\
\end{longtable}
\normalsize

\section{Precise limits of the conclusion}\label{J15:sec:scope}

The counting and analytic identities in Theorems~\ref{J15:thm:discriminant}
and \ref{J15:thm:reduced} have no remaining mathematical gap. Theorem
\ref{J15:thm:branch} assembles the existing branch-word and singular-point
inputs with these identities and preserves their chosen word and link
representative. The reduced positive-word conclusion uses the specified
J14 construction, as required by J15(b). The exact source spelling of the
resolution-level structures is not supplied in the attached native listings,
so the resolution-level assembly is stated using the existing declarations
as quoted in the supplied audit.

Theorem~\ref{J15:thm:all-fulltwist-words} closes the broader M407 quantifier
for positive full-twist words known only through their ambiently isotopic
closures. Its additional input is the finite matrix and cyclic-covering
argument proved in Section~\ref{J15:sec:fulltwist-invariant}. The analytic
word length still comes from the argument principle, discriminant,
resultant, and Teissier identity. No fibre-surface theorem is used.
The native sphere identification in the reduced case retains the
M403(b) geometric input assigned to J13 by the audit. The definition of
its proposed closure predicate belongs to that assembly. This is the
same explicit geometric dependency as in J14, rather than an additional
word-length gap.

Similarly, the displayed J14 API with only regularity in the fibre variable
must be distinguished from its narrative hypothesis of multiplicity $S$
and a transverse projection. Corollary~\ref{J15:cor:fulltwist} uses the latter
hypotheses, which are explicitly required in the user's reduced-germ task.
The tangent smooth example above explains why those hypotheses cannot be
silently inferred from projection degree alone.




\chapter{J14: Multi-branch positivity and the full-twist prefix}
\label{chapter:J14}



\section{Statements and the distinction between the two hypotheses}

There is a hypothesis error in the displayed proposed API of M405.
Regularity of order $S$ in the first coordinate means that the restriction
$f(\zeta,0)$ has order $S$. It does not say that the two-variable germ has
multiplicity $S$. The API, including its smooth sphere-link conclusion, is
false with only its displayed hypotheses. Proposition~\ref{J14:prop:false}
gives a counterexample of projection degree two.

The user's geometric statement includes the extra information that the
multiplicity is $S$ and that the projection is transverse. That information
is sufficient for the full-twist argument. The distinction is retained
throughout this document. No assertion with the corrected hypotheses is
presented as a proof of the displayed API.

Write $\sigma_i$ for the positive adjacent generator at native index
$i\in\{0,\ldots,S-2\}$. Put
\begin{equation}\label{J14:eq:fullword}
 c_S=\sigma_0\sigma_1\cdots\sigma_{S-2},\qquad d_S=c_S^S.
\end{equation}
For $S=1$, both words are empty. These are the precise lists proposed as
\decl{coxeterWord S} and \decl{fullTwistWord S} in the audit. The supplied
listings define \decl{PositiveAdjacentWord}, its conversion to signed words,
its permutation, and its cycle count. They do not include a completed
definition of the proposed multi-component smooth model
\decl{positiveWordClosedLinkModel}.

Here is the exact analytic statement proved below.

\begin{theorem}[New lemma: multi-sheet positivity]\label{J14:thm:multisheet}
Let $q,S$ be positive integers. Let $Y_1,\ldots,Y_S$ be convergent holomorphic
functions on a disc about zero. Assume that their germs are pairwise
distinct. Let $c\in\C^*$ satisfy
\[
 \Ree\bigl(c(Y_i(0)-Y_j(0))\bigr)\ne0
 \quad\text{whenever }Y_i(0)\ne Y_j(0).
\]
There is $r_0>0$ such that, for every $0<r<r_0$, every projected crossing of
the strands
\[
 K_i(\theta)=cY_i\bigl(r e^{\ii\theta/q}\bigr)
\]
is positive. Every pair with $Y_i(0)=Y_j(0)$ has strictly increasing argument
of its difference along the parameter interval. Pairs with distinct constant
terms have no projected crossing. Every compact parameter interval has
finitely many crossing events, with multiple events counted pairwise.

If, in addition, the unordered set of functions is invariant under
$t\mapsto e^{2\pi\ii/q}t$, these strands define a closed $S$-braid. On a
period with regular endpoints that braid is represented by a word in
\decl{PositiveAdjacentWord S}.
\end{theorem}

The monodromy-invariance condition in the last paragraph is necessary for a
closed braid. A list consisting only of $\xi^{1/2}$ does not close after one
turn in $\xi$. The root sheets of a reduced equation satisfy this condition.
The schematic wording of M1153 must be read as a complete set of root sheets,
or must include this condition explicitly.

\begin{theorem}[New lemma: full twist for explicit transverse branch data]
\label{J14:thm:factorization}
Let $b\geq1$, let $a_1,\ldots,a_b$ be positive integers with
$\sum_j a_j=S$, and let
\[
 F_j(s)=(h_j(s),s^{a_j})
\]
be convergent injective parametrizations of distinct branch germs, whose
punctured images are disjoint in a common neighborhood. Assume
$\ord h_j\geq a_j$, where the zero function is allowed. For all sufficiently
small $\eta>0$, let $B_\eta$ be their root braid over $\xi=\eta e^{\ii\theta}$.
There is a fibre rotation $c\in\C^*$ and one positive word
$w\in\decl{PositiveAdjacentWord S}$ such that $B_\eta$ is isotopic through
closed braids to the word $d_Sw$. The word $w$ may be chosen as a positive
event word of the strict-transform sheets $\zeta/\xi$.
Its native permutation has exactly $b$ cycles, with the cycle of branch $j$
of length $a_j$. The same holds for $d_Sw$.
\end{theorem}

The theorem is a statement about the same root braid, through a
collision-free braid isotopy. It supplies the braid identification needed by
the discriminant argument of J15. It is not merely an equality of link types.
Section~\ref{J14:sec:curve} constructs the needed finite data directly from
the prepared polynomial and its finite root covering. It proves injectivity,
the exact kernel ideals, punctured disjointness, and the degree sum without
an additional normalization assumption. The sphere conclusion is explicitly conditional on M403(b) and
the realization in M419(B), as permitted by audit Section 11.1.

\section{The displayed API is false}

\begin{lemma}[New lemma: positive two-strand closures]
\label{J14:lem:two}
For every integer $r\geq2$, the closure of the positive two-strand word
$\sigma_0^r$ is not the unknot.
\end{lemma}
\begin{proof}
The permutation is the $r$th power of the transposition of two elements.
For even $r$, following the two strands around the closure gives two
components, so the link is not a knot.

For odd $r$, the closure is the curve of slope $(2,r)$ on the torus separating
the two standard solid tori of $\mathbb S^3$. This follows directly from the
sheet formula: each generator rotates the difference of the two fibre
points counterclockwise through $\pi$, and closure gives the parametrization
\[
 t\longmapsto (a e^{2\ii t},b e^{r\ii t}),\qquad
 a,b>0,\quad a^2+b^2=1.
\]
It is embedded because $r$ is odd. The complement group has presentation
\begin{equation}\label{J14:eq:torusgroup}
 \langle u,v\mid u^2=v^r\rangle.
\end{equation}
For completeness, remove an annular neighborhood of this primitive slope
curve from the separating torus. The remainder of the torus is an annulus.
Thicken the two solid tori slightly into the knot complement. Each thickened
piece retracts to its solid-torus core, and their intersection retracts to
that annulus. Its core winds twice around one solid-torus core and $r$ times
around the other. Applying the van Kampen gluing relation gives exactly
\eqref{J14:eq:torusgroup}.

This group maps onto the dihedral group generated by a reflection $s$ and a
rotation $v_0$ of order $r$, by $u\mapsto s$, $v\mapsto v_0$.
The relation is satisfied since $s^2=v_0^r=1$. For odd $r\geq3$, this quotient
is nonabelian: $sv_0s^{-1}=v_0^{-1}\ne v_0$.
The unknot complement is a solid torus and has group $\mathbb Z$.
An ambient diffeomorphism induces an isomorphism of complement groups, which
excludes an unknot closure in the odd case as well.
\end{proof}

\begin{proposition}[New lemma: counterexample to the exact M405 API]
\label{J14:prop:false}
The proposed assertion
\decl{Germ.exists_reduced_link_fullTwist_positive_word}
in M405 is false under its displayed hypotheses.
\end{proposition}
\begin{proof}
Take
\[
 H(\zeta,\xi)=\zeta^2-\xi,\qquad S=2.
\]
The function is analytic and its germ is squarefree, since it is a smooth
coordinate germ. Its restriction $H(\zeta,0)=\zeta^2$ has order two, so it
satisfies the displayed first-coordinate regularity hypothesis. It has
two-variable multiplicity one.

Its small sphere link is the unknot. Indeed, its parametrization is
$t\mapsto(t,t^2)$. On the sphere of radius $\epsilon>0$, the parameter has
the unique positive modulus $r$ solving $r^2+r^4=\epsilon^2$. After
normalization the link is
\[
 \theta\longmapsto(ae^{\ii\theta},be^{2\ii\theta}),
 \qquad a=r/\epsilon>0,\quad b=r^2/\epsilon,
 \quad a^2+b^2=1.
\]
Replace $b$ by $(1-s)b$ and $a$ by
$\sqrt{1-(1-s)^2b^2}$ for $0\leq s\leq1$.
This is a smooth family of embedded graphs over the first circle, ending
at $(e^{\ii\theta},0)$. A compactly supported fibre vector field extending
the velocity of these graphs gives an ambient isotopy, by the explicit
construction in Lemma~\ref{J14:lem:ambient} below.

Every native positive two-strand word is $\sigma_0^\ell$ for some
$\ell\geq0$. The proposed prefix $d_2$ is $\sigma_0^2$.
Thus every candidate model has closure $\sigma_0^{2+\ell}$, which cannot be
the unknot by Lemma~\ref{J14:lem:two}. This contradicts the displayed sphere
conclusion for every possible $w$.
\end{proof}

The same example gives a word-level obstruction with no knot invariant:
the root sheets $\pm\sqrt\xi$ give one positive crossing, whereas $d_2w$
has exponent sum $2+\ell$. The discriminant computation of J15 makes that
obstruction explicit. The geometric counterexample above also addresses
the weaker sphere-link conclusion of the displayed API.
On the common cover $\xi=t^2$, the proposed strict-transform sheets would
be $\pm t^{-1}$. Their negative order records the failed transverse
hypothesis directly. They are outside the nonnegative-order statement of
M1153.

\section{Positive rotation for a finite collection of analytic differences}

\begin{lemma}[New lemma: analytic difference estimate]\label{J14:lem:log}
Let $D$ be a nonzero holomorphic germ at zero. Write
$D(t)=t^eV(t)$ with $e\in\mathbb N$ and $V(0)\ne0$.
On a sufficiently small punctured disc, $D$ has no zeros and
\[
 \frac{tD'(t)}{D(t)}=e+\frac{tV'(t)}{V(t)}.
\]
If $e>0$, its real part is greater than $e/2$ on a sufficiently small
punctured disc.
\end{lemma}
\begin{proof}
Choose the first nonzero coefficient of the convergent power series of $D$.
Factoring the corresponding power of $t$ produces the holomorphic $V$ with
$V(0)\ne0$. Continuity gives a disc on which $V$ is nonzero.
Differentiating the factorization gives the displayed identity there.
The continuous function $tV'(t)/V(t)$ tends to zero. Taking its modulus less
than $e/2$ proves the last assertion.
\end{proof}

The limiting calculation is also the existing declaration
\decl{tendsto_mul_deriv_div_of_analytic_factor} in the supplied
\decl{PlaneBranchSheetRotation.lean}. The new use is to every difference in
a finite collection, including differences from distinct branches.

\begin{proof}[Proof of Theorem~\ref{J14:thm:multisheet}]
For each unordered pair let $D_{ij}=Y_i-Y_j$. Its germ is nonzero by
hypothesis. If $D_{ij}(0)=0$, apply Lemma~\ref{J14:lem:log} with $e_{ij}>0$.
If $D_{ij}(0)\ne0$, continuity gives
\[
 \left|\Ree(cD_{ij}(t))-\Ree(cD_{ij}(0))\right|
 <\frac12\left|\Ree(cD_{ij}(0))\right|
\]
on a sufficiently small disc. There are only finitely many pairs, so a common
disc works for all estimates and all pairwise nonvanishing conclusions.

Let $t=r e^{\ii\theta/q}$ and $d(\theta)=cD_{ij}(t)$.
For a positive-order difference,
\begin{equation}\label{J14:eq:arg}
 \frac{d}{d\theta}\arg d(\theta)
 =\Imm\frac{d'(\theta)}{d(\theta)}
 =\frac1q\Ree\frac{tD_{ij}'(t)}{D_{ij}(t)}>0.
\end{equation}
The argument on any real interval is defined by continuous lifting since
$d$ is nonzero. At a projected crossing write $d(\theta_0)=\ii v$ with
$v\in\mathbb R\setminus\{0\}$. If $a=\Ree d'(\theta_0)$, then
\begin{equation}\label{J14:eq:sign}
 \frac{d}{d\theta}\arg d(\theta_0)=-\frac{a}{v},
 \qquad -a v>0.
\end{equation}
This is the native positive projected-crossing determinant convention
recorded by
\decl{exists_radius_planeBranchAngularDifference_projected_crossings_pos}.
In particular $a\ne0$, so each pairwise projected crossing is simple.
The order-zero pairs have no crossing by the preceding real-part estimate.

Each pair has finitely many zeros of its projected real difference on a
compact interval. Otherwise a convergent subsequence would accumulate at a
zero, contradicting the nonzero derivative there. A finite union over the
pairs is finite. Choose the initial angle outside these events whenever a
closed period is required.

At an event where several strands have the same real part, their imaginary
parts are distinct. Equation~\eqref{J14:eq:sign} says that the order of their
real velocities is the reverse of the order of their imaginary coordinates.
Consequently every coincident pair reverses its real order exactly once in
a sufficiently small event interval, and every such crossing is positive.
Lemma~\ref{J14:lem:events} below resolves the event into positive adjacent
letters, preserving the braid. Concatenating the event words gives the
required positive adjacent word. Monodromy invariance identifies the endpoint
configurations, and therefore makes this a closed braid.
\end{proof}

\begin{lemma}[New lemma: resolution of finite positive events]
\label{J14:lem:events}
Suppose finitely many smooth strands on a compact interval have distinct
complex values at each time and distinct projected real coordinates at the
endpoints. Suppose their projected real coincidences occur at finitely many
times, each pairwise zero is simple, and every crossing satisfies
\eqref{J14:eq:sign}. Their braid is isotopic relative to the endpoints to a
concatenation of positive adjacent generators. An event at which a group of
$r$ strands has common real coordinate contributes exactly $r(r-1)/2$
letters, one for each pair in that group.
\end{lemma}
\begin{proof}
Choose pairwise disjoint small intervals around the event times.
On the complement, the minimum separation of projected real coordinates is
positive. Within an event interval, shorten the interval until every
coincident pair has real derivative of a fixed nonzero sign and its imaginary
difference stays nonzero with fixed sign. All other pairs stay separated in
real coordinate. The implicit function theorem preserves one simple crossing
of the same sign under sufficiently small perturbations of real coordinates
supported inside the interval.

Choose such perturbations with different small constant offsets for the
strands near the event, and multiply them by a bump function supported in
the interval. The offsets can be chosen so that no two pairwise crossing
times coincide. Here is the finite avoidance argument. Each simple crossing
time depends smoothly on the offsets. For two distinct pairs, vary an offset
belonging to one pair and not the other. The first crossing time changes with
nonzero derivative and the second does not. Thus the equation asserting
equality of those two times has empty interior in the finite-dimensional
offset space. Avoid the finite union of these sets. The perturbation and its
linear interpolation from zero are small enough to preserve complex
separation, because the original strands have positive minimum separation.

After perturbation there is one pairwise crossing at each crossing time.
That pair is adjacent in real order, since all other projected coordinates
are distinct. The sign remains positive by the fixed-sign estimates.
Every originally coincident pair crosses once, and all other pairs cross
zero times in that interval. This proves the asserted count.

Finally identify the braid with the standard generator model on each small
crossing interval. The two crossing strands are separated from all other
strands in real coordinate. Their difference travels through one open
half-plane to the other with the positive orientation in
\eqref{J14:eq:sign}. Contract its radial coordinate to a fixed positive value
and its angular coordinate to a strictly increasing half-turn, keeping its
endpoints fixed by short paths inside the initial and final real-order
chambers. More explicitly, first straighten the endpoint configurations to
distinct real rank positions while keeping their real order. On the small
crossing interval the imaginary difference has one fixed nonzero sign.
After adjoining the endpoint straightening paths, its argument runs from
one real ray to the other through the prescribed upper or lower half-plane.
Interpolate that lifted argument with the corresponding monotone half-turn
and interpolate the positive modulus with a constant. Both interpolations
avoid zero. The midpoint of the crossing pair is put at the midpoint of its
two rank positions. All other strands can be kept outside the intervening
real rank gap. This realizes the pair homotopy without a collision with a
third strand. On a gap with no crossing, the ordered
configuration chamber is convex, so a straight homotopy connects to the
constant rank model. Choosing collars at interval endpoints makes these
homotopies compatible. This constructs the claimed relative braid isotopy.
\end{proof}

\section{Several tangent points and strict-transform sheets}

\begin{lemma}[New lemma: a common Puiseux parameter]\label{J14:lem:common}
Under the hypotheses of Theorem~\ref{J14:thm:factorization}, let
$q=\operatorname{lcm}(a_1,\ldots,a_b)$ and
$\omega_a=e^{2\pi\ii/a}$. Index sheets by pairs $(j,l)$ with
$0\leq l<a_j$. The functions
\begin{equation}\label{J14:eq:sheets}
 Z_{jl}(t)=h_j(\omega_{a_j}^l t^{q/a_j}),\qquad
 U_{jl}(t)=\frac{Z_{jl}(t)}{t^q}
\end{equation}
extend holomorphically across zero and have pairwise distinct germs.
Their set is preserved by monodromy, which acts by
$(j,l)\mapsto(j,l+1\bmod a_j)$. For fixed $j$, all constants $U_{jl}(0)$
are equal to the tangent slope $\lambda_j$ of that branch.
\end{lemma}
\begin{proof}
Write $h_j(s)=s^{a_j}v_j(s)$ with $v_j$ holomorphic, using the assumed order
bound. Since $(\omega_{a_j}^l)^{a_j}=1$,
$U_{jl}(t)=v_j(\omega_{a_j}^l t^{q/a_j})$, proving extension and
$U_{jl}(0)=v_j(0)=\lambda_j$.

If two different root labels of the same branch gave equal $Z$ on a
punctured neighborhood, their parametrizations would give the same point
with different nonzero parameter values, contrary to injectivity. For
different branches, equality would contradict disjoint punctured images.
Thus the differences are nonzero germs. Dividing by $t^q$ preserves this
property. Substitution of $e^{2\pi\ii/q}t$ increments $l$ by one in
\eqref{J14:eq:sheets}, which proves the monodromy assertion.
\end{proof}

\begin{lemma}[New lemma: one projection separates all tangent clusters]
\label{J14:lem:projection}
For any finite set of distinct complex numbers $\lambda_1,\ldots,\lambda_r$,
there is a unit complex number $c$ such that
$\Ree(c(\lambda_i-\lambda_j))\ne0$ for every distinct pair.
For the functions in Lemma~\ref{J14:lem:common}, sufficiently small circles have
no projected crossings between sheets with distinct tangent slopes.
\end{lemma}
\begin{proof}
For each nonzero difference $\lambda_i-\lambda_j$, the equation
$\Ree(c(\lambda_i-\lambda_j))=0$ excludes exactly two points of the unit
circle. A finite set cannot exhaust that circle. Choose $c$ outside it.
The projected constants are then distinct. Continuity and finiteness give
disjoint open real intervals containing, respectively, the projected values
of all sheets in each tangent cluster on a common sufficiently small disc.
Sheets from distinct clusters therefore never have the same projected real
coordinate.
\end{proof}

These lemmas address several strict-transform points simultaneously.
The points on the exceptional line are $(\lambda_j,0)$. No hypothesis says
that these points coincide. Within a cluster the differences have positive
order, so Theorem~\ref{J14:thm:multisheet} applies even when the strict transform
is tangent to the exceptional line. Across clusters the differences have
order zero and Lemma~\ref{J14:lem:projection} excludes crossings. Positivity of a
logarithmic derivative is not asserted for an order-zero difference.

\section{The full turn is the native full-twist word}

Let $\Conf_S(\C)$ be the space of unordered sets of $S$ distinct complex
points. A closed braid is a smooth loop in this space. A collision-free
homotopy of loops is a braid isotopy. When a regular projection is chosen,
the rank ordering of the real parts identifies a positive crossing with the
native adjacent generator at that rank.

\begin{lemma}[New lemma: rigid rotation gives the prescribed word]
\label{J14:lem:rotation}
Let $S\geq1$ and let $Q$ be a configuration of $S$ distinct points.
The full counterclockwise rotation
$R_Q(t)=e^{2\pi\ii t}Q$, $0\leq t\leq1$, represents $d_S$ of
\eqref{J14:eq:fullword}, using real-rank identifications at a regular basepoint.
In particular, its chirality is positive, its word length is $S(S-1)$,
and its native permutation is the identity.
\end{lemma}
\begin{proof}
For $S=1$ the assertion concerns the trivial one-strand braid. Assume $S\geq2$.
First place the points on the upper unit semicircle, including $1$ and $-1$,
with strictly increasing arguments
$0=\beta_0<\cdots<\beta_{S-1}=\pi$.
Choose a smooth strictly increasing lift $\phi:\mathbb R\to\mathbb R$
with $\phi(x+2\pi)=\phi(x)+2\pi$ and
$\phi(2\pi j/S)=\beta_j$.

The cyclic-shift path
\[
 Q(t)=\left\{\exp\bigl(\ii\phi(2\pi(j+t)/S)\bigr):0\leq j<S\right\},
 \qquad 0\leq t\leq1,
\]
is a loop of unordered configurations. One point travels from $-1$ to $1$
along the lower semicircle. Every other point moves along the upper
semicircle to the next point. The lower point has strictly increasing real
coordinate, and every upper point has strictly decreasing real coordinate.
It therefore crosses each upper point exactly once, in order from the
leftmost to the rightmost rank. Its imaginary coordinate is negative and
the upper point's is positive at each crossing. Thus the difference has
positive real derivative and negative imaginary part. Equation~\eqref{J14:eq:sign}
gives a positive crossing. The word is exactly
$\sigma_0\sigma_1\cdots\sigma_{S-2}=c_S$.

Its $S$fold iterate has angles $\phi(\alpha_j+2\pi t)$ with
$\alpha_j=2\pi j/S$. The following is a based collision-free homotopy from
that iterate to rigid rotation of the same initial configuration:
\begin{equation}\label{J14:eq:rotatehomotopy}
 \left\{
 \exp\left(\ii\bigl[(1-s)\phi(\alpha_j+2\pi t)
       +s(\phi(\alpha_j)+2\pi t)\bigr]\right):0\leq j<S
 \right\},\quad 0\leq s,t\leq1.
\end{equation}
For $j<l$, each of the two argument differences being combined lies
strictly between $0$ and $2\pi$. Their convex combination does too, so the
points never collide. The endpoints at $t=0,1$ are fixed as unordered sets.
This proves that rigid rotation at this configuration is $c_S^S$.

For an arbitrary configuration with distinct projected real coordinates,
connect it to the semicircle configuration with the same real-rank labels
by a straight homotopy. Every real-coordinate difference keeps its strict
sign, so this path has no crossing. Rotating this path through $2\pi$
gives a two-parameter homotopy transporting the rotation loop to the one
just computed. The transport has no generator letters and preserves the
rank identification. A preliminary arbitrarily small change of projection
handles a base configuration with a projected coincidence.

There are $S$ copies of a word of length $S-1$, proving the length formula.
Under the native product convention, applying the adjacent swaps in $c_S$
to each index gives the permutation
\[
 0\mapsto1\mapsto\cdots\mapsto S-1\mapsto0.
\]
Its $S$th power is the identity. This verifies the
permutation assertion directly from \decl{SignedAdjacentWord.permutation}.
\end{proof}

\begin{lemma}[New lemma: separating rotation from a sheet loop]
\label{J14:lem:product}
Let $U:[0,1]\to\Conf_S(\C)$ be a smooth loop. The loop
$t\mapsto e^{2\pi\ii t}U(t)$ is isotopic through configuration loops to the
concatenation of $R_{U(0)}$ followed by $U$.
\end{lemma}
\begin{proof}
Consider $F(a,b)=e^{2\pi\ii a}U(b)$ on the unit square.
Every value is a configuration. The diagonal path from $(0,0)$ to $(1,1)$
is homotopic relative to its endpoints to the path along the bottom and
right sides. For example, interpolate linearly between
$(t,t)$ and
$(\min(2t,1),\max(2t-1,0))$.
Composing with $F$ gives the desired homotopy. The bottom side is rotation
of $U(0)$, and the right side is $U$ because $e^{2\pi\ii}=1$.
Replace the piecewise linear parameter functions by smooth nondecreasing
ones with endpoint collars. This preserves the square, the endpoints, and
collision-freeness, and gives a smooth braid isotopy.
\end{proof}

\begin{lemma}[New lemma: a braid isotopy has an ambient realization]
\label{J14:lem:ambient}
Every smooth collision-free homotopy of finite-sheet closed braids whose
images stay in a compact subset of the open solid torus extends to an ambient
isotopy supported in that solid torus. The assertion allows any endpoint
permutation and any number of components.
\end{lemma}
\begin{proof}
Use the coordinates $(\theta,z)$ of $\mathbb S^1\times\C$ on a neighborhood
of the compact union of the braids. Locally label the fibre points as
$Z_j(s,\theta)$. Their minimum pairwise distance over the compact parameter
set is positive. Choose a radius $a>0$ smaller than one third of that
distance and smaller than the distance to the edge of the chosen fibre
chart. Let $\chi$ be a smooth bump function equal to one near zero and zero
outside $[0,1)$. Define a vertical vector field by
\[
 V_s(\theta,z)=
 \sum_j\chi\left(\frac{|z-Z_j(s,\theta)|^2}{a^2}\right)
       \partial_s Z_j(s,\theta).
\]
For one strand there is no pairwise-distance restriction, and a radius
smaller than the distance to the edge of the fibre chart suffices.
The discs are disjoint, so this vector field agrees with the prescribed
velocity on each strand. Relabeling permutes the summands. Thus it is smooth
and periodic in $\theta$, including around the seam of a multi-component
closure. It has compact support. Its time-dependent flow exists for the
whole compact $s$ interval and takes each initial point along its prescribed
strand, by uniqueness for the differential equation. Extending the vector
field by zero gives the claimed ambient isotopy. The same construction
works in the standard solid-torus chart
$(\theta,z)\mapsto(\sqrt{1-|z|^2}e^{\ii\theta},z)$ of $\mathbb S^3$.
\end{proof}

\section{Proof of the factorization and the branch cycles}

\begin{proof}[Proof of Theorem~\ref{J14:thm:factorization}]
Use the common parameter of Lemma~\ref{J14:lem:common} and choose $c$ by
Lemma~\ref{J14:lem:projection}. The strict-transform functions $cU_{jl}$ satisfy
Theorem~\ref{J14:thm:multisheet}. Their closed braid therefore has a positive
event word $w$ at a fixed sufficiently small reference radius.
The same word works at every smaller positive radius in the common disc.
Indeed, interpolate two such radii through positive radii. The sheet
formulas in \eqref{J14:eq:sheets} give a smooth family of configurations at
every angle, and their differences never vanish in the common punctured
disc. This is a closed-braid isotopy from the braid at either radius to the
one at the reference radius. Thus the chosen word is independent of the
radius as a representative, even if a particular event decomposition
changes. The construction includes every tangent cluster and has no
crossings between different clusters.

On the circle $\xi=\eta e^{\ii\theta}$ the original sheets satisfy
\[
 c\zeta_{jl}(\theta)=\eta e^{\ii\theta}\,c u_{jl}(\theta).
\]
Multiplication by the fixed positive number $\eta$ is a configuration
isotopy and creates no crossings. Lemma~\ref{J14:lem:product} splits the remaining
loop into a rigid full rotation followed by the strict-transform loop.
Lemma~\ref{J14:lem:rotation} identifies the first loop with the exact word $d_S$.
The event resolution in Lemma~\ref{J14:lem:events} identifies the second with
$w$. Thus the original root braid is isotopic to $d_Sw$, with this literal
prefix. All homotopies stay in configuration space. After a uniform fibre
rescaling their images lie in a compact part of the prescribed solid torus,
so Lemma~\ref{J14:lem:ambient} also realizes them as ambient braid isotopies.

Let $\tau$ be the permutation of sheet labels given by
$\tau(j,l)=(j,l+1\bmod a_j)$. Following $\xi$ once counterclockwise has
exactly this monodromy, by Lemma~\ref{J14:lem:common}. In native real-rank
coordinates, let $\sigma_-$ and $\sigma_+$ enumerate the labels before and
after the period. The native convention is
\[
 \sigma_-\, w^{\rm perm}=\sigma_+,
 \qquad \tau\sigma_+=\sigma_-.
\]
The first identity is obtained by composing the adjacent swaps read in
time order, as in the existing event-word declarations. The second follows
by identifying the endpoint sheets. Solving gives
\begin{equation}\label{J14:eq:cycles}
 w^{\rm perm}=\sigma_-^{-1}\tau^{-1}\sigma_-.
\end{equation}
This is the same inverse-conjugacy convention as
\decl{braid_permutation_eq_conjugate_inverse_rotation} in the branch case.
The permutation $\tau$ is a disjoint union of $b$ cycles of respective
lengths $a_j$. Taking the inverse and conjugating preserves each orbit and
its cardinality. Since \decl{SignedAdjacentWord.cycleCount} counts these
orbits, $w$ has cycle count $b$. Lemma~\ref{J14:lem:rotation} gives
$d_S^{\rm perm}=1$, so $d_Sw$ has the same permutation and cycle count.
\end{proof}

\begin{corollary}[New lemma: the word-level conclusion M406]
\label{J14:cor:cycles}
For the words constructed in Theorem~\ref{J14:thm:factorization},
\[
 (d_Sw)^{\rm perm}=w^{\rm perm},\qquad
 \decl{cycleCount}(d_Sw)=\decl{cycleCount}(w)=b.
\]
A smooth branch transverse to the projection contributes one fixed point
when its covering degree is one.
\end{corollary}
\begin{proof}
The first two assertions were proved in \eqref{J14:eq:cycles} and
Lemma~\ref{J14:lem:rotation}. A one-element block is a one-element orbit and is
counted by the native \decl{SameCycle} quotient, which includes fixed points.
\end{proof}

\section{Application to every transverse reduced germ}
\label{J14:sec:curve}

Let $f$ be a nonzero reduced germ of exact multiplicity $S\geq1$ at the
origin. Write its first nonzero homogeneous part as $f_S(\zeta,\xi)$.
Transversality of the projection to $\xi$ is the condition
$f_S(1,0)\ne0$. The existing declaration
\decl{Germ.exists_linear_equiv_regularInFirstCoordinateOfOrder}, recorded
as L430, chooses such coordinates for a germ of exact order $S$.
Indeed, only the finitely many tangent directions given by the roots of
$f_S$ are forbidden as the kernel of the projection. This is a choice of
base projection, distinct from the subsequent choice of the real fibre
projection $\Ree(c\zeta)$.

The following elementary order argument explains exactly where
transversality enters.

\begin{lemma}[New lemma: transverse parametrizations have bounded slope]
\label{J14:lem:slope}
Suppose $f$ has multiplicity $S$, $f_S(1,0)\ne0$, and
$f(h(s),s^a)=0$ for a holomorphic $h$ with $h(0)=0$ and an integer $a>0$.
Then $\ord h\geq a$.
\end{lemma}
\begin{proof}
The conclusion holds if $h=0$. Otherwise put $r=\ord h$ and suppose $r<a$.
For a monomial $\zeta^p\xi^q$ of $f$, its order after substitution is
$pr+qa\geq r(p+q)\geq rS$.
Equality holds only when $q=0$ and $p=S$.
The coefficient of $\zeta^S$ is $f_S(1,0)\ne0$.
Thus the substituted series has a unique nonzero term of lowest order
$rS$, which contradicts its vanishing. Hence $r\geq a$.
\end{proof}

\begin{lemma}[New lemma: roots of a holomorphic function persist]
\label{J14:lem:persistence}
Let $g(\zeta,\xi)$ be holomorphic near $(0,0)$, and suppose
$g(\zeta,0)$ has finite positive order $d$ at $\zeta=0$.
There are positive radii $R,\eta$ such that, for every $|\xi|<\eta$,
the function $\zeta\mapsto g(\zeta,\xi)$ has exactly $d$ zeros counted with
multiplicity in $|\zeta|<R$ and no zero on $|\zeta|=R$.
\end{lemma}
\begin{proof}
Choose $R$ so that $g(\zeta,0)$ has no zero other than zero on
$|\zeta|\leq R$. Its modulus has a positive minimum on the boundary.
Uniform continuity gives $\eta$ such that the straight interpolation from
$g(\zeta,0)$ to $g(\zeta,\xi)$ is nonzero on that boundary whenever
$|\xi|<\eta$. Its boundary winding is therefore constant and equals $d$.
The local argument principle, proved as the new lemma
Lemma~\ref{J15:lem:local-principle}, identifies that
winding with the sum of the interior zero orders. Its proof factors the
finitely many zeros and writes the remaining nonzero function as the
exponential of a holomorphic logarithm. It applies on a neighborhood of the
closed disc chosen here. This proves the assertion.
\end{proof}

\begin{lemma}[New lemma: finite branch data from kernel parametrizations]
\label{J14:lem:packaging}
Let $f_1,\ldots,f_b$ be pairwise nonassociated irreducible pointed analytic
germs in two variables, with $b\geq1$, and let
$f=u\prod_{j=1}^b f_j$ for an analytic unit $u$.
For every $j$, let $a_j\geq1$ and let
$F_j(s)=(h_j(s),s^{a_j})$ be an injective analytic map on a disc about zero,
with $F_j(0)=0$ and pullback kernel exactly $(f_j)$.
Then there is a distinguished monic polynomial
$Q_j(\zeta,\xi)$ of degree $a_j$ in $\zeta$, with coefficients holomorphic
in $\xi$, such that $(Q_j)=(f_j)$. Moreover, after shrinking neighborhoods,
the following conclusions hold:
\begin{enumerate}
\item The zero set of $f$ is the union of the images of these parametrizations.
\item The images are pairwise disjoint away from the origin.
\item For each sufficiently small $\xi\ne0$, the roots belonging to branch
$j$ are exactly $h_j(s)$ with $s^{a_j}=\xi$, each occurring once.
\item The order of $f(\zeta,0)$ is $\sum_j a_j$.
\end{enumerate}
If $f$ has multiplicity $S$ and $f_S(1,0)\ne0$, then
$\sum_j a_j=S$ and $\ord h_j\geq a_j$ for every $j$.
\end{lemma}
\begin{proof}
Fix $j$ temporarily and write $a=a_j$, $h=h_j$.
Form the product
\begin{equation}\label{J14:eq:branchproduct}
 Q(\zeta,s^a)=\prod_{l=0}^{a-1}
       \bigl(\zeta-h(\omega_a^l s)\bigr).
\end{equation}
Its coefficients are convergent power series in $s$ invariant under
$s\mapsto\omega_a s$. If a coefficient is $\sum_{k\geq0}b_k s^k$,
this invariance implies $(\omega_a^k-1)b_k=0$.
Thus only exponents divisible by $a$ occur, and the coefficient is a
convergent power series in $\xi=s^a$. The radius of convergence in $\xi$
may be taken smaller than the $a$th power of a common radius in $s$.
Every $h(0)$ is zero, so $Q(\zeta,0)=\zeta^a$.
Hence $Q$ is distinguished monic of degree $a$.

The pullback of $Q$ along $F_j$ vanishes because one factor in
\eqref{J14:eq:branchproduct} is zero. The kernel assumption gives
$Q=f_jv$ for some analytic germ $v$. We claim that $v$ is a unit.
If it were not, then $v(0,0)=0$.
Restriction of the product identity to $\xi=0$ gives
$\zeta^a=f_j(\zeta,0)v(\zeta,0)$, so $v(\zeta,0)$ is a nonzero germ of
positive finite order. Lemma~\ref{J14:lem:persistence} gives a zero
$\zeta_\xi$ of $v(\cdot,\xi)$ for every sufficiently small nonzero $\xi$.
Choose all representatives and the radii so that $Q=f_jv$ holds at these
zeros and all parameters with $s^a=\xi$ lie in the injectivity disc.

At $\zeta_\xi$ the polynomial $Q$ is zero. Formula
\eqref{J14:eq:branchproduct} implies $\zeta_\xi=h(\omega_a^l s)$ for some
$s$ with $s^a=\xi$. Consequently $f_j(\zeta_\xi,\xi)=0$, since $f_j$
belongs to the pullback kernel. Both factors in $Q=f_jv$ therefore vanish
at this point, which gives $\partial_\zeta Q(\zeta_\xi,\xi)=0$ by the
product rule. On the other hand, the $a$ parameters $\omega_a^l s$ are
distinct. Injectivity of $F_j$, together with their common projection
coordinate, makes their $h$ values distinct. Formula
\eqref{J14:eq:branchproduct} therefore gives only simple roots. This is a
contradiction. Hence $v$ is a unit and $(Q)=(f_j)$.

Apply this construction to all branches. Absorbing the finitely many units
gives $f=v_0\prod_jQ_j$ for a unit $v_0$. At a small point of its zero
set, some $Q_j$ vanishes. If $\xi\ne0$, its product formula identifies that
point with $F_j(s)$ for a root $s^{a_j}=\xi$. If $\xi=0$, the identity
$Q_j(\zeta,0)=\zeta^{a_j}$ forces $\zeta=0$. Conversely, each
parametrization lies on $f=0$ because its pullback annihilates $f_j$.
This proves the local zero-set coverage and the asserted fibre description.

For $j\ne k$, the germ $f_k\circ F_j$ is nonzero. Otherwise the kernel
condition would imply $f_j\mid f_k$, and irreducibility would make them
associated. A nonzero one-variable holomorphic germ has no punctured zero
after shrinking its disc, by its factorization into a power and a unit.
Choose a common radius for the finitely many pairs. Then a nonzero point
of the image of $F_j$ cannot lie on $f_k=0$, and therefore cannot lie in
the image of $F_k$. This proves disjointness.

Finally, restriction of $f=v_0\prod_jQ_j$ to $\xi=0$ gives a unit times
$\zeta^{\sum_j a_j}$. The order is exactly $\sum_j a_j$.
If $f_S(1,0)\ne0$, that same restriction has exact order $S$, proving
$\sum_j a_j=S$. Each $F_j$ annihilates $f$, so Lemma~\ref{J14:lem:slope}
proves $\ord h_j\geq a_j$.
\end{proof}

\begin{corollary}[New lemma: the transverse germ under the precise
per-branch normalization input]\label{J14:cor:germ}
Let $f$ be a reduced germ of exact multiplicity $S\geq1$ with
$f_S(1,0)\ne0$. Assume, for each irreducible factor $f_j$ of $f$, a pointed
analytic map germ $\beta_j$ with an injective representative and
$\ker\beta_j^*=(f_j)$. Then the word-level conclusions of
Theorem~\ref{J14:thm:factorization} hold for the root braid of $f$, with
$b=\decl{PlaneCurve.irreducibleBranchCount}\ f$.
\end{corollary}
\begin{proof}
Choose pairwise nonassociated irreducible factors of the reduced germ.
There is at least one factor because the positive multiplicity makes $f$
a nonunit.
The native finite factorization is recorded in the supplied Milnor--delta
source by \decl{associated_branchProduct_branchSet} and
\decl{irreducible_branchEquation_of_mem}. Its number of factors is the
native irreducible branch count.

The pullback of $\xi$ by each $\beta_j$ is nonzero. Otherwise
$\xi\in(f_j)$, so the irreducible germ $f_j$ would be associated to the
coordinate $\xi$. This would make $\xi$ divide $f$, contrary to
$f(\zeta,0)$ having finite order $S$. Since $\beta_j$ is pointed, its
projection coordinate has finite positive order $a_j$.
Apply the supplied
\decl{PointedMapGerm.exists_injective_power_representative_sheetRotation}
to $\beta_j$, this coordinate, and this positive order.
Its conclusions provide an injective representative with projection
$s^{a_j}$. A source reparametrization by an analytic equivalence preserves
the kernel ideal, since its pullback on the one-variable germ ring is an
automorphism. Thus these representatives have exactly the hypotheses of
Lemma~\ref{J14:lem:packaging}.
That lemma supplies coverage, disjointness, $\sum_j a_j=S$, and all order
bounds. Theorem~\ref{J14:thm:factorization} and Corollary~\ref{J14:cor:cycles}
now give the claimed word and the native branch-cycle count.
\end{proof}

\begin{lemma}[New lemma: parametrizations and branch cycles from a finite root covering]
\label{J14:lem:root-cover}
Let $f\in\C\{\zeta,\xi\}$ be a reduced nonunit germ with
$f(\zeta,0)$ nonzero of order $S\geq1$. Let
$f_1,\ldots,f_b$ be pairwise nonassociated irreducible factors of $f$,
each occurring once. There are positive integers $a_1,\ldots,a_b$
with $\sum_j a_j=S$ and convergent holomorphic functions $h_j$ such
that the following assertions hold after shrinking all neighborhoods.
\begin{enumerate}
\item The map $F_j(s)=(h_j(s),s^{a_j})$ is pointed and injective.
Its pullback kernel is exactly $(f_j)$.
\item The zero set of $f$ is the union of these images, and their
punctured images are pairwise disjoint.
\item Over each sufficiently small $\xi\ne0$, the $j$th branch has
exactly the $a_j$ distinct sheets $h_j(s)$ with $s^{a_j}=\xi$.
Positive continuation once around $\xi=0$ permutes them in one
$a_j$-cycle. Consequently the complete root permutation has exactly
$b$ cycles.
\end{enumerate}
If, in addition, $f$ has multiplicity $S$ and $f_S(1,0)\ne0$, then
$\ord h_j\geq a_j$ for every $j$, with the zero function allowed.
\end{lemma}
\begin{proof}
Lemma~\ref{J15:lem:preparation}, with the coordinates named
$(\zeta,\xi)$ here, writes
\[
 f=uP,\qquad P(\zeta,\xi)=\zeta^S+
       A_1(\xi)\zeta^{S-1}+\cdots+A_S(\xi),
 \qquad A_l(0)=0,
\]
where $u$ is a unit. It proves also that $(f,\xi)$ has finite
colength. The existing declaration
\decl{PlaneCurve.finite_first_partial_polar_of_squarefree_axis_finite}
therefore makes $(f,f_\zeta)$ a finite intersection.
The product rule and invertibility of $u$ give
$(f,f_\zeta)=(P,P_\zeta)$, as proved in Lemma~\ref{J15:lem:unit}.
The supplied declaration
\decl{Germ.resultant_ne_zero_of_hasFiniteIntersection}
and the monic identity \decl{Polynomial.resultant_deriv} show that
the discriminant $D(\xi)$ of $P$ is not the zero germ. In the resultant
declaration take the degree parameter $S-1$, since $P$ is distinguished
monic of positive degree $S$.
Thus $D(\xi)=\xi^d v(\xi)$ with $v(0)\ne0$ for some $d\geq0$.
There is a disc $|\xi|<\eta$ on which the roots are distinct whenever
$\xi\ne0$.

Each restriction $f_j(\zeta,0)$ is nonzero, because its product with
the restrictions of the other factors and a unit is $f(\zeta,0)$.
It has positive order $a_j$, since $f_j$ is pointed. Orders in this
one-variable product add, so $\sum_j a_j=S$.
Apply the same proved preparation lemma to each factor to obtain
$f_j=u_jQ_j$, where $Q_j$ is distinguished monic of degree $a_j$.
For each fixed small $\xi$, the zeros of $P$ and of $\prod_jQ_j$
in the preparation disc agree with multiplicity. Both polynomials are
monic of degree $S$, so
\begin{equation}\label{J14:eq:root-cover-factor}
 P=\prod_{j=1}^bQ_j.
\end{equation}
The equality holds coefficientwise as an identity of holomorphic
functions. Every $Q_j$ is irreducible as an analytic germ, since it
differs from $f_j$ by a unit.

We give the continuation and extension arguments explicitly.
A simple root of a holomorphic polynomial has a unique local
holomorphic continuation. To prove this assertion near a root
$(\zeta_0,\xi_0)$, subtract its nonzero $\zeta$ derivative from
the equation and rewrite it as a fixed point equation for
$\zeta-\zeta_0$. On a small closed polydisc its derivative in that
variable has norm less than one. Iteration from zero converges
uniformly on a smaller polydisc to a holomorphic fixed point, which
is unique. All roots remain in a fixed disc, by the preparation
construction. Around any punctured base point, the finitely many
distinct roots have disjoint such neighborhoods, each mapping
bijectively to a smaller base neighborhood. Thus they form a finite
covering of the punctured base disc.

Continuation along a path is obtained by subdividing its compact
parameter interval into these neighborhoods. Uniqueness on overlaps
makes the result independent of the subdivision. Continuation along
homotopic paths with fixed endpoints has the same endpoint. Indeed,
subdivide the parameter square of a homotopy into sufficiently small
rectangles whose images lie in covering neighborhoods. The unique
local lift around each rectangle has the same values at its corners,
and induction over the rectangles proves equality along the two
boundary paths. A loop in a punctured disc is homotopic to the
corresponding integer number of turns: divide it by its positive
radius and use a continuous argument on its parameter interval.
Interpolation of that argument with the linear argument of the same
total increase supplies the homotopy. Consequently continuation is
specified by one permutation $\pi$ of the roots, namely continuation
around one positive turn.

For each orbit $E$ of $\pi$, form the monic polynomial
\[
 Q_E(\zeta,\xi)=\prod_{y\in E(\xi)}(\zeta-y).
\]
Its coefficients are single-valued holomorphic functions on the
punctured disc. Local holomorphy follows from the simple-root
continuations, and orbit invariance makes the local definitions agree
after any loop. They are bounded because all roots are bounded.
A bounded holomorphic function on a punctured disc extends across
zero. To verify the needed removable-singularity statement, its
Laurent coefficient at degree $-k$, for $k\geq1$, is
$\frac1{2\pi\ii}\int_{|s|=r}h(s)s^{k-1}\,ds$.
If $|h|\leq M$, its modulus is at most $Mr^k$, and the coefficient
is independent of $r$. Letting $r$ decrease to zero makes it zero.
The Laurent series therefore has no negative terms.

All roots tend to zero with $\xi$. If this were false, a sequence of
bounded roots over base points tending to zero would have a subsequence
converging to a nonzero root of $P(\zeta,0)=\zeta^S$.
This is impossible. Hence the extended polynomial $Q_E$ is
distinguished, and every orbit of positive size contributes a nonunit
analytic factor of $P$.

The roots of each $Q_j$ form a union of $\pi$-orbits. In fact, a
continuation of a root of $Q_j$ remains a root of $Q_j$ locally, and
uniqueness continues this property along each path.
If it contained two or more orbits, the preceding construction would
factor $Q_j$ into two distinguished polynomials of positive degree.
Their equality with $Q_j$ follows on the punctured disc by comparing
their roots and leading coefficients, then at zero by continuity.
Both factors are nonunits, contradicting irreducibility of $Q_j$.
Therefore its $a_j$ roots form precisely one orbit. Different $j$
give disjoint orbits by \eqref{J14:eq:root-cover-factor} and the
nonvanishing discriminant. This proves the asserted orbit count
without assuming a normalization theorem.

Fix one root in the $j$th orbit. Pull the punctured covering back by
$\xi=s^{a_j}$. A positive turn of $s$ induces $a_j$ turns of $\xi$,
whose permutation is the identity on this orbit. The continuation
argument above therefore gives a single-valued holomorphic root
$h_j(s)$ on the punctured $s$ disc. It is bounded, so the same
Laurent argument extends it holomorphically across zero, with
$h_j(0)=0$. Replacing $s$ by $e^{2\pi\ii/a_j}s$ continues that root
once around $\xi=0$. Thus the values $h_j(\omega_{a_j}^l s)$, for
$0\leq l<a_j$, are all the distinct roots of $Q_j$ over $s^{a_j}$.
It follows that $F_j$ is injective on its punctured disc.
No punctured point maps to the origin because its second coordinate
is nonzero, so it is injective on the whole disc.

Equation~\eqref{J14:eq:root-cover-factor} now proves local coverage and
punctured disjointness. Indeed a zero over $\xi\ne0$ is one of the
listed roots, and a zero over $\xi=0$ has $\zeta=0$.
Distinct factors have no common root over the punctured disc.
The parametrizations annihilate their respective factors.

We also verify the exact pullback kernel. Let $g$ be a holomorphic
germ whose pullback by $F_j$ is zero. Shrink a product neighborhood
so that all representatives are defined and choose a circle
$|v|=R$ surrounding all roots of $Q_j$ with no root on the circle.
For $|\zeta|<R$, define
\[
 H(\zeta,\xi)=\frac1{2\pi\ii}\int_{|v|=R}
       \frac{g(v,\xi)}{Q_j(v,\xi)(v-\zeta)}\,dv
\]
and
\[
 R_j(\zeta,\xi)=\frac1{2\pi\ii}\int_{|v|=R}
 \frac{g(v,\xi)\bigl(Q_j(v,\xi)-Q_j(\zeta,\xi)\bigr)}
      {Q_j(v,\xi)(v-\zeta)}\,dv.
\]
The contour denominators are nonzero, so the functions are jointly
holomorphic by differentiation under the compact integral.
The quotient
$(Q_j(v,\xi)-Q_j(\zeta,\xi))/(v-\zeta)$ is a polynomial in
$\zeta$ of degree at most $a_j-1$. Thus $R_j$ has that degree bound.
The power-series form of Cauchy's formula, proved in the preparation
lemma of J15, gives
$g=Q_jH+R_j$.
For each punctured $\xi$, the hypothesis on $g\circ F_j$ says that
$g$ vanishes at all $a_j$ roots of $Q_j$. So does $R_j$.
A polynomial of degree at most $a_j-1$ with $a_j$ distinct roots is
zero. Hence all coefficients of $R_j$ vanish on the punctured disc
and then at zero. It follows that $g\in(Q_j)=(f_j)$.
The reverse inclusion was already proved by the root equation.
This establishes the kernel assertion.

Under the final transversality hypotheses, each $F_j$ annihilates $f$.
Lemma~\ref{J14:lem:slope} therefore gives $\ord h_j\geq a_j$.
\end{proof}

\begin{theorem}[New lemma: full J14 word statement for a transverse reduced germ]
\label{J14:thm:transverse-germ}
Let $f\in\C\{\zeta,\xi\}$ be reduced of exact multiplicity
$S\geq1$, and suppose $f_S(1,0)\ne0$. Let $b$ denote the number
\decl{PlaneCurve.irreducibleBranchCount} $f$.
For all sufficiently small positive base radii, the closed root braid
of $f=0$ over the positively oriented $\xi$ circle is represented,
through closed braids, by
\[
 \decl{fullTwistWord}(S)\,w,
 \qquad w\in\decl{PositiveAdjacentWord}(S).
\]
One may choose $w$ as the positive event word of the strict-transform
sheets $\zeta/\xi$, with one fibre projection chosen for all tangent
clusters. The native permutations of $w$ and of the whole word each
have exactly $b$ cycles. The cycle belonging to $f_j$ has length
$i_0(f_j,\xi)$.
\end{theorem}
\begin{proof}
The native finite reduced factorization is indexed by the irreducible
branch set, by \decl{associated_branchProduct_branchSet} and
\decl{irreducible_branchEquation_of_mem}, used in the supplied
Milnor--delta source. Its number of factors is the native branch
count. Transversality makes $f(\zeta,0)$ have exact order $S$.
Lemma~\ref{J14:lem:root-cover} supplies precisely the injective,
punctured-disjoint finite branch data of
Theorem~\ref{J14:thm:factorization}, with
$\sum_j a_j=S$ and $\ord h_j\geq a_j$.
Apply that theorem and Corollary~\ref{J14:cor:cycles}.
They give the positive strict-transform word, the full-twist prefix,
the identification through the same closed braid, and the stated
native cycle counts. The preparation calculation for $f_j$ gives
$\C\{\zeta,\xi\}/(f_j,\xi)\cong
\C\{\zeta\}/(\zeta^{a_j})$, whose dimension is $a_j$.
Consequently the cycle length equals $i_0(f_j,\xi)$.
\end{proof}

Theorem~\ref{J14:thm:transverse-germ} has no extra normalization input.
It proves the user's transverse, multiplicity-$S$ word statement for
every reduced germ with those hypotheses. The earlier
Corollary~\ref{J14:cor:germ} remains a useful direct application of an
already supplied parametrization. Its kernel and injectivity hypotheses
are now constructed in Lemma~\ref{J14:lem:root-cover}, so the omitted return
type of \decl{exists_intersectionParametrization_of_irreducible} is not
needed for this conclusion. The counterexample to the uncorrected M405
API remains unchanged.


For the smooth sphere conclusion, assume additionally M403(b): the normalized
small sphere link is smoothly ambiently isotopic, with its component
orientations, to this same root braid in the native solid-torus chart.
Assume M419(B) realizes the given native word by the standard collared
generator sheets used in Lemma~\ref{J14:lem:events}.
Composing that sphere isotopy with the ambient isotopies of
Lemma~\ref{J14:lem:ambient} proves the corrected sphere conclusion for $d_Sw$.
The identities specifying \decl{planeSolidTorusMap 0} must be checked in
those inputs. No identification of its index with the chosen base coordinate
is silently assumed here. Those two geometric obligations remain explicit
hypotheses and are not replaced by an assertion that an arbitrary path
homotopy is an ambient diffeotopy.

\section{Exceptional cases and source ledger}

If $S=1$ with transverse projection, the only sheet has no pairwise
crossing, both words $d_1,w$ are empty, and the permutation has one fixed
point. If $S=2$ and the curve is an ordinary node with distinct slopes,
the strict-transform word is empty and the original word is
$d_2=\sigma_0^2$. For the transverse cusp
$\zeta^2=\xi^3$, the strict-transform sheets are
$\pm\xi^{1/2}$, so $w=\sigma_0$ and the original word is
$\sigma_0^3$. Thus the sign agrees with the native positive convention
already in the smallest singular examples.

For an ordinary $S$-fold point, every branch is smooth with degree one and
all slopes are distinct. Lemma~\ref{J14:lem:projection} removes every
strict-transform crossing, so $w$ is empty and $b=S$.
Several branches with a common tangent require no additional argument:
their strict-transform differences have positive order in the common
parameter. Tangency of a strict-transform branch to the exceptional line
does not affect Lemma~\ref{J14:lem:log}.
Events with three or more sheets are treated by Lemma~\ref{J14:lem:events},
which counts each pair once. Choosing a regular initial angle prevents an
event from lying at the seam of the closed period.

The assertions involve complex analytic functions, integer covering degrees,
and native finite words. No Floer chain coefficients or homology theory are
introduced. In particular, no holomorphic Heegaard Floer result or
fibre-surface argument is a proof step.

\small
\begin{longtable}{>{\raggedright\arraybackslash}p{0.43\linewidth}
                  >{\raggedright\arraybackslash}p{0.15\linewidth}
                  >{\raggedright\arraybackslash}p{0.34\linewidth}}
Declaration or statement & Status & Exact role\\\hline\endhead
\decl{PositiveAdjacentWord}, \decl{PositiveAdjacentWord.toSigned}
& Existing & Finite lists and positive signs.\\
\decl{adjacentSwap}, \decl{SignedAdjacentWord.permutation},
\decl{SignedAdjacentWord.cycleCount}
& Existing & Native permutation multiplication and orbit count.\\
\decl{tendsto_mul_deriv_div_of_analytic_factor}
& Existing & The analytic logarithmic-derivative identity behind Lemma~\ref{J14:lem:log}.\\
\decl{exists_radius_planeBranchAngularDifference_projected_crossings_pos}
& Existing, recorded & Native crossing-sign convention, not a multibranch theorem.\\
Lemma~\ref{J14:lem:log}, Theorem~\ref{J14:thm:multisheet}
& New & Uniform finite-pair estimates, order-zero pairs, and positivity.\\
Lemma~\ref{J14:lem:events}
& New & Multiple-event resolution to positive adjacent words.\\
Lemmas~\ref{J14:lem:common} and \ref{J14:lem:projection}
& New & Common ramification, several tangent points, and choice of real projection.\\
Lemma~\ref{J14:lem:rotation}
& New & Literal word $d_S$, length $S(S-1)$, positive chirality, identity permutation.\\
Lemmas~\ref{J14:lem:product} and \ref{J14:lem:ambient}
& New & Explicit braid homotopy and its ambient realization.\\
\decl{braid_permutation_eq_conjugate_inverse_rotation}
& Existing & Verification of the inverse-conjugacy convention used in \eqref{J14:eq:cycles}.\\
Theorem~\ref{J14:thm:factorization}, Corollary~\ref{J14:cor:cycles}
& New & Full twist and exact branch cycles for the displayed finite branch data.\\
\decl{Germ.exists_linear_equiv_regularInFirstCoordinateOfOrder}
& Existing, recorded & Choice of transverse coordinates for a germ of exact order $S$.\\
Lemma~\ref{J14:lem:slope}
& New & Transversality makes division by $\xi$ holomorphic on the common cover.\\
Lemma~\ref{J15:lem:local-principle}
& New, proved in J15 & The local argument principle used in Lemma~\ref{J14:lem:persistence}.\\
Lemmas~\ref{J14:lem:persistence} and \ref{J14:lem:packaging}
& New & Kernel ideals give the branch polynomials, full local coverage, punctured disjointness, and exact degree sum.\\
\decl{associated_branchProduct_branchSet},
\decl{irreducible_branchEquation_of_mem}
& Existing, used in supplied source & Finite reduced factorization indexed by native branches.\\
\decl{PointedMapGerm.exists_injective_power_representative_sheetRotation}
& Existing & Power normalization preserves injectivity and the kernel after source equivalence.\\
\decl{exists_intersectionParametrization_of_irreducible}
& Existing, recorded & Optional alternate source of parametrizations. It is not needed for Theorem~\ref{J14:thm:transverse-germ}.\\
Corollary~\ref{J14:cor:germ}
& New & Direct application when a kernel parametrization is already supplied.\\
Lemma~\ref{J14:lem:root-cover}
& New & Constructs all parametrizations from the finite root covering, proves their exact kernels, and identifies monodromy orbits with irreducible factors.\\
Theorem~\ref{J14:thm:transverse-germ}
& New & The complete transverse reduced-germ word statement, with no extra normalization assumption.\\
\decl{PlaneCurve.finite_first_partial_polar_of_squarefree_axis_finite},
\decl{Germ.resultant_ne_zero_of_hasFiniteIntersection},
\decl{Polynomial.resultant_deriv}
& Existing & The prepared polynomial has nonzero discriminant.\\
Preparation and unit lemmas of \decl{J15-braid-length.tex}
& New, proved in J15 & Produce distinguished polynomials for the equation and its factors, and identify their polar pair ideals.\\
M403(b), M419(B)
& Unproved input & Required only for the stated conditional native sphere conclusion.\\
Lemma~\ref{J14:lem:two}, Proposition~\ref{J14:prop:false}
& New & Counterexample to the uncorrected exact API, including its sphere formulation.\\
\end{longtable}
\normalsize

The supplied reference by Fomin, Pylyavskyy, Shustin, and Thurston,
\emph{Morsifications and mutations}, Lemma 13.5, motivates the decomposition
by a full twist. Its one-sentence proof is not used as a replacement for
Lemmas~\ref{J14:lem:common}--\ref{J14:lem:product}. The source conventions and
acceptance statement are those of the attached \decl{combined_01-10.md},
J14 and records M405, M406, M1153, together with the supplied native word and
analytic listings. No Lean code is written. No finite computation is used
as a proof.


\renewcommand{\bibname}{References}
\begin{thebibliography}{9}

\bibitem{OSSbook}
P.~Ozsv\'ath, A.~I.~Stipsicz, and Z.~Szab\'o,
\emph{Grid Homology for Knots and Links}, preliminary version supplied
with the request. Definitions~3.1.6 and 3.1.7 and Section~3.2 are used
only for the elementary grid-move conventions.


\bibitem{MortonBraid}
H.~R.~Morton,
\emph{The multivariable Alexander polynomial for a closed braid},
in H.~Nencka, editor, \emph{Low-dimensional topology},
Contemporary Mathematics 233, American Mathematical Society (1999),
167--172. Author's manuscript:
\url{https://www.liverpool.ac.uk/~su14/papers/multiburau.pdf}.
Accessed 1 November 2026. The Introduction records the classical
determinant formula. Lemma~\ref{J15:lem:alexander-determinant} proves
its precise coefficient-field version used here.

\end{thebibliography}

\end{document}
